\section*{第 1 章 集合与逻辑(必修 1)}

\section*{知识梳理*}

\section*{一、集合及其运算}

1.集合的特性:集合中的元素具有确定性、无序性、互异性;

2. 列举法:将集合中的元素一一列举出来写在大括号内表示集合的方法;

3. 描述法:将集合中元素的通性描述出来写在大括号内表示集合的方法；通式: $\{ x \mid  x$ 满足 $P\}$ ；读法:满足条件 $P$ 的 $x$ 的集合；

4. 空集(《》)是不含任何元素的集合；它是任何集合的子集，是任何非空集合的真子集；

5. “ $\in  , \notin$ ”表示元素与集合间的从属关系; “ $\subset  , \subseteq$ ”表示集合与集合间的包含关系.

6. 给出下列条件:

(1)集合 $A$ 中任何一个元素都是集合 $B$ 中的元素;

(2)集合 $B$ 至少存在一个元素不在集合 $A$ 中；

(3)集合 $B$ 中任何一个元素都是集合 $A$ 中的元素.

如果集合 $A$ 、 $B$ 满足 \ding{172} ，则 $A$ 是 $B$ 的子集( $A \subseteq  B$ )；

如果集合 $A$ 、 $B$ 满足 \ding{172}  \ding{173} ，则 $A$ 是 $B$ 的真子集( $A \subset  B$ )；

如果集合 $A$ 、 $B$ 满足 \ding{172}  \ding{174} ，则 $A$ 与 $B$ 是相等的集合 $\left( {A = B}\right)$ .

注意: \ding{172}  $A \subseteq  B$ ，讨论时别忘 $A = \varnothing$ 的情况；考察集合的关系借助韦恩图；

 \ding{173} 几个特殊数集之间的包含关系为: ${N} \subset  \mathbf{Z} \subset  \mathbf{Q} \subset  \mathbf{R} \subset  \mathbf{C}$ (本书中 ${\mathbf{N}}^{ * }$ 表示正整数集);

 \ding{174} 正确区分 $\varnothing \text{ 、 }\{ 0\} \text{ 、 }\{ \varnothing \}  : \varnothing$ 是不含任何元素的集合，即空集. \{0\}是含有一个元素 0 的集合，它不是空集， 因为它有一个元素,这个元素是 $0.\{ \varnothing \}$ 是含有一个元素 $\varnothing$ 的集合,其关系为: $\varnothing  \subset  \{ 0\} ,\varnothing  \subset  \{ \varnothing \} ,\varnothing  \in  \{ \varnothing \}$ .

7. 集合的含义:

(1) $A = \{ x \mid  y = f\left( x\right) \}$ 表示函数的定义域；

(2) $B = \{ y \mid  y = f\left( x\right) \}$ 表示函数的值域；

(3) $C = \{ \left( {x,y}\right)  \mid  y = f\left( x\right) \}$ 表示方程 $y = f\left( x\right)$ 的解的集合，或表示曲线上的点的集合……

8. 集合的运算: $A \cap  B = \{ x \mid  x \in  A$ 且 $x \in  B\}$ ； $A \cup  B = \{ x \mid  x \in  A$ 或 $x \in  B\}$ ； $\overline{A} = \{ x \mid  x \in  U$ ，且 $x \notin  A\}$ (其中 $U$ 表示全集).

9. 运算性质:

$A \subseteq  B$ ,且 $B \subseteq  C \Rightarrow  A \subseteq  C$ ;

$\varnothing  \subseteq  A \cap  B = B \cap  A \subseteq  A$ (或 $B$ ) $\subseteq  A \cup  B = B \cup  A$ ;

$\bar{A} \cap  \bar{B} = \overline{A \cup  B};\bar{A} \cup  \bar{B} = \overline{A \cap  B};\overline{\varnothing } = U;\bar{U} = \varnothing$ .

10.几个结论:

$B \cup  \bar{A} = U \Leftrightarrow  A \subseteq  B;\bar{B} \subseteq  \bar{A} \Leftrightarrow  A \cap  \bar{B} = \varnothing$ ;

$A \subseteq  B \Leftrightarrow  A \cap  B = A \Leftrightarrow  A \cup  B = B \Leftrightarrow  A \cap  B \supseteq  A \Leftrightarrow  A \cup  B \subseteq  B$ .

11. 正整数 (整数) 分类: 被 2 整除与否可分为 ${2k} - 1,{2k}\left( {k \in  \mathbf{N}\text{ (或 }\mathbf{Z}}\right) ,k \geq  1$ );

被 3 整除与否可分为 ${3k} - 2,{3k} - 1,{3k}\left( {k \in  \mathbf{N}\text{ (或 }\mathbf{Z}}\right) ,k \geq  1$ );

被 4 整除与否可分为 ${4k} - 3,{4k} - 2,{4k} - 1,{4k}\left( {k \in  \mathbf{N}\text{ (或 }\mathbf{Z}}\right) ,k \geq  1$ ); 其余依此类推.

\HRule

* 注 本章节所属知识模块是“预备知识”.

\HRule

12. $n$ 个元素的子集有 ${{C}}_{n}^{0} + {{C}}_{n}^{1} + {{C}}_{n}^{2} + \cdots  + {{C}}_{n}^{n} = {2}^{n}$ 个;

真子集有 ${{C}}_{n}^{0} + {{C}}_{n}^{1} + {{C}}_{n}^{2} + \cdots  + {{C}}_{n}^{n - 1} = {2}^{n} - 1$ 个;

非空子集有 ${{C}}_{n}^{1} + {{C}}_{n}^{2} + {{C}}_{n}^{3} + \cdots  + {{C}}_{n}^{n} = {2}^{n} - 1$ 个;

非空真子集有 ${{C}}_{n}^{1} + {{C}}_{n}^{2} + {{C}}_{n}^{3} + \cdots  + {{C}}_{n}^{n - 1} = {2}^{n} - 2$ 个.

13.* 容斥原理:记集合 $M$ 中的元素个数为 $\operatorname{Card}\left( M\right)$ ，则

$\operatorname{Card}\left( {A \cup  B}\right)  = \operatorname{Card}\left( A\right)  + \operatorname{Card}\left( B\right)  - \operatorname{Card}\left( {A \cap  B}\right)$ (可推广到有限个集合).

\section*{二、命题与充要条件}

1. 可以判断真假的语句叫命题. 其一般形式是:“如果……，那么 ······”(或“若 $\alpha$ ，则 $\beta$ ”)，命题有真命题和假命题两类.

要判断命题“若 $\alpha$ ，则 $\beta$ ”是假命题，只要存在一个满足条件 $\alpha$ 但不满足结论 $\beta$ 的对象就可以了. 要判断命题“若 $\alpha$ ,则 $\beta$ ”是真命题,就需要证明所有满足条件 $\alpha$ 的对象都满足结论 $\beta$ ; 若要否定命题,切记举反例.

2. 证明命题成立的方法:

综合法:从已知的条件出发，根据已知的定义、公理、定理和推论进行严密的逻辑推理，逐步得到结论的方法. 即 “由因导果”. 连接词: “ “ …… ， ， …… ... ”。

分析法:从需要的结论出发，一步步寻找结论成立的充分条件，直到找到的条件是已知的定义、公理、定理和推论或题设条件的方法. 即“执果索因”. 连接词:“要证……，只要证……”.

反证法:当证明“若 $\alpha$ ，则 $\beta$ ”感到困难时，可从它的反面入手证明.

证题步骤: \ding{172} 假设命题的结论 $\beta$ 不成立； \ding{173} 从这个假设出发，经过正确的推理论证，得出矛盾； \ding{174} 由矛盾判断假设不成立，从而肯定结论正确.

矛盾的来源: \ding{172}  与原命题的条件； \ding{173}  与假设； \ding{174}  恒真命题.

适用于待证命题的结论涉及“不可能”“不是”“至少”“至多”“唯一”等否定词.

\begin{center}
\begin{tabular}{|c|c|c|c|c|c|c|c|c|}
\hline
正面词语 & 等于 & \multicolumn{2}{|c|}{大于} & \multicolumn{2}{|c|}{小于} & \multicolumn{2}{|c|}{是} & 都是 \\
\cline{1-9}
否定 & 不等于 & \multicolumn{2}{|c|}{小于或等于} & \multicolumn{2}{|c|}{大于或等于} & \multicolumn{2}{|c|}{不是} & 不都是 \\
\cline{1-9}
正面词语 & \multicolumn{2}{|c|}{至少一个} & \multicolumn{2}{|c|}{任意的} & \multicolumn{2}{|c|}{所有的} & \multicolumn{2}{|c|}{至多有 $n$ 个} \\
\cline{1-9}
否定 & \multicolumn{2}{|c|}{一个都没有} & \multicolumn{2}{|c|}{存在一个} & \multicolumn{2}{|c|}{存在一个} & \multicolumn{2}{|c|}{至少 $n + 1$ 个} \\
\cline{1-9}
\hline
\end{tabular}
\end{center}

3. 充要条件及其判定 $\left( {A\text{ 、 }B\text{ 分别可以是命题，也可以是一个“开语句”): }}\right)$

若 $A \Leftrightarrow  B$ (或 $\bar{B} \Rightarrow  \bar{A}$ ),则 $A$ 是 $B$ 的充分条件;

若 $A \Leftarrow  B$ (或 $\overline{\overline{A \Rightarrow  B}}$ ),则 $A$ 是 $B$ 的必要条件;

若 $A \Leftrightarrow  B$ (或 $\bar{A} \Leftrightarrow  \bar{B}$ ),则 $A$ 是 $B$ 的充分必要条件,简称充要条件.

4. 充要条件的判断中,分清“条件”和“结论”是关键,其次是看推出符号的箭头指向.

5. 若命题 $\alpha$ 成立的范围是集合 $A$ ，命题 $\beta$ 成立的范围是集合 $B\texttt{ ( }A\text{ 、 }B\texttt{ 非空 }\texttt{ ),则 }{}^{\alpha }A\text{ 、 }B$ 之间的包含关系” 与“充要条件”之间有如下的关系:

若 $A \subseteq  B$ ,则 $\alpha$ 是 $\beta$ 的充分条件;

若 $A \subset  B$ ,则 $\left\{  \begin{array}{l} \alpha  \Rightarrow  \beta , \\  \beta  \nRightarrow  \alpha , \end{array}\right.$ 即 $\alpha$ 是 $\beta$ 的充分非必要条件;

若 $B \subseteq  A$ ,则 $\alpha$ 是 $\beta$ 的必要条件;

\begin{center}
\includegraphics[width=0.2\textwidth]{images/bo_d7tk7oc91nqc73875fu0_1_1166_1790_304_248_0.jpg}
\end{center}
 

若 $B \subset  A$ ,则 $\left\{  \begin{array}{l} \alpha  \Leftarrow  \beta , \\  \alpha  \nRightarrow  \beta , \end{array}\right.$ 即 $\alpha$ 是 $\beta$ 的必要非充分条件;

若 $A = B$ ,则 $\alpha$ 是 $\beta$ 的充要条件.

6. 一般地, 充要条件的范围具有如右图的包含关系.

7. 集合的运算是“交、并、补”; 命题的运算是“或、且、非”，两者都是逻辑运算.

\section*{$\text{ § }{1.1}$ 集合及其运算}

\section*{典型例题精讲}

【例 1】设集合 $A = \{ x \mid   - 2 \leq  x \leq  5\} ,B = \{ x \mid  m + 1 \leq  x \leq  {2m} - 1\}$ .

(1)若 $A \cap  B = B$ ，求实数 $m$ 的取值范围；

(2)当 $x \in  \mathbf{R}$ 时，不存在 $x$ 使得 $x \in  A$ 与 $x \in  B$ 同时成立，求实数 $m$ 的取值范围.

【例 2】已知 $a\text{ 、 }m \in  \mathbf{R}$ ,集合 $A = \left\{  {x \mid  {x}^{2} - {3x} + 2 = 0}\right\}  ,B\left\{  {x \mid  {x}^{2} + 2\left( {m + 1}\right) x + {m}^{2} - 5 = 0,x \in  \mathbf{R}}\right\}  ,C =  \left\{  {x \mid  {x}^{2} - {ax} + a - 1 = 0,x \in  \mathbf{R}}\right\}  .$

(1)若 $A \cup  C = A$ ，求实数 $a$ 的值；

(2)若 $A \cap  B = B$ ，求实数 $m$ 的取值范围.

【例 3】设 $f\left( x\right)  = {x}^{2} + {ax} + b\left( {a\text{ 、 }b \in  \mathbf{R}}\right) ,A = \{ x \mid  f\left( x\right)  = x\} ,B = \left\{  {x\left| {f\left\lbrack  {f\left( x\right) }\right\rbrack   = x}\right. }\right\}$ .

(1)证明 $A \subseteq  B$ ；

(2)若 $A = \{  - 1,3\}$ ，求集合 $B$ .

【例 4】已知实数集 $A = \left\{  {{a}_{1},{a}_{2},\cdots ,{a}_{n}}\right\}  \left( {n \in  \mathbf{N},n \geq  3}\right)$ ，定义 $\varphi \left( A\right)  = \left\{  {{a}_{i}{a}_{j} \mid  {a}_{i}\text{ 、 }{a}_{j} \in  A\left( {i \neq  j}\right) }\right.$ .

(1)若 $A = \{  - 2,0,1,2\}$ ，求 $\varphi \left( A\right)$ ；

(2)若 $\varphi \left( A\right)  = \{ 0, - 6, - 8, - {12},{12},{18},{24}\}$ ，求集合 $A$ .

【例 5】已知函数 $f\left( x\right)  = {x}^{2} - {2x} - {a}^{2} + {2a}\left( {a \in  \mathbf{R}}\right)$ ,集合 $A = \{ x \mid  f\left( x\right)  \leq  0\}$ .

(1)若集合 $A$ 中有且仅有 3 个整数，求实数 $a$ 的取值范围；

(2)集合 $B = \{ x \mid  f\left( {f\left( x\right)  + b}\right)  \leq  0\}$ ，若存在实数 $a \leq  1$ ，使得 $A \subseteq  B$ ，求实数 $b$ 的取值范围.

\section*{四基达标精练}

\section*{一、填空题}

1. 已知集合 $A = \{ 1,2,4\}$ ，集合 $B = \{ 2,4,5\}$ ，则 $A \cap  B =$ \_\_\_\_\_.

2. 集合 $A = \{ \left( {x,y}\right)  \mid  x + y = 4$ 且 $x - y = 2\}$ ，用列举法表示为\_\_\_\_\_.

3. 已知集合 $A = \{ x \mid  x \leq  1\} ,B = \{ x \mid  x \geq  a\}$ ，且 $A \cup  B = \mathbf{R}$ ，则实数 $a$ 的取值范围是\_\_\_\_\_.

\begin{center}
\includegraphics[width=0.2\textwidth]{images/bo_d7tk7oc91nqc73875fu0_4_1168_465_271_175_0.jpg}
\end{center}
 

(第 5 题图)

4. 全集 $U = \{ 2,3,5\} ,A = \{ \left| {a - 5}\right| ,2\} ,\overline{A} = \{ 5\}$ ,则 $a =$ \_\_\_\_\_.

5. 已知全集 $U = \mathbf{R}$ ,集合 $A = \{ x \mid  x <  - 1$ 或 $x > 4),B = \{ x \mid   - 2 \leq  x \leq  3)$ ,那么阴影部分表示的集合为\_\_\_\_\_.

6. 已知 $M = \{ \left( {x,y}\right)  \mid  {2x} + {3y} = {4320},x\text{ 、 }y \in  \mathbf{N}\} ,P = \{ \left( {x,y}\right)  \mid  {4x} - {3y} = 1$ , $x$ 、 $y \in  \mathbf{N}\}$ ，则其中的有限集是\_\_\_\_\_.

7. 已知集合 $A = \left\{  {x\left| {a{x}^{2} - {3x} + 2 = 0}\right. }\right\}$ 中有且仅有一个元素，则 $a$ 的取值组成的集合为\_\_\_\_\_.

8. 集合 $A = \{ \left( {x,y}\right) \left| {y = a}\right| x \mid  \} ,B = \{ \left( {x,y}\right)  \mid  y = x + 1\}$ ,若 $A \cap  B$ 为单元素集，实数 $a$ 的取值范围为\_\_\_\_\_.

9. 设数集 $M = \left\{  {x\left| {m \leq  x \leq  m + \frac{3}{4}}\right. }\right\}  ,N = \left\{  {x\left| {n - \frac{1}{3} \leq  x \leq  n}\right. }\right\}$ ,且 $M\text{ 、 }N$ 都是集合 $\{ x \mid  0 \leq  x \leq  1\}$ 的子集,如果把 $b - a$ 叫做集合 $\{ x \mid  a \leq  x \leq  b\}$ 的“长度”,那么集合 $M \cap  N$ 的长度的最小值是\_\_\_\_\_.

\section*{二、选择题}

10. 设集合 $P = \{ x \mid  0 \leq  x \leq  \sqrt{2}\} ,m = \sqrt{3}$ ，则下列关系中正确的是 ( )

A. $m \subseteq  P$ . B. $m \text{ ⊄ } P$ . C. $m \in  P$ . D. $m \notin  P$ .

11. 由 ${a}^{2} + 1\text{ 、 }a + 3\text{ 、 }a$ 组成的集合含有元素 2,则实数 $a$ 的可能取值的个数是 ( )

A. 1 个. B. 2 个. C. 3 个. D. 4 个.

12. 设 $a$ 、 $b$ 、 $c \in  \mathbf{R}$ ,函数 $f\left( x\right)  = \left( {x + a}\right) \left( {{x}^{2} + {bx} + c}\right) ,g\left( x\right)  = \left( {{ax} + 1}\right) \left( {c{x}^{2} + {bx} + 1}\right)$ . 记集合 $S = \{ x \mid  f\left( x\right)  = 0$ , $x \in  \mathbf{R}\} ,T = \{ x \mid  g\left( x\right)  = 0,x \in  \mathbf{R}\}$ . 若 $\left| S\right| \text{ 、 }\left| T\right|$ 分别为集合 $S\text{ 、 }T$ 的元素个数,则下列结论不可能的是 ( )

A. $\left| S\right|  = 1,\left| T\right|  = 0$ . B. $\left| S\right|  = 1,\left| T\right|  = 1$ . C. $\left| S\right|  = 2,\left| T\right|  = 2$ . D. $\left| S\right|  = 2,\left| T\right|  = 3$ .

\section*{三、解答题}

13. 已知集合 $A = \{ x \mid  x \geq  3\} ,B = \{ x \mid  1 \leq  x \leq  7\} ,C = \{ x \mid  x \geq  a - 1\}$ .

(1)求 $A \cap  B$ 和 $A \cup  B$ ；

(2)若 $C \cup  A = A$ ，求实数 $a$ 的取值范围.

14. 已知 $a\text{ 、 }x \in  \mathbf{R}$ ,集合 $A = \{ 2,4,{x}^{2} - {5x} + 9\} ,B = \left\{  {3,{x}^{2} + {ax} + a}\right\}  ,C = \left\{  {{x}^{2} + \left( {a + 1}\right) x - 3,1}\right\}$ .

(1)若 $2 \in  B$ ，且 $B \subseteq  A$ ，求 $a$ 和 $x$ 的值；

(2)若集合 $B = C$ ，求 $a$ 和 $x$ 的值.

15. 已知全集是 $\mathbf{R}$ ,函数 $f\left( x\right)  = \sqrt{x - 3} - \frac{1}{\sqrt{7 - x}}$ 的定义域为集合 $A,B = \{ x \mid  2 < x < {10},x \in  \mathbf{Z}\} ,C = \{ x \mid  x < a$ 或 $x > a + 1\}$ .

(1)求 $A$ 和 $\overline{A} \cap  B$ ；

(2)若 $A \cup  C = \mathbf{R}$ ，求实数 $a$ 的取值范围.

16. 设 $A$ 为集合 $M$ 的子集，且 $A = \left\{  {{a}_{1},{a}_{2},\cdots ,{a}_{n}}\right\}  \left( {n \in  \mathbf{N},n \geq  2}\right)$ ，若 ${a}_{1} + {a}_{2} + \cdots  + {a}_{n} = {a}_{1} \cdot  {a}_{2} \cdot  \cdots  \cdot  {a}_{n}$ ，则， 称 $A$ 为集合 $M$ 的 $n$ 元“好集”.

(1)写出实数集 $\mathbf{R}$ 的一个二元“好集”；

(2)求出正整数集的所有三元“好集”.

\section*{$\text{ § }{1.2}$ 充要条件与反证法}

\section*{典型例题精讲}

【例 1】已知集合 $A = \left\{  {y\left| {y = {x}^{2} - \frac{3}{2}x + 1}\right. ,x \in  \left\lbrack  {\frac{3}{4},2}\right\rbrack  }\right\}  ,B = \left\{  {x\left| {x + {m}^{2} \geq  1}\right. }\right\}$ .

(1)求集合 $A$ ；

(2)若 $p : x \in  A;q : x \in  B$ 且 $p$ 是 $q$ 的充分条件，求实数 $m$ 的取值范围.

【例 2】设 $p$ : 实数 $x$ 满足 ${x}^{2} - {4ax} + 3{a}^{2} \leq  0\left( {a > 0}\right) ,q$ : 实数 $x$ 满足 $\frac{x - 3}{x - 2} < 0$ .

(1)若 $a = 1,p$ 且 $q$ 是真命题，求实数 $x$ 的取值范围；

(2)若 $p$ 是 $q$ 的必要非充分条件，求实数 $a$ 的取值范围.

【例 3】已知函数 $f\left( x\right)  =  - {x}^{2} + {bx} + c$ ,则 “ $f\left( {f\left( \frac{b}{2}\right) }\right)  > 0$ ” 是 “方程 $f\left( x\right)  = 0$ 有两个不同实数解且方程 $f\left( {f\left( x\right) }\right)  = 0$ 恰有两个不同实数解”的 ( )

A. 充要条件. B. 充分非必要条件. C. 必要不充分条件. D. 既非充分也非必要条件.

【例 4】已知定义域为 $D$ 的函数 $f\left( x\right)$ ,若存在实数 $a$ ,使得对任意 ${x}_{1} \in  D$ ,都存在 ${x}_{2} \in  D$ 满足 $\frac{{x}_{1} + f\left( {x}_{2}\right) }{2} = a$ ,则称函数 $f\left( x\right)$ 具有性质 $P\left( a\right)$ .

(1)判断下列函数 $f\left( x\right)  = {2}^{x}$ 和 $g\left( x\right)  = {\log }_{2}x\left( {0 < x < 1}\right)$ 是否具有性质 $P\left( 0\right)$ ，说明理由；

(2)若函数 $f\left( x\right)$ 的定义域为 $D$ ，且具有性质 $P\left( 1\right)$ ，则 “ $f\left( x\right)$ 存在零点” 是 “ $2 \in  D$ ” 的什么条件，说明理由.

【例 5】证明下列命题是真命题:

(1)已知 $n \in  \mathbf{Z}$ . 用反证法证明:若 ${n}^{3}$ 是奇数，则 $n$ 是奇数；

(2)已知 $a$ 、 $b \in  \left( {0, + \infty }\right)$ ， $a + b > 2$ . 证明: $\frac{1 + a}{b}$ 与 $\frac{1 + b}{a}$ 中至少有一个小于 2 .

\section*{四基达标精练}

\section*{一、填空题}

1. 命题:“若 $a < b$ ，则 $a{c}^{2} < b{c}^{2}$ ”是\_\_\_\_\_命题(填“真”或“假”).

2. 设 $m\text{ 、 }n$ 是整数，则 “ $m\text{ 、 }n$ 均为偶数” 的否定是\_\_\_\_\_.

3. 已知 “ $p : x > 1$ ” 是 “ $q : {x}^{2} > x$ ”，则 $p\text{ 、 }q$ 的推出关系用 “ $\Rightarrow$ 、 $\Leftarrow$ 、 $\Leftrightarrow$ ”中的一个表示为\_\_\_\_\_.

4. 一元二次方程 $a{x}^{2} + {bx} + c = 0$ 有一正一负根的充要条件是\_\_\_\_\_.

5. 否定命题 “若 $a > b$ ，则 $\frac{1}{a} < \frac{1}{b}$ ” 的一个反例可以是\_\_\_\_\_.

6. 若 “ $a$ 满足 $\frac{{2a} - 1}{a + 1} \leq  1$ ” 为真命题，则实数 $a$ 的取值范围是\_\_\_\_\_.

7. 设条件 $p : \left| x\right|  \leq  m\left( {m > 0}\right) ,q :  - 1 \leq  x \leq  4$ ，若 $p$ 是 $q$ 的充分条件，则 $m$ 的最大值为\_\_\_\_\_.

8. 已知 $A = \{ x \mid  \left| {x - a}\right|  < 4\} ,B = \{ x \mid  \left( {x - 1}\right) \left( {2 - x}\right)  > 0\}$ ,若 $x \in  A$ 是 $x \in  B$ 的必要非充分条件,则实数 $a$ 的取值范围是\_\_\_\_\_.

9. 已知命题 $p :$ 对任意 $x \in  \left\lbrack  {1,2}\right\rbrack$ ， ${x}^{2} - a \geq  0$ 都成立；命题 $q :$ 存在 $x \in  \mathbf{R}$ ，使得成立 ${x}^{2} + {2ax} + 2 - a = 0$ ， 若命题“ $p$ 且 $q$ ”是真命题，则实数 $a$ 的取值范围为\_\_\_\_\_.

\section*{二、选择题}

10. 设 $a\text{ 、 }b\text{ 、 }c \in  {\mathbf{R}}^{n}{b}^{2} - {4ac} > 0$ ” 是 “方程 $a{x}^{2} + {bx} + c = 0$ 有实根” 的 ( )

A. 充分非必要条件. B. 必要非充分条件. C. 充要条件. D. 既非充分也非必要条件.

11. " $x + y \neq  3$ " 是 " $x \neq  1$ 或 $y \neq  2$ "的 ( )

A. 充分非必要条件. B. 必要非充分条件. C. 充要条件. D. 既非充分也非必要条件.

12. ” ${\sin }^{2}\alpha  + {\cos }^{2}\beta  = 1$ ” 是 “ $\alpha  = \beta$ ” 的充要条件,则 $\alpha$ 、 $\beta$ 满足的条件是 ( )

A. $\alpha  = \beta$ . B. $\alpha  - \beta  = {2k\pi }\left( {k \in  \mathbf{Z}}\right)$ . C. $\alpha  - \beta  = {k\pi }\left( {k \in  \mathbf{Z}}\right)$ . D. $\alpha  - \beta  = \frac{k\pi }{2}\left( {k \in  \mathbf{Z}}\right)$ .

\section*{三、解答题}

13. 已知命题 $p$ : 关于 $x$ 的方程 ${x}^{2} + \left( {m - 2}\right) x + 1 = 0$ 无实根; 命题 $q$ : 关于 $x$ 的方程 ${x}^{2} + {mx} + 1 = 0$ 有两个不等的负实根.

(1)若命题 $p$ 为假命题，求实数 $m$ 的取值范围；

(2)若命题 $p$ 、 $q$ 中有且仅有一个是真命题，求实数 $m$ 的取值范围.

14. 证明下列命题是真命题:

(1)已知 $x$ 为正数， $a =  - x + \frac{10}{x}$ ， $b = {5x} - \frac{1}{x}$ ，则 $a$ 、 $b$ 中至少有一个不小于 6；

(2)已知实数 $a$ 、 $b$ 满足 ${a}^{3} + {b}^{3} = 2$ ，则 $a + b \leq  2$ .

15. 设集合 $A = \left\{  {x \mid  {x}^{2} + \left( {p + 2}\right) x + 1 = 0,x \in  \mathbf{R}}\right\}$ .

(1)求 $A = \varnothing$ 的充要条件；

(2)求 $A \cap  {\mathbf{R}}^{ + } = \varnothing$ 的充要条件，其中 ${\mathbf{R}}^{ + }$ 表示正实数集.

16. 设 $a$ 、 $b$ 、 $c$ 为 $\bigtriangleup {ABC}$ 的三条边长,方程 ${x}^{2} + {2ax} + {b}^{2} = 0$ 与 ${x}^{2} + {2cx} - {b}^{2} = 0$ .

(1)求方程 ${x}^{2} + {2ax} + {b}^{2} = 0$ 有实数根的充要条件；

(2)求这两个方程有公共根的充要条件.

\section*{第 2 章 等式与不等式(必修 1)}

\section*{知识梳理*}

1. 不等式的依据:实数的有序性，即 $a > b \Leftrightarrow  a - b > 0;a < b \Leftrightarrow  a - b < 0;a = b \Leftrightarrow  a - b = 0$ ；

2. 等式与不等式的性质:

\begin{center}
\begin{tabular}{|c|c|c|}
\hline
性质 & 等式 & 不等式 \\
\cline{1-3}
意义 & $a = b \Leftrightarrow  a - b = 0$ & $a > b \Leftrightarrow  a - b > 0;a < b \Leftrightarrow  a - b < 0$ \\
\cline{1-3}
传递性 & $a = b,b = c \Leftrightarrow  a = c$ & $a > b,b > c \Rightarrow  a > c$ \\
\cline{1-3}
加法性质 & $a = b \Leftrightarrow  a + c = b + c$ & $a > b \Leftrightarrow  a + c > b + c$ \\
\cline{1-3}
推论 & $a = b,c = d \Rightarrow  a + c = b + d$ & $a > b,c > d \Rightarrow  a + c > b + d$ \\
\cline{1-3}
推广 & ${a}_{i} = {b}_{i}\left( {i \in  \mathbf{N},i \geq  2}\right)  \Rightarrow  \mathop{\sum }\limits_{{i = 1}}^{n}{a}_{i} = \mathop{\sum }\limits_{{i = 1}}^{n}{b}_{i}$ & ${a}_{i} > {b}_{i}\left( {i \in  \mathbf{N},i \geq  2}\right)  \Rightarrow  \mathop{\sum }\limits_{{i = 1}}^{n}{a}_{i} > \mathop{\sum }\limits_{{i = 1}}^{n}{b}_{i}$ \\
\cline{1-3}
乘法性质 & $a = b \Rightarrow  {ac} = {bc}$ & $a > b,c > 0 \Rightarrow  {ac} > {bc};a > b,c < 0 \Rightarrow  {ac} < {bc}$ \\
\cline{1-3}
推论 & $a = b,c = d \Rightarrow  {ac} = {bd}$ & $a > b > 0,c > d > 0 \Rightarrow  {ac} > {bd}$ \\
\cline{1-3}
推广 & \makecell{ ${a}_{i} = {b}_{i}\left( {i \in  \mathbf{N},i \geq  2}\right)$ \\ $\Rightarrow  {a}_{1}{a}_{2}\cdots {a}_{n} = {b}_{1}{b}_{2}\cdots {b}_{n}$ } & \makecell{ ${a}_{i} > {b}_{i} > 0\left( {i \in  \mathbf{N},i \geq  2}\right)$ \\ $\Rightarrow  {a}_{1}{a}_{2}\cdots {a}_{n} > {b}_{1}{b}_{2}\cdots {b}_{n}$ } \\
\cline{1-3}
乘方性质 & $a = b,n \in  \mathbf{N},n \geq  1 \Rightarrow  {a}^{n} = {b}^{n}$ & $a > b > 0,n \in  \mathbf{N},n \geq  1 \Rightarrow  {a}^{n} > {b}^{n} > 0$ \\
\cline{1-3}
开方性质 & $a = b > 0,n \in  \mathbf{N},n \geq  1 \Rightarrow  \sqrt[n]{a} = \sqrt[n]{b} > 0$ & $a > b > 0,n \in  \mathbf{N},n \geq  2 \Rightarrow  \sqrt[n]{a} > \sqrt[n]{b} > 0$ \\
\cline{1-3}
倒数性质 & $a = b \neq  0 \Rightarrow  \frac{1}{a} = \frac{1}{b}$ & $a > b,{ab} > 0 \Rightarrow  \frac{1}{a} < \frac{1}{b}$ \\
\cline{1-3}
\hline
\end{tabular}
\end{center}

3. 方程 (一元二次方程) 与方程的解集:

 \ding{172}  方程的解与数域有关；

 \ding{173} 一元二次方程的解集与韦达定理:

若一元二次方程 $a{x}^{2} + {bx} + c = 0\left( {a \neq  0}\right)$ 的解为 ${x}_{1}\text{ 、 }{x}_{2}$ ,则 ${x}_{1}{x}_{2} =  - \frac{b}{a},{x}_{1}{x}_{2} = \frac{c}{a}$ .

4. 不等式的解法:高次不等式低次化；分式不等式整式化；绝对值不等式常规化；无理不等式有理化；含参数的不等式分类讨论; 数形结合、换元法、根据函数的单调性等是转化的重要途径.

5. 均值不等式:

 \ding{172} 常用不等式: $a$ 、 $b \in  \mathbf{R} \Rightarrow  {a}^{2} + {b}^{2} \geq  {2ab}$ (当且仅当 $a = b$ 时，等号成立)；

 \ding{173} 平均值不等式: $a \geq  0,b \geq  0 \Rightarrow  \frac{a + b}{2} \geq  \sqrt{ab}$ (当且仅当 $a = b$ 时，等号成立)；

\HRule

* 注:本章节所属知识模块是“预备知识”.

\HRule

 \ding{174}  $a\text{ 、 }b \in  \mathbf{R} \Rightarrow  {\left( \frac{a + b}{2}\right) }^{2} \geq  {ab}$ (当且仅当 $a = b$ 时，等号成立)；

 \ding{175}  $a\text{ 、 }b \in  \mathbf{R} \Rightarrow  \frac{{a}^{2} + {b}^{2}}{2} \geq  {\left( \frac{a + b}{2}\right) }^{2} \geq  {ab}$ (当且仅当 $a = b$ 时，等号成立)；

 \circled{5}  ${ab} > 0 \Rightarrow  \frac{b}{a} + \frac{a}{b} \geq  2$ (当且仅当 $a = b$ 时，等号成立)； ${ab} < 0 \Rightarrow  \frac{b}{a} + \frac{a}{b} \leq   - 2$ (当且仅当 $a =  - b$ 时，等号成立).

注意: ${a}^{2} + {b}^{2}\text{ 、 }a + b\text{ 、 }a - b\text{ 、 }{ab}$ 之间的关系.

重要结论: 设 $x\text{ 、 }y$ 为正数,

 \ding{172} 如果积 ${xy}$ 为定值 $P$ ，那么当 $x = y$ 时，和 $x + y$ 有最小值 $2\sqrt{P}$ ；

 \ding{173} 如果和 $x + y$ 为定值 $S$ ，那么当 $x = y$ 时，积 $x \cdot  y$ 有最大值 $\frac{1}{4}{S}^{2}$ .

注意:利用平均值不等式求最值时要注意“一正、二定、三取等”.

6. 三角不等式:

 \ding{172}  $\left| {a + b}\right|  \leq  \left| a\right|  + \left| b\right|$ (当且仅当 ${ab} \geq  0$ 时，等号成立)

 \ding{173}  $\left| {a - b}\right|  \leq  \left| a\right|  + \left| b\right|$ (当且仅当 ${ab} \leq  0$ 时，等号成立)

 \ding{174}  $\left| \right| a\left| -\right| b\left| \right|  \leq  \left| {a + b}\right|$ (当且仅当 ${ab} \leq  0$ 时，等号成立)

 \ding{175}  $\left| \right| a\left| -\right| b\left| \right|  \leq  \left| {a - b}\right|$ (当且仅当 ${ab} \geq  0$ 时，等号成立)

即 $a\text{ 、 }b \in  \mathbf{R} \Rightarrow  \left| \right| a\left| -\right| b\left| \right|  \leq  \left| {a \pm  b}\right|  \leq  \left| a\right|  + \left| b\right|$ .

7. 不等式的证明:

有比较法(比差、比商)、综合法、反证法、分析法等.

(1)比较法:作差 $\rightarrow$ 变形 (配方或分解因式) $\rightarrow$ 判断符号.

两边均正时也可比商，判断与 1 的大小.

(2)综合法:从已知出发，依据性质和基本不等式，直接推出所需的不等式(连接词:: ···，，....).

(3)反证法:首先假设结论 $\beta$ 不成立( $\beta$ 为假)，然后经过正确的逻辑推理得出与已知条件或(已学)定理等相矛盾的结论,从而说明“ $\beta$ 为假”是不可能发生的,即结论 $\beta$ 是正确的. 这样的证明方法叫反证法.

(4)分析法:从欲证的不等式出发，寻找不等式成立的充分条件 (连接词:要证...，只要证...).

除此以外, 还有放缩法、数学归纳法等.

8. 不等式的恒成立和有解问题:

(1) $a{x}^{2} + {bx} + c > \left( { \geq  0}\right)$ 0 恒成立 $\Leftrightarrow  \left\{  \begin{array}{l} a > 0, \\  \Delta  < \left(  \leq  \right) 0 \end{array}\right.$ 或 $\left\{  \begin{array}{l} a = 0,b = 0, \\  c > \left(  \geq  \right) 0; \end{array}\right.$

(2)若 $m \leq  f\left( x\right)  \leq  M$ ，则 $a > f\left( x\right)$ 恒成立 $\Leftrightarrow  a > M;a > f\left( x\right)$ 有解 $\Leftrightarrow  a > m;a \geq  f\left( x\right)$ 恒成立 $\Leftrightarrow  a \geq  M$ ； $a \geq  f\left( x\right)$ 有解 $\Leftrightarrow  a \geq  m$ .

※若 $m < f\left( x\right)  < M$ ( $m$ 、 $M$ 分别是 $f\left( x\right)$ 的最大下界和最小上界),则 $a > f\left( x\right)$ 恒成立 $\Leftrightarrow  a \geq  M;a >  f\left( x\right)$ 有解 $\Leftrightarrow  a > m;a \geq  f\left( x\right)$ 恒成立 $\Leftrightarrow  a \geq  M;a \geq  f\left( x\right)$ 有解 $\Leftrightarrow  a > m$ . 对于 $a < f\left( x\right)$ 和 $a \leq  f\left( x\right)$ 类似考虑.

注意区别 “恒成立” 和 “有解”: “恒成立” 中的 $x$ 具有任意性,“有解” 中的 $x$ 只要具有存在性.

9. 常用解题策略:认清变量和参数，恰当分离参数，灵活转化问题是解决不等式问题的重要保证.

\section*{$\text{ § }{2.1}$ 等式与方程}

\section*{典型例题精讲}

【例 1】已知全集 $U = \{ 1,2,3,4,5\}$ ,集合 $A = \left\{  {x \mid  {x}^{2} - {5x} + m = 0}\right\}  ,B = \left\{  {x \mid  {x}^{2} + {nx} + {12} = 0}\right\}$ ,且 $\bar{A} \cup  B  = \{ 1,3,4,5\}$ ,求 $m + n$ 的值.

【例 2】已知一元二次函数 $f\left( x\right)  = a{x}^{2} + {bx} + c\left( {a > 0,c > 0}\right)$ 的图像与 $x$ 轴有两个不同的公共点,其中一个公共点的坐标为 $\left( {c,0}\right)$ ,且当 $0 < x < c$ 时,恒有 $f\left( x\right)  > 0$ .

(1)当 $a = 1\text{ 、 }c = \frac{1}{2}$ 时，求出不等式 $f\left( x\right)  < 0$ 的解;

(2)求不等式 $f\left( x\right)  < 0$ 的解 (用 $a$ 、 $c$ 表示)；

(3)若不等式 ${m}^{2} - {2km} + 1 + b + {ac} \geq  0$ 对所有 $k \in  \left\lbrack  {-1,1}\right\rbrack$ 恒成立，求实数 $m$ 的取值范围.

【例 3】已知二次函数 $f\left( x\right)  = a{x}^{2} + {bx}$ ( $a$ 、 $b$ 为常数，且 $a \neq  0$ )满足条件: $f\left( {x - 1}\right)  = f\left( {3 - x}\right)$ ，且方程 $f\left( x\right)  = {2x}$ 有两等根.

(1)求 $f\left( x\right)$ 的解析式；

(2)是否存在实数 $m$ 、 $n$ ( $m < n$ )，使 $f\left( x\right)$ 的定义域和值域分别为 $\left\lbrack  {m,n}\right\rbrack$ 和 $\left\lbrack  {{4m},{4n}}\right\rbrack$ ？如果存在，求出 $m\text{ 、 }n$ 的值；如果不存在，说明理由.

【例 4】函数 $y = f\left( x\right)$ 的定义域为 $D$ ,若存在 ${x}_{0} \in  D$ ,使得 $f\left( {x}_{0}\right)  = {x}_{0}$ 成立,则称 ${x}_{0}$ 为函数 $y = f\left( x\right)$ 的“不动点”.

(1)若 $f\left( x\right)  = a{x}^{2} + {bx} + c\left( {a \neq  0}\right)$ 有两个不动点 -1 和 3，求 ${bc}$ 的最小值；

(2)若 $a > 0$ ，且 $f\left( x\right)  = a{x}^{2} + {bx} + c$ 有两个不动点 $m$ 、 $n$ ，且满足: $0 < m < n < \frac{1}{a}$ ，求证:当 $x \in  \left( {0,m}\right)$ 时， $f\left( x\right)  < m$ .

【例 5】已知 $a > b > c,a + b + c = 0$ ,方程 $a{x}^{2} + {bx} + c = 0$ 的两个实根为 ${x}_{1}\text{ 、 }{x}_{2}$ .

(1)证明: $- \frac{1}{2} < \frac{b}{a} < 1$ ；

( 2 )若 ${x}_{1}^{2} + {x}_{1}{x}_{2} + {x}_{2}^{2} = 1$ ，求 ${x}_{1}^{2} - {x}_{1}{x}_{2} + {x}_{2}^{2}$ 的值；

(3)求 $\left| {{x}_{1}^{2} - {x}_{2}^{2}}\right|$ 的取值范围.

\section*{四基达标精练}

\section*{一、填空题}

1. 已知一元二次方程 $a{x}^{2} + {bx} + 2 = 0$ 的解是 -2 $\frac{1}{2}$ 和 $\frac{1}{3}$ ，则 $a + b =$ \_\_\_\_\_.

2. 若关于 $x$ 的不等式 $a{x}^{2} + {6x} - {a}^{2} < 0$ 的解集是 $\left( {-\infty ,1}\right)  \cup  \left( {m, + \infty }\right)$ ，则实数 $m =$ \_\_\_\_\_.

3. 已知实数 $a$ 、 $b$ 满足 ${\left( a + 2\right) }^{2} + {\left( {b}^{2} - 2b - 3\right) }^{2} = 0$ ，则 $a + b$ 的值为\_\_\_\_\_.

4. 已知方程 ${x}^{2} + {2mx} + m + 2 = 0$ 有两个不等正根，则实数 $m$ 的取值范围是\_\_\_\_\_.

5. 关于 $x$ 的方程 ${x}^{2} - {2ax} - 8{a}^{2} = 0$ 的解为 ${x}_{1}$ 和 ${x}_{2}$ ，且 ${x}_{1}^{2} - {x}_{2}^{2} = {15}$ ，则实数 $a =$ \_\_\_\_\_.

6. 二次函数 $y = a{x}^{2} + {bx} + c\left( {x \in  \mathbf{R}}\right)$ 的部分对应值如表:

\begin{center}
\begin{tabular}{|c|c|c|c|c|c|c|c|c|}
\hline
$x$ & -3 & -2 & -1 & 0 & 1 & 2 & 3 & 4 \\
\cline{1-9}
$y$ & 6 & 0 & -4 & -6 & -6 & -4 & 0 & 6 \\
\cline{1-9}
\hline
\end{tabular}
\end{center}

则不等式 $a{x}^{2} + {bx} + c > 0$ 的解集是\_\_\_\_\_.

7. 设 ${x}_{1}$ 、 ${x}_{2}$ 是关于 $x$ 的方程 $2{x}^{2} - {4mx} + 2{m}^{2} + {3m} - 2 = 0$ 的两个实数根，则 ${x}_{1}^{2} + {x}_{2}^{2}$ 的最小值为 .

8. 在 $\mathbf{R}$ 上定义运算 $\odot   : a \odot  b = {ab} + {2a} + b$ ，则不等式 $x \odot  \left( {x - 2}\right)  < 0$ 的解集是\_\_\_\_\_.

9. 已知 $\alpha$ 、 $\beta \left( {\alpha  < \beta }\right)$ 是方程 $\left( {x - a}\right) \left( {x - b}\right)  =  - 2\left( {a < b}\right)$ 的两个根，则 $\alpha$ 、 $\beta$ 、 $a$ 、 $b$ 的大小关系是\_\_\_\_\_.

\section*{二、选择题}

10. “ $- 2 < a < 2$ ” 是 “不等式 $\left( {a - 2}\right) {x}^{2} + 2\left( {a - 2}\right) x - 4 < 0$ 的解集为 $\mathbf{R}$ ” 的 ( )

A. 充分非必要条件. B. 必要非充分条件. C. 充要条件. D. 既非充分也非必要条件.



\begin{center}
\includegraphics[width=0.2\textwidth]{images/bo_d7tk7oc91nqc73875fu0_15_1268_157_163_345_0.jpg}
\end{center}
 

(第 12 题图)

11. 已知关于 $x$ 的不等式 ${x}^{2} - \left( {a + 1}\right) x + a < 0$ 的解集中恰有两个正整数,则实数 $a$ 的取值范围是 ( )

A. $\lbrack 2,4)$ . B. $\left\lbrack  {3,4}\right\rbrack$ . C. $(3,4\rbrack$ . D. $\left( {3,4}\right)$ .

12. 古希腊时期, 人们认为最美人体的头顶至肚脐的长度与肚脐至足底的长度之比是 $\frac{\sqrt{5} - 1}{2}\left( {\frac{\sqrt{5} - 1}{2} \approx  {0.618}}\right.$ ,称为黄金分割比例 $)$ ,著名的“断臂维纳斯”便是如此. 此外,最美人体的头顶至咽喉的长度与咽喉至肚脐的长度之比也是 $\frac{\sqrt{5} - 1}{2}$ . 若某人满足上述两个黄金分割比例, 且腿长为 ${105}{{cm}}$ ，头顶至脖子下端的长度为 ${26}{{cm}}$ ，则其身高可能是 ( )

A. ${165}{{cm}}$ . B. ${175}{{cm}}$ . C. ${185}{{cm}}$ . D. ${190}{{cm}}$ .

\section*{三、解答题}

13. 已知 $f\left( x\right)  = {x}^{2} + {mx} - 6\left( {m > 0}\right)$ ，实数 ${x}_{1}$ 、 ${x}_{2}$ 满足 $f\left( {x}_{1}\right)  = f\left( {x}_{2}\right)  = 0$ ，且 ${x}_{2} - {x}_{1} = 5$ .

(1)求函数 $y = f\left( x\right)$ 的解析式；

(2)解关于 $x$ 的不等式 $y < 4 - {2x}$ .

14. 已知集合 $\{ x \mid  {x}^{2} - \left( {k + 2}\right) x - {3k} + 1 \geq  0\}  = \left( {-\infty , - 1\rbrack \cup \lbrack 5, + \infty }\right)$ .

(1)求实数 $k$ 的值；

(2)已知 $t \in  \left( {-\infty ,2}\right)$ ，若不等式 ${x}^{2} - \left( {k + 2}\right) x - {3k} - {m}^{2} + {4m} + {15} \geq  0$ 在 $t \leq  x \leq  4$ 上恒成立，求实数 $m$ 的取值范围.

15. 已知函数 $f\left( x\right)  = {x}^{2} + {2x} + 1$ .

(1)求关于 $x$ 的不等式 $f\left( x\right)  \geq  {b}^{2}\left( {b \geq  0}\right)$ 的解集；

(2)若不等式 ${\left\lbrack  f\left( x\right) \right\rbrack  }^{2} - {2mf}\left( x\right)  + {m}^{2} - 1 \geq  0$ 对于任意 $x \in  \left\lbrack  {-3,2}\right\rbrack$ 都成立，求 $m$ 的取值范围.

16. 已知关于 $x$ 的一元二次方程 ${x}^{2} - {6x} + {2m} - 1 = 0$ 有两实数根 ${x}_{1}$ 、 ${x}_{2}$ .

(1)若 ${x}_{1} = 1$ ，求 ${x}_{2}$ 和 $m$ 的值；

(2)是否存在实数 $m$ ，满足 $\left( {{x}_{1} - 1}\right) \left( {{x}_{2} - 1}\right)  = \frac{6}{m - 5}$ ？若存在，求出实数 $m$ 的值；若不存在，请说明理由.

\section*{$\text{ § }{2.2}$ 不等式的性质与解法}

\section*{典型例题精讲}

【例 1】若不等式 $f\left( x\right)  \geq  0$ 的解集是 $\left\lbrack  {-1,2}\right\rbrack$ ,不等式 $g\left( x\right)  \geq  0$ 的解集为空集,且函数 $f\left( x\right) \text{ 、 }g\left( x\right)$ 的定义域都是 $\mathbf{R}$ ，则不等式 $\frac{f\left( x\right) }{g\left( x\right) } > 0$ 的解集为\_\_\_\_\_.

【例 2】解不等式组 $\left\{  \begin{array}{l} a\left( {x - 2}\right)  + 1 > 0, \\  {\left( x - 1\right) }^{2} > a\left( {x - 2}\right)  + 1, \end{array}\right.$ 其中 $a > 1$ .

【例 3】已知关于 $x$ 的不等式组 $\left\{  \begin{array}{l} {x}^{2} - x - 2 > 0, \\  2{x}^{2} + \left( {{2k} + 5}\right) x + {5k} < 0 \end{array}\right.$ 的解集为 $A,\mathbf{Z}$ 为整数集.

(1)若 $A \cap  \mathbf{Z} = \{  - 2\}$ ，求实数 $k$ 的取值范围；

(2)若 $A \cap  \mathbf{Z}$ 中恰有 $n$ ( $n$ 为正整数)个整数，求实数 $k$ 的取值范围(用 $n$ 表示).

【例 4】已知不等式: ${0x}^{2} - {4x} + 3 < 0,{0x}^{2} - {6x} + 8 < 0,{0x}^{3}{x}^{2} - {9x} + m < 0$ .

(1)若同时满足 \ding{172}  \ding{173} 的 $x$ 也满足 \ding{174} ，求实数 $m$ 的取值范围；

(2)若 \ding{174} 的解集非空，且满足  \ding{174}  的 $x$ 至少满足  \ding{172}  \ding{173}  中的一个，求实数 $m$ 的取值范围.

【例 5】已知关于 $x$ 的不等式 $\left( {{kx} - {k}^{2} - 4}\right) \left( {x - 4}\right)  > 0$ ，其中 $k \in  \mathbf{R}$ .

(1)试求不等式的解集 $A$ ；

(2)对于不等式的解集 $A$ ,若满足 $A \cap  \mathbf{Z} = B$ (其中 $\mathbf{Z}$ 为整数集). 试探究集合 $B$ 能否为有限集? 若能,求出使得集合 $B$ 中元素个数最少的 $k$ 的所有取值,并用列举法表示集合 $B$ ; 若不能,请说明理由.

\section*{四基达标精练}

\section*{一、填空题}

1. 不等式 $- \frac{1}{2} < \frac{1}{x} \leq  \frac{1}{3}$ 的解集是\_\_\_\_\_.

2. 已知 $x \in  \mathbf{R},M = 2{x}^{2} + {5x} + 3,N = {x}^{2} + {4x} + 2$ ，则 $M$ 、 $N$ 的大小关系是\_\_\_\_\_.

3. 已知 $- 1 < a < 2, - 1 < b < 2$ ,则 $a - b$ 的范围是\_\_\_\_\_.

4. 不等式 $- {x}^{2} + x + 6 < 0$ 的解集是\_\_\_\_\_.

5. 已知函数 $f\left( x\right)  = \left\{  \begin{array}{l} x + 2,x \leq  0, \\   - x + 2,x > 0, \end{array}\right.$ 则不等式 $f\left( x\right)  \geq  {x}^{2}$ 的解集为\_\_\_\_\_.

6. 若关于 $x$ 的不等式 ${ax} - b > 0$ 的解集是 $\left( {1, + \infty }\right)$ ,则关于 $x$ 的不等式 $\frac{{ax} + b}{x - 2} > 0$ 的解集是\_\_\_\_\_.

7. 已知不等式 $a{x}^{2} + {bx} + c > 0$ 的解集为 $\left( {-\infty ,\frac{1}{4}}\right)  \cup  \left( {\frac{1}{2}, + \infty }\right)$ ,则不等式 $c{x}^{2} + {bx} + a < 0$ 的解集为\_\_\_\_\_.

8. 若关于 $x$ 的不等式 $1 \leq  k{x}^{2} + {2x} + k \leq  2$ 有唯一实数解,则实数 $k =$ \_\_\_\_\_.

9. 对于问题:“已知关于 $x$ 的不等式 $a{x}^{2} + {bx} + c > 0$ 的解集为 $\left( {-1,2}\right)$ ，解关于 $x$ 的不等式 $a{x}^{2} - {bx} + c > 0$ ”， 给出如下一种解法:

解:由 $a{x}^{2} + {bx} + c > 0$ 的解集为 $\left( {-1,2}\right)$ ,得 $a{\left( -x\right) }^{2} + b\left( {-x}\right)  + c > 0$ 的解集为 $\left( {-2,1}\right)$ ,即关于 $x$ 的不等式 ${ax} - {bx} + c > 0$ 的解集为 $\left( {-2,1}\right)$ .

参考上述解法,若关于 $x$ 的不等式 $\frac{k}{x + a} + \frac{x + b}{x + c} < 0$ 的解集为 $\left( {-1, - \frac{1}{3}}\right)  \cup  \left( {\frac{1}{2},1}\right)$ ,则关于 $x$ 的不等式 $\frac{kx}{{ax} + 1} + \frac{{bx} + 1}{{cx} + 1} < 0$ 的解集为\_\_\_\_\_.

\section*{二、选择题}

10. 已知 $x > 0,P = \sqrt{x + 1} - \sqrt{x},Q = \sqrt{x + 3} - \sqrt{x + 2}$ ,则 $P$ 与 $Q$ 的大小关系为 ( )

A. $P > Q$ . B. $P < Q$ . C. $P = Q$ . D. 不确定.

11.16 世纪中叶, 英国数学家雷科德在《砺智石》一书中首先把 “ = ”作为等号使用，后来英国数学家哈利奥特首次使用“<”和“>”符号，并逐步被数学界接受，不等号的引入对不等式的发展影响深远. 已知 $b < a < 0$ ,则下列选项正确的是 ( )

A. ${a}^{2} > {b}^{2}$ . B. $a + b < {ab}$ . C. $\left| b\right|  < \left| a\right|$ . D. ${ab} > {b}^{2}$ .

12. 下列各组中, $M$ 是 $N$ 的充要条件的是 ( )

A. $M : \left| x\right|  + \left| y\right|  \leq  1,N : {x}^{2} + {y}^{2} \leq  1$ .

B. $M$ : 实数 $a$ 、 $b$ 满足 $a + b > 2$ 且 ${ab} > 1,N : a > 1$ 且 $b > 1$ .

C. 集合 $E\text{ 、 }F\text{ 、 }P,M : P \subset  E$ 且 $P \subset  F,N : P \subset  E \cap  F$ .

D. $M : \left| a\right|  \geq  1,N$ : 曲线 $y = \left| x\right|$ 与直线 $y = {ax} + 2$ 有且只有一个交点.

\section*{三、解答题}

13. 设 $x\text{ 、 }y\text{ 、 }a\text{ 、 }b \in  \mathbf{R}$ ,且 $a\text{ 、 }b > 0,a \neq  b$ ,分别比较下列两个数的大小.

(1) $m = {x}^{2} + {y}^{2} - {2x} + {2y},n =  - 3$ ，比较 $m$ 、 $n$ 的大小；

(2) $S = {a}^{3} + {b}^{3},T = {a}^{2}b + a{b}^{2}$ ，比较 $S$ 、 $T$ 的大小.

14. 解下列不等式:

(1) $3{x}^{2} - {2x} - 1 \geq  0$ ； (2) $\frac{x + 2}{x - 1} \leq  0$ ； (3) $\left| {x + 3}\right|  - \left| {{2x} - 1}\right|  < \frac{x}{2} + 1$ .

15. 已知 $m \in  \mathbf{R},f\left( x\right)  = \frac{x - m}{x + 1}$ .

(1)若 $m = 2$ ，解不等式 $f\left( x\right)  > 0$ ；

(2)讨论 $f\left( x\right)  > 0$ 的解集.

16. 如图,在梯形 ${PMNQ}$ 中, ${PQ}//{MN}$ ,对角线 ${PN}$ 和 ${MQ}$ 相交于点 $O$ ,并把梯形分成四部分,记这四部分的面积分别为 ${S}_{1}$ 、 ${S}_{2}$ 、 ${S}_{3}$ 、 ${S}_{4}$ . 试判断 ${S}_{1} + {S}_{2}$ 和 ${S}_{3} + {S}_{4}$ 的大小关系,并证明你的结论.

\begin{center}
\includegraphics[width=0.2\textwidth]{images/bo_d7tk7oc91nqc73875fu0_20_1107_1371_323_248_0.jpg}
\end{center}
 

(第 16 题图)

§ 2.3 基本不等式

\section*{典型例题精讲}

【例 1】解答下列问题:

(1)若 $x$ 、 $y > 0$ ，且 $x + {2y} = 1$ ，求 $\frac{3}{x} + \frac{2}{y}$ 的最小值；

(2)若 $x$ 、 $y > 0$ ，且 $\frac{1}{x} + \frac{2}{y} = 1$ ，求 ${3x} + {2y}$ 的最小值；

(3)若 $0 < x < 1$ ，求函数 $f\left( x\right)  = \frac{3}{x} + \frac{4}{1 - x}$ 的最小值.

【例 2】已知函数 $f\left( x\right)  = \left| {x + 1}\right|  + \left| x\right|$ .

(1)求不等式 $f\left( x\right)  \leq  3$ 的解集；

(2)若 $x \in  \left\lbrack  {a - 2,a + 2}\right\rbrack$ 时， $f\left( x\right)  < {3a} - 2$ ，求实数 $a$ 的取值范围.

【例 3】设函数 $f\left( x\right)  = \left| {x - 1}\right|  + \left| {{2x} + a}\right| \left( {a \in  \mathbf{R}}\right)$ .

(1)若 $a = 2$ ，求 $f\left( x\right)  \leq  8$ 的解集；

(2)若 $f\left( x\right)  \geq  3 - \left| {x - 1}\right|$ 对任意 $x \in  \mathbf{R}$ 恒成立，求实数 $a$ 的取值范围.

【例 4】已知 $a$ 、 $b$ 、 $c > 0$ ，且 $a + b + c = 1$ . 证明下列不等式:

(1) $\frac{1}{a} + \frac{1}{b + c} \geq  4$ ；

(2) $\frac{1}{a + {2b}} + \frac{1}{b + {2c}} + \frac{1}{c + {2a}} \geq  3$ .

【例 5】已知正实数 $a\text{ 、 }b$ 满足 ${\left( a + 1\right) }^{3} + {\left( b + 1\right) }^{3} = {16}$ ,设 $a + b$ 的最大值为 $m$ .

(1)求 $m$ 的值；

(2)若 $\left| {{2x} - a}\right|  \leq  m,\left| {{2y} - a}\right|  \leq  m$ ，求证: $\left| {{2x} - {4y} + a}\right|  \leq  6$ .

\section*{四基达标精练}

\section*{一、填空题}

1.若正数 $x\text{ 、 }y$ 满足 ${xy} = 1$ ，则 $x + {2y}$ 的最小值为\_\_\_\_\_.

2. 已知 $m > 0,n > 0,m + n = 1$ ,则 $\frac{1}{m} + \frac{2}{n}$ 的最小值为\_\_\_\_\_.

3. 若 $\left| {a - 1}\right|  - \left| a\right|  = 1$ ，则实数 $a$ 的取值范围是\_\_\_\_\_.

4. 若 $x$ 、 $y > 0$ ，且 $\frac{1}{x} + {2y} = 3$ ，则 $\frac{y}{x}$ 的最大值为\_\_\_\_\_.

5. 已知 $a$ 、 $b \in  \mathbf{R}$ ，且 $a - {3b} + 6 = 0$ ，则 ${2}^{a} + \frac{1}{{8}^{b}}$ 的最小值为\_\_\_\_\_.

6. 若不等式 $\left| {x - 4}\right|  - \left| {x - 3}\right|  \leq  a$ 对一切 $x \in  \mathbf{R}$ 恒成立，则实数 $a$ 的取值范围是\_\_\_\_\_.

7. 已知 $x > y > 0$ ，则 $\frac{{x}^{2} + {y}^{2}}{{xy} - {y}^{2}}$ 的最小值是\_\_\_\_\_.

8．正项等比数列 $\left\{  {a}_{n}\right\}$ 中， ${a}_{2023} = {a}_{2022} + 2{a}_{2021}$ ，若 ${a}_{m}{a}_{n} = {16}{a}_{1}^{2}$ ，则 $\frac{1}{m} + \frac{4}{n}$ 的最小值等于\_\_\_\_\_.

9. 已知定义在 $\mathbf{R}$ 上的偶函数 $f\left( x\right)  = \left| {x - m + 1}\right|  - 2$ ,若正实数 $a\text{ 、 }b$ 满足 $f\left( a\right)  + f\left( {2b}\right)  = m$ ,则 $\frac{1}{a} + \frac{2}{b}$ 的最小值为\_\_\_\_\_.

\section*{二、选择题}

10. 下列不等式恒成立的是 ( )

A. ${a}^{2} + {b}^{2} \leq  {2ab}$ . B. ${a}^{2} + {b}^{2} \geq   - {2ab}$ . C. $a + b \geq  2\sqrt{\left| ab\right| }$ . D. $a + b \geq   - 2\sqrt{\left| ab\right| }$ .

11. 等式 $\left| {a - b}\right|  = \left| {a - 1}\right|  + \left| {b - 1}\right|$ 成立的条件是 ( )

A. $\left( {a - 1}\right) \left( {b - 1}\right)  < 0$ . B. $\left( {a - 1}\right) \left( {b - 1}\right)  \leq  0$ . C. $\left( {a - 1}\right) \left( {b - 1}\right)  > 0$ . D. $\left( {a - 1}\right) \left( {b - 1}\right)  \geq  0$ .

12. 已知 ${2}^{m} = {3}^{n} = 6$ ,则 $m\text{ 、 }n$ 不可能满足的关系是 ( )

A. $m + n > 4$ . B. ${mn} > 4$ . C. ${m}^{2} + {n}^{2} < 8$ . D. ${\left( m - 1\right) }^{2} + {\left( n - 1\right) }^{2} > 2$ .

\section*{三、解答题}

13. 某化工厂生产某种产品,当年产量在 150 吨至 250 吨 (包括 150 吨与 250 吨) 时,每年的生产成本 $y$ 万元与年产量 $x$ 吨之间的关系可近似地表示为 $y = \frac{1}{10}{x}^{2} - {30x} + {4000}$ .

(1)若每年的生产总成本不超过 2000 万元，求年产量 $x$ 的取值范围；

(2)求年产量为多少吨时, 每吨的平均成本最低, 并求每吨的最低成本.

14. 已知 $a\text{ 、 }b$ 是正常数， $a \neq  b,x\text{ 、 }y \in  \left( {0, + \infty }\right)$ .

(1)求证: $\frac{{a}^{2}}{x} + \frac{{b}^{2}}{y} \geq  \frac{{\left( a + b\right) }^{2}}{x + y}$ ，指出等号成立的条件；

(2)利用(1)的结论求函数 $f\left( x\right)  = \frac{2}{x} + \frac{9}{1 - {2x}}$ ( $0 < x < \frac{1}{2}$ )的最小值，并指出取最小值时 $x$ 的值.

15. 设 $a > 0$ ，实数 $x$ 、 $y$ 满足 $\left| {x - 1}\right|  < \frac{a}{3}$ ， $\left| {y - 2}\right|  < \frac{a}{3}$ .

(1)若 $y = {3x} - 2$ ，且 $\left\{  {x\left| \right| x - 1\left| { < \frac{a}{3}\} \cap \{ y}\right| \left| {y - 2}\right|  < \frac{a}{3}}\right\}   \neq  \varnothing$ ，求 $a$ 的取值范围；

(2)求证: $\left| {{{2x} + y} - 4}\right|  < a$ .

16. 已知函数 $f\left( x\right)  = \left| {x - 2}\right| ,g\left( x\right)  =  - \left| {x + 3}\right|  + m\left( {m \in  \mathbf{R}}\right)$ .

(1)解关于 $x$ 的不等式 $f\left( x\right)  + a - 2 > 0\left( {a \in  \mathbf{R}}\right)$ ；

(2)若函数 $f\left( x\right)$ 的图像恒在函数 $g\left( x\right)$ 图像的上方，求 $m$ 的取值范围.

\section*{§ 2.4 不等式综合}

\section*{典型例题精讲}

【例 1】已知函数 $f\left( x\right)  = {ax} - 1,g\left( x\right)  = {x}^{2} - a$ .

(1)若不等式 $\frac{f\left( x\right) }{g\left( x\right) } < 0$ 的解集为 $S$ ,且 $3 \in  S,4 \notin  S$ ,求实数 $a$ 的取值范围；

(2)若不等式 $f\left( x\right)  \leq  g\left( x\right)$ 对任意 $x \in  \left\lbrack  {0,1}\right\rbrack$ 恒成立，求实数 $a$ 的取值范围.

【例 2】已知 $a\text{ 、 }b \in  \left( {0, + \infty }\right)$ ,证明下列结论:

(1) $\frac{a}{\sqrt{b}} + \frac{b}{\sqrt{a}} \geq  \sqrt{a} + \sqrt{b}$ ；

(2)若 $a + b = {10}$ ，求证: $\sqrt{1 + {3a}} + \sqrt{1 + {3b}} \leq  8$ .

【例 3】已知 $m + {2n} = 2$ ，且 $m >  - 1$ ， $n > 0$ .

(1)求 $\frac{1}{m + 1} + \frac{2}{n}$ 的最小值；

(2)求 $\frac{{m}^{2}}{{2n} + 2} + \frac{4{n}^{2}}{m + 1}$ 的最小值.

【例 4】已知正数 $a$ 、 $b$ 和实数 $t$ 满足 ${a}^{2} + {tab} + {b}^{2} = 1$ .

(1)若 $t = 1$ ，求 $a + b$ 的最大值；

(2)若 $a + b$ 存在最大值，求 $t$ 的取值范围.

【例 5】我们知道 $\left| {AB}\right|  = \left| {{x}_{A} - {x}_{B}}\right|$ 表示数轴上 $A$ 、 $B$ 两点的距离，它也可以看作满足一定条件的一种运算. 这样,我们也可以将满足下列三个条件的一个 $x$ 与 $y$ 间的运算 $p\left( {x,y}\right)$ 叫做 $x\text{ 、 }y$ 之间的距离: 条件 \ding{172} : 非负性 $p\left( {x,y}\right)  \geq  0$ ,当且仅当 $x = y$ 等号成立; 条件 \ding{173} : 交换律 $p\left( {x,y}\right)  = p\left( {y,x}\right)$ ; 条件 \ding{174} : $p\left( {x,z}\right)  \leq  p\left( {x,y}\right)  + p\left( {y,z}\right)$ 成立. 试判断运算 $s\left( {x,y}\right)  = \frac{\left| x - y\right| }{1 + \left| {x - y}\right| }$ 是否为一个距离? 如果是,请加以证明; 如果不是, 请举出反例.

\section*{四基达标精练}

\section*{一、填空题}

1. 若函数 $f\left( x\right)  = x + \frac{1}{x - 2}\left( {x > 2}\right)$ 在 $x = a$ 处取得最小值，则 $a =$ \_\_\_\_\_.

2. 已知 ${\log }_{5}x + {\log }_{5}y = 2$ ,则 $x + {4y}$ 的最小值为\_\_\_\_\_.

3. 若 $x = k$ 时，函数 $f\left( x\right)  = \left| {x - 2}\right|  + \left| {x + 1}\right|$ 取得最小值时，则实数 $k$ 的取值范围是\_\_\_\_\_.

4. 若不等式 $\left| {x - 1}\right|  - \left| {x + 1}\right|  > {2a}$ 的解集非空，则实数 $a$ 的取值范围是\_\_\_\_\_.

5. 已知 $a\text{ 、 }b > 0,a \geq  \frac{1}{a} + \frac{2}{b},b \geq  \frac{1}{b} + \frac{2}{a}$ ,则 $a + b$ 的最小值为\_\_\_\_\_.

6. 关于 $x$ 的不等式 $a{x}^{2} - \left| x\right|  + {2a} \geq  0$ 的解集是 $\mathbf{R}$ ，则实数 $a$ 的取值范围为\_\_\_\_\_.

7. 已知 $x + \frac{4}{x} - a - \frac{4}{a} \geq  0$ 对任意 $x > 0$ 恒成立，则实数 $a$ 的取值范围是\_\_\_\_\_. 8. 已知 $x$ 、 $y > 0$ ， $x + y + \sqrt{xy} = 2$ ，则 $x + y$ 的取值范围是\_\_\_\_\_.

9. 设 $A$ 、 $B$ 是非空数集，若对任意 $x \in  A$ 、 $y \in  B$ ，都有唯一确定的 $f\left( {x,y}\right)$ 与之对应，则称 $f\left( {x,y}\right)$ 为关于 $x\text{ 、 }y$ 的二元函数,现定义满足下列性质的二元函数 $f\left( {x,y}\right)$ 为关于实数 $x\text{ 、 }y$ 的广义“距离”.

(1)非负性: $f\left( {x,y}\right)  \geq  0$ ，当且仅当 $x = y$ 时取等号；

(2)对称性: $f\left( {x,y}\right)  = f\left( {y,x}\right)$ ；

(3)三角形不等式: $f\left( {x,y}\right)  \leq  f\left( {x,z}\right)  + f\left( {y,z}\right)$ 对任意的实数 $z$ 均成立.

给出下列三个二元函数: \ding{172}  $f\left( {x,y}\right)  = \sqrt{x - y}$ ;  \ding{173}  $f\left( {x,y}\right)  = {\left( x - y\right) }^{2}$ ;  \ding{174}  $f\left( {x,y}\right)  = \left| {x - y}\right|$ .

其中能够成为关于 $x$ 、 $y$ 的广义“距离”的二元函数的序号是\_\_\_\_\_.

\section*{二、选择题}

10. 已知 $x \in  \mathbf{R}$ ,下列不等式中正确的是 ( )

A. $\frac{1}{{2}^{x}} > \frac{1}{{3}^{3}}$ . B. $\frac{1}{{x}^{2} - x + 1} > \frac{1}{{x}^{2} + x + 1}$ .

C. $\frac{1}{2\left| x\right| } > \frac{1}{{x}^{2} + 1}$ . D. $\frac{1}{{x}^{2} + 1} > \frac{1}{{x}^{2} + 2}$ .

11. 在 $\mathbf{R}$ 上定义运算 $\otimes   : x \otimes  y = \frac{x}{2 - y}$ ,若关于 $x$ 的不等式 $\left( {x - a}\right)  \otimes  \left( {x - 1 - a}\right)  \geq  0$ 的解集是集合 $\{ x \mid   - 2 < x \leq  4\}$ 的子集,则实数 $a$ 的取值范围为 ( )

A. $- 2 < a < 1$ . B. $- 2 \leq  a < 1$ . C. $- 2 < a \leq  1$ . D. $- 2 \leq  a \leq  1$ .

12. 已知 $a\text{ 、 }b > 0$ ,有下列不等式:  \ding{172}  $\left| {a - 1}\right|  + \left| a\right|  \geq  1$ ;  \ding{173}  ${ab} + \frac{2}{ab} > 2$ ;  \ding{174}  ${a}^{2} + {b}^{2} > {4ab} - 4{b}^{2}$ ;  \ding{175}  $\frac{b}{a} > \frac{b + 1}{a + 1}$ , 其中恒成立的有 ( )

A.  \ding{172}  \ding{173} . B.  \ding{173}  \ding{175} . C.  \ding{172}  \ding{173}  \ding{174} . D.  \ding{172}  \ding{175} .

\section*{三、解答题}

13. 已知 $x > 0$ ,求 $x - 1 + \frac{2}{x}$ 的最小值. 某同学的解答如下:

因为 $x > 0$ ,所以 $x - 1 + \frac{2}{x} \geq  2\sqrt{\left( {x - 1}\right)  \cdot  \frac{2}{x}}$ ,当且仅当 $x - 1 = \frac{2}{x}$ 上式中等号成立,

即 ${x}^{2} - x - 2 = 0$ ,解得 ${x}_{1} = 2,{x}_{2} =  - 1$ (舍)

当 $x = 2$ 时, $2\sqrt{\left( {x - 1}\right)  \cdot  \frac{2}{x}} = 2$ ,所以当 $x = 2$ 时, $x - 1 + \frac{2}{x}$ 的最小值为 2 .

请判断该同学的结论的正误，如果你认为正确，则解答完毕；如果你认为错误，请指出所有的错误，并给出正确的解答.

14.某商品计划两次提价,有甲、乙、丙三种方案,其中 $p > q > 0$ .

\begin{center}
\begin{tabular}{|c|c|c|}
\hline
\makecell{ 次 \\ 方案 } & 第一次提价 & 第二次提价 \\
\cline{1-3}
甲 & p\% & q\% \\
\cline{1-3}
乙 & q\% & p\% \\
\cline{1-3}
丙 & $\frac{p + q}{2}\%$ & $\frac{p + q}{2}\%$ \\
\cline{1-3}
\hline
\end{tabular}
\end{center}

经两次提价后, 哪一种方案的提价幅度最大? 为什么?

15. 已知函数 $f\left( x\right)  = \left| {x - 1}\right|  + \left| {x + 2}\right| ,a > 0,b > 0$ .

(1)求不等式 $f\left( x\right)  \leq  7$ 的解集；

(2)若 $f\left( x\right)$ 的最小值为 $a + {2b}$ ，且 $\frac{1}{a} + \frac{2}{b} \geq  {m}^{2}$ 恒成立，求实数 $m$ 的取值范围.

16. 已知函数 $f\left( x\right)  = {kx} + b$ 的图像与 $x$ 、 $y$ 轴分别相交于点 $A$ 、 $B$ ， $\overrightarrow{AB} = 2\overrightarrow{i} + 2\overrightarrow{j}$ ( $\overrightarrow{i}$ 、 $\overrightarrow{j}$ 分别是与 $x$ 、 $y$ 轴正半轴同方向的单位向量),函数 $g\left( x\right)  = {x}^{2} - x - 6$ .

(1)求 $k$ 、 $b$ 的值；

(2)当 $x$ 满足 $f\left( x\right)  > g\left( x\right)$ 时,求函数 $y = \frac{g\left( x\right)  + 1}{f\left( x\right) } + \left| {x - 4}\right|  + \left| {x + 2}\right|$ 的最小值.

\section*{第 3 章 函数的概念与性质(必修 1)}

\section*{知识梳理*}

\section*{一、函数的概念}

1. 函数的两要素:定义域，对应法则；注意定义域优先的原则.

2. 函数定义域的求法:列不等式组，解不等式组.

原则: \ding{172}  分式分母不为零； \ding{173}  偶次方根被开方数非负； \ding{174}  零次方底数非零； \ding{175}  对数的底数为正且不等于 1，真数大于零； \circled{5} 正切函数的角的终边不能落在 $y$ 轴上； \circled{6} 实际问题的定义域要根据实际意义来确定； \circled{7} 复合函数定义域,若已知 $f\left( x\right)$ 的定义域是 $\left\lbrack  {a,b}\right\rbrack$ ,则 $f\left\lbrack  {g\left( x\right) }\right\rbrack$ 的定义域由 $a \leq  g\left( x\right)  \leq  b$ 的解集，且 $g\left( x\right)$ 有意义的 $x$ 的范围确定.

3. 注意函数的定义域与不等式恒成立之间的关系.

\section*{二、函数的基本性质:}

1. 函数的奇偶性 (函数定义域内的整体性质、定义域关于原点对称是必要条件):

(1)若 $f\left( x\right)$ 是偶函数，那么 $f\left( x\right)  = f\left( {-x}\right)  = f\left( \left| x\right| \right)$ (图像关于 $y$ 轴对称)；

(2)若 $f\left( x\right)$ 是奇函数，那么 $f\left( x\right)  =  - f\left( {-x}\right)$ (图像关于原点中心对称)；

(3)判断函数奇偶性可用定义的等价形式: $f\left( x\right)  \pm  f\left( {-x}\right)  = 0$ 或 $\frac{f\left( {-x}\right) }{f\left( x\right) } =  \pm  1\left( {f\left( x\right)  \neq  0}\right)$ ;

(4)判断函数奇偶性的实质是判断两个互为相反数的自变量的函数值相等或互为相反数;

(5)确定一个函数既不是奇函数也不是偶函数,通过举反例:如 $f\left( {-1}\right)  \neq  f\left( 1\right)$ 且 $f\left( {-1}\right)  \neq   - f\left( 1\right)$ 等, 注意此时不能用连续的不等号写成 $f\left( {-1}\right)  \neq  f\left( 1\right)  \neq   - f\left( 1\right)$ 的形式;

(6)分段函数的奇偶性分段求值和判断，复杂形式先化简后判断；

(7)原点处有意义的奇函数图像必过原点 (可用于求参数，但要给予一般性证明)；但 $f\left( 0\right)  = 0$ 是一个函数为奇函数的既不充分也不必要条件, 因为有定义域的限制;

(8)奇(偶)函数在关于原点对称的区间内单调性相同(反);

(9)奇(偶)函数在关于原点对称的区间内符号相反(同).

2. 函数图像 (或方程曲线) 的对称性 (中心对称关键:中点坐标公式；轴对称关键:垂直、平分):

(1) $f\left( x\right)$ 满足 $f\left( {a + x}\right)  = f\left( {a - x}\right)$ (或 $f\left( x\right)  = f\left( {{2a} - x}\right)$ )，则 $f\left( x\right)$ 图像关于 $x = a$ 对称；

(2) $f\left( x\right)$ 满足 $f\left( {a + x}\right)  =  - f\left( {a - x}\right)$ (或 $f\left( x\right)  =  - f\left( {{2a} - x}\right)$ )，则 $f\left( x\right)$ 图像关于(a，0)对称；

(3) $f\left( x\right)$ 满足 $f\left( {a + x}\right)  = {2b} - f\left( {a - x}\right)$ (或 $f\left( x\right)  = {2b} - f\left( {{2a} - x}\right)$ )，则 $f\left( x\right)$ 图像关于 $\left( {a,b}\right)$ 对称；

(4)特别的, $f\left( {x,y}\right)  = 0$ 关于 $y = x + a\left( {y =  - x + a}\right)$ 对称的曲线方程为: $f\left( {y - a,x + a}\right)  = 0(f(a - y$ , $a - x) = 0)$ ;

(5) 曲线 $f\left( {x,y}\right)  = 0$ 关于点 $\left( {a,b}\right)$ 对称的曲线方程为: $f\left( {{2a} - x,{2b} - y}\right)  = 0$ ;

(6) 函数 $y = f\left( {x - a}\right)$ 与 $y = f\left( {b - x}\right)$ 的图像关于直线 $x = \frac{a + b}{2}$ 对称.

说明:

 \ding{172} 证明函数图像的对称性，即证图像上任意点关于对称中心(对称轴)的对称点仍在图像上； \ding{173} 证明图像 ${C}_{1}\text{ 、 }{C}_{2}$ 的对称性,即证 ${C}_{1}$ 上任意点关于中心(轴)的对称点仍在 ${C}_{2}$ 上,反之亦然.

\HRule

* 注:本章节所属知识模块是 “函数与分析”.

\HRule

3. 函数的单调性 (函数定义域内的局部性质,单调性相同但间断的区间一般不能并在一起书写):

(1)单调性的定义:设 ${x}_{1}$ 、 ${x}_{2}$ 是函数 $y = f\left( x\right)$ 定义域内某区间 $I$ 上的任意两个自变量，

若 ${x}_{1} < {x}_{2}$ ,都有 $f\left( {x}_{1}\right)  \leq  f\left( {x}_{2}\right) \left( {f\left( {x}_{1}\right)  \geq  f\left( {x}_{2}\right) }\right)$ ,则 $y = f\left( x\right)$ 在区间 $I$ 上是增 (减) 函数;

若 ${x}_{1} < {x}_{2}$ ,都有 $f\left( {x}_{1}\right)  < f\left( {x}_{2}\right) \left( {f\left( {x}_{1}\right)  > f\left( {x}_{2}\right) }\right)$ ,则 $y = f\left( x\right)$ 在区间 $I$ 上是严格增 (减) 函数;

以上 “若……，都有……” 可改写为: “若 ${x}_{1} \neq  {x}_{2}$ ，都有 $\frac{f\left( {x}_{1}\right)  - f\left( {x}_{2}\right) }{{x}_{1} - {x}_{2}} > 0\left( {\frac{f\left( {x}_{1}\right)  - f\left( {x}_{2}\right) }{{x}_{1} - {x}_{2}} < 0}\right)$ ”，或改写为:“若 ${x}_{1} \neq  {x}_{2}$ ，都有 $\left( {{x}_{1} - {x}_{2}}\right) \left\lbrack  {f\left( {x}_{1}\right)  - f\left( {x}_{2}\right) }\right\rbrack   > 0$ ( $\left( {{x}_{1} - {x}_{2}}\right) \left\lbrack  {f\left( {x}_{1}\right)  - f\left( {x}_{2}\right) }\right\rbrack   < 0$ )”；

(2)规范证明步骤:任意取 ${x}_{1}\text{ 、 }{x}_{2} \in  I \subseteq  D$ 且 ${x}_{1} < {x}_{2}$ ，则 $f\left( {x}_{1}\right)  - f\left( {x}_{2}\right)  = \cdots \cdots ;\because {x}_{1}\text{ 、 }{x}_{2} \in  I \subseteq  D$ 且 ${x}_{1} < {x}_{2},\therefore \cdots \cdots ,\therefore f\left( {x}_{1}\right)  > f\left( {x}_{2}\right)$ (或 $f\left( {x}_{1}\right)  < f\left( {x}_{2}\right) ),\therefore f\left( x\right)$ 在 $I$ 上是减 (增函数);

(3)求单调区间的方法:定义法、图像法；(注意:极限思想的应用、单调区间一定是定义域的子集.

(4)复合函数 $y = f\left\lbrack  {g\left( x\right) }\right\rbrack$ 的单调性(同增异减): $y = f\left\lbrack  {g\left( x\right) }\right\rbrack$ 由 $y = f\left( u\right)$ 和 $u = g\left( x\right)$ 复合而成，若 $x \in  \left( {a,b}\right)$ 时, $u \in  \left( {c,d}\right)$ ,当 $u = g\left( x\right)$ 在 $x \in  \left( {a,b}\right)$ 上的单调性与 $y = f\left( u\right)$ 在 $u \in  \left( {c,d}\right)$ 上的单调性相同 (反) 时 $y = f\left\lbrack  {g\left( x\right) }\right\rbrack$ 在 $x \in  \langle a,b\rangle$ 上单调递增 (减); (其中 “ ( ) ” 随题目的变化确定区间的开闭)

(5)函数单调性的简要运算性质:在公共定义域内，增十增是增函数；减十减是减函数；增一减是增函数; 减一增是减函数.

注意:(4)(5)可作为判断单调性的依据，但不可作为解答题的证明过程.

4. 函数的最值(函数定义域内的整体性质):

(1)最值的定义:

对于函数 $y = f\left( x\right) \left( {x \in  D}\right)$ ,若存在 ${x}_{0} \in  D$ ,使得 $f\left( x\right)  \geq  \left(  \leq  \right) f\left( {x}_{0}\right)$ 对任意 $x \in  D$ 恒成立,则 $f\left( {x}_{0}\right)$ 称为函数 $f\left( x\right)$ 的最小 (大) 值. (本书用 $f{\left( x\right) }_{\max }$ 表示 $f\left( x\right)$ 的最大值， $f{\left( x\right) }_{\min }$ 表示最小值)

(2)函数值域(和最值)的求法:

 \ding{172} 配方法:适用于二次函数(或可转化为二次函数)及区间上的二次函数问题;

 \ding{173} 不等式法:适用于“耐克函数”；注意“一正二定三取等”；

 \ding{174} 换元法:通常用于对函数式进行变形以获得更为熟悉的函数；

 \ding{175} 有界性法:利用已知函数的有界性，借助解不等式求解；

 \circled{5} 单调性法:当所产生的函数单调性较明显或其他方法不太奏效时，考察函数的单调性有时很有效；

 \circled{6} 三角函数的特殊方法:辅助角公式；

 \circled{7} 数列的最大项和最小项的特殊性和特殊方法:最常用的是单调性.

注意:

 \ding{172} 数形结合是解决区间上的函数的最值和值域问题的重要手段，是求解这类问题时避免发生错误的根本保证.  \ding{173}  注意“最值”与“极值”的区别:最值是函数定义域内的整体性质，而极值是局部的.

5. 函数的零点 (函数定义域内的整体性质):

(1)零点的意义:

对于函数 $y = f\left( x\right) \left( {x \in  D}\right)$ ,若存在 ${x}_{0} \in  D$ ,使得 $f\left( {x}_{0}\right)  = 0$ 成立,则 ${x}_{0}$ 称为函数 $f\left( x\right)$ 的零点; 注意函数的零点是函数定义域内的一个实数，而不是有序实数对；

(2)零点存在定理:如果在区间 $\left\lbrack  {a,b}\right\rbrack$ 上，函数 $y = f\left( x\right)$ 的图像满足一段连续曲线，并且 $f\left( a\right)  \cdot  f\left( b\right)  < 0$ ，那么 $y = f\left( x\right)$ 在区间 $\left( {a,b}\right)$ 上一定有零点;

(3)二分法求零点:在确定了区间 $\left( {a,b}\right)$ 上一定有零点的前提下，将区间一分为二，这两部分中总有一个含有零点，而含有零点的区间的长度变为原先的一半，反复执行这种“一分为二”的操作，就能将零点限制在一个足够小的区间中，从而容易求得其近似值了.

6. 函数的周期性(函数定义域内的整体性质):

(1)对于函数 $y = f\left( x\right) \left( {x \in  D}\right)$ ,存在常数 $T \neq  0$ ,使 $f\left( {x + T}\right)  = f\left( x\right)$ 对任意 $x \in  D$ 恒成立,则 $T$ 为 $f\left( x\right)$ 的一个周期;

(2)若偶函数 $f\left( x\right)$ 的图像关于直线 $x = a\left( {a \neq  0}\right)$ 对称，则 $f\left( x\right)$ 是周期函数， $2\left| a\right|$ 是其一个周期；

二若函数 $f\left( x\right)$ 的图像关于直线 $x = a$ 和 $x = b\left( {a \neq  b}\right)$ 对称,则 $f\left( x\right)$ 是周期函数， $2\left| {a - b}\right|$ 是其周期； (3)若奇函数 $f\left( x\right)$ 的图像关于直线 $x = a\left( {a \neq  0}\right)$ 对称，则 $f\left( x\right)$ 是周期函数， $4\left| a\right|$ 是其一个周期；

* 若函数 $f\left( x\right)$ 的图像关于 $\left( {a,0}\right)$ 和 $x = b\left( {a \neq  b}\right)$ 对称,则 $f\left( x\right)$ 是周期函数， $4\left| {a - b}\right|$ 是其一个周期； (4)若奇函数 $f\left( x\right)$ 的图像关于点 $\left( {b,0}\right) \left( {b \neq  0}\right)$ 对称，则 $f\left( x\right)$ 是周期函数， $2\left| b\right|$ 是其一个周期；

二若函数 $f\left( x\right)$ 的图像关于 $\left( {a,0}\right)$ 和 $\left( {b,0}\right) \left( {a \neq  b}\right)$ 对称,则 $f\left( x\right)$ 是周期函数, $2\left| {a - b}\right|$ 是其一个周期.

7. 处理二次函数的问题勿忘数形结合；二次函数在闭区间上必有最值；求解最值问题用“两看”:一看开口方向；二看对称轴与所给区间的相对位置关系.

8.二次函数涉及到方程根的分布问题看:开口、对称轴、端点的函数值与零的关系，特别提醒:端点特别考虑.

\section*{三、函数图像的变换:}

1. 函数图像平移变换 (不改变形状、只改变位置):

(1)将 $y = f\left( x\right)$ 图像沿 $x$ 轴向左 $\left( {a > 0}\right)$ (或右 $\left( {a < 0}\right)$ )平移 $\left| a\right|$ 个单位得到函数 $y = f\left( {x + a}\right)$ 的图像；

(2)将 $y = f\left( x\right)$ 图像沿 $y$ 轴向上 $\left( {b > 0}\right)$ (或下 $\left( {b < 0}\right)$ 〕平移 $\left| b\right|$ 个单位得到函数 $y = f\left( x\right)  + b$ 的图像.

2. 几种特殊的对称变换:

(1)函数 $y = f\left( x\right)$ 与函数 $y = f\left( {-x}\right)$ 的图像关于 $y$ 轴 ( $x = 0$ ) 对称;

(2)函数 $y = f\left( x\right)$ 与函数 $y =  - f\left( x\right)$ 的图像关于 $x$ 轴( $y = 0$ )对称；

(3) 函数 $y = f\left( x\right)$ 与函数 $y =  - f\left( {-x}\right)$ 的图像关于坐标原点对称.

3. 伸缩变换:

(1)将 $y = f\left( x\right)$ 图像上各点的纵坐标伸长 $\left( {a > 1}\right)$ 或缩短 $\left( {0 < a < 1}\right)$ 到原来的 $a$ 倍(横坐标不变)，可得 $y  = {af}\left( x\right) \left( {a > 0}\right)$ 图像;

(2) 将 $y = f\left( x\right)$ 的图像上各点的横坐标伸长 $\left( {0 < a < 1}\right)$ 或缩短 $\left( {a > 1}\right)$ 到原来的 $\frac{1}{a}$ 倍 (纵坐标不变),可得 $y = f\left( {ax}\right) \left( {a > 0}\right)$ 的图像.

4. 翻折变换:

(1) $y = f\left( \left| x\right| \right)$ 的图像作法:保持 $y = f\left( x\right)$ 在 $y$ 轴右侧的图像不变，同时翻折到 $y$ 轴左侧；

(2) $y = \left| {f\left( x\right) }\right|$ 的图像作法:保持 $y = f\left( x\right)$ 在 $x$ 轴上方的图像不变，同时将下方的图像翻折到 $x$ 轴上方.

5. 对于函数 $y = f\left( x\right) \left( {x \in  \mathbf{R}}\right) ,f\left( {x + a}\right)  = f\left( {b - x}\right)$ 恒成立，则函数 $f\left( x\right)$ 的对称轴是 $x = \frac{a + b}{2}$ ；两个函数 $y = f\left( {x + a}\right)$ 与 $y = f\left( {b - x}\right)$ 的图像关于直线 $x = \frac{b - a}{2}$ 对称.

\section*{$\text{ § }{3.1}$ 函数的概念及其表示法}

\section*{典型例题精讲}

【例 1】已知函数 $f\left( x\right)  = \left\{  \begin{array}{l} {x}^{2},x < 0, \\  4 - {2x},x > 0. \end{array}\right.$

\begin{center}
\includegraphics[width=0.2\textwidth]{images/bo_d7tk7oc91nqc73875fu0_32_1141_538_336_339_0.jpg}
\end{center}
 

(例 1 图)

(1)画出函数 $f\left( x\right)$ 的图像；

(2)若 $x \in  \left\lbrack  {-2,0)\cup (0,3}\right\rbrack$ ，求其值域；

(3)当 $f\left( x\right)  \geq  2$ 时,求实数 $x$ 的取值范围.

【例 2】已知函数 $f\left( {{\log }_{3}x}\right)  = \frac{{2x} - 3}{x - 1}$ .

(1)求函数 $f\left( x\right)$ 的解析式及定义域；

(2)求函数 $f\left( x\right)$ 在 $x \in  \left( {-\infty ,0}\right)$ ∪ $\left( {0,2}\right)$ 时的值域.

【例 3】( 1 )已知函数 $f\left( x\right)$ 的定义域是 $\lbrack 1,7)$ ，求函数 $f\left( {{2x} - 3}\right)$ 的定义域；

(2)若函数 $f\left( \left| {x - 2}\right| \right)$ 的定义域是 $\lbrack 1,7)$ ，求 $f\left( x\right)$ 的定义域；

(3)若函数 $f\left( {{3x} - 2}\right)$ 的定义域是 $\lbrack 1,7)$ ，求 $f\left( {{2x} - 1}\right)$ 的定义域.

【例 4】已知函数 $g\left( x\right)  = \sqrt{x} + 1,h\left( x\right)  = \frac{1}{x + 3}\left( {-3 < x \leq  a}\right)$ ,其中 $a$ 为常数,且 $a > 0$ ,令函数 $f\left( x\right)  =  g\left( x\right) h\left( x\right)$ .

(1)求函数 $f\left( x\right)$ 的解析式，指出其定义域，并求当 $a = \frac{1}{4}$ 时，函数 $f\left( x\right)$ 的值域；

(2)是否存在正整数 $a$ ，使得函数 $f\left( x\right)$ 的值域恰为 $\left\lbrack  {\frac{1}{3},\frac{1}{2}}\right\rbrack$ ? 若存在，试写出所有满足条件的正整数 $a$ 所构成的集合; 若不存在,试说明理由.

【例 5】如图, $\bigtriangleup {OAB}$ 是边长为 2 的正三角形,记 $\bigtriangleup {OAB}$ 位于直线 $x = t\left( {t > 0}\right)$ 左侧的图形的面积为 $f\left( t\right)$ .

(1)求函数 $f\left( t\right)$ 解析式；

(2)当函数 $g\left( t\right)  = f\left( t\right)  - {at}$ 有且只有一个零点时，求实数 $a$ 的值.

\begin{center}
\includegraphics[width=0.2\textwidth]{images/bo_d7tk7oc91nqc73875fu0_33_1139_1055_298_267_0.jpg}
\end{center}
 

(例 5 图)

\section*{四基达标精练}

\section*{一、填空题}

1. 函数 $f\left( x\right)  = \frac{1}{x - 1} + \ln x$ 的定义域是\_\_\_\_\_.

2. 已知函数 $f\left( x\right)  = \left\{  \begin{array}{l} {\log }_{2}x,x > 0, \\  {3}^{x},x \leq  0, \end{array}\right.$ 则 $f\left\lbrack  {f\left( \frac{1}{4}\right) }\right\rbrack   =$ \_\_\_\_\_.

3. 若函数 $f\left( {{2x} + 1}\right)  = {4}^{x} - 2$ ，则函数 $f\left( x\right)$ 的解析式为\_\_\_\_\_.

4. 已知 $f\left( x\right)  = \left\{  {\begin{array}{l} {x}^{2},x \geq  0, \\   - x,x < 0, \end{array}g\left( x\right)  = f\left( {-x}\right) }\right.$ ，则 $g\left( x\right)$ 的解析式是\_\_\_\_\_.

5. 已知 $f\left( {x + 1}\right)$ 的定义域为 $\left\lbrack  {2,3}\right\rbrack$ ，则 $f\left( {{2x} - 1}\right)$ 的定义域为\_\_\_\_\_.313 (2).

6. 设 $f\left( x\right)  = \left\{  \begin{array}{l} {\left( \frac{1}{2}\right) }^{x} - 3,x \leq  0, \\  {x}^{\frac{1}{2}},x > 0, \end{array}\right.$ 若 $f\left( a\right)  < 1$ ，则实数 $a$ 的范围为\_\_\_\_\_.

7. 若 $y = \sqrt{\left| {x + 2}\right|  - \left| {x - m}\right|  + 1}$ 的定义域为 $\mathbf{R}$ ，则实数 $m$ 的取值范围是\_\_\_\_\_.

8. 用长为 $l$ 的铁丝弯成下部为矩形,上部为半圆形的框架,若矩形底边长为 ${2x}$ ,则此框架围成的面积 $y$ 关于 $x$ 的函数解析式为\_\_\_\_\_.

9. 已知 $a \in  \mathbf{R}$ ，若存在定义域为 $\mathbf{R}$ 的函数 $f\left( x\right)$ 同时满足下列两个条件: \ding{172} 对任意 ${x}_{0} \in  \mathbf{R}$ ， $f\left( {x}_{0}\right)$ 的值为 ${x}_{0}$ 或 ${x}_{0}^{2}$ ;  \ding{173} 关于 $x$ 的方程 $f\left( x\right)  = a$ 无实数解; 则 $a$ 的取值范围是\_\_\_\_\_.

\section*{二、选择题}

10. 某校有一班级,设变量 $x$ 是该班同学的姓名,变量 $y$ 是该班同学的学号,变量 $z$ 是该班同学的身高, 变量 $w$ 是该班同学的数学考试成绩,则下列选项中正确的是 ( )

A. $y$ 是 $x$ 的函数. B. $w$ 是 $y$ 的函数. C. $w$ 是 $z$ 的函数. D. $w$ 是 $x$ 的函数.

11. 存在函数 $y = f\left( x\right)$ ,满足对任意 $x \in  \mathbf{R}$ 都有 ( )

A. $f\left( {\left| x\right|  - 1}\right)  = \left| {x + 1}\right|$ . B. $f\left( \left| {x - 1}\right| \right)  = \left| {x + 1}\right|$ .

C. $f\left( {{x}^{2} + {2x}}\right)  = \left| {x + 1}\right|$ . D. $f\left( {{x}^{2} - 1}\right)  = \left| {x + 1}\right|$ .

12. 四参数方程的拟合函数表达式为 $y = \frac{a - d}{1 + {\left( \frac{x}{c}\right) }^{b}} + d\left( {x > 0}\right)$ ,常用于竞争系统和免疫检测,它的图像是一个递增 (或递减) 的类似指数曲线或对数曲线或双曲线 (如 $y = {x}^{-1}$ ),还可以是一条 ${S}$ 形曲线. 当 $a = 4$ , $b =  - 1,c = 1,d = 1$ 时,该拟合函数图像是 ( )

A. 类似递增的双曲线. B. 类似递增的对数曲线.

C. 类似递减的指数曲线. D. 是一条 S 形曲线.

\section*{三、解答题}

13. 已知函数 $f\left( x\right)  = \left\{  \begin{array}{l} \frac{2}{x},x < 0, \\   - x,0 \leq  x \leq  2, \\  \frac{1}{2}x - 3,x > 2. \end{array}\right.$

(1)求 $f\left( 0\right)$ 和 $f\left\lbrack  {f\left( 1\right) }\right\rbrack$ 的值；

\begin{center}
\includegraphics[width=0.3\textwidth]{images/bo_d7tk7oc91nqc73875fu0_34_1055_1435_422_378_0.jpg}
\end{center}
 

(第 13 题图)

(2)若 $f\left( m\right)  =  - 1$ ，求 $m$ 的值；

(3)在给定的坐标系中，作出函数 $y = f\left( x\right)$ 的图像.

14. 函数 $f\left( x\right)  = \sqrt{2 - \frac{x + 3}{x + 1}}$ 的定义域为 $A,g\left( x\right)  = \lg \left\lbrack  {\left( {x - a - 1}\right) \left( {{2a} - x}\right) }\right\rbrack  \left( {a < 1}\right)$ 的定义域为 $B$ .

(1)求 $A$ ；

(2)若 $B \subseteq  A$ ，求实数 $a$ 的取值范围.

15. 已知函数 $f\left( {1 - {2x}}\right)$ 的定义域为 $A = \left\lbrack  {\frac{1}{2},1}\right\rbrack$ .

(1)求 $f\left( x\right)$ 的定义域 $B$ ；

(2)对于(1)中的集合 $B$ ，若存在 $x \in  B$ ，使得 $a > {x}^{2} - x + 1$ 成立，求实数 $a$ 的取值范围.

16. 已知函数 $f\left( x\right)  = \frac{{x}^{2} - {3x}}{\sqrt{{2x} + 1}}$ 和 $g\left( x\right)  = \frac{\sqrt{{2x} + 1}}{x - 3}$ ,设 $h\left( x\right)  = f\left( x\right)  \cdot  g\left( x\right)$ .

(1)若函数 $H\left( x\right)  = x$ ，试判断 $y = h\left( x\right)$ 与 $y = H\left( x\right)$ 是否为同一函数，并说明理由；

(2)求 $G\left( x\right)  = h\left( x\right)  - \sqrt{{2h}\left( x\right)  + 1}$ 的值域.

\section*{$\text{ § }{3.2}$ 函数的奇偶性}

\section*{典型例题精讲}

【例 1】已知函数 $f\left( x\right)  = a{x}^{2} + \frac{b}{x}\left( {x \neq  0}\right.$ ,常数 $\left. {a\text{ 、 }b \in  \mathbf{R}}\right)$ .

(1)当 $a = 0,b = 1$ 时，解关于 $x$ 的不等式: $f\left( x\right)  > k\left( {k \in  \mathbf{R}}\right)$ ；

(2)讨论函数 $f\left( x\right)$ 的奇偶性,并说明理由.

【例 2】已知函数 $f\left( x\right)  = {x}^{2} + 2\left| {x - a}\right| \left( {a \in  \mathbf{R}}\right)$ .

( 1 )若函数 $f\left( x\right)$ 为偶函数，求 $a$ 的值；

( 2 )若函数 $g\left( x\right)  = {af}\left( x\right)  + 2$ 的最小值为 8，求 $a$ 的值.

【例 3】已知函数 $f\left( x\right)  = \left\{  {\begin{array}{l} {x}^{2} + {ax},x \geq  0, \\  {x}^{2} + {2x},x < 0 \end{array}\left( {a \in  \mathbf{R}}\right) }\right.$ .

(1)若 $f\left( x\right)$ 是偶函数，求 $a$ 的值；

(2)若 $f\left( x\right)$ 在 $x =  - 1$ 时取到最小值，求 $a$ 的取值范围.

【例 4】已知 $f\left( x\right)$ 是定义在 $\mathbf{R}$ 上的函数,对任意 $x$ 、 $y \in  \mathbf{R}$ ，都有 $f\left( {x + y}\right)  = f\left( x\right)  + f\left( y\right)$ ，且当 $x > 0$ 时, $f\left( x\right)  < 0,f\left( {-1}\right)  = 2$ .

(1)求证: $f\left( x\right)$ 为奇函数；

(2)求证: $f\left( x\right)$ 是 $\mathbf{R}$ 上的减函数；

(3)求 $f\left( x\right)$ 在 $\left\lbrack  {-2,4}\right\rbrack$ 上的最值.

【例 5】对于函数 $f\left( x\right)$ ,若在定义域内存在实数 ${x}_{0}$ ,满足 $f\left( {-{x}_{0}}\right)  =  - f\left( {x}_{0}\right)$ ,则称 $f\left( x\right)$ 为“ $M$ 类函数”. (1)设 $f\left( x\right)  = {2}^{x} + m$ 是定义在 $\left\lbrack  {-1,1}\right\rbrack$ 上的“ $M$ 类函数”，求实数 $m$ 的最小值；

(2)若 $f\left( x\right)  = \left\{  \begin{array}{l} {\log }_{2}\left( {{x}^{2} - {2mx}}\right) ,x \geq  2, \\   - 3,x < 2 \end{array}\right.$ 为定义域上的“ $M$ 类函数”，求实数 $m$ 的取值范围.

\section*{四基达标精练}

\section*{一、填空题}

1. 已知 $y = f\left( x\right)$ 是奇函数，当 $x \geq  0$ 时， $f\left( x\right)  =  - {x}^{2}$ ，则 $f\left( {-1}\right)$ 的值是\_\_\_\_\_.

2. 若定义在 $\left( {-4,{\log }_{2}a}\right)$ 上的函数 $y = f\left( x\right)$ 是偶函数，则实数 $a =$ \_\_\_\_\_.

3. 若函数 $y = {x}^{2} + \left( {a - 3}\right) x + 4$ 为偶函数，则 $a =$ \_\_\_\_\_.

4. 已知函数 $f\left( x\right)$ 是定义在 $\mathbf{R}$ 上的偶函数. 当 $x \in  \left( {-\infty ,0}\right)$ 时, $f\left( x\right)  = {2}^{x} - {x}^{2}$ ,则当 $x \in  \left( {0, + \infty }\right)$ 时, $f\left( x\right)  =$

5. 函数 $f\left( x\right)  = a{x}^{3} - {bx} - \tan x + 2$ ，若 $f\left( m\right)  = 1$ ，则 $f\left( {-m}\right)  =$ \_\_\_\_\_.

6. 设 $y = f\left( x\right)$ 为定义在 $\mathbf{R}$ 上的奇函数，当 $x \geq  0$ 时， $f\left( x\right)  = {2}^{x} + {2x} + b$ ( $b$ 为常数),则 $f\left( {-1}\right)  =$

7. 设 $f\left( x\right)$ 是定义域为 $\mathbf{R}$ 的奇函数, $g\left( x\right)$ 是定义域为 $\mathbf{R}$ 的偶函数,若函数 $f\left( x\right)  + g\left( x\right)$ 的值域为 $\lbrack 1,3)$ , 则函数 $f\left( x\right)  - g\left( x\right)$ 的值域为\_\_\_\_\_.

8. 函数 $f\left( x\right)  = \frac{\sqrt{a - {x}^{2}}}{\left| {x + 1}\right|  - 1}$ 是奇函数，则实数 $a$ 的取值范围是\_\_\_\_\_.

9. 已知函数 $f\left( x\right)  = {x}^{2} - \cos x\left( {-\frac{\pi }{2} \leq  x \leq  \frac{\pi }{2}}\right)$ ,如果 $f\left( x\right)  < f\left( \frac{\pi }{3}\right)$ ,那么实数 $x$ 的取值范围是\_\_\_\_\_.

\section*{二、选择题}

10. 函数 $f\left( x\right)  =  - x + \frac{1}{2x}$ 的图像关于 ( )

A. $y$ 轴对称. B. 直线 $y =  - x$ 对称. C. 坐标原点对称. D. 直线 $y = x$ 对称.

11. 若 $f\left( x\right)$ 为奇函数，且 ${x}_{0}$ 是 $y = f\left( x\right)  - {{e}}^{x}$ 的一个零点，则 $- {x}_{0}$ 一定是下列哪个函数的零点( )

D. $y = f\left( {-x}\right) {{e}}^{x} + 1$ .



\begin{center}
\includegraphics[width=0.2\textwidth]{images/bo_d7tk7oc91nqc73875fu0_39_1176_155_266_243_0.jpg}
\end{center}
 

(第 12 题图)

12. 如图，在平面直角坐标系内有一个边长为 $a$ 、中心在原点 $O$ 的正六边形 ${ABCDEF},{AB}//{Ox}$ . 直线 $l : y = {kx} + t\left( {k\text{ 为常数 }}\right)$ 与正六边形的边交于 $M\text{ 、 }N$ 两点, 记 $\bigtriangleup {OMN}$ 的面积为 $S$ ,则函数 $S = f\left( t\right)$ 的奇偶性为 ( )

A. 偶函数.

B. 奇函数.

C. 不是奇函数，也不是偶函数.

D. 奇偶性与 $k$ 有关.

\section*{三、解答题}

13. 已知 $y = f\left( x\right)$ 是定义在 $\mathbf{R}$ 上的奇函数,当 $x \geq  0$ 时, $f\left( x\right)  =  - {x}^{2} + {2x}$ .

(1)求函数 $f\left( x\right)$ 的解析式；

(2)若函数 $f\left( x\right)$ 在 $\left\lbrack  {-1,a - 2}\right\rbrack$ 上单调递增，求实数 $a$ 的取值范围.

14. 已知函数 $f\left( x\right)  = \ln \left( {3 + x}\right)  + \ln \left( {3 - x}\right)$ .

(1)判断函数 $f\left( x\right)$ 的奇偶性,并证明你的结论;

(2)若 $f\left( {{2m} - 1}\right)  < f\left( m\right)$ ，求 $m$ 的取值范围.

15. 已知函数 $f\left( x\right)  = \frac{\left( {x + 1}\right) \left( {x + a}\right) }{{x}^{2}}$ 为偶函数.

(1)求实数 $a$ 的值；

(2)当 $x \in  \left\lbrack  {\frac{1}{m},\frac{1}{n}}\right\rbrack  \left( {m\text{ 、 }n > 0}\right)$ 时,若函数 $f\left( x\right)$ 的值域为 $\left\lbrack  {2 - {3m},2 - {3n}}\right\rbrack$ ，求 $m\text{ 、 }n$ 的值.

16. 已知函数 $f\left( x\right)  = \frac{-{2}^{x} + 1}{{2}^{x + 1} + 2}$ .

(1)判断函数 $f\left( x\right)$ 的奇偶性,并证明;

(2)若不等式 $f\left( x\right)  > {\log }_{9}\left( {{2c} - 1}\right)$ 有解，求 $c$ 的取值范围.

\section*{§ 3.3 函数的单调性}

\section*{典型例题精讲}

【例 1】已知函数 $f\left( x\right)  = \frac{2x}{x - 1}$ .

(1)求 $f\left( x\right)$ 的定义域、值域并写出其单调区间及单调性(不要求证明)；

(2)判断并用定义证明函数 $g\left( x\right)  = {xf}\left( x\right)$ 在区间 $\left( {0,1}\right)$ 上的单调性.

【例 2】设 $a$ 为实数，函数 $f\left( x\right)  = {\left( x - a\right) }^{2} + \left| {x - a}\right|  - a\left( {a - 1}\right)$ .

(1)若 $f\left( 0\right)  \leq  1$ ，求 $a$ 的取值范围；

(2)讨论 $f\left( x\right)$ 的单调性.

【例 3】已知函数 $f\left( x\right)  = \left| {x - a}\right|  - \frac{a}{2}\ln x\left( {a \in  \mathbf{R},x > 0}\right) ,g\left( x\right)  = \left| {x - 2}\right|$ .

(1)若 $a = 0$ ，求函数 $F\left( x\right)  = f\left( x\right)  - g\left( x\right)$ 的值域；

(2)求函数 $f\left( x\right)$ 的单调区间.

【例 4】已知函数 $f\left( x\right)  = a{x}^{3} + x - a\left( {a \in  \mathbf{R}}\right) ,g\left( x\right)  = \frac{x}{1 - {x}^{3}}$ .

(1)判断 $g\left( x\right)$ 在 $\left( {-1, - \frac{\sqrt[3]{4}}{2}}\right\rbrack$ 和 $\left\lbrack  {-\frac{\sqrt[3]{4}}{2},1}\right)$ 的单调性,并说明理由;

(2)证明:函数 $f\left( x\right)$ 存在零点 $q$ ，使得 $a = q + {q}^{4} + {q}^{7} + \cdots  + {q}^{{3n} - 2} + \cdots$ 成立的充要条件是 $a \geq  \frac{-\sqrt[3]{4}}{3}$ .

【例 5】已知函数 $f\left( x\right)  = a\left( {x + \frac{1}{x}}\right)  - \left| {x - \frac{1}{x}}\right| \left( {x > 0}\right) ,a \in  \mathbf{R}$ .

(1)若 $a = \frac{1}{2}$ ，求 $y = f\left( x\right)$ 的单调区间；

(2)若关于 $x$ 的方程 $f\left( x\right)  = t$ 有四个不同的解,求实数 $a\text{ 、 }t$ 应满足的条件.

\section*{四基达标精练}

\section*{一、填空题}

1. 函数 $f\left( x\right)  = {x}^{2} - {3x} - 1$ 的严格递增区间是\_\_\_\_\_.

2. 函数 $y = \sqrt{-{x}^{2} - {2x} + 8}$ 的单调递减区间为\_\_\_\_\_.

3. 若函数 $f\left( x\right)  = \frac{5}{{3x} - 2}$ 的严格递减区间是\_\_\_\_\_.

4. 若函数 $f\left( x\right)  = {x}^{2} + 2\left( {a - 1}\right) x + 2$ 在 $( - \infty ,1\rbrack$ 上递减，则实数 $a$ 的取值范围是\_\_\_\_\_.

5. 若函数 $f\left( x\right)  = {x}^{2} + 2\left( {a - 1}\right) x + 2$ 在 $\left\lbrack  {-1,1}\right\rbrack$ 上是增函数，则实数 $a$ 的取值范围是\_\_\_\_\_.

6. 能说明“若 $f\left( x\right)  > f\left( 0\right)$ 对任意的 $x \in  (0,2\rbrack$ 都成立，则 $f\left( x\right)$ 在 $\left\lbrack  {0,2}\right\rbrack$ 上是增函数”为假命题的一个函数是\_\_\_\_\_.

7. 设 $a$ 为常数，函数 $f\left( x\right)  = {x}^{2} - {4x} + 3$ ，若 $f\left( {x + a}\right)$ 在 $\lbrack 0, + \infty )$ 上是增函数，则 $a$ 的取值范围是\_\_\_\_\_.

8. 若函数 $f\left( x\right)  = {x}^{2} + \left( {{2a} + 1}\right) \left| x\right|  + 1$ 的定义域被分成四个不同的单调区间，则实数 $a$ 的取值范围是\_\_\_\_\_.

9. 设函数 $y = f\left( x\right)$ 定义在 $\mathbf{R}$ 上,对于给定的正数 $K$ ,定义函数 ${f}_{K}\left( x\right)  = \left\{  \begin{array}{l} f\left( x\right) ,f\left( x\right)  \leq  K, \\  K,f\left( x\right)  > K. \end{array}\right.$ 若函数 $f\left( x\right)  = {2}^{-\left| x\right| }$ ,当 $K = \frac{1}{2}$ 时,函数 ${f}_{K}\left( x\right)$ 的严格递增区间为\_\_\_\_\_.

\section*{二、选择题}

10. 已知函数 $f\left( x\right)  = \frac{{ax} + 1}{x + 2}\left( {a \in  \mathbf{R}}\right)$ ,则 “ $f\left( 2\right)  < f\left( 3\right)$ ” 是 “ $f\left( x\right)$ 在区间 $\left( {-2, + \infty }\right)$ 上单调递增” 的

( )

B. 必要非充分条件. C. 充要条件. D. 既非充分也非必要条件.

11. 若函数 $y = f\left( x\right)$ 在区间 $I$ 上是增函数,且函数 $y = \frac{f\left( x\right) }{x}$ 在区间 $I$ 上是减函数,则称函数 $f\left( x\right)$ 是区间 $I$ 上的“ $H$ 函数”,对于命题: \ding{172} 函数 $f\left( x\right)  =  - x + 2\sqrt{x}$ 是 $\left( {0,1}\right)$ 上的“ $H$ 函数”;  \ding{173} 函数 $g\left( x\right)  = \frac{2x}{1 - {x}^{2}}$ 是 $\left( {0,1}\right)$ 上的“ $H$ 函数”; 下列判断正确的是 ( )

A.  \ding{172} 和 \ding{173} 均为真命题. B.  \ding{172} 为真命题， \ding{173} 为假命题.

C.  \ding{172} 为假命题， \ding{173} 为真命题. D.  \ding{172} 和 \ding{173} 均为假命题.

12. 命题 $p$ : 存在 $a \in  \mathbf{R}$ 且 $a \neq  0$ ,对于任意的 $x \in  \mathbf{R}$ ,使得 $f\left( {x + a}\right)  < f\left( x\right)  + f\left( a\right)$ ; 命题 ${q}_{1} : f\left( x\right)$ 单调递减且 $f\left( x\right)  > 0$ 恒成立; 命题 ${q}_{2} : f\left( x\right)$ 单调递增,存在 ${x}_{0} < 0$ 使得 $f\left( {x}_{0}\right)  = 0$ ,则下列说法正确的是 ( )

A. 只有 ${q}_{1}$ 是 $p$ 的充分条件. B. 只有 ${q}_{2}$ 是 $p$ 的充分条件.

C. ${q}_{1}\text{ 、 }{q}_{2}$ 都是 $p$ 的充分条件. D. ${q}_{1}\text{ 、 }{q}_{2}$ 都不是 $p$ 的充分条件.

\section*{三、解答题}

13. 已知函数 $f\left( x\right)  = 2 + \frac{1}{a} - \frac{1}{{a}^{2}x}\left( {a \in  \mathbf{R},a \neq  0}\right)$ .

(1)设 ${mn} > 0$ ，判断函数 $f\left( x\right)$ 在 $\left\lbrack  {m,n}\right\rbrack$ 上的单调性，并加以证明；

(2)若不等式 ${a}^{2}f\left( x\right)  \leq  {2x}$ 对 $x \geq  1$ 恒成立，求实数 $a$ 的取值范围.

14. 已知函数 $f\left( x\right)  = \frac{1 + {x}^{2}}{x}$ .

(1)求函数的定义域和值域；

(2)求函数的单调区间.



15. 某群体的人均通勤时间，是指单日内该群体中成员从居住地到工作地的平均用时. 某地上班族 $S$ 群体中的成员仅以自驾或公交方式通勤. 分析显示: 当 $S$ 中 $x\% \left( {0 < x < {100}}\right)$ 的成员自驾时,自驾群体的人均通勤时间为 $f\left( x\right)  = \left\{  \begin{array}{l} {30},0 < x \leq  {30}, \\  {2x} + \frac{1800}{x} - {90},{30} < x < {100} \end{array}\right.$ (单位: 分钟),而公交群体的人均通勤时间不受 $x$ 影响,恒为 40 分钟. 试根据上述分析结果回答下列问题:

(1)当 $x$ 在什么范围内时,公交群体的人均通勤时间少于自驾群体的通勤时间？

(2)求该地上班族 $S$ 群体的人均通勤时间 $g\left( x\right)$ 的表达式:讨论 $g\left( x\right)$ 的单调性，并说明其实际意义.

16. 已知函数 $f\left( x\right)  = \frac{{2}^{x} + a}{{2}^{x} + 1},x \in  \mathbf{R}$ .

(1)证明:当 $a > 1$ 时，函数 $y = f\left( x\right)$ 在 $\mathbf{R}$ 上是严格减函数；

(2)求函数 $y = f\left( x\right)$ 的值域.

\section*{§ 3.4 函数的最值和值域}

\section*{典型例题精讲}

【例 1】已知函数 $f\left( x\right)  = x + b$ 在区间 $\left\lbrack  {1,2}\right\rbrack$ 上的最大值为 2,且函数 $g\left( x\right)  = \frac{x + b}{{x}^{2} + 1}$ .

(1)求出 $b$ 的值，判断 $g\left( x\right)$ 的奇偶性，并证明你的结论；

(2)设 $h\left( x\right)  =  - x - {2c}$ ，对任意的 ${x}_{1} \in  \mathbf{R}$ ，总存在 ${x}_{2} \in  \left\lbrack  {-2,2}\right\rbrack$ ，使得 $g\left( {x}_{1}\right)  = h\left( {x}_{2}\right)$ 成立，求实数 $c$ 的取值范围.

【例 2】已知函数 $f\left( x\right)$ 的定义域为 $D$ ,对于给定的正整数 $k$ ,若存在 $\left\lbrack  {a,b}\right\rbrack   \subseteq  D$ ,使得函数 $f\left( x\right)$ 满足: 函数 $f\left( x\right)$ 在 $\left\lbrack  {a,b}\right\rbrack$ 上是单调函数,且 $f\left( x\right)$ 的最小值为 ${ka}$ ,最大值为 ${kb}$ ,则称函数 $f\left( x\right)$ 是 “倍缩函数”,区间 $\left\lbrack  {a,b}\right\rbrack$ 是函数 $f\left( x\right)$ 的“ $k$ 倍值区间”.

(1)判断函数 $f\left( x\right)  = {x}^{3}$ 是否是“倍缩函数”? 请说明理由;

(2)设函数 $g\left( x\right)  = \frac{8x}{4{x}^{2} + 1},x \in  \left\lbrack  {0,\frac{1}{2}}\right\rbrack$ ，若函数 $g\left( x\right)$ 存在 “ $k$ 倍值区间”，求 $k$ 的值.



【例 3】已知 $f\left( x\right)  = {\log }_{a}\left( {{ax} - 2}\right)$ (其中 $a > 0$ 且 $a \neq  1$ ).

(1)若 $a = 2$ ， $f\left( x\right)  < 2$ ，求实数 $x$ 的取值范围；

(2)若 $x \in  \left\lbrack  {4,6}\right\rbrack$ ， $f\left( x\right)$ 的最大值大于 1，求 $a$ 的取值范围.

【例 4】已知函数 $f\left( x\right)  = \sqrt{\left| {x + 2}\right|  + \left| {x - 4}\right|  - m}$ 的定义域为 $\mathbf{R}$ .

(1)求实数 $m$ 的范围；

(2)若 $m$ 的最大值为 $n$ ，当正数 $a$ 、 $b$ 满足 $\frac{4}{a + {5b}} + \frac{1}{{3a} + {2b}} = n$ 时，求 ${4a} + {7b}$ 的最小值.

【例 5】已知 $a \in  \mathbf{R}$ ,函数 $f\left( x\right)  = x\left| {x - a}\right|$ .

(1)判断函数 $f\left( x\right)$ 的奇偶性,请说明理由;

(2)设 $a > 0$ ，求函数 $f\left( x\right)$ 在区间 $\left\lbrack  {1,3}\right\rbrack$ 上的最小值.

\section*{四基达标精练}

\section*{一、填空题}

1. 函数 $f\left( x\right)  = \left| {x - 3}\right|  + \left| {x + 1}\right|$ 的最小值是\_\_\_\_\_.

2. 函数 $y = \frac{{x}^{2} + 5}{\sqrt{{x}^{2} + 4}}$ 的最小值为\_\_\_\_\_.

3. 函数 $y = \sqrt{x\left( {3 - {2x}}\right) }$ 的值域是\_\_\_\_\_.

4. 若函数 $f\left( x\right)  = {\left( \frac{1}{3}\right) }^{{ax} + {a}^{2}}$ 在 $\lbrack 1, + \infty )$ 上有最大值 $\frac{1}{9}$ ，则实数 $a =$ \_\_\_\_\_.

5. 定义 $\min \{ a,b\}  = \left\{  \begin{array}{l} a,a \leq  b, \\  b,a > b, \end{array}\right.$ 则函数 $y = \min \{ {3x} - 1, - x + 2\}$ 的最大值为\_\_\_\_\_.

6. 函数 $y = \sqrt{{2x} + 1} - \sqrt{2 - x}$ 的值域是\_\_\_\_\_.

7. 若函数 $f\left( x\right)  = x + \frac{m}{x}\left( {1 < x < 2}\right)$ 存在最大值或最小值，则实数 $m$ 的取值范围是\_\_\_\_\_.

8. 函数 $f\left( x\right)  = \left| {x - 1}\right|  - \ln x$ 的最小值为\_\_\_\_\_.

9. 设 ${M}_{I}$ 表示函数 $f\left( x\right)  = \left| {{x}^{2} - {4x} + 2}\right|$ 在闭区间 $I$ 上的最大值. 若 $a > 0$ ，且 ${M}_{\left\lbrack  0,a\right\rbrack  } \geq  2{M}_{\left\lbrack  a,2a\right\rbrack  }$ ，则正实数 $a$ 的取值范围是\_\_\_\_\_.

\section*{二、选择题}

10. 已知 $f\left( x\right)$ 是定义在 $\left\lbrack  {0,1}\right\rbrack$ 上的函数，那么 “函数 $f\left( x\right)$ 在 $\left\lbrack  {0,1}\right\rbrack$ 上单调递增” 是 “函数 $f\left( x\right)$ 在 $\left\lbrack  {0,1}\right\rbrack$ 上的最大值为 $f\left( 1\right)$ ”的 ( )

A. 充分非必要条件. B. 必要非充分条件. C. 充要条件. D. 既非充分也非必要条件.

11. 关于 “函数 $f\left( x\right)  = \frac{x - \sqrt{13}}{x - \sqrt{14}},x \in  \left( {-\infty ,\sqrt{14}}\right)  \cup  \left( {\sqrt{14}, + \infty }\right)$ 的最大、最小值与函数 $g\left( x\right)  =  \frac{x - \sqrt{13}}{x - \sqrt{14}},x \in  \mathbf{Z}$ 的最大、最小值”,下列说法中正确的是 ( )

A. $f\left( x\right)$ 有最大、最小值, $g\left( x\right)$ 有最大、最小值. B. $f\left( x\right)$ 有最大、最小值, $g\left( x\right)$ 无最大、最小值.

C. $f\left( x\right)$ 无最大、最小值, $g\left( x\right)$ 有最大、最小值. D. $f\left( x\right)$ 无最大、最小值, $g\left( x\right)$ 无最大、最小值.



12. 若函数 $f\left( x\right)  = \left| {x - a}\right|  + b\left( {-1 \leq  x \leq  1}\right)$ 的最大值是 $M$ ,最小值是 $m$ ,则 $m - n$ ( )

A. 与 $a$ 有关,且与 $b$ 有关. B. 与 $a$ 有关,但与 $b$ 无关.

C. 与 $a$ 无关，且与 $b$ 无关. D. 与 $a$ 无关，但与 $b$ 有关.

\section*{三、解答题}

13. 设函数 $f\left( x\right)  = a{x}^{2} + \left( {b - 2}\right) x + 3\left( {a \neq  0}\right)$ .

(1)若关于 $x$ 的不等式 $f\left( x\right)  > 0$ 的解集为 $\left( {-1,1}\right)$ ，求实数 $a$ 和 $b$ 的值；

(2)若 $f\left( 1\right)  = 3$ ，求函数 $f\left( x\right)$ 在区间 $\left\lbrack  {1,3}\right\rbrack$ 上的最小值.

14. 已知二次函数 $f\left( x\right)$ 满足 $f\left( x\right)  = f\left( {4 - x}\right) ,f\left( 2\right)  = f\left( 1\right)  - 1$ ,若不等式 $f\left( x\right)  \leq   - {2x} + 2$ 有唯一实数解.

(1)求函数 $f\left( x\right)$ 的解析式；

(2)若函数 $f\left( x\right)$ 在 $\left\lbrack  {t,t + 1}\right\rbrack$ 上的最小值为 $g\left( t\right)$ . 求 $g\left( t\right)$ 的解析式.

15. 已知函数 $f\left( x\right)  = \left| {a{x}^{2} - b}\right|  + a{x}^{2} + {bx}\left( {a > 0}\right)$ .

(1)若 $a = b = 1$ ，求函数 $f\left( x\right)$ 的最小值；

(2)若 $b \leq  0$ ，函数 $f\left( x\right)$ 存在两个不同的零点 ${x}_{1}$ 和 ${x}_{2}$ ，求 $\frac{{x}_{2}}{{x}_{1}} + \frac{{x}_{1}}{{x}_{2}}$ 的取值范围.

16. 已知函数 $f\left( x\right)  = \left| {x - 1}\right|  + {\left( x - 1\right) }^{2} + \frac{\sqrt{3}}{3}$ .

(1)求 $g\left( x\right)  = f\left( x\right)  + \left| {f\left( x\right)  - \frac{\sqrt{3}}{3}}\right|$ 的最小值；

(2)若存在实数 $x$ ，使得 $\frac{\sqrt{3}}{{4f}\left( x\right) } + \frac{{4f}\left( x\right) }{\sqrt{3}} \leq  \left| {{m}^{2} - {2m} - 1}\right|$ 成立，求 $m$ 的取值范围.

\section*{$\text{ § }{3.5}$ 函数的零点}

\section*{典型例题精讲}

【例 1】已知函数 $f\left( x\right)  = \frac{a}{x - 1},a$ 为常数.

(1)当 $a > 0$ 时，根据定义证明函数 $f\left( x\right)$ 在区间 $\left( {1, + \infty }\right)$ 上单调递减；

(2)设 $g\left( x\right)  = f\left( x\right)  + {ax} - 3$ ,若 $g\left( x\right)$ 在 $\left( {1, + \infty }\right)$ 上存在两个零点,求实数 $a$ 的取值范围.

【例 2】已知函数 $f\left( x\right)  = {x}^{3} + {3a}{x}^{2} + {bx} + {a}^{2}$ 在 $x =  - 1$ 时有极值 0 .

(1)求函数 $f\left( x\right)$ 的解析式；

(2)记 $h\left( x\right)  = f\left( x\right)  - m + 1$ ，若函数 $h\left( x\right)$ 有三个零点，求实数 $m$ 的取值范围.

【例 3】设函数 $f\left( x\right)  = \left( {x - a}\right) \left( {x - b}\right) \left( {x - c}\right) ,a\text{ 、 }b\text{ 、 }c \in  \mathbf{R},{f}^{\prime }\left( x\right)$ 为 $f\left( x\right)$ 的导函数.

(1)若 $a = b = c$ ， $f\left( 4\right)  = 8$ ，求 $a$ 的值；

(2)若 $a \neq  b$ ， $b = c$ ，且 $f\left( x\right)$ 和 ${f}^{\prime }\left( x\right)$ 的零点均在集合 $\{  - 3,1,3\}$ 中，求 $f\left( x\right)$ 的极小值.

【例 4】已知函数 $f\left( x\right)  = \frac{{{e}}^{x} - {{e}}^{-x}}{2},g\left( x\right)  = \frac{{{e}}^{x} + {{e}}^{-x}}{2}$ ，其中 ${e}$ 是自然对数的底数.

(1)求证: $g\left( {2x}\right)  = {\left\lbrack  f\left( x\right) \right\rbrack  }^{2} + {\left\lbrack  g\left( x\right) \right\rbrack  }^{2}$ ；

(2)求函数 $h\left( x\right)  = g\left( {2x}\right)  - \frac{7}{2}g\left( x\right)$ 的零点.

【例 5】已知函数 $f\left( x\right)  = {\log }_{3}\left( {{9}^{x} + 1}\right)  - {kx}$ (其中 $k \in  \mathbf{R}$ ).

(1)若 $k = 2$ 且方程 $f\left( x\right)  - a + 1 = 0$ 有解,求实数 $a$ 的取值范围；

(2)若 $f\left( x\right)$ 是偶函数，讨论函数 $g\left( x\right)  = {\left( 3k\right) }^{f\left( x\right) } - \left( {m \cdot  {3}^{x} - \frac{4}{3}m}\right) \left( {m > 0}\right)$ 的零点情况.

\section*{四基达标精练}

\section*{一、填空题}

1. 函数 $f\left( x\right)  = {{e}}^{x} + \ln x$ 的零点个数为\_\_\_\_\_.

2. 已知函数 $f\left( x\right)  = x + {2}^{x}$ 的零点在区间 $\left( {n,n + 1}\right)$ 内, $n \in  \mathbf{Z}$ ,则 $n =$ \_\_\_\_\_.

3. 若函数 $f\left( x\right)  = {2x} + m$ 在 $\left( {1,2}\right)$ 上存在一个零点，则实数 $m$ 的取值范围是\_\_\_\_\_.

4. 已知函数 $h\left( x\right)$ 是奇函数，且 $f\left( x\right)  = h\left( x\right)  + 2$ ，若 $x = 2$ 是函数 $y = f\left( x\right)$ 的一个零点，则 $f\left( {-2}\right)  =$

5. 已知二次函数 $y =  - {x}^{2} + {bx} + c$ 的零点为 -2 和 1，则关于 $x$ 的不等式 ${x}^{2} + {bx} - c > 0$ 的解集为\_\_\_\_\_.

6. 已知函数 $f\left( x\right)  = {x}^{3} - 7{x}^{2} + {ax} - 8$ 有 3 个不同的零点分别为 ${x}_{1}$ 、 ${x}_{2}$ 、 ${x}_{3}$ ，且 ${x}_{1}$ ， ${x}_{2}$ ， ${x}_{3}$ 依次成等比数列，则实数 $a$ 的值为\_\_\_\_\_.

7. 若函数 $f\left( x\right)  = x - a\ln x$ 在定义域上恰有两个零点，则实数 $a$ 的取值范围是\_\_\_\_\_.

8. 已知定义域为 $\mathbf{R}$ 的函数 $f\left( x\right)  = \left\{  \begin{array}{l} {x}^{2} - {2x} + 2,0 < x \leq  2, \\  {\log }_{\frac{1}{2}}\left( {{2x} - \frac{15}{4}}\right) ,x > 2, \end{array}\right.$ 且满足 $f\left( {-x}\right)  =  - f\left( x\right) \left( {x \in  \mathbf{R}}\right)$ ,若方程 $f\left( x\right)  = m\left( {m > 0}\right)$ 的解为 ${x}_{1}\text{ 、 }{x}_{2}\text{ 、 }{x}_{3}\text{ 、 }{x}_{4}$ ,则 ${x}_{1} + {x}_{2} + {x}_{3} + {x}_{4}$ 的取值范围为\_\_\_\_\_.

9. 设 $a > 0$ ，函数 $f\left( x\right)  = \sqrt{\left| x - a\right| } + \sqrt{\left| x + a\right| }$ ，若 $g\left( x\right)  = f\left( x\right)  - b$ 恰有三个不同的零点，且 $b$ 是其中的一个零点，则实数 $b$ 的值为\_\_\_\_\_.

\section*{二、选择题}

10. 已知 $f\left( x\right)  = {x}^{2} - {ax} + a$ ，则 “ $a > 4$ ” 是 “ $f\left( x\right)$ 有两个不同的零点” 的 ( )

A. 充分非必要条件. B. 必要非充分条件. C. 充要条件. D. 既非充分也非必要条件.

11. 下列函数中不能用二分法求零点的是 ( )

A. $y = {3x} - 1$ . B. $y = {x}^{3}$ . C. $y = \left| x\right|$ . D. $y = \ln x$ .

12. 关于函数 $f\left( x\right)  = \frac{1}{4}\left( {x - \frac{1}{x}}\right)  - \ln x$ 的零点,下列说法正确的是 ( )

A. 函数 $f\left( x\right)$ 有两个零点 ${x}_{1}\text{ 、 }{x}_{2}$ ,且 ${x}_{1}{x}_{2} = 1$ .

B. 函数 $f\left( x\right)$ 有两个零点 ${x}_{1}\text{ 、 }{x}_{2}$ ,且 ${x}_{1}{x}_{2} \neq  1$ .

C. 函数 $f\left( x\right)$ 有三个零点 ${x}_{1}\text{ 、 }{x}_{2}\text{ 、 }{x}_{3}$ ,且 ${x}_{1}{x}_{2}{x}_{3} = 1$ .

D. 函数 $f\left( x\right)$ 有三个零点 ${x}_{1}\text{ 、 }{x}_{2}\text{ 、 }{x}_{3}$ ,且 ${x}_{1}{x}_{2}{x}_{3} \neq  1$ .

\section*{三、解答题}

13. 已知 $p :$ 存在 $x \in  \left\lbrack  {0,2}\right\rbrack$ ,使得 ${{e}}^{x} - 1 \leq  m;q :$ 函数 $y = {x}^{2} - {2mx} + 1$ 有两个零点.

(1)若 $p$ 或 $q$ 为假命题，求实数 $m$ 的取值范围；

(2)若 $p$ 或 $q$ 为真命题，而 $p$ 且 $q$ 为假命题，求实数 $m$ 的取值范围.

14. 已知奇函数 $f\left( x\right)  = {x}^{3} + b{x}^{2} + {cx} + d$ 的三个零点 ${x}_{1}$ 、 ${x}_{2}$ 、 ${x}_{3}$ 满足 ${x}_{1}{x}_{2} + {x}_{2}{x}_{3} + {x}_{3}{x}_{1} =  - 2$ .

(1)求函数 $f\left( x\right)$ 的解析式;

(2)求出所有零点.



15. 设函数 $f\left( x\right)  = {x}^{3} - {3a}{x}^{2} + {3bx}$ 的图像与直线 ${12x} + y - 1 = 0$ 相切于点 $\left( {1, - {11}}\right)$ .

(1)求a、b的值；

(2)若函数 $g\left( x\right)  = f\left( x\right)  - c$ 有 3 个零点，求实数 $c$ 的取值范围.

16. 已知函数 $f\left( x\right)  = \frac{1 - {x}^{2}}{1 + {x}^{2}}$ .

(1)求函数 $f\left( x\right)$ 的零点；

(2)设函数 $f\left( x\right)$ 的值域为 $A$ ,若至少有两个不同的 $x \in  \left\lbrack  {\frac{\pi }{2},\pi }\right\rbrack$ ，使得 $\sin \left( {{\omega x} + \frac{\pi }{6}}\right)  \notin  A$ ，求正数 $\omega$ 的取值范围.

\section*{$\text{ § }{3.6}$ 函数的周期性与对称性}

\section*{典型例题精讲}

【例 1】我们知道,函数 $y = f\left( x\right)$ 的图像关于坐标原点成中心对称图形的充要条件是函数 $y = f\left( x\right)$ 为奇函数,有同学发现可以将其推广为: 函数 $y = f\left( x\right)$ 的图像关于点 $P\left( {a,b}\right)$ 成中心对称图形的充要条件是函数 $y = f\left( {x + a}\right)  - b$ 为奇函数.

(1)利用上述材料，求函数 $f\left( x\right)  = {x}^{3} - 3{x}^{2} + {6x} - 2$ 图像的对称中心；

(2)利用函数单调性的定义，证明函数 $f\left( x\right)  = {x}^{3} - 3{x}^{2} + {6x} - 2$ 在 $\mathbf{R}$ 上是严格增函数.

【例 2】定义在 $\mathbf{R}$ 上的奇函数 $f\left( x\right)$ 有最小正周期 4,且 $x \in  \left( {0,2}\right)$ 时, $f\left( x\right)  = \frac{{2}^{x}}{{4}^{x} + 1}$ .

(1)判断 $f\left( x\right)$ 在 $\left( {0,2}\right)$ 上的单调性,并求 $f\left( x\right)$ 在 $\left\lbrack  {-2,2}\right\rbrack$ 上的解析式;

(2)当λ为何值时，关于 $x$ 的方程 $f\left( x\right)  =$ λ在 $\left\lbrack  {2,6}\right\rbrack$ 上有实数解?

【例 3】记函数 $f\left( x\right)$ 的定义域为 $D$ . 如果存在实数 $a\text{ 、 }b$ 使得 $f\left( {a - x}\right)  + f\left( {a + x}\right)  = b$ 对任意满足 $a - x \in  D$ 且 $a + x \in  D$ 的 $x$ 恒成立,则称 $f\left( x\right)$ 为 $\Psi$ 函数.

(1)设函数 $f\left( x\right)  = \frac{1}{x} - 1$ ，试判断 $f\left( x\right)$ 是否为 $\Psi$ 函数，并说明理由；

(2)若 $h\left( x\right)$ 是定义在 $\mathbf{R}$ 上的 $\Psi$ 函数，且函数 $h\left( x\right)$ 的图像关于直线 $x = m$ ( $m$ 为常数，且 $m \neq  a$ )对称，试判断 $h\left( x\right)$ 是否为周期函数? 并证明你的结论.

【例 4】对于定义域为 $\mathbf{R}$ 的函数 $g\left( x\right)$ ,若存在正常数 $T$ ,使得 $\cos \left\lbrack  {g\left( x\right) }\right\rbrack$ 是以 $T$ 为周期的函数,则称 $g\left( x\right)$ 为余弦周期函数,且称 $T$ 为其余弦周期. 已知 $f\left( x\right)$ 是以 $T$ 为余弦周期的余弦周期函数.

(1)证明: $h\left( x\right)  = x + \sin \frac{x}{3}$ 是以 ${6\pi }$ 为周期的余弦周期函数;

(2)证明: “ ${u}_{0}$ 为方程 $\cos \left\lbrack  {f\left( x\right) }\right\rbrack   = 1$ 在 $\left\lbrack  {0,T}\right\rbrack$ 上的解”的充要条件是 “ ${u}_{0} + T$ 为方程 $\cos \left\lbrack  {f\left( x\right) }\right\rbrack   = 1$ 在 $\left\lbrack  {T,{2T}}\right\rbrack$ 上有解”.

【例 5】已知函数 $f\left( x\right)$ 的定义域为 $\mathbf{R}$ ,若存在常数 $T \neq  0$ ,使得 $f\left( x\right)  = {Tf}\left( {x + T}\right)$ 对任意的 $x \in  \mathbf{R}$ 成立,则称函数 $f\left( x\right)$ 是 $\Omega$ 函数.

(1)判断函数 $F\left( x\right)  = x,h\left( x\right)  = \sin {\pi x}$ 是否是 $\Omega$ 函数，请说明理由;

(2)若函数 $f\left( x\right)$ 是 $\Omega$ 函数，且 $f\left( x\right)$ 是偶函数，求证:函数 $f\left( x\right)$ 是周期函数；

(3)若函数 $f\left( x\right)  = \sin {kx}$ 是 $\Omega$ 函数. 求实数 $k$ 的取值范围.

\section*{四基达标精练}

\section*{一、填空题}

1. 函数 $f\left( x\right)  = {2}^{\left| x - 1\right| }$ 的对称轴是直线\_\_\_\_\_.

2. 函数 $f\left( x\right)  = \frac{1}{{3x} + 1} + 1$ 的对称中心是\_\_\_\_\_.

3. 函数 $f\left( x\right)  = \left| {x - 3}\right|  - \left| {x + 1}\right|$ 的对称中心是\_\_\_\_\_.

4.若函数 $f\left( x\right)$ 满足 $f\left( {4 - x}\right)  = f\left( x\right)$ 对任意 $x \in  \mathbf{R}$ 恒成立,当 $x \geq  2$ 时, $f\left( x\right)  = {2x} - 1$ ,则 $f\left( 0\right)  =$

5. 若 $f\left( x\right)$ 是定义在 $\mathbf{R}$ 上的奇函数，且 $f\left( {x + 1}\right)  =  - f\left( x\right)$ ，则 $f\left( {10}\right)  =$ \_\_\_\_\_.

6. 若函数 $y = f\left( x\right)$ 是偶函数，且 $f\left( {x + 3}\right)  = f\left( x\right)$ 对任意 $x \in  \mathbf{R}$ 恒成立，则除 $x = 0$ 外，函数必有一条对称轴是\_\_\_\_\_(写出一条即可).

7. 已知函数 $f\left( x\right)  = \ln \frac{ex}{2 - x}$ ,则 $f\left( x\right)  + f\left( {2 - x}\right)  =$ \_\_\_\_\_；进一步有 $\mathop{\sum }\limits_{{k = 1}}^{{{2n} - 1}}f\left( \frac{k}{n}\right)  =$ \_\_\_\_\_.

8. 已知 $y = f\left( x\right)$ 是 $\mathbf{R}$ 上的奇函数，且 $f\left( x\right)  + f\left( {2 - x}\right)  = 0$ ，当 $- 1 < x < 0$ 时， $f\left( x\right)  = \sqrt{-x}$ ，则 $f\left( {2023}\right)  + f\left( \frac{2023}{2}\right)  =$ \_\_\_\_\_.

9. 已知 $f\left( x\right)$ 是定义在 $\mathbf{R}$ 上的函数,且满足 $f\left( {x + 2}\right)  + f\left( {x + 2}\right) f\left( x\right)  + f\left( x\right)  = 1,f\left( 1\right)  = \frac{1}{2}$ ,则 $f\left( 1\right)  +  f\left( 3\right)  + f\left( 5\right)  + \cdots  + f\left( {2023}\right)  =$

\section*{二、选择题}

10. 已知 $a$ 、 $b$ 、 $c$ 是常数，则函数 $f\left( x\right)  = \frac{{2x} + b}{x + a}$ 和 $g\left( x\right)  = \frac{{2x} + c}{x + a}\left( {c \neq  b}\right)$ 有相同的 ( )

A. 对应法则. B. 奇偶性. C. 对称中心. D. 单调区间和单调性.



11. 已知 $f\left( x\right)$ 是定义在 $\mathbf{R}$ 上的偶函数，且以 2 为周期，则 “ $f\left( x\right)$ 在 $\left\lbrack  {0,1}\right\rbrack$ 上是增函数” 是 “ $f\left( x\right)$ 在 $\left\lbrack  {3,4}\right\rbrack$ 上是减函数”的 ( )

A. 充分非必要条件. B. 必要非充分条件. C. 充要条件. D. 既非充分也非必要条件.

12. 定义在 $\mathbf{R}$ 上的函数 $f\left( x\right)$ 满足 $f\left( {x + 3}\right)  + f\left( x\right)  = 0$ ,且函数 $f\left( {x - \frac{3}{2}}\right)$ 为奇函数. 给出以下 3 个命题  \ding{172} 函数 $f\left( x\right)$ 的周期是 6;  \ding{173} 函数 $f\left( x\right)$ 的图像关于点 $\left( {-\frac{3}{2},0}\right)$ 对称;  \ding{174} 函数 $f\left( x\right)$ 的图像关于 $y$ 轴对称,则其中真命题的个数是 ( )

A.3. B.2. C.1. D.0.

\section*{三、解答题}

13. 已知函数 $f\left( x\right)  = \left\{  \begin{array}{l}  - {2x}, - 1 \leq  x < 0, \\  {x}^{2} - 1,0 \leq  x \leq  1. \end{array}\right.$

(1)求函数 $f\left( x\right)$ 的值域；

(2)函数 $f\left( x\right)$ 的图像上是否存在关于坐标原点对称的点？若存在，求出这些点的坐标；若不存在，说明理由.

14. 对于函数 $f\left( x\right)$ ,若在定义域 $D$ 内存在实数 ${x}_{0}$ 满足 $f\left( {2 - {x}_{0}}\right)  =  - f\left( {x}_{0}\right)$ ,则称函数 $y = f\left( x\right)$ 为 “类对称函数”.

(1)判断函数 $g\left( x\right)  = {x}^{2} - {2x} + 1$ 是否为 “类对称函数” ？若是，求出所有满足条件的 ${x}_{0}$ 的值；若不是， 请说明理由;

(2)若函数 $h\left( x\right)  = {3}^{x} + t$ 为定义在 $\lbrack  - 1,3)$ 上的“类对称函数”,求实数 $t$ 的取值范围.

15. 已知真命题: “函数 $y = f\left( x\right)$ 的图像关于点 $P\left( {a\text{ 、 }b}\right)$ 成中心对称图形” 的充要条件为 “函数 $y = f(x +  a) - b$ 是奇函数”.

(1)将函数 $f\left( x\right)  = {x}^{3} - 3{x}^{2}$ 图像向左平移 1 个单位，再向上平移 2 个单位，所得图像对应的函数是 $y =  g\left( x\right)$ ,请利用题设中的真命题求函数 $f\left( x\right)$ 图像对称中心的坐标;

(2)求函数 $h\left( x\right)  = {\log }_{2}\frac{2x}{4 - x}$ 图像对称中心的坐标.

16. 已知函数 $f\left( x\right)$ 定义在 $\mathbf{R}$ 上，且对任意 $x \in  \mathbf{R}$ ，都有 $f\left( {x + 1}\right)  = f\left( {x - 1}\right) ,f\left( {1 - x}\right)  = f\left( {1 + x}\right)$ ，当 $x \in  \left\lbrack  {0,1}\right\rbrack$ 时， $f\left( x\right)  = {3}^{x} - {3}^{-x}$ .

(1)请指出 $f\left( x\right)$ 在区间 $\left\lbrack  {-1,1}\right\rbrack$ 上的奇偶性、单调区间、最大(小)值和零点;

(2)试证明 $f\left( x\right)$ 是周期函数，并求其在区间 $\left\lbrack  {{2k} - 1,{2k}}\right\rbrack  \left( {k \in  \mathbf{Z}}\right)$ 上的解析式.

\section*{第 4 章 幂函数、指数函数和对数函数(必修 1)}

\section*{知识梳理*}

\section*{一、幂、指数与对数}

1. 指数幂的拓展:正整数指数幂、整数指数幂、有理数指数幂、实数指数幂.

2. 幂的运算性质:对任意给定的正数 $a$ 、 $b$ 及实数 $s$ 、 $t$ ，

(1) ${a}^{s}{a}^{t} = {a}^{s + t};\left( 2\right) {\left( {a}^{s}\right) }^{t} = {a}^{st};\left( 3\right) {\left( ab\right) }^{t} = {a}^{t}{b}^{t}$ ;

${\left( 4\right) }^{ * }{\left( \sqrt[n]{a}\right) }^{n} = a\left( {n \in  \mathbf{N},n \geq  2}\right) ;{\left( 5\right) }^{ * }\sqrt[n]{{a}^{n}} = \left\{  {\begin{array}{l} a,n\text{ 为奇数, } \\  \left| a\right| ,n\text{ 为偶数. } \end{array}\left( {n \in  \mathbf{N},n \geq  2,a \in  \mathbf{R}}\right) }\right.$

3. 指数基本不等式:当 $a > 1$ 时，若 $s > 0$ ，则 ${a}^{s} > 1$ ；若 $s < 0$ ，则 $0 < {a}^{s} < 1$ .

(当 $0 < a < 1$ 时，若 $s > 0$ ，则 $0 < {a}^{s} < 1$ ；若 $s < 0$ ，则 ${a}^{s} > 1$ )

4. 对数的定义: 若 ${a}^{x} = N$ ,则 $x = {\log }_{a}N$ ,其中 $a > 0$ ,且 $a \neq  1,N > 0$ .

5. 对数的基本性质: 设 $a > 0$ ,且 $a \neq  1,M\text{ 、 }N$ 是任意给定的正数, $c$ 是任意给定的实数.

(1) ${\log }_{a}1 = 0$ ;

(2) ${\log }_{a}a = 1$ ；

(3) ${\log }_{a}\left( {MN}\right)  = {\log }_{a}M + {\log }_{a}N$ ；

(4) ${\log }_{a}\frac{M}{N} = {\log }_{a}M - {\log }_{a}N$ ；

(5) ${\log }_{a}{N}^{c} = c{\log }_{a}N$ ;

(6) (换底公式) 如果 $b > 0$ ,且 $b \neq  1$ ,那么 ${\log }_{b}N = \frac{{\log }_{a}N}{{\log }_{a}b}$ .

6. 重要结论: ${a}^{{\log }_{a}N} = N\left( {N > 0}\right) ;{\log }_{a}{a}^{N} = N\left( {N \in  \mathbf{R}}\right) ;{\log }_{{a}^{n}}{N}^{m} = \frac{m}{n}{\log }_{a}N$ .

7. 常用对数和自然对数: 以 10 为底的对数叫常用对数,记作 $\lg N = {\log }_{10}N$ ; 以 ${e}$ 为底的对数叫自然对数,记作 $\ln N = {\log }_{{e}}N$ ,其中 ${e} = {2.71828}\cdots$ 是欧拉常数, ${e}$ 是一个无理数.

\section*{二、几个重要的初等函数的图像和性质}

1. 幂函数、指数函数、对数函数(属于基本初等函数，概念和性质任何时候可以直接应用)的图像与性质:

\begin{center}
\begin{tabular}{|c|c|c|c|c|}
\hline
幂函数 $y = {x}^{a}$ & \multicolumn{2}{|c|}{指数函数 $y = {a}^{x}\left( {a > 0,a \neq  1}\right)$} & \multicolumn{2}{|c|}{对数函数 $y = {\log }_{a}x\left( {a > 0,a \neq  1}\right)$} \\
\cline{1-5}
第一象限内图像 & $a > 1$ & $0 < a < 1$ & $a > 1$ & $0 < a < 1$ \\
\cline{1-5}
\includegraphics[width=0.2\textwidth]{images/bo_d7tk7oc91nqc73875fu0_61_185_1598_201_186_0.jpg} & \includegraphics[width=0.2\textwidth]{images/bo_d7tk7oc91nqc73875fu0_61_446_1594_187_184_0.jpg} & \includegraphics[width=0.2\textwidth]{images/bo_d7tk7oc91nqc73875fu0_61_672_1593_191_185_0.jpg} & \includegraphics[width=0.2\textwidth]{images/bo_d7tk7oc91nqc73875fu0_61_891_1596_220_180_0.jpg} & \includegraphics[width=0.2\textwidth]{images/bo_d7tk7oc91nqc73875fu0_61_1141_1594_223_180_0.jpg} \\
\cline{1-5}
定点 $\left( {1,1}\right)$ & 定点 $\left( {0,1}\right)$ & 定点 $\left( {0,1}\right)$ & 定点 $\left( {1,0}\right)$ & 定点 $\left( {1,0}\right)$ \\
\cline{1-5}
\multirow{2}{*}{\makecell{ $a > 0$ 时 $\lbrack 0, + \infty )$ \textbackslash \\ $a < 0$ 时 $\left( {0, + \infty }\right)  \downarrow$ }} & $\mathbf{R} \uparrow$ & R ↓ & $\left( {0, + \infty }\right)  \uparrow$ & $\left( {0, + \infty }\right)  \downarrow$ \\
\cline{2-5}
 & $x > 0 \Rightarrow  y > 1  x < 0 \Rightarrow  0 < y < 1$ & \makecell{ $x > 0 \Rightarrow  0 < y < 1$ \\ $x < 0 \Rightarrow  y > 1$ } & \makecell{ $x > 1 \Rightarrow  y > 0$ \\ $0 < x < 1 \Rightarrow  y < 0$ } & \makecell{ $x > 1 \Rightarrow  y < 0$ \\ $0 < x < 1 \Rightarrow  y > 0$ } \\
\cline{1-5}
\hline
\end{tabular}
\end{center}

\HRule

*注:本章节所属知识模块是“函数与分析”.

\HRule

2. “耐克 (双勾) 函数”:形如的 $f\left( x\right)  = {ax} + \frac{b}{x}\left( {{ab} \neq  0}\right)$ 函数，不属于基本初等函数，所以单调性用于解答题时应予以证明.

“耐克(双勾)函数”在各种情况下的图像(双曲线):

\begin{center}
\includegraphics[width=0.9\textwidth]{images/bo_d7tk7oc91nqc73875fu0_62_300_343_1081_302_0.jpg}
\end{center}
 

渐近线 $y$ 轴和 $y = {ax}\left( {a \neq  0}\right)$

性质:

\begin{center}
\begin{tabular}{|c|c|c|c|}
\hline
函数 & 参数范围 & 单调递增区间 & 单调递减区间 \\
\cline{1-4}
\multirow{4}{*}{$y = {ax} + \frac{b}{x}$} & $a > 0,b > 0$ & $\left( {-\infty , - \frac{\sqrt{ab}}{a}}\right\rbrack$ 和 $\left\lbrack  {\frac{\sqrt{ab}}{a}, + \infty }\right)$ & $\left\lbrack  {-\frac{\sqrt{ab}}{a},0)\text{ 和 }(0,\frac{\sqrt{ab}}{a}}\right\rbrack$ \\
\cline{2-4}
 & $a \geq  0,b < 0$ & $\left( {-\infty ,0}\right)$ 和 $\left( {0, + \infty }\right)$ & 无 \\
\cline{2-4}
 & $a < 0,b < 0$ & $\left\lbrack  {\frac{\sqrt{ab}}{a},0}\right)$ 和 $\left( {0, - \frac{\sqrt{ab}}{a}}\right\rbrack$ & $\left( {-\infty ,\frac{\sqrt{ab}}{a}}\right\rbrack$ 和 $\left\lbrack  {-\frac{\sqrt{ab}}{a}, + \infty }\right)$ \\
\cline{2-4}
 & $a \leq  0,b > 0$ & 无 & $\left( {-\infty ,0}\right)$ 和 $\left( {0, + \infty }\right)$ \\
\cline{1-4}
\hline
\end{tabular}
\end{center}

3. 由具体函数抽象出来的函数性质的特征表达式:

(1)正比例函数: $f\left( {x + y}\right)  = f\left( x\right)  + f\left( y\right)$ 或 $f\left( {x - y}\right)  = f\left( x\right)  - f\left( y\right)$ 或 $f\left( {nx}\right)  = {nf}\left( x\right) \left( {n \in  {\mathbf{N}}^{ * }}\right)$ ；

(2)幂函数: $f\left( {x \cdot  y}\right)  = f\left( x\right)  \cdot  f\left( y\right)$ 或 $f\left( \frac{x}{y}\right)  = \frac{f\left( x\right) }{f\left( y\right) }$ 或 $f\left( {x}^{n}\right)  = {f}^{n}\left( x\right) \left( {n \in  {\mathbf{N}}^{ * }}\right)$ ；

(3)指数函数: $f\left( {x + y}\right)  = f\left( x\right)  \cdot  f\left( y\right)$ 或 $f\left( {x - y}\right)  = \frac{f\left( x\right) }{f\left( y\right) }$ 或 $f\left( {nx}\right)  = {f}^{n}\left( x\right) \left( {n \in  {\mathbf{N}}^{ * }}\right)$ ；

(4)对数函数: $f\left( {x \cdot  y}\right)  = f\left( x\right)  + f\left( y\right)$ 或 $f\left( \frac{x}{y}\right)  = f\left( x\right)  - f\left( y\right)$ 或 $f\left( {x}^{n}\right)  = {nf}\left( x\right) \left( {n \in  {\mathbf{N}}^{ * }}\right)$ .

\section*{§ 4.1 幂函数、二次函数与“对勾”函数}

\section*{典型例题精讲}

【例 1】已知幂函数 $f\left( x\right)  = {x}^{\frac{1}{{m}^{2} + m}}\left( {m \in  {\mathbf{N}}^{ * }}\right)$ .

(1)试确定该函数的定义域，并指明该函数在其定义域上的单调性；

(2)若函数 $f\left( x\right)$ 的图像经过点 $\left( {2,\sqrt{2}}\right)$ ，试确定 $m$ 的值，并求满足条件 $f\left( {2 - a}\right)  > f\left( {a - 1}\right)$ 的实数 $a$ 的取值范围.

【例 2】已知幂函数 $f\left( x\right)  = {\left( m - 1\right) }^{2}{x}^{m{}^{2} - {4m} + 3}$ 在 $\left( {0, + \infty }\right)$ 上是严格增函数.

(1)求实数 $m$ 的值；

(2)若函数 $h\left( x\right)  =  - \sqrt[3]{{\left\lbrack  f\left( x\right) \right\rbrack  }^{2}} + {2bx} + 1 - b$ 在 $\left\lbrack  {0,2}\right\rbrack$ 上的最大值为 3，求实数 $b$ 的值.

【例 3】因函数 $y = x + \frac{t}{x}\left( {t > 0}\right)$ 的图像形状像对勾,我们称形如 “ $y = x + \frac{t}{x}\left( {t > 0}\right)$ ” 的函数为 “对勾函数”.

(1)已知 $f\left( x\right)  = {2x} + \frac{4}{{2x} - 1} - 5,x \in  \left\lbrack  {1,3}\right\rbrack$ ，利用对勾函数性质，求函数 $f\left( x\right)$ 的单调区间和值域；

(2)对于(1)中的函数 $f\left( x\right)$ 和函数 $g\left( x\right)  = {x}^{2} - {mx} + 4$ ，若对任意 ${x}_{1} \in  \left\lbrack  {1,3}\right\rbrack$ ，总存在 ${x}_{2} \in  \left\lbrack  {1,3}\right\rbrack$ ，使得 $g\left( {x}_{2}\right)  < f\left( {x}_{1}\right)$ 成立，求实数 $m$ 的取值范围.

【例 4】已知函数 $f\left( x\right)  = m\left( {x + \frac{1}{x}}\right)  + n$ 的图像与 $h\left( x\right)  = \frac{1}{4}\left( {x + \frac{1}{x}}\right)$ 的图像关于 $\left( {0, - 1}\right)$ 对称.

(1)求实数 $m$ 、 $n$ 的值；

(2)若 $g\left( x\right)  = f\left( x\right)  + \frac{a}{4x}$ ，求 $g\left( x\right)$ 在区间 $\left\lbrack  {1,2}\right\rbrack$ 上的最小值.

【例 5】设函数 $y = f\left( x\right)$ 的定义域为 $D$ ,若存在 $\left\lbrack  {a,b}\right\rbrack   \subseteq  D$ ,使得 $y = f\left( x\right)$ 在 $\left\lbrack  {a,b}\right\rbrack$ 上的值域也为 $\left\lbrack  {a,b}\right\rbrack$ ,则称函数 $y = f\left( x\right)$ 具有性质 $P$ ,已知幂函数 $f\left( x\right)  = \left( {{p}^{2} + p - 1}\right) {x}^{p - \frac{1}{2}}$ 在 $\left( {0, + \infty }\right)$ 上是严格增函数.

(1) $g\left( x\right)  = f\left( x\right)  - \frac{2}{9}$ 是否具有性质 $P$ ? 若是，请指出区间 $\left\lbrack  {a,b}\right\rbrack$ ；若不是，请说明理由；

(2)若函数 $h\left( x\right)  = n - f\left( {x + 1}\right)$ ，且 $h\left( x\right)$ 具有性质 $P$ ，求实数 $n$ 的取值范围.

\section*{四基达标精练}

\section*{一、填空题}

1. 已知幂函数 $y = f\left( x\right)$ 的图像过点 $P\left( {2,8}\right)$ ，则函数 $y = f\left( x\right)  - x$ 的零点为\_\_\_\_\_.

2. 若幂函数 $y = \left( {{m}^{2} - m - 1}\right) {x}^{{m}^{2} - {2m} - 1}$ 在 $\left( {0, + \infty }\right)$ 上是增函数，则实数 $m =$ \_\_\_\_\_.

3. 若 ${\left( \frac{1}{2}\right) }^{-\frac{1}{2}} > {a}^{-\frac{1}{2}}$ ，则实数 $a$ 的取值范围是\_\_\_\_\_.

4. 一元二次方程 ${x}^{2} + \left( {{a}^{2} - 1}\right) x + \left( {a - 2}\right)  = 0$ 的一根比 1 大，另一根比 -1 小，则实数 $a$ 的取值范围是\_\_\_\_\_.

5. 设定点 $A\left( {0,1}\right)$ ，若动点 $P$ 在函数 $y = \frac{x + 2}{x}\left( {x > 0}\right)$ 图像上，则 $\left| {PA}\right|$ 的最小值为\_\_\_\_\_.

6. 已知函数 $y = {x}^{2} + x + \frac{1}{2}\left( {n \leq  x \leq  n + 1}\right)$ ，其中 $n$ 为正整数的值域中共含有 ${a}_{n}$ 个整数，则 ${a}_{n} =$ \_\_\_\_\_.

7. 已知四个函数: \ding{172}  $y =  - x$ ;  \ding{173}  $y =  - \frac{1}{x}$ ;  \ding{174}  $y = {x}^{3}$ ； \ding{175}  $y = {x}^{\frac{1}{2}}$ . 从中任选 2 个，则事件“所选 2 个函数的图像有且仅有一个公共点”的概率为\_\_\_\_\_.

8. 已知函数 $f\left( x\right)  = \left| {x + \frac{9}{x} - a}\right|  + a$ 在区间 $\left\lbrack  {1,9}\right\rbrack$ 上的最大值是 10,则实数 $a$ 的取值范围是\_\_\_\_\_.

9. 设函数 $F\left( x\right)$ 和 $f\left( x\right)$ 都在区间 $D$ 上有定义,若对 $D$ 的任意子区间 $\left\lbrack  {u,v}\right\rbrack$ ,总有 $\left\lbrack  {u,v}\right\rbrack$ 上的实数 $p\text{ 、 }q$ , 使得不等式 $f\left( p\right)  \leq  \frac{F\left( u\right)  - F\left( v\right) }{u - v} \leq  f\left( q\right)$ 成立,则称 $F\left( x\right)$ 是 $f\left( x\right)$ 在区间 $D$ 上的甲函数, $f\left( x\right)$ 是 $F\left( x\right)$ 在区间 $D$ 上的乙函数. 已知 $F\left( x\right)  = {x}^{2} - {3x}\left( {x \in  \mathbf{R}}\right)$ ，那么 $F\left( x\right)$ 的乙函数 $f\left( x\right)  =$ \_\_\_\_\_.

\section*{二、选择题}

10. 设 $f\left( x\right)  = {x}^{2} - x + a$ ,若 $f\left( {-m}\right)  < 0$ ,则 $f\left( {m + 1}\right)$ 的值是 ( )

A. 正数. B. 负数. C. 非负数. D. 正负与 $m$ 和 $a$ 有关.

11. 已知幂函数的图像过 $\left( {\frac{1}{2},\frac{\sqrt{2}}{4}}\right) ,P\left( {{x}_{1},{y}_{1}}\right) \text{ 、 }Q\left( {{x}_{2},{y}_{2}}\right) \left( {{x}_{1} < {x}_{2}}\right)$ 是函数图像上的任意不同两点,则下列结论中正确的是 ( )

A. ${x}_{1}f\left( {x}_{1}\right)  > {x}_{2}f\left( {x}_{2}\right)$ . B. ${x}_{1}f\left( {x}_{2}\right)  < {x}_{2}f\left( {x}_{1}\right)$ . C. $\frac{f\left( {x}_{1}\right) }{{x}_{2}} > \frac{f\left( {x}_{2}\right) }{{x}_{1}}$ . D. $\frac{f\left( {x}_{1}\right) }{{x}_{1}} < \frac{f\left( {x}_{2}\right) }{{x}_{2}}$ .

12. 函数 $f\left( x\right)  = a{x}^{2} + {bx} + c\left( {a \neq  0}\right)$ 的图像关于直线 $x =  - \frac{b}{2a}$ 对称. 据此可推测,对任意的非零实数 $a$ 、 $b\text{ 、 }c\text{ 、 }m\text{ 、 }n\text{ 、 }p$ ,关于 $x$ 的方程 $m{\left\lbrack  f\left( x\right) \right\rbrack  }^{2} + {nf}\left( x\right)  + p = 0$ 的解集都不可能是 ( )

A. $\{ 1,2\}$ . B. $\{ 1,3,5\}$ . C. $\{ 1,4,7,{10}\}$ . D. $\{ 1,4,{16},{64}\}$ .

\section*{三、解答题}

13. 已知幂函数 $f\left( x\right)  = {x}^{9 - {3m}}\left( {m \in  {\mathbf{N}}^{ * }}\right)$ 的图像关于原点对称,且在 $\mathbf{R}$ 上严格递增.

(1)求 $f\left( x\right)$ 表达式；

(2)求满足 $f\left( {a + 1}\right)  + f\left( {{3a} - 4}\right)  < 0$ 实数 a 的取值范围.

14. 已知幂函数 $f\left( x\right)  = \left( {{m}^{2} - {5m} + 7}\right) {x}^{m - 1}$ 为偶函数.

(1)求 $f\left( x\right)$ 的解析式;

(2)若 $g\left( x\right)  = f\left( x\right)  - {ax} - 3$ 在 $\left\lbrack  {1,3}\right\rbrack$ 上不是单调函数，求实数 $a$ 的取值范围.

15. 已知二次函数 $f\left( x\right)  = a{x}^{2} + {bx}$ 对任意 $x \in  \mathbf{R}$ 均有 $f\left( {x - 1}\right)  = f\left( {-1 - x}\right)$ 成立,且函数的图像过点 $A\left( {1,\frac{3}{2}}\right)$ .

(1)求函数 $y = f\left( x\right)$ 的解析式;

(2)若不等式 $f\left( {x - t}\right)  \leq  x$ 的解集为 $\left\lbrack  {4,m}\right\rbrack$ ，求实数 $t$ 、 $m$ 的值.

16. 某通信公司需要在三角形地带 ${OAC}$ 区域内建造甲、乙两种通信信号加强中转站,甲中转站建在区域 ${BOC}$ 内,乙中转站建在区域 ${AOB}$ 内. 分界线 ${OB}$ 固定,且 ${OB} = \left( {1 + \sqrt{3}}\right)$ 百米,边界线 ${AC}$ 始终过点 $B$ ,边界线 ${OA}\text{ 、 }{OC}$ 满足 $\angle {AOC} = {75}^{ \circ  },\angle {AOB} = {30}^{ \circ  }$ ,设 ${OA} = x\left( {3 \leq  x \leq  6}\right)$ 百米, ${OC} = y$ 百米.

(1)试将 $y$ 表示成 $x$ 的函数，并求出函数 $y$ 的解析式；

(2)当 $x$ 取何值时， $\bigtriangleup {OAC}$ 的面积 ${S}_{\bigtriangleup {OAC}}$ 最小？求出其最小值.

\begin{center}
\includegraphics[width=0.2\textwidth]{images/bo_d7tk7oc91nqc73875fu0_67_1167_1110_268_278_0.jpg}
\end{center}
 

(第 16 题图)

\section*{$\text{ § }{4.2}$ 指数运算与指数函数}

\section*{典型例题精讲}

【例 1】求下列不等式的解集:

(1) ${3}^{x - 1} > {9}^{x}$ ；

(2) ${a}^{-{5x}} > {a}^{x + 7}\left( {a > 0,a \neq  1}\right)$ .

【例 2】已知函数 $g\left( x\right)  = {\left( a + 1\right) }^{x - 2} + 1\left( {a > 0}\right)$ 的图像恒过定点 $A$ ,且点 $A$ 又在函数 $f\left( x\right)  = {\log }_{\sqrt{3}}\left( {x + a}\right)$ 的图像上.

(1)求实数 $a$ 的值；

(2)当 $\left| {g\left( {x + 2}\right)  - 2}\right|  = {2b}$ 有两个不等实根时,求实数 $b$ 的取值范围.

【例 3】已知奇函数 $f\left( x\right)$ 与偶函数 $g\left( x\right)$ 的定义域都是 $\mathbf{R}$ ,且满足 $f\left( x\right)  + g\left( x\right)  = {2}^{x}$ .

(1)求 $f\left( x\right)$ 的解析式；

(2)设函数 $h\left( x\right)  = f\left( x\right)  + x$ .

 \ding{172} 判断 $h\left( x\right)$ 的单调性，并用定义法证明；

 \ding{173}  若 $f\left( {{\log }_{2}m}\right)  + f\left( {2{\log }_{2}m - 1}\right)  \leq  1 - 3{\log }_{2}m$ ，求实数 $m$ 的取值范围.

【例 4】已知函数 $f\left( x\right)  = {a}^{x} + b\left( {a > 0,a \neq  1}\right)$ .

(1)若 $f\left( x\right)$ 的图像如图(1)所示，求 $a$ 、 $b$ 的值；

(2)若 $f\left( x\right)$ 的图像如图(2)所示，求 $a$ 、 $b$ 的取值范围；

(3)在(1)中，若 $\left| {f\left( x\right) }\right|  = m$ 有且仅有一个实数解，求出 $m$ 的范围.

\begin{center}
\includegraphics[width=0.3\textwidth]{images/bo_d7tk7oc91nqc73875fu0_69_1003_1211_447_239_0.jpg}
\end{center}
 

(例 4 图)

【例 5】已知函数 $f\left( x\right)  = {4}^{x} - a \cdot  {2}^{x + 1} - 2$ .

(1)若 $f\left( x\right)$ 的最小值为 -3，求实数 $a$ 的值；

(2)若 $f\left( x\right)  < 0$ 对任意的 $x \in  \lbrack 0,1)$ ，恒成立，求实数 $a$ 的取值范围.

\section*{四基达标精练}

\section*{一、填空题}

1．化简 ${\left( {2}^{n + 1}\right) }^{2} \cdot  {2}^{-{2n} - 1} \div  {4}^{n}$ 的结果是\_\_\_\_\_.

2. 化简 $\frac{\sqrt{-{x}^{3}}}{x}$ 的结果是\_\_\_\_\_.

3. 用分数指数幂形式表示: $\sqrt{a\sqrt{a\sqrt{a\sqrt{a}}}} =$ \_\_\_\_\_.

4. 若指数函数 $y = {\left( 2a - 1\right) }^{x}$ 在 $\mathbf{R}$ 上严格递减，则实数 $a$ 的取值范围是\_\_\_\_\_.

5. 已知函数 $f\left( x\right)  = {a}^{x - 3} + 2$ 的图像恒过定点 $A$ ，则 $A$ 的坐标为\_\_\_\_\_.

6. 若 ${3}^{2x} + 9 = {10} \cdot  {3}^{x}$ ，则 ${x}^{2} + 1 =$ \_\_\_\_\_.

7. 函数 $y = {\left( \frac{1}{2}\right) }^{{x}^{2} + {2x} - 1}$ 的值域是\_\_\_\_\_.

8. 方程 $\left( {{4}^{x} + {4}^{-x}}\right)  - 2\left( {{2}^{x} + {2}^{-x}}\right)  + 2 = 0$ 的解集是\_\_\_\_\_.

9. 已知 $f\left( x\right)  + g\left( x\right)  = {a}^{x} - {a}^{-x} + 2\left( {a > 0,a \neq  1}\right)$ ,其中 $f\left( x\right)$ 是定义在 $\mathbf{R}$ 上的奇函数, $g\left( x\right)$ 是定义在 $\mathbf{R}$ 上的偶函数. 若 $g\left( a\right)  = a$ ,则 $f\left( a\right)  =$ \_\_\_\_\_.

\section*{二、选择题}

10. 已知函数 $g\left( x\right)  = {3}^{x} + t$ 的图像不经过第二象限,则 $t$ 的取值范围是 ( )

A. $t \leq   - 1$ . B. $t <  - 1$ . C. $t \leq   - 3$ . D. $t \geq   - 3$ .

11. 若 $0 < b < a < 1$ ,则在 ${a}^{b}$ 、 ${b}^{a}$ 、 ${a}^{a}$ 、 ${b}^{b}$ 中的最大值是 ( )

A. ${b}^{a}$ . B. ${a}^{a}$ . C. ${a}^{b}$ . D. ${b}^{b}$ .

12. 已知 $n \in  {\mathbf{N}}^{ * }$ ,则函数 $y = {x}^{n}\left( {x \in  \mathbf{R}}\right)$ 与 $y = {n}^{x}\left( {x \in  \mathbf{R}}\right)$ 图像的交点不可能 ( )

A. 只有 $\left( {n,{n}^{n}}\right)$ . B. 在直线 $y = {nx}$ 上. C. 多于三个. D. 在第二象限.

\section*{三、解答题}

13. 已知函数 $f\left( x\right)  = {a}^{x - 1}\left( {x \geq  0}\right)$ 的图像经过点 $\left( {2,\frac{1}{2}}\right)$ ，其中 $a > 0$ 且 $a \neq  1$ .

(1)求 $a$ 的值；

(2)求函数 $y = f\left( x\right)$ 的值域.

14. 解下列方程:

(1) ${3}^{{2x} + 5} = 5 \cdot  {3}^{x + 2} + 2$ ； (2) ${\left( \sqrt{3 + 2\sqrt{2}}\right) }^{x} + {\left( \sqrt{3 - 2\sqrt{2}}\right) }^{x} = 6$ .

15. 已知 $f\left( x\right)  = {{e}}^{x} - {{e}}^{-x},g\left( x\right)  = {{e}}^{x} + {{e}}^{-x}$ ,其中 ${e} = {2.71828}\cdots$ .

(1)求 ${\left\lbrack  f\left( x\right) \right\rbrack  }^{2} - {\left\lbrack  g\left( x\right) \right\rbrack  }^{2}$ 的值；

(2)设 $f\left( x\right)$ 、 $f\left( y\right)  = 4,g\left( x\right)  \cdot  g\left( y\right)  = 8$ ，求 $\frac{g\left( {x + y}\right) }{g\left( {x - y}\right) }$ 的值.

16. 已知 $f\left( x\right)  = \frac{{4}^{x} + 1}{{4}^{x} + a}$ 是奇函数,其中 $a$ 是常数.

(1)求 $f\left( x\right)$ 的定义域和值域；

(2)讨论 $f\left( x\right)$ 的单调性并证明.

\section*{§ 4.3 对数运算与对数函数}

\section*{典型例题精讲}

【例 1】已知函数 $f\left( x\right)  = {\log }_{a}x + b$ (其中 $a\text{ 、 }b$ 均为常数, $a > 0$ 且 $a \neq  1$ ) 的图像经过点 $\left( {2,5}\right)$ 和点 $\left( {8,7}\right)$ .

(1)求 $a$ 、 $b$ 的值；

(2)设函数 $g\left( x\right)  = {b}^{x} - {a}^{x + 2}$ ，若对任意的 ${x}_{1} \in  \left\lbrack  {1,4}\right\rbrack$ ，总存在 ${x}_{2} \in  \left\lbrack  {0,{\log }_{2}5}\right\rbrack$ ，使得 $f\left( {x}_{1}\right)  = g\left( {x}_{2}\right)  + m$ 成立,求实数 $m$ 的取值范围.

【例 2】如图,已知函数 $y = {kx}\left( {k > 0}\right)$ 的图像与函数 $y = {\log }_{2}x$ 的图像交于 $A\text{ 、 }B$ 两点.过 $A\text{ 、 }B$ 分别作 $x$ 轴的垂线,垂足分别为点 ${A}^{\prime }$ 、点 ${B}^{\prime }$ ，并且 $A{A}^{\prime }$ 、 $B{B}^{\prime }$ 分别交函数 $y = {\log }_{4}x$ 的图像于 $M$ 、 $N$ 两点.

(1)探究线段 ${AM}$ 与 $M{A}^{\prime }$ 长度的大小关系；

(2)若 ${AN}$ 平行于 $x$ 轴,求四边形 ${AMNB}$ 的面积.

\begin{center}
\includegraphics[width=0.2\textwidth]{images/bo_d7tk7oc91nqc73875fu0_73_1149_1317_288_216_0.jpg}
\end{center}
 

(例 2 图)

【例 3】( 1 )若方程 ${\log }_{2}\left( {x + 3}\right)  - {\log }_{4}{x}^{2} = a$ 的根在( 3,4 )内，求实数 $a$ 的取值范围；

(2)已知 $a > 0,a \neq  1$ ，试求使方程 ${\log }_{a}\left( {x - {ak}}\right)  = {\log }_{{a}^{2}}\left( {{x}^{2} - {a}^{2}}\right)$ 有解的实数 $k$ 的取值范围.

【例 4】已知函数 $f\left( x\right)  = {\log }_{\frac{1}{2}}\frac{1 - {ax}}{x - 1}$ 的图像关于原点对称,其中 $a$ 为常数.

(1)求 $a$ 的值；

(2)当 $x \in  \left( {1, + \infty }\right)$ 时， $f\left( x\right)  + {\log }_{\frac{1}{2}}\left( {x - 1}\right)  < m$ 恒成立，求实数 $m$ 的取值范围；

(3)若关于 $x$ 的方程 $f\left( x\right)  = {\log }_{\frac{1}{2}}\left( {x + k}\right)$ 在 $\left\lbrack  {2,3}\right\rbrack$ 上有解,求 $k$ 的取值范围.

【例 5】已知函数 $f\left( x\right)  = {\log }_{a}x\left( {a > 0\text{ 且 }a \neq  1}\right)$ .

(1)若 $f\left( {a + 4}\right)  \leq  f\left( {3a}\right)$ ，求实数 $a$ 的取值范围；

(2)设 $a = 2$ ，函数 $g\left( x\right)  =  - {f}^{2}\left( x\right)  + \left( {3 - {2m}}\right) f\left( x\right)  + m + 2\left( {0 < m \leq  1}\right)$ . 若 $x \in  \left\lbrack  {1,{2}^{m}}\right\rbrack$ ，证明: $g\left( x\right)  \leq  \frac{10}{3}$ .

\section*{四基达标精练}

\section*{一、填空题}

1. 若 $a > 1,b > 1,p = \frac{{\log }_{b}\left( {{\log }_{b}a}\right) }{{\log }_{b}a}$ ，则 ${a}^{p} =$ \_\_\_\_\_(用 $a$ 、 $b$ 表示).

2. 已知 $2\lg \left( {x - {2y}}\right)  = \lg x + \lg y$ ，则 $\frac{x}{y}$ 的值为\_\_\_\_\_.

3. 设 ${2}^{m} = {5}^{n} = {10}$ ，则 $\frac{1}{m} + \frac{1}{n} =$ \_\_\_\_\_.

4. 已知 ${\log }_{6}{27} = a$ ，试用 $a$ 表示 ${\log }_{18}{16} =$ \_\_\_\_\_.

5. 已知对数函数 $y = {\log }_{a}x$ 的图像过点 $P\left( {8,3}\right)$ ，则 $f\left( \frac{1}{32}\right)  =$ \_\_\_\_\_.

6．如果方程 ${\lg }^{2}x + \left( {\lg 7 + \lg 5}\right) \lg x + \lg 7 \cdot  \lg 5 = 0$ 的两根为 $\alpha$ 、 $\beta$ ，则 $\alpha  \cdot  \beta  =$ \_\_\_\_\_.

7. 方程 ${x}^{\lg x} = {10}$ 的所有实数根之积是\_\_\_\_\_.

8. 已知 ${\log }_{2}x = {\log }_{3}y = {\log }_{5}z =  - 2$ ，则 ${x}^{\frac{1}{2}}$ 、 ${y}^{\frac{1}{3}}$ 、 ${z}^{\frac{1}{5}}$ 由小到大的排列顺序是\_\_\_\_\_.

9. 已知函数 $f\left( x\right)$ 满足:对于任意 ${x}_{1}$ 、 ${x}_{2} \in  \mathbf{R}$ ，都有 $f\left( {{x}_{1} + {x}_{2}}\right)  = f\left( {x}_{1}\right)  + f\left( {x}_{2}\right)  - {2016}$ ，且当 $x > 0$ 时， $f\left( x\right)  > {2016}$ ，若 $f\left( x\right)$ 在区间 $\left\lbrack  {-{2023},{2023}}\right\rbrack$ 上的最大值、最小值分别为 $M$ 、 $N$ ，则 $M + N =$ \_\_\_\_\_.

\section*{二、选择题}

10. 设 $a$ 、 $b$ 、 $c$ 都是正数，且 ${3}^{a} = {4}^{b} = {6}^{c}$ ，则 ( )

A. $\frac{1}{c} = \frac{1}{a} + \frac{1}{b}$ . B. $\frac{2}{c} = \frac{2}{a} + \frac{1}{b}$ . C. $\frac{1}{c} = \frac{2}{a} + \frac{2}{b}$ . D. $\frac{2}{c} = \frac{1}{a} + \frac{2}{b}$ .

11. 已知函数 $f\left( x\right)  = \left| {{\log }_{2}x}\right|$ ，正实数 $m$ 、 $n\left( {m < n}\right)$ 满足 $f\left( m\right)  = f\left( n\right)$ ，若 $f\left( x\right)$ 在区间 $\left\lbrack  {{m}^{2},n}\right\rbrack$ 上的最大值为 2,则 $m\text{ 、 }n$ 的值分别为 ( )

A. $\frac{\sqrt{2}}{2},\sqrt{2}$ . B. $\frac{1}{4}\text{ 、 }2$ . C. $\frac{1}{2}\text{ 、 }2$ . D. $\frac{1}{4}$ 、4.

12. 已知函数 $f\left( x\right)$ 满足 $f\left( {x + 2}\right)  = f\left( x\right)$ 对任意 $x \in  \mathbf{R}$ 恒成立，且当 $x \in  \left\lbrack  {1,3}\right\rbrack$ 时， $f\left( x\right)  = \left| {x - 2}\right|  - 1$ ，若函数 $y  = f\left( x\right)  - {\log }_{a}\left( {x + 1}\right)$ 至少有三个零点，则实数 $a$ 的取值范围为 ( )

A. $\left( {0,\frac{\sqrt{3}}{3}}\right)$ . B. $\left( {0,\frac{1}{3}}\right)$ . C. $\left( {0,\frac{1}{2}}\right)$ . D. $\left( {\frac{1}{3},1}\right)$ .

\section*{三、解答题}

13. 已知函数 $f\left( x\right)  = \lg \left( {x - 1}\right) ,g\left( x\right)  = \lg \left( {4 - x}\right)$ .

(1)求函数 $h\left( x\right)  = f\left( x\right)  - g\left( x\right)$ 的定义域；

(2)当不等式 $f\left( x\right)  > g\left( x\right)$ 成立时,求实数 $x$ 的取值范围.

14. 解关于 $x$ 的方程.

(1) ${\log }_{4}\left( {3 - x}\right)  + {\log }_{0.25}\left( {3 + x}\right)  = {\log }_{4}\left( {1 - x}\right)  + {\log }_{0.25}\left( {{2x} + 1}\right)$ ;

(2) ${\log }_{\left( x + a\right) }{2x} = 2$ .

15. 函数 $f\left( x\right)  = b + {\log }_{a}x\left( {a > 0\text{ 且 }a \neq  1}\right)$ 的图像经过点 $\left( {8,2}\right)$ 和 $\left( {1, - 1}\right)$ .

(1)求函数 $f\left( x\right)$ 的解析式；

(2)函数 $g\left( x\right)  = {f}^{2}\left( x\right)  - f\left( x\right)$ ，求函数 $g\left( x\right)$ 的最小值.

16. 已知函数 $f\left( x\right)  = {a}^{x - 1}\left( {a > 0\text{ 且 }a \neq  1}\right)$ .

(1)若函数 $y = f\left( x\right)$ 的图像经过 $P\left( {3,4}\right)$ 点，求 $a$ 的值；

(2)若 $f\left( {\lg a}\right)  = {100}$ ，求 $a$ 的值.


\section*{第 5 章 导数及其应用(选择性必修 2)}

\section*{知识梳理*}

\section*{一、平均变化率与导数的概念:}

对于函数 $y = f\left( x\right)$ ,将 $\frac{f\left( {{x}_{0} + h}\right)  - f\left( {x}_{0}\right) }{h}$ 称为函数在以 ${x}_{0}$ 和 ${x}_{0} + h$ 为端点的区间上的平均变化率, 如果当 $h$ 趋近于 0 时, $\frac{f\left( {{x}_{0} + h}\right)  - f\left( {x}_{0}\right) }{h}$ 趋近于一个稳定值,则将这一稳定值称为 $h$ 趋于 0 时, $\frac{f\left( {{x}_{0} + h}\right)  - f\left( {x}_{0}\right) }{h}$ 的极限,记作 $\mathop{\lim }\limits_{{h \rightarrow  0}}\frac{f\left( {{x}_{0} + h}\right)  - f\left( {x}_{0}\right) }{h},\mathop{\lim }\limits_{{h \rightarrow  0}}\frac{f\left( {{x}_{0} + h}\right)  - f\left( {x}_{0}\right) }{h}$ 是函数在 $x = {x}_{0}$ 处的瞬时变化率,也称为函数 $y = f\left( x\right)$ 在 $x = {x}_{0}$ 处的导数,记作 ${f}^{\prime }\left( {x}_{0}\right)$ .

即 ${f}^{\prime }\left( {x}_{0}\right)  = \mathop{\lim }\limits_{{h \rightarrow  0}}\frac{f\left( {{x}_{0} + h}\right)  - f\left( {x}_{0}\right) }{h}$ ; 也就是说,导数值 $=$ 瞬时变化率 $=$ 平均变化率的极限.

\section*{二、平均变化率和导数值的几何意义}

定点 $P\left( {{x}_{0},f\left( {x}_{0}\right) }\right)$ 和动点 $Q\left( {{x}_{0} + h,f\left( {{x}_{0} + h}\right) }\right)$ 在函数 $y = f\left( x\right)$ 上,如果 ${f}^{\prime }\left( {x}_{0}\right)$ 存在,则 $\frac{f\left( {{x}_{0} + h}\right)  - f\left( {x}_{0}\right) }{h} =$ 割线 ${PQ}$ 的斜率,当 $h$ 趋于 0 时,割线 ${PQ}$ 逐渐趋近于 $y = f\left( x\right)$ 在点 $P$ 处的切线,因此 ${f}^{\prime }\left( {x}_{0}\right)$ 是 $y = f\left( x\right)$ 在点 $P$ 处切线的斜率.

\section*{三、导函数}

1. 对于函数 $y = f\left( x\right)$ ，在导数值存在的前提下，对不同的 ${x}_{0}$ ，总有一个确定的导数值 $y = f\left( {x}_{0}\right)$ 对应，换句话说，如果用 $x$ 表示自变量，那么 ${y}^{\prime } = {f}^{\prime }\left( x\right)$ 也是关于 $x$ 的一个函数，称为 $y = f\left( x\right)$ 的导函数.

注意:导数值和导函数的区别可类比函数值和函数的区别.

2. 驻点:将导数为零的点称为函数的驻点.

\section*{四、导数的运算}

1.基本初等函数的导数公式和运算法则

\begin{center}
\begin{tabular}{|c|c|}
\hline
基本初等函数 & 导函数 \\
\cline{1-2}
$f\left( x\right)  = C\left( {C\text{ 为常数 }}\right)$ & ${f}^{\prime }\left( x\right)  = 0$ \\
\cline{1-2}
$f\left( x\right)  = {x}^{a}\left( {a \in  Q,a \neq  0}\right)$ & ${f}^{\prime }\left( x\right)  = a{x}^{a - 1}$ \\
\cline{1-2}
$f\left( x\right)  = \sin x$ & ${f}^{\prime }\left( x\right)  = \cos x$ \\
\cline{1-2}
$f\left( x\right)  = \cos x$ & ${f}^{\prime }\left( x\right)  =  - \sin x$ \\
\cline{1-2}
$f\left( x\right)  = {e}^{x}$ & ${f}^{\prime }\left( x\right)  = {e}^{x}$ \\
\cline{1-2}
$f\left( x\right)  = {a}^{x}\left( {a > 0,a \neq  1}\right)$ & $f\left( x\right)  = {a}^{x}\ln a$ \\
\cline{1-2}
$f\left( x\right)  = \ln x$ & $f\left( x\right)  = \frac{1}{x}$ \\
\cline{1-2}
$f\left( x\right)  = {\log }_{a}x\left( {a > 0,a \neq  1}\right)$ & $f\left( x\right)  = \frac{1}{x\ln a}$ \\
\cline{1-2}
\hline
\end{tabular}
\end{center}

\begin{center}
\begin{tabular}{|c|}
\hline
运算法则 \\
\cline{1-1}
设 ${f}^{\prime }\left( x\right)$ 和 ${g}^{\prime }\left( x\right)$ 都存在,则 \\
\cline{1-1}
\makecell{  \ding{172}  ${\left\lbrack  f\left( x\right)  + g\left( x\right) \right\rbrack  }^{\prime } = {f}^{\prime }\left( x\right)  + {g}^{\prime }\left( x\right)$ \\  \ding{173}  ${\left\lbrack  f\left( x\right)  - g\left( x\right) \right\rbrack  }^{\prime } = {f}^{\prime }\left( x\right)  - {g}^{\prime }\left( x\right)$ \\  \ding{174}  ${\left\lbrack  f\left( x\right)  \cdot  g\left( x\right) \right\rbrack  }^{\prime } = {f}^{\prime }\left( x\right)  \cdot  g\left( x\right)  + f\left( x\right)  \cdot  {g}^{\prime }\left( x\right)$ \\  \ding{175}  ${\left\lbrack  \frac{f\left( x\right) }{g\left( x\right) }\right\rbrack  }^{\prime } = \frac{{f}^{\prime }\left( x\right)  \cdot  g\left( x\right)  - f\left( x\right)  \cdot  {g}^{\prime }\left( x\right) }{{g}^{2}\left( x\right) }$ \\ 其中 $g\left( x\right)  \neq  0$ } \\
\cline{1-1}
\hline
\end{tabular}
\end{center}

\HRule

* 注:本章节所属知识模块是“函数与分析”.

\HRule

2. 简单复合函数的求导法则:

简单复合函数 $y = f\left( {{ax} + b}\right)$ 的导数和函数 $y = f\left( u\right)$ 的关系为 ${\left\lbrack  f\left( ax + b\right) \right\rbrack  }^{\prime } = {f}^{\prime }\left( {{ax} + b}\right)  \cdot  a$ ,其中 ${\left\lbrack  f\left( ax + b\right) \right\rbrack  }^{\prime }$ 是指复合函数 $f\left( {{ax} + b}\right)$ 关于 $x$ 的导数, ${f}^{\prime }\left( {{ax} + b}\right)$ 是指将 ${f}^{\prime }\left( u\right)$ 中的 $u$ 用 ${ax} + b$ 代替.

注: \ding{172}  求函数的导数要准确地把函数拆分成基本初等函数的和、差、积、商，再利用运算法则求导.

 \ding{173} 抽象函数求导，恰当赋值是关键，然后活用方程思想求解.

 \ding{174} 复合函数求导，应由外到内逐层求导，必要时要进行换元.

\section*{五、切线方程}

函数 $y = f\left( x\right)$ 在 $x = {x}_{0}$ 处的导数的几何意义就是曲线 $y = f\left( x\right)$ 在点 $P\left( {{x}_{0},f\left( {x}_{0}\right) }\right)$ 处的切线的斜率, 相应的切线方程为 $y - f\left( {x}_{0}\right)  = {f}^{\prime }\left( {x}_{0}\right) \left( {x - {x}_{0}}\right)$ .

1．求曲线的切线方程要分清“在点处”与“过点处”的切线方程的不同.

求曲线在点 $P\left( {{x}_{0},{y}_{0}}\right)$ 处的切线,则表明点 $P$ 是切点,可直接求出函数在 $P$ 处的导数,然后利用点斜式写出切线方程,若在该点 $P$ 处的导数不存在,则切线垂直于 $x$ 轴,切线方程为 $x = {x}_{0}$ .

求过点 $P\left( {{x}_{0},{y}_{0}}\right)$ 的切线,点 $P$ 不一定是切点,要设出切点坐标,根据斜率相等建立方程 (组) 求解,求出切点坐标是解题的关键.

2. 处理与切线有关的参数问题,通常利用曲线、切线、切点的三个关系列出参数的方程 (组) 并解出参数:

(1)切点处的导数是切线的斜率；

(2)切点在切线上，故满足切线方程；

(3)切点在曲线上，故满足曲线方程。

3. 公切线问题: 两条曲线的公切线, 一般是把两条曲线分开考虑: 分别求两条曲线的切线, 化为斜截式方程，利用方程相等，列式求参数. 若其中一条曲线涉及二次曲线，可先求另一条曲线的切线，与二次曲线联立、 消元, 用判别式法.

4. 利用导数的几何意义求参数问题时, 注意利用数形结合, 化归与转化的思想方法.

六、函数的单调性与导数.

1. 函数的单调性与导数的关系.

\begin{center}
\begin{tabular}{|c|c|c|}
\hline
条件 & 恒有 & 结论 \\
\cline{1-3}
\multirow{3}{*}{函数 $y = f\left( x\right)$ 在区间 $\left( {a,b}\right)$ 上可导} & ${f}^{\prime }\left( x\right)  > 0$ & $f\left( x\right)$ 在 $\left( {a,b}\right)$ 上单调增 \\
\cline{2-3}
 & ${f}^{\prime }\left( x\right)  < 0$ & $f\left( x\right)$ 在 $\left( {a,b}\right)$ 上单调减 \\
\cline{2-3}
 & ${f}^{\prime }\left( x\right)  = 0$ & $f\left( x\right)$ 在 $\left( {a,b}\right)$ 上是常数函数 \\
\cline{1-3}
\hline
\end{tabular}
\end{center}

2. 利用导数判断函数单调性的步骤:

第 1 步，确定函数的定义域；

第 2 步，求出导函数 ${f}^{\prime }\left( x\right)$ 的零点；

第 3 步,用 ${f}^{\prime }\left( x\right)$ 的零点将 $f\left( x\right)$ 的定义域划分为若干个区间,列表给出 ${f}^{\prime }\left( x\right)$ 在各区间上的正负,由此得出函数 $y = f\left( x\right)$ 在定义域内的单调性.

(1)若函数 $f\left( x\right)$ 在区间 $\left( {a,b}\right)$ 上递增，则 ${f}^{\prime }\left( x\right)  \geq  0$ ，所以 “ ${f}^{\prime }\left( x\right)  > 0$ 在 $\left( {a,b}\right)$ 上成立” 是 “ $f\left( x\right)$ 在 $\left( {a,b}\right)$ 上单调增”的充分不必要条件.

(2)研究含参数的函数的单调性，要依据参数对不等式解集的影响进行分类讨论.

若导函数为二次函数式，首先看能否因式分解，再讨论二次项系数的正负及两根的大小；若不能因式分解,则需讨论判别式 $\Delta$ 的正负,二次项系数的正负,两根的大小及根是否在定义域内.

(3)个别导数为 0 的点不影响所在区间的单调性，如 $f\left( x\right)  = {x}^{3}$ ， ${f}^{\prime }\left( x\right)  = 3{x}^{2} \geq  0$ ( ${f}^{\prime }\left( x\right)  = 0$ 在 $x = 0$ 时取到), $f\left( x\right)$ 在 $\mathbf{R}$ 上是严格增函数.

3. 根据函数单调性求参数的一般思路:

(1)利用集合间的包含关系处理:

$y = f\left( x\right)$ 在 $\left( {a,b}\right)$ 上单调,则区间 $\left( {a,b}\right)$ 是相应单调区间的子集.

(2) $f\left( x\right)$ 为增 (减) 函数的充要条件是对任意的 $x \in  \left( {a,b}\right)$ 都有 ${f}^{\prime }\left( x\right)  \geq  0\left( {{f}^{\prime }\left( x\right)  \leq  0}\right)$ ，且在 $\left( {a,b}\right)$ 内的任一非空子区间上， ${f}^{\prime }\left( x\right)$ 不恒为零，应注意此时式子中的等号不能省略，否则会漏解.

(3)函数在某个区间上存在单调区间可转化为不等式有解问题.

\section*{七、函数的极值与最值}

1. 函数的极值

(1)对于函数 $y = f\left( x\right)$ ，如果在 $x = {x}_{1}$ 附近存在一个小区间，使得在该区间内， $f\left( {x}_{1}\right)$ 是最大值，那么， 就称 $y = f\left( x\right)$ 有极大值 $f\left( {x}_{1}\right) ,x = {x}_{1}$ 是极大值点. 从图像上说，极大值点附近是 “山峰”，极大值实际上是 “局部的最大值”，但不一定是整体的最大值.

类似地,如果在 $x = {x}_{2}$ 附近存在一个小区间,使得在该区间内, $f\left( {x}_{2}\right)$ 是最小值,那么,就称 $y = f\left( x\right)$ 有极小值 $f\left( {x}_{2}\right) ,x = {x}_{2}$ 是极小值点. 从图像上说,极小值点附近是 “山谷”,极小值实际上是 “局部的最小值”, 但不一定是整体的最小值.

(2)极大值和极小值统称为极值，极大值点和极小值点统称为极值点.

(3)在导数存在的前提下，“ $x = {x}_{0}$ 是函数的驻点”是 “ $x = {x}_{0}$ 是函数的极值点”的必要非充分条件.

2. 运用导数求函数 $f\left( x\right)$ 极值的一般步骤:

(1)确定函数 $f\left( x\right)$ 的定义域;

(2)求导数 ${f}^{\prime }\left( x\right)$ ；

(3)解方程 ${f}^{\prime }\left( x\right)  = 0$ ，求出函数定义域内的所有根；

(4)列表检验 ${f}^{\prime }\left( x\right)$ 在 ${f}^{\prime }\left( x\right)  = 0$ 的根 ${x}_{0}$ 左右两侧值的符号;

(5)求出极值.

注: \ding{172} 已知函数极值，确定函数解析式中的参数时，要注意:根据极值点的导数为 0 和极值这两个条件列方程组，利用待定系数法求解.  \ding{173} 导数值为 0 不是此点为极值点的充要条件，所以用待定系数法求解后必须检验.

3. 函数的最大 (小) 值

函数 $f\left( x\right)$ 在区间 $\left\lbrack  {a,b}\right\rbrack$ 上有最值的条件:

如果在区间 $\left\lbrack  {a,b}\right\rbrack$ 上函数 $y = f\left( x\right)$ 的图像是一条连续不断的曲线,那么它必有最大值和最小值.

4. 利用导数求函数 $f\left( x\right)$ 在 $\left\lbrack  {a,b}\right\rbrack$ 上的最值的一般步骤:

(1)求函数在 $\left( {a,b}\right)$ 内的极值.

(2)求函数在区间端点处的函数值 $f\left( a\right) ,f\left( b\right)$ .

(3)将函数 $f\left( x\right)$ 的各极值与 $f\left( a\right) ,f\left( b\right)$ 比较,其中最大的一个为最大值,最小的一个为最小值.

\section*{八、不等式证明的常用思路}

1. 移项构造函数法: 证明不等式 $f\left( x\right)  > g\left( x\right)$ (或 $f\left( x\right)  < g\left( x\right)$ ) 转化为证明 $f\left( x\right)  - g\left( x\right)  > 0$ (或 $f\left( x\right)  - g\left( x\right)  < 0)$ ,进而构造辅助函数 $h\left( x\right)  = f\left( x\right)  - g\left( x\right)$ ;

2. 适当放缩构造法:一是根据已知条件适当放缩；二是利用常见放缩结论；

3. 构造“形似”函数，稍作变形再构造，对原不等式同解变形，根据相似结构构造辅助函数；

4. 双变量不等式的处理策略:含两个变量的不等式，基本的思路是将之转化为一元的不等式，具体转化方法主要有三种:整体代换，分离变量，选取主元.

\section*{九、不等式成立问题常用的转化规则(注意等号的相对性)}

单变量不等式成立问题:一般利用参变分离法求解函数不等式恒 (能) 成立.

1. 对任意 $x \in  D,m \leq  f\left( x\right)  \Leftrightarrow  m \leq  f{\left( x\right) }_{\min }$ ;

2. 对任意 $x \in  D,m \geq  f\left( x\right)  \Leftrightarrow  m \geq  f{\left( x\right) }_{\min }$ ;

3. 存在 $x \in  D,m \leq  f\left( x\right)  \Leftrightarrow  m \leq  f{\left( x\right) }_{\min }$ ;

4. 存在 $x \in  D,m \geq  f\left( x\right)  \Leftrightarrow  m \geq  f{\left( x\right) }_{\min }$ .

\section*{$\text{ § }{5.1}$ 导数的概念}

\section*{典型例题精讲}

【例 1】求下列函数的导数.

(1) $y = {x}^{2}\sin x$ ; (2) $y = \ln x + \frac{1}{x}$ ;

(3) $y = \sin \left( {{2x} + \frac{\pi }{3}}\right)$ ； (4) $y = \ln \left( {{2x} - 5}\right)$ .

【例 2】已知某质点的运动方程为 $s\left( t\right)  = 3{t}^{2} + {2t} + 1$ ( $s$ 的单位为 ${m},t$ 的单位为 ${s}$ ).

(1)求从 $t = {2{s}}$ 到 $t = \left( {2 + {\Delta t}}\right) {s}$ 的平均速度；

(2)求当 $t = {2{s}}$ 时的瞬时速度.

【例 3】设 $f\left( 5\right)  = 5,{f}^{\prime }\left( 5\right)  = 3,g\left( 5\right)  = 4,{g}^{\prime }\left( 5\right)  = 1,h\left( x\right)  = \frac{f\left( x\right)  + 2}{g\left( x\right) }$ .

(1)求 $h\left( 5\right)$ 及 ${h}^{\prime }\left( 5\right)$ ；

(2)求曲线 $y = h\left( x\right)  + \sin \frac{\pi }{6}$ 在 $x = 5$ 处的切线方程.

【例 4】已知曲线 $S : y = {2x} - {x}^{3}$ .

(1)求曲线 $S$ 在点 $A\left( {1,1}\right)$ 处的切线方程；

(2)求过点 $B\left( {2,0}\right)$ 并与曲线 $S$ 相切的直线方程.


【例 5】记 ${f}^{\prime }\left( x\right) \text{ 、 }{g}^{\prime }\left( x\right)$ 分别为函数 $f\left( x\right) \text{ 、 }g\left( x\right)$ 的导函数. 若存在 ${x}_{0} \in  \mathbf{R}$ ,满足 $f\left( {x}_{0}\right)  = g\left( {x}_{0}\right)$ 且 ${f}^{\prime }\left( {x}_{0}\right)  = {g}^{\prime }\left( {x}_{0}\right)$ ,则称 ${x}_{0}$ 为函数 $f\left( x\right)$ 与 $g\left( x\right)$ 的一个“ $S$ 点”.

(1)证明:函数 $f\left( x\right)  = x$ 与 $g\left( x\right)  = {x}^{2} + {2x} - 2$ 不存在 “ $S$ 点”；

(2)若函数 $f\left( x\right)  = a{x}^{2} - 1$ 与 $g\left( x\right)  = \ln x$ 存在 “ $S$ 点”，求实数 $a$ 的值.

\section*{四基达标精练}

\section*{一、填空题}

1. 已知 $f\left( x\right)  = {2}^{x}$ ,则 ${f}^{\prime }\left( 1\right)  =$ \_\_\_\_\_.

2. 已知函数 $y = f\left( x\right)$ ，满足 $f\left( x\right)  = \ln \left( {2 - {3x}}\right)$ ，则它的导函数 ${y}^{\prime } =$ \_\_\_\_\_.

3. 函数 $f\left( x\right)  = x + \frac{1}{{x}^{2}}$ 的驻点为 $x =$ \_\_\_\_\_.

4. 已知直线 $y = x + 2$ 是曲线 $y = f\left( x\right)$ 在点 $\left( {0,2}\right)$ 处的切线,则 $\mathop{\lim }\limits_{{h \rightarrow  0}}\frac{f\left( {0 + h}\right)  - f\left( 0\right) }{h} =$ \_\_\_\_\_.

5. 已知函数 $f\left( x\right)$ 为定义在 $\mathbf{R}$ 上的偶函数,且当 $x \geq  0$ 时, $f\left( x\right)  = x + 2{e}^{x}$ ,则函数 $f\left( x\right)$ 在 $x =  - 1$ 处的切线斜率为\_\_\_\_\_.

6. 已知函数 $f\left( x\right)  = a{x}^{3} + {bx} + 1$ 的图像在点 $\left( {1,f\left( 1\right) }\right)$ 处的切线方程为 $y = {4x} - 1$ ，则 ${ab} =$ \_\_\_\_\_.

7. 若 $f\left( x\right)  = \ln x$ 与 $g\left( x\right)  = {x}^{2} + {3x} + a$ 两个函数的图像有一条与直线 $y = x$ 平行的公共切线，则实数 $a =$

\begin{center}
\includegraphics[width=0.2\textwidth]{images/bo_d7tk7oc91nqc73875fu0_83_1213_1667_257_298_0.jpg}
\end{center}
 

(第 8 题图)

8. 某高台跳水运动员在运动过程中的重心相对于水面的高度 $h$ (单位: ${m}$ ) 与跳起后的时间 $t$ (单位: ${s}$ ) 存在函数关系 $h\left( t\right)  =  - {4.9}{t}^{2} + {4.8t} + {11},h\left( t\right)$ 的图像如图所示,已知曲线 $h\left( t\right)$ 在 $t = {t}_{0}$ 处的切线 ${l}_{0}$ 平行于 $t$ 轴,根据图像,给出下列四个结论:

 \ding{172}  在 $t = {t}_{0}$ 时高度 $h$ 关于 $t$ 时间的瞬时变化率为 0 ;

 \ding{173} 曲线 $h\left( t\right)$ 在 $t = {t}_{2}$ 附近比在 $t = {t}_{1}$ 附近下降得慢;

 \ding{174} 曲线 $h\left( t\right)$ 在 $t = {t}_{3}$ 附近比在 $t = {t}_{4}$ 附近上升得快;

 \ding{175} 设在 $t = {t}_{2}$ 和 $t = {t}_{4}$ 时该运动员的瞬时速度分别为 ${v}_{2}$ 、 ${v}_{4}\left( {{m}/{s}}\right)$ ，

则 $\left| {v}_{2}\right|  > \left| {v}_{4}\right|$ . 其中所有正确结论的序号是\_\_\_\_\_.

\begin{center}
\includegraphics[width=0.2\textwidth]{images/bo_d7tk7oc91nqc73875fu0_84_1132_176_300_312_0.jpg}
\end{center}
 

(第 9 题图)

9. 如图,在墙角有一根长 $1{m}$ 的直木棒 ${AB}$ 紧贴墙面,墙面与底面垂直. 在 $t =  0{s}$ 时,木棒的端点 $A$ 从离地 $1{m}$ 的高度以 $1{m}/{s}$ 的速度竖直向下匀速运动,端点 $B$ 向右沿直线运动，则端点 $B$ 在 $t = \frac{\sqrt{2}}{2}{s}$ 这一时刻的瞬时速度为\_\_\_\_\_ ${m}/{s}$ .

\section*{二、选择题}

10. 下列函数的求导正确的是 ( )

A. ${\left( \ln {e}\right) }^{\prime } = \frac{1}{{e}}$ . B. ${\left( x\cos x\right) }^{\prime } = \cos x - x\sin x$ .

C. ${\left( {x}^{-2}\right) }^{\prime } =  - {2x}$ . D. ${\left( 2{{e}}^{x}\right) }^{\prime } = {{e}}^{x}$ .

11. 若函数 $y = f\left( x\right)$ 在 $x = {x}_{0}$ 处导数为 ${f}^{\prime }\left( {x}_{0}\right)$ ,则 $\mathop{\lim }\limits_{{h \rightarrow  0}}\frac{f\left( {{x}_{0} - {3h}}\right)  - f\left( {{x}_{0} + h}\right) }{h}$ 等于 ( )

A. $- {f}^{\prime }\left( {x}_{0}\right)$ . B. $3{f}^{\prime }\left( {x}_{0}\right)$ . C. $- 3{f}^{\prime }\left( {x}_{0}\right)$ . D. $- 4{f}^{\prime }\left( {x}_{0}\right)$ .

12. 已知函数 $y = f\left( x\right)$ ,其导函数 $y = {f}^{\prime }\left( x\right)$ 的图像如图所示. 如图四个选项中,可能表示函数 $y = f\left( x\right)$ 图像的是 ( )

yng

$y = {f}^{\prime }\left( x\right)$

-1 0 1 $x$

\begin{center}
\includegraphics[width=0.2\textwidth]{images/bo_d7tk7oc91nqc73875fu0_84_150_818_217_203_0.jpg}
\end{center}
 

A.

y

-1 0 1 $x$

B.

y

-1 0 1 $x$

C.

(第 12 题图)

$y$

-1 o $x$

D.

\section*{三、解答题}

13. 下列函数的导数.

(1) $y = {\left( 2x - 1\right) }^{3}$ ; (2) $y = x - \sin x\cos x$ ； (3) $y = \frac{\cos x}{{{e}}^{x}}$ .

14. 求下列各式的值.

(1) $f\left( x\right)  = {x}^{2} + 2\cos x$ ，求 ${f}^{\prime }\left( \frac{\pi }{2}\right)$ ；

(2)已知函数 $f\left( x\right)  = {{e}}^{x} + {2x}$ ，求 $\mathop{\lim }\limits_{{h \rightarrow  0}}\frac{{{e}}^{h} + {2h} - 1}{h}$ 的值.

15. 已知函数 $f\left( x\right)  = {x}^{3} - {x}^{2} - \frac{2}{3}{f}^{\prime }\left( {-1}\right) x$ ,其中 ${f}^{\prime }\left( x\right)$ 是 $f\left( x\right)$ 的导函数.

(1)求 ${f}^{\prime }\left( {-1}\right)$ ；

(2)求过原点与曲线 $y = f\left( x\right)$ 相切的切线方程.

16. 如图，酒杯的形状为倒立的圆锥，杯深 $8{{cm}}$ ，上口宽 $6{{cm}}$ ，水以 ${20}{{{cm}}}^{3}/{s}$ 的流量倒入杯中.

(1)给出水深 $h\left( {{cm}}\right)$ 和时间 $t\left( {s}\right)$ 的函数关系式;

(2)当水深为 $4{{cm}}$ 时,求水面升高的瞬时变化率.

\begin{center}
\includegraphics[width=0.2\textwidth]{images/bo_d7tk7oc91nqc73875fu0_85_1293_1099_182_239_0.jpg}
\end{center}
 

(第 16 题图)

\section*{§ 5.2 导数的应用一}

\section*{典型例题精讲}

【例 1】已知函数 $f\left( x\right)  = x - \ln x$ .

(1)求定义域及单调区间；

(2)求 $f\left( x\right)$ 的极值.

【例 2】已知函数 $f\left( x\right)  = x{\left( x - c\right) }^{2}$ 在 $x =  - 1$ 处取得极小值.

(1)求 $c$ 的值；

(2)求 $f\left( x\right)$ 在区间 $\left\lbrack  {-4,0}\right\rbrack$ 上的最值.

【例 3】设函数 $f\left( x\right)  = x - {x}^{3}{{e}}^{{ax} + b}$ ,曲线 $y = f\left( x\right)$ 在点 $\left( {1,f\left( 1\right) }\right)$ 处的切线方程为 $y =  - x + 1$ .

(1)设 $g\left( x\right)  = {f}^{\prime }\left( x\right)$ ，求 $g\left( x\right)$ 的单调区间；

(2)求 $f\left( x\right)$ 的极值点的个数.

【例 4】已知函数 $f\left( x\right)  = \ln x - a{x}^{2} + 1\left( {a \in  \mathbf{R}}\right)$ .

(1)求 $f\left( x\right)$ 的单调区间；

(2)若 $f\left( x\right)$ 有两个零点，求 $a$ 的取值范围.

【例 5】已知函数 $f\left( x\right)  = \left( {x + a}\right)  \cdot  {{e}}^{x}$ .

(1)当 $a = 1$ 时，求函数 $f\left( x\right)$ 的单调区间；

(2)过点 $O\left( {0,0}\right)$ 可作曲线 $y = f\left( x\right)$ 的两条切线,求实数 $a$ 的取值范围.

\section*{四基达标精练}

\section*{一、填空题}

1. 函数 $y = \frac{1}{x} - \ln x$ 的单调递减区间为\_\_\_\_\_.

2. 已函数 $f\left( x\right)  = \ln x - x + 1$ 的最大值为\_\_\_\_\_.

3. 已知函数 $f\left( x\right)  = {ax} - a - \ln x$ ，且 $f\left( x\right)$ 的最小值为 0，则实数 $a =$ \_\_\_\_\_.

4. 函数 $f\left( x\right)  = x - \ln x$ 在区间 $\left( {0,{e}}\right)$ 上的极小值为\_\_\_\_\_.

5. 函数 $f\left( x\right)  = \ln x - \left| {x - 2}\right|$ 的最大值为\_\_\_\_\_.

6. 已知函数 $f\left( x\right)  = \frac{1}{2}{x}^{2} - \left( {a + 2}\right) x + {2a}\ln x + 1$ 在 $\left( {4,6}\right)$ 上存在极值点,则实数 $a$ 的取值范围是\_\_\_\_\_.

\begin{center}
\includegraphics[width=0.2\textwidth]{images/bo_d7tk7oc91nqc73875fu0_88_1176_1667_287_222_0.jpg}
\end{center}
 

(第 8 题图)

7. 函数 $f\left( x\right)  = \frac{1}{3}{x}^{3} + a{x}^{2} + \left( {a + 2}\right) x - 1$ 有极大值又有极小值,则实数 $a$ 的范围是\_\_\_\_\_.

8. 已知直线 ${AB}$ 是函数 $y = f\left( x\right)$ 图像的一条割线 (如图所示), ${f}^{\prime }\left( x\right)$ 是函数 $f\left( x\right)$ 的导函数,若 $a = 2{f}^{\prime }\left( 3\right) \text{ 、 }b = 2{f}^{\prime }\left( 5\right) \text{ 、 }c = f\left( 5\right)  - f\left( 3\right)$ ,则关于 $a\text{ 、 }b\text{ 、 }c$ 从小到大的顺序是\_\_\_\_\_.

9. 已知函数 $f\left( x\right)  = \sin x - 1,g\left( x\right)  = a\ln x - x$ ,若对任意 ${x}_{1} \in  \mathbf{R}$ 都存在 ${x}_{2} \in  \left( {1,{e}}\right)$ ,使 $f\left( {x}_{1}\right)  < g\left( {x}_{2}\right)$ 成立，则实数 $a$ 的取值范围是\_\_\_\_\_.

\section*{二、选择题}

10. 函数 $f\left( x\right)  = \frac{x - 2}{{{e}}^{x}}$ 的单调递增区间为 ( )

A. $\left( {-\infty ,3}\right)$ . B. $\left( {0,3}\right)$ . C. $\left( {3, + \infty }\right)$ . D. $\left( {-\infty ,2}\right)$ .

11. 函数 $y = \frac{x}{{{e}}^{x}}$ 的极值点是 ( )

A.0. B.1. C. $\left( {0,1}\right)$ . D. $\left( {1,0}\right)$ .

12. 函数 $f\left( x\right)  = a{x}^{3} - {bx} - \ln x\left( {a > 0,b \in  \mathbf{R}}\right)$ 的一个极值点是 1,则 $\ln a - b + 1$ 的值 ( )

A. 恒大于 0 . C. 恒等于 0 . D. 不确定.

\section*{三、解答题}

13. 给出函数 $f\left( x\right)  = {x}^{3} - 4{x}^{2} + {4x}$ .

(1)求 $f\left( x\right)$ 的单调区间；

(2)求 $f\left( x\right)$ 在 $\left\lbrack  {0,3}\right\rbrack$ 上的最值.

14. 已知函数 $f\left( x\right)  = a{x}^{2} + b\ln x$ 在 $x = 1$ 处有极值 $\frac{1}{2}$ .

(1)求a、b的值；

(2)求函数 $y = f\left( x\right)$ 的单调性.

15. 已知函数 $f\left( x\right)  = \ln x - {ax} + 1\left( {a \in  \mathbf{R}}\right)$ .

(1)讨论 $f\left( x\right)$ 的单调性;

(2)若 $f\left( x\right)  \leq  0$ 恒成立，求 $a$ 的取值范围.

16. 设 $a > 0$ ,已知函数 $f\left( x\right)  = {\left( x - 2\right) }^{3} - {ax}$ .

(1)求函数 $y = f\left( x\right)$ 的单调区间；

(2)对于函数 $y = f\left( x\right)$ 的极值点 ${x}_{0}$ ，存在 ${x}_{1}\left( {{x}_{1} \neq  {x}_{0}}\right)$ ，使得 $f\left( {x}_{1}\right)  = f\left( {x}_{0}\right)$ ，试问对任意的正数 $a$ ， ${x}_{1} + 2{x}_{0}$ 是否为定值？若是，求出这个定值；若不是，请说明理由.

§ 5.3 导数的应用二

\section*{典型例题精讲}

【例 1】已知函数 $f\left( x\right)  = \left( {x + a}\right) {{e}}^{x}$ .

(1)若 $f\left( x\right)$ 在 $x = 1$ 处取得极小值，求实数 $a$ 的值；

(2)若 $f\left( x\right)$ 在 $\left( {-1,1}\right)$ 上单调递增，求实数 $a$ 的取值范围.

【例 2】已知函数 $f\left( x\right)  = \frac{1}{3}{x}^{3} - {x}^{2} + a$ .

(1)当 $a = 0$ 时，求函数 $f\left( x\right)$ 的极大值与极小值；

(2)若函数 $f\left( x\right)$ 在 $\left\lbrack  {1,3}\right\rbrack$ 上的最大值是最小值的 3 倍，求 $a$ 的值.

【例 3】已知 $x = 3$ 是函数 $f\left( x\right)  = a\ln x + {x}^{2} - {10x}$ 的一个极值点.

(1)求函数 $f\left( x\right)$ 的单调区间；

(2)若直线 $y = b$ 与 $y = f\left( x\right)$ 图像有 3 个交点，求 $b$ 的取值范围.

【例 4】已知函数 $f\left( x\right)  = {{e}}^{x} + m{{e}}^{-x} + \left( {m - 1}\right) x\left( {m \leq  0}\right)$ .

(1)讨论函数 $f\left( x\right)$ 的单调性;

(2)若对任意的 $x \in  \left( {0, + \infty }\right)$ ， $f\left( x\right)  \geq   - 2$ 恒成立，求 $m$ 的取值范围.


【例 5】已知函数 $f\left( x\right)  = \ln x + \left( {a - 1}\right) x + a + 1\left( {a \in  \mathbf{R}}\right)$ .

(1)讨论函数 $f\left( x\right)$ 的极值情况；

(2)证明:当 $a \leq  1$ 时， ${{e}}^{x} - f\left( x\right)  > 0$ .

\section*{四基达标精练}

\section*{一、填空题}

1. 已知函数 $f\left( x\right)  = \ln x + k$ 在 $\left\lbrack  {1,{e}}\right\rbrack$ 上的最大值为 2，则 $f\left( k\right)  =$ \_\_\_\_\_.

2. 已知函数 $f\left( x\right)  = {x}^{2} + \frac{a}{x}$ 在区间 $\left\lbrack  {1,3}\right\rbrack$ 上单调增，则实数 $a$ 的取值范围是\_\_\_\_\_.

3. 若函数 $f\left( x\right)  = 2{x}^{3} - 6{x}^{2} + m$ 在区间 $\left\lbrack  {-2,2}\right\rbrack$ 上有最大值 3,则 $f\left( x\right)$ 在区间 $\left\lbrack  {-2,2}\right\rbrack$ 上的最小值为\_\_\_\_\_.

4. 若函数 $f\left( x\right)  = {ax} - \ln x,x \in  \left( {1, + \infty }\right)$ 是严格增函数，则实数 $a$ 的最小值是\_\_\_\_\_.

5. 已知函数 $f\left( x\right)  = {x}^{3} - x$ 有 2 个极值点 ${x}_{1}\text{ 、 }{x}_{2}$ ,则 ${x}_{1} + {x}_{2} + f\left( {x}_{1}\right)  + f\left( {x}_{2}\right)  =$ \_\_\_\_\_.

6. 已知 $f\left( x\right)  = \ln x + {x}^{2} - {bx} + b$ 在区间 $\left( {1,2}\right)$ 上不单调，则实数 $b$ 的取值范围是\_\_\_\_\_.

7. 已知函数 $f\left( x\right)  = \ln x - \frac{1}{2}{x}^{2} - a$ ,如果存在 $x > 0$ ,使得 $f\left( x\right)  \geq  0$ ,则实数 $a$ 的取值范围

是

8. 已知函数 $f\left( x\right)  = {{e}}^{x} + m\ln x$ ( $m \in  \mathbf{R}$ , e 为自然对数的底数),若对任意正数 ${x}_{1}\text{ 、 }{x}_{2}$ ,当 ${x}_{1} > {x}_{2}$ 时都有 $f  \left( {x}_{1}\right)  - f\left( {x}_{2}\right)  > {x}_{1} - {x}_{2}$ 成立，则实数 $m$ 的取值范围是\_\_\_\_\_.

9. 已知点 $P$ 是曲线 $y = {x}^{2} - \ln x$ 上一点,求点 $P$ 到直线 $y = x - 3$ 的最小距离为\_\_\_\_\_.

\section*{二、选择题}

10. 设 $a \in  \mathbf{R}$ ,若函数 $f\left( x\right)  = {{e}}^{x} + {ax},x \in  \mathbf{R}$ 有大于零的极值点,则 $a$ 的取值范围为 ( )

A. $a <  - 1$ . B. $a <  - \frac{1}{{e}}$ C. $a >  - \frac{1}{{e}}$ . D. $a >  - 1$ .

11. 已知函数 $f\left( x\right)  = \left( {x - 1}\right) \ln x$ ,下列说法中正确的是 ( )

A. 1 既是 $y = f\left( x\right)$ 的一个零点,又是 $y = f\left( x\right)$ 的一个极小值点.

B. 1 既是 $y = f\left( x\right)$ 的一个零点,又是 $y = f\left( x\right)$ 的一个极大值点.

C. 1 是 $y = f\left( x\right)$ 的一个零点,不是 $y = f\left( x\right)$ 的极值点.

D. 1既不是 $y = f\left( x\right)$ 的一个零点,也不是 $y = f\left( x\right)$ 的极值点.

12. 函数 $f\left( x\right)$ 定义域为 $\left( {0, + \infty }\right)$ ,其导函数为 ${f}^{\prime }\left( x\right)$ ,若 $\forall x \in  \left( {0, + \infty }\right) ,x{f}^{\prime }\left( x\right)  + f\left( x\right)  < 1$ ,且 $f\left( 1\right)  =$ 2,则不等式 ${xf}\left( x\right)  - x < 1$ 的解集为 ( )

A. $\left( {1,2}\right)$ . B. $\left( {2, + \infty }\right)$ . C. $\left( {0,1}\right)$ . D. $\left( {1, + \infty }\right)$ .

\section*{三、解答题}

13. 已知函数 $y = f\left( x\right)  = a{x}^{3} - {x}^{2} + x - 5$ .

(1)若 $a =  - 1$ ，求 $y = f\left( x\right)$ 的极小值；

(2)若 $y = f\left( x\right)$ 在 $\mathbf{R}$ 上无极值，求 $a$ 的取值范围.

14. 设 $f\left( x\right)  =  - \frac{1}{3}{x}^{3} + \frac{1}{2}{x}^{2} + {ax}$ .

(1)若 $f\left( x\right)$ 在 $\left( {\frac{2}{3}, + \infty }\right)$ 上存在单调增区间，求 $a$ 的取值范围；

(2)若 $f\left( x\right)$ 在 $\left( {\frac{2}{3}, + \infty }\right)$ 上单调减，求 $a$ 的取值范围.

15. 已知函数 $f\left( x\right)  = \frac{1}{3}{x}^{3} - b{x}^{2} + {2x} - a,x = 2$ 是 $f\left( x\right)$ 的一个极值点.

(1)求 $f\left( x\right)$ 的单调增区间；

(2)当 $a > 0$ 时,求方程 $f\left( x\right)  = 0$ 的解的个数.

16. 已知函数 $f\left( x\right)  = \frac{1}{2}{x}^{2} + a\ln x$ .

(1)讨论 $f\left( x\right)$ 的单调性；

(2)若 $a = 1$ ，函数 $g\left( x\right)  = f\left( x\right)  - {mx}$ 是其定义域上的增函数，求实数 $m$ 的取值范围.

\section*{第 6 章 三角与三角函数(必修 2)}

\section*{知识梳理，}

\section*{一、任意角及其三角比}

1. 与 $\alpha$ 角终边相同、关于 $x$ 轴、 $y$ 轴、原点对称的角的集合分别是 $\{ \beta  \mid  \beta  = {2k\pi } + \alpha ,k \in  \mathbf{Z}\}$ 、 $\left\{  {\beta  \mid  \beta  = {2k\pi } - \alpha }\right\}$ ， $k \in  \mathbf{Z}\} \text{ 、 }\{ \beta  \mid  \beta  = {2k\pi } + \pi  - \alpha ,k \in  \mathbf{Z}\} \text{ 、 }\{ \beta  \mid  \beta  = {2k\pi } + \pi  + \alpha ,k \in  \mathbf{Z}\} .$

2. 弧度制: $\left| \alpha \right|  = \frac{l}{r}$ ( $\left| \alpha \right|$ 是将 $\left| \alpha \right|$ 看成圆心角时角 $\left| \alpha \right|$ 的弧度数， $r$ 是半径， $l$ 是弧长).

3. 半径为 $r$ 、圆心角为 $\alpha$ (弧度) (或 $\left. {n}^{ \circ  }\right)$ 的扇形弧长 $l = {r\alpha } = \frac{n\pi r}{180}$ ，面积 ${S}_{\text{ 扇 }} = \frac{1}{2}{lr} = \frac{1}{2}{r}^{2}\alpha  = \frac{{n\pi }{r}^{2}}{360}$ .

4. 设 $P\left( {x,y}\right)$ 是角 $\alpha$ 终边上任意一点 (异于原点) $\left| {OP}\right|  = \sqrt{{x}^{2} + {y}^{2}}$ ,则 $\sin \alpha  = \frac{y}{r};\cos \alpha  = \frac{x}{r};\tan \alpha  = \frac{y}{x}$ ; $\cot \alpha  = \frac{x}{y}$ ; 各三角比的符号由点 $P$ 所在的象限确定.

\section*{二、三角恒等变形}

1. 同角三角比基本关系: ${\sin }^{2}\theta  + {\cos }^{2}\theta  = 1,1 + {\tan }^{2}\theta  = \frac{1}{{\cos }^{2}\theta },1 + {\cot }^{2}\theta  = \frac{1}{{\sin }^{2}\theta },\sin \theta  = \tan \theta \cos \theta ,\cos \theta  =  \cot \theta \sin \theta ,\tan \theta  \cdot  \cot \theta  = 1$ .

2. 诱导公式: $n \cdot  \frac{\pi }{2} + \alpha \left( {n \in  \mathbf{Z}}\right)$ 变 $\alpha$ 的三角比遵循 “ $\left( n\right)$ 奇变偶不变，符号看象限(原象限原三角比)”.

3. 和角与差角三角比公式: $\sin \left( {\alpha  \pm  \beta }\right)  = \sin \alpha \cos \beta  \pm  \cos \alpha \sin \beta$ ;

$\cos \left( {\alpha  \pm  \beta }\right)  = \cos \alpha \cos \beta  \mp  \sin \alpha \sin \beta ,\tan \left( {\alpha  \pm  \beta }\right)  = \frac{\tan \alpha  \pm  \tan \beta }{1 \mp  \tan \alpha \tan \beta }.$

4. 辅助角公式 $a\sin \alpha  + b\cos \alpha  = \sqrt{{a}^{2} + {b}^{2}}\sin \left( {\alpha  + \varphi }\right)$ (其中 $\cos \varphi  = \frac{a}{\sqrt{{a}^{2} + {b}^{2}}},\sin \varphi  = \frac{b}{\sqrt{{a}^{2} + {b}^{2}}}$ ).

5. 二倍角公式: $\sin {2\alpha } = 2\sin \alpha \cos \alpha ;\cos {2\alpha } = {\cos }^{2}\alpha  - {\sin }^{2}\alpha  = 2{\cos }^{2}\alpha  - 1 = 1 - 2{\sin }^{2}\alpha$ .

6. 降幂公式: ${\sin }^{2}\alpha  = \frac{1 - \cos {2\alpha }}{2}$ ; ${\cos }^{2}\alpha  = \frac{1 + \cos {2\alpha }}{2}$ .

7. 万能公式: $\sin {2\alpha } = \frac{2\tan \alpha }{1 + {\tan }^{2}\alpha };\cos {2\alpha } = \frac{1 - {\tan }^{2}\alpha }{1 + {\tan }^{2}\alpha };\tan {2\alpha } = \frac{2\tan \alpha }{1 - {\tan }^{2}\alpha }$ .

\section*{三、解三角形}

1. $\bigtriangleup {ABC}$ 中,正弦定理: $\frac{a}{\sin A} = \frac{b}{\sin B} = \frac{c}{\sin C} = {2R}$ . 其中 $R$ 为 $\bigtriangleup {ABC}$ 外接圆半径;

2. 余弦定理: ${a}^{2} = {b}^{2} + {c}^{2} - {2bc}\cos A;{b}^{2} = {c}^{2} + {a}^{2} - {2ca}\cos B;{c}^{2} = {a}^{2} + {b}^{2} - {2ab}\cos C$ .

$\cos A = \frac{{b}^{2} + {c}^{2} - {a}^{2}}{2bc};\cos B = \frac{{c}^{2} + {a}^{2} - {b}^{2}}{2ca};\cos C = \frac{{a}^{2} + {b}^{2} - {c}^{2}}{2ab}$ .

3. 三角形的面积公式: \ding{172}  $S = \frac{1}{2}a{h}_{a} = \frac{1}{2}b{h}_{b} = \frac{1}{2}c{h}_{c}$ ( ${h}_{a}$ 、 ${h}_{b}$ 、 ${h}_{c}$ 分别表示 $a$ 、 $b$ 、 $c$ 边上的高)；

\HRule

* 注 本章节所属知识模块是 “函数与分析”.

\HRule

 \ding{173}  $S = \frac{1}{2}{ab}\sin C = \frac{1}{2}{bc}\sin A = \frac{1}{2}{ca}\sin B$ ； \ding{174}  ${S}_{\bigtriangleup {OAB}} = \frac{1}{2}\sqrt{{\left( \left| \overrightarrow{OA}\right|  \cdot  \left| \overrightarrow{OB}\right| \right) }^{2} - {\left( \overrightarrow{OA} \cdot  \overrightarrow{OB}\right) }^{2}}$ ；\_\_\_\_\_

公式的变式: $S = 2{R}^{2}\sin A\sin B\sin C = \frac{abc}{4R}(R$ 是 $\bigtriangleup {ABC}$ 的外接圆半径);

$S = \frac{1}{2}\left( {a + b + c}\right) r$ ( $r$ 是 $\bigtriangleup {ABC}$ 的内切圆半径).

\section*{四、三角函数的图像和性质}

\begin{center}
\begin{tabular}{|c|c|c|c|}
\hline
\phantom{X} & 正弦函数 & 余弦函数 & 正切函数 \\
\cline{1-4}
记 号 & $y = \sin x$ & $y = \cos x$ & $y = \tan x$ \\
\cline{1-4}
图像 & \includegraphics[width=0.2\textwidth]{images/bo_d7tk7oc91nqc73875fu0_97_289_557_326_267_0.jpg} & \includegraphics[width=0.2\textwidth]{images/bo_d7tk7oc91nqc73875fu0_97_674_557_328_266_0.jpg} & \includegraphics[width=0.2\textwidth]{images/bo_d7tk7oc91nqc73875fu0_97_1092_548_261_274_0.jpg} \\
\cline{1-4}
定义域 & R & $\mathbf{R}$ & $\{ x \mid  x \neq  {k\pi } + \frac{\pi }{2},k \in  \mathbf{Z}\}$ \\
\cline{1-4}
值域(最值) & $\left\lbrack  {-1,1}\right\rbrack  \left( {k \in  \mathbf{Z}}\right)  x = {2k\pi } - \frac{\pi }{2}$ 时, ${y}_{\min } =  - 1  x = {2k\pi } + \frac{\pi }{2}$ 时, ${y}_{\max } = 1$ & $\left\lbrack  {-1,1}\right\rbrack  \left( {k \in  \mathbf{Z}}\right)  x = {2k\pi } + \pi ,{y}_{\min } =  - 1  x = {2k\pi }$ 时， ${y}_{\max } = 1$ & R(无最值) \\
\cline{1-4}
单调性 & $\left\lbrack  {{2k\pi } - \frac{\pi }{2},{2k\pi } + \frac{\pi }{2}}\right\rbrack   \uparrow  \left( {k \in  \mathbf{Z}}\right)  \left\lbrack  {{2k\pi } + \frac{\pi }{2},{2k\pi } + \frac{3\pi }{2}}\right\rbrack   \downarrow$ & $\left\lbrack  {{2k\pi } - \pi ,{2k\pi }}\right\rbrack   \uparrow  \left( {k \in  \mathbf{Z}}\right)  \left\lbrack  {{2k\pi },{2k\pi } + \pi }\right\rbrack   \downarrow$ & $\left( {{k\pi } - \frac{\pi }{2},{k\pi } + \frac{\pi }{2}}\right)  \uparrow  \left( {k \in  \mathbf{Z}}\right)$ \\
\cline{1-4}
奇偶性 & 奇函数 & 偶函数 & 奇函数 \\
\cline{1-4}
对称轴 & $x = {k\pi } + \frac{\pi }{2}\left( {k \in  \mathbf{Z}}\right)$ & $x = {k\pi }\left( {k \in  \mathbf{Z}}\right)$ & 无 \\
\cline{1-4}
对称中心 & $\left( {{k\pi },0}\right) \left( {k \in  \mathbf{Z}}\right)$ & $\left( {{k\pi } + \frac{\pi }{2},0}\right) \left( {k \in  \mathbf{Z}}\right)$ & $\left( {\frac{k\pi }{2},0}\right) \left( {k \in  \mathbf{Z}}\right)$ \\
\cline{1-4}
周期性 & ${2\pi }$ & ${2\pi }$ & $\pi$ \\
\cline{1-4}
\hline
\end{tabular}
\end{center}

\section*{五、三角函数图像变换}

1. 变换法作函数 $y = A\sin \left( {{\omega x} + \varphi }\right) \left( {\omega  > 0,A > 0}\right)$ 的步骤:

(1)作函数 $y = \sin x$ 的图像；

(2)将 $y = \sin x$ 的图像右移 $\left( {\varphi  < 0}\right)$ 或左移 $\left( {\varphi  > 0}\right)$ 或 $\left| \varphi \right|$ 个单位得 $y = \sin \left( {x + \varphi }\right)$ 的图像；

(3)将 $y = \sin \left( {x + \varphi }\right)$ 的图像纵标不变，横标伸长(0 $< \omega  < 1$ )或缩短( $\omega  > 1$ )到原来的 $\frac{1}{\omega }$ 得到 $y = \sin \left( {{\omega x} + \varphi }\right)$ 的图像;

(4)将 $y = \sin \left( {{\omega x} + \varphi }\right)$ 的图像横标不变，纵标伸长 (A > 1) 或缩短 (0 $< A < 1$ ) 到原来的 $A$ 倍得 $y =  A\sin \left( {{\omega x} + \varphi }\right)$ 图像;

注意: 由 $y = \sin {\omega x}$ 的图像得到 $y = \sin \left( {{\omega x} + \varphi }\right)$ 的图像时是将 $y = \sin {\omega x}$ 的图像左 (或右) 移 $\left| \frac{\varphi }{\omega }\right|$ 个单

位. (要注意反向平移问题)

2. 已知图像求 $y = A\sin \left( {{\omega x} + \varphi }\right) \left( {A > 0,\omega  > 0}\right)$ 解析式的问题:

(1)根据图像确定 $\omega$ :

方法一:相邻最大值和最小值对应的横坐标之差的绝对值是周期的一半；

方法二:相邻两个零点对应的横坐标之差的绝对值是周期的一半；

方法三:最值与相邻的零点的横坐标之差的绝对值是周期的四分之一.

(2)根据图像确定 $A$ :最大值与最小值之差的一半的绝对值是 $\left| A\right|$ .

(3)根据图像确定 $\varphi$ :递增区间中零点的横坐标 ${x}_{0}$ 对应于 $\omega {x}_{0} + \varphi  = {2k\pi }\left( {k \in  \mathbf{Z}}\right)$ ；递减区间中零点的横坐标 ${x}_{0}$ 对应于 $\omega {x}_{0} + \varphi  = {2k\pi } + \pi \left( {k \in  \mathbf{Z}}\right)$ .

\section*{六、反三角和三角方程}

反三角函数的本质:

$\arcsin x\left( {x \in  \left\lbrack  {0,1}\right\rbrack  }\right)$ 表示 $\left\lbrack  {0,\frac{\pi }{2}}\right\rbrack$ 内正弦值为 $x$ 的角;

$\arccos x\left( {x \in  \left\lbrack  {0,1}\right\rbrack  }\right)$ 表示 $\left\lbrack  {0,\frac{\pi }{2}}\right\rbrack$ 内余弦值为 $x$ 的角;

$\arctan x\left( {x \geq  0}\right)$ 表示 $\left\lbrack  {0,\frac{\pi }{2}}\right)$ 内正切值为 $x$ 的角.

\section*{$\text{ § }{6.1}$ 三角及其恒等变形}

\section*{典型例题精讲}

【例 1】已知角 $\alpha$ 的终边过点 $P\left( {{7m}, - {3m}}\right) \left( {m \neq  0}\right)$ .

(1) 求 $\frac{\cos \left( {\frac{\pi }{2} + \alpha }\right) \sin \left( {-\pi  - \alpha }\right) }{\cos \left( {\frac{11\pi }{2} - \alpha }\right) \sin \left( {\frac{9\pi }{2} + \alpha }\right) }$ 的值；

(2) 求 $\sin \left( {\pi  - \alpha }\right) \cos \left( {{2\pi } - \alpha }\right)  - {\sin }^{2}\left( {\frac{\pi }{2} + \alpha }\right)  + 2$ 的值.

\section*{高分钥匙高考数学复习一轮通}

【例 2】已知是关于 $x$ 的方程 $k{x}^{2} - {kx} + k - 3 = 0$ 的两个实根.

(1)若 ${x}_{1} = \sin \alpha$ ， ${x}_{2} = \cos \alpha \left( {\alpha  \in  \mathbf{R}}\right)$ ，求实数 $k$ 的值；

(2)若 ${x}_{1} = \sin \alpha$ ， ${x}_{2} = \cos \beta \left( {\alpha \text{ 、 }\beta  \in  \mathbf{R}}\right)$ ，求实数 $k$ 的取值范围.

【例 3】已知角 $\theta$ 的顶点为坐标原点 $O$ ,始边为 $x$ 轴的非负半轴,终边与单位圆相交于点 $P\left( {x,y}\right)$ ,若点 $P$ 位于 $x$ 轴上方且 $x + y = \frac{1}{2}$ .

(1)求 $\sin \theta  - \cos \theta$ 的值；

(2)求 ${\sin }^{4}\theta  + {\cos }^{4}\theta$ 的值.

【例 4】已知 $0 < \alpha  < \frac{\pi }{3}, - \frac{\pi }{2} < \beta  <  - \frac{\pi }{3},\cos \left( {\alpha  + \frac{\pi }{6}}\right)  = \frac{\sqrt{3}}{3},\sin \left( {\alpha  - \beta  + \frac{\pi }{6}}\right)  = \frac{2}{3}$ .

(1)求 $\sin \alpha$ 的值；

(2)求 $\cos \left( {{3\alpha } - \beta }\right)$ 的值.

【例 5】已知向量 $\overrightarrow{a} = \left( {\cos \alpha ,\sin \alpha }\right) ,\overrightarrow{b} = \left( {\cos \beta , - \sin \beta }\right) ,\overrightarrow{c} = \left( {2,1}\right) ,\overrightarrow{a}//\overrightarrow{c},\overrightarrow{a} \cdot  \overrightarrow{b} = \frac{\sqrt{5}}{5}$ .

(1)求 $\frac{{\sin }^{2}\alpha  + 2\cos {2\alpha }}{1 + \tan \alpha }$ 的值；

(2)若 $\alpha$ 、 $\beta$ 均为锐角,求 $\tan \left( {\alpha  - \beta }\right)$ 的值.

\section*{四基达标精练}

\section*{一、填空题}

1. 角 $\alpha$ 的终边过点 $P\left( {-{6a}, - {8a}}\right) \left( {a > 0}\right)$ ，则 $\sin \alpha  - \cos \alpha  =$ \_\_\_\_\_.

2. 关于 $x$ 的不等式 $\sin x \leq   - \frac{\sqrt{3}}{2}$ 的解集为\_\_\_\_\_.

3. 若一个扇形的半径是 2,圆心角 $\alpha$ 所对的弦长为 $2\sqrt{3}$ ,则 $\alpha  =$ \_\_\_\_\_.

4. 若 $\cos \alpha  = \frac{2}{3}$ ，且 $\alpha$ 在第四象限，则 $\sin \left( {\alpha  - {2\pi }}\right)  + \sin \left( {-\alpha  - {3\pi }}\right) \cos \left( {\alpha  - {3\pi }}\right)  =$ \_\_\_\_\_.

5. 已知 $\tan \theta  = \frac{1}{3}$ ,则 $\frac{\tan \left( {{2\pi } - \theta }\right) }{\cos \left( {\frac{\pi }{2} - \theta }\right) \sin \left( {\frac{3\pi }{2} + \theta }\right) } =$ \_\_\_\_\_.

6. 已知点 $A$ 的坐标为 $\left( {-3,4}\right)$ ，将 ${OA}$ 绕坐标原点 $O$ 顺时针旋转 $\frac{\pi }{3}$ 至 ${OB}$ . 则点 $B$ 的坐标为\_\_\_\_\_.

7. 已知 $A\left( {{x}_{1},{y}_{1}}\right)$ 、 $B\left( {{x}_{2},{y}_{2}}\right)$ 是角 $\alpha$ 、 $\beta$ 终边与单位圆的两个不同交点，且 ${x}_{1}{y}_{2} = {x}_{2}{y}_{1}$ ，则 $2{x}_{1} - {x}_{2} +  2{y}_{1} - {y}_{2}$ 的最大值为\_\_\_\_\_.

8. 已知长方形的四个顶点: $A\left( {0,0}\right) \text{ 、 }B\left( {2,0}\right) \text{ 、 }C\left( {2,1}\right) \text{ 、 }D\left( {0,1}\right)$ ,一质点从点 $A$ 出发,沿与 ${AB}$ 夹角为 $\theta$ 的方向射到 ${BC}$ 上的点 ${P}_{1}$ 后,依次反射到 ${CD}\text{ 、 }{DA}$ 和 ${AB}$ 上的点 ${P}_{2}\text{ 、 }{P}_{3}\text{ 、 }{P}_{4}$ (入射角等于反射角). 设 ${P}_{4}$ 的坐标为 $\left( {{x}_{4},0}\right)$ ,若 $1 < {x}_{4} < 2$ ,则 $\tan \theta$ 的范围是\_\_\_\_\_.

9. 已知 $a$ 、 $b$ 、 $x \in  \mathbf{R},{ab} \neq  0$ ，且 $\frac{{\sin }^{4}x}{{a}^{2}} + \frac{{\cos }^{4}x}{{b}^{2}} = \frac{1}{{a}^{2} + {b}^{2}}$ ，则 $\frac{{\sin }^{2022}x}{{a}^{2020}} + \frac{{\cos }^{2022}x}{{b}^{2020}} =$ \_\_\_\_\_.

\section*{二、选择题}

10. 在平面直角坐标系 ${xOy}$ 中,角 $\alpha  = {120}^{ \circ  }$ ,其顶点为坐标原点 $O$ ,始边与 $x$ 轴的非负半轴重合,终边过点 $P\left( {-4,m}\right)$ ,则实数 $m$ 的值是 ( )

A.2. $\sqrt{3}$ B. $- 2\sqrt{3}$ C. $4\sqrt{3}$ D. $- 4\sqrt{3}$

11. 在直角坐标平面上,角 $\sin \frac{\pi }{6}$ 所在的象限是 ( )

A. 一 B. 二 C. 三 D. 四

12. 已知 $\sin \frac{q + p}{2} < 0,\tan \frac{q + p}{2} < 0,\cos \frac{p - q}{2} > 0$ ,则下列关系中必定不成立的是 ( )

A. $\sin q < 0,\cos q > 0$ B. $\sin q > 0,\cos q < 0$ C. $\sin q > 0,\cos q > 0$ D. $\sin q < 0,\cos q < 0$

\section*{三、解答题}

13. 已知 $\sin x + \cos x = \frac{1}{5},x \in  \left( {0,\pi }\right)$ .

(1) 求 $\sin x - \cos x$ 的值；

(2)求 $\cos {2x}$ 的值.

14. 已知关于 $x$ 的方程 ${13}{x}^{2} - {7x} + m = 0$ 的两根为 $\sin \theta \text{ 、 }\cos \theta ,\theta  \in  \left( {0,\pi }\right)$ .

(1)求 $m$ 的值；

(2)求 $\tan \theta$ 的值.

15. 已知 $\overrightarrow{a} = \left( {\cos \alpha ,\sin \alpha }\right) ,\overrightarrow{b} = \left( {\cos \beta ,\sin \beta }\right) ,0 < \beta  < \alpha  < \pi$ .

(1)若 $\left| {\overrightarrow{a} - \overrightarrow{b}}\right|  = \sqrt{2}$ ，求证: $\overrightarrow{a}\bot \overrightarrow{b}$ ；

(2)设 $\overrightarrow{c} = \left( {0,1}\right)$ ，若 $\overrightarrow{a} + \overrightarrow{b} = \overrightarrow{c}$ ，求 $\cos \left( {\alpha  - \beta }\right)$ 的值.

16. 如图，在平面直角坐标系中，角 $\alpha$ 和 $\beta$ 的终边与单位圆分别交于 $P$ 、 $Q$ 两点.

(1)若 $\overrightarrow{OP} + \overrightarrow{OQ} = \left( {\frac{1}{2},\frac{3}{2}}\right)$ ，求 $\cos \left( {\alpha  - \beta }\right)$ 的值；

(2)若 $\alpha  = \frac{\pi }{6},\left| {\overrightarrow{OP} - \overrightarrow{OQ}}\right|  = \frac{\sqrt{2}}{3}$ ，求 $\cos \left( {{2\beta } + \frac{2}{3}\pi }\right)$ 的值.

\begin{center}
\includegraphics[width=0.2\textwidth]{images/bo_d7tk7oc91nqc73875fu0_102_1159_1180_281_245_0.jpg}
\end{center}
 

(第 16 题图)

\section*{§ 6.2 解三角形}

\section*{典型例题精讲}

【例 1】在 $\bigtriangleup {ABC}$ 中,角 $A\text{ 、 }B\text{ 、 }C$ 所对的边分别为 $a\text{ 、 }b\text{ 、 }c$ ,已知 ${c}^{2} = {a}^{2} + {b}^{2} + {ab}$ .

(1)求角 $C$ 的大小；

(2)若 $c = 4\sqrt{7},\bigtriangleup {ABC}$ 的周长为 ${12} + 4\sqrt{7}$ ，求 $\sin A + \sin B$ .

【例 2】如图所示,某人为了测量出到河对岸铁塔的距离与铁塔的高,选与塔底 $B$ 同在水平面内的两个观测点 $C$ 与 $D$ ,在 $C$ 点测得塔底 $B$ 在北偏东 ${45}^{ \circ  }$ 方向,然后向正东方向前进 ${10}{m}$ 到达 $D$ ,测得此时塔底 $B$ 在北偏东 ${15}^{ \circ  }$ 方向.

(1)求点 $D$ 到塔底 $B$ 的距离 ${BD}$ ；

(2)若在点 $C$ 测得塔顶 $A$ 的仰角为 ${60}^{ \circ  }$ ，求铁塔 ${AB}$ 的高.

(注:结果保留根号)

\begin{center}
\includegraphics[width=0.2\textwidth]{images/bo_d7tk7oc91nqc73875fu0_103_1184_1447_268_377_0.jpg}
\end{center}
 

(例 2 图)

【例 3】已知 $\bigtriangleup  {ABC}$ 的内角和边满足 $a\cos C + \sqrt{3}a\sin C - b - c = 0$ ，其中角 $A$ 、 $B$ 、 $C$ 所对的边分别是 $a$ 、 $b$ 、 $c$ . (1)求角 $A$ 的大小；

(2)若 $a = \sqrt{3}$ ，求 $\bigtriangleup  {ABC}$ 面积的最大值.

【例 4】如图,已知点 $D$ 在 $\bigtriangleup {ABC}$ 的边 ${BC}$ 上,且 $\angle {DAC} = {90}^{ \circ  },\cos \angle {DAB} = \frac{2\sqrt{2}}{3},{AB} = 6$ .

(1)若 $\sin C = \frac{\sqrt{3}}{3}$ ，求线段 ${BC}$ 的长；

(2)若点 $E$ 是 ${BC}$ 的中点， ${AE} = \sqrt{17}$ ，求线段 ${AC}$ 的长.

\begin{center}
\includegraphics[width=0.3\textwidth]{images/bo_d7tk7oc91nqc73875fu0_104_1076_1215_363_180_0.jpg}
\end{center}
 

(例 4 图)

【例 5】已知 $\bigtriangleup {ABC}$ 的三个内角 $A\text{ 、 }B\text{ 、 }C$ 所对的边分别是 $a\text{ 、 }b\text{ 、 }c$ ,且满足 $3\overrightarrow{AB} \cdot  \overrightarrow{AC} + 4\overrightarrow{CA} \cdot  \overrightarrow{CB} =  \overrightarrow{BA} \cdot  \overrightarrow{BC}$ .

(1)求 $\frac{c}{b}$ 的值；

(2)若 $\overrightarrow{AD} = \frac{2}{3}\overrightarrow{AB} + \frac{1}{3}\overrightarrow{AC}$ ， $\left| \overrightarrow{AD}\right|  = 2$ ，求 $\bigtriangleup {ABC}$ 面积的最大值.

\section*{四基达标精练}

\section*{一、填空题}

1. 在 $\bigtriangleup {ABC}$ 中， ${a}^{2} + {bc} = {b}^{2} + {c}^{2}$ ，则 $A =$ \_\_\_\_\_.

2. 已知 $\bigtriangleup  {ABC}$ 的内角 $A$ 、 $B$ 、 $C$ 所对的边分别是 $a$ 、 $b$ 、 $c$ ，若 $A = {30}^{ \circ  }$ ， $B = {105}^{ \circ  }$ ， $a = 4$ ，则 $c =$ \_\_\_\_\_.

3. 在 $\bigtriangleup  {ABC}$ 中， $a = 2\sqrt{6}$ ， $b = {2c}$ ， $\cos A =  - \frac{1}{4}$ ，则 $\bigtriangleup  {ABC}$ 的面积 ${S}_{\bigtriangleup {ABC}} =$ \_\_\_\_\_.

4. 边长分别为 ${10}\text{ 、 }{14}\text{ 、 }{16}$ 的三角形中,最大角与最小角的和为\_\_\_\_\_.

5. 在 $\bigtriangleup  {ABC}$ 中， ${BC} = 6$ ， ${AC} = 8$ ， $\angle A = {40}^{ \circ  }$ ，则 $\angle B$ 的解的个数是\_\_\_\_\_.

6. 在 $\bigtriangleup {ABC}$ 中,角 $A\text{ 、 }B\text{ 、 }C$ 所对的边分别为 $a\text{ 、 }b\text{ 、 }c,\bigtriangleup {ABC}$ 的面积 $S\bigtriangleup {ABC} = \frac{\sqrt{3}}{4}\left( {{a}^{2} + {b}^{2} - {c}^{2}}\right)$ ,则 $\frac{c}{a + b}$ 的取值范围是\_\_\_\_\_.

7. 在 $\bigtriangleup  {ABC}$ 中，若 $\left( {\sin A + \sin B}\right) \left( {\sin A - \sin B}\right)  = {\sin }^{2}C$ ，则 $A =$ \_\_\_\_\_.

8. 在 $\bigtriangleup  {ABC}$ 中，角 $A$ 、 $B$ 、 $C$ 所对的边分别是 $a$ 、 $b$ 、 $c$ ，若 $\frac{a}{\sin B} = \frac{b}{\sin C} = \frac{c}{\sin A}$ ，则 $\bigtriangleup  {ABC}$ 的形状是

9. 已知 $\bigtriangleup {ABC}$ 的三个内角 $A$ 、 $B$ 、 $C$ 成等差数列，且 ${AB} = 1$ ， ${BC} = 4$ ，则边 ${AC}$ 上的中线 ${BD}$ 的长为\_\_\_\_\_.

\section*{二、选择题}

10. 在 $\bigtriangleup  {ABC}$ ，角 $A$ 、 $B$ 、 $C$ 所对的边分别为 $a$ 、 $b$ 、 $c$ ，若 $A : B : C = 3 : 2 : 1$ ，则 $a$ : $b : c$ 等于\_\_\_\_\_。

A. $1 : 2 : 3$ . B. $3 : 2 : 1$ . C. $1 : \sqrt{3} : 2$ . D. $2 : \sqrt{3} : 1$ .

11. 在 $\bigtriangleup  {ABC}$ 中，角 $A$ 、 $B$ 、 $C$ 所对的边分别为 $a$ 、 $b$ 、 $c$ ，若 $c = 4$ ， $a = 4\sqrt{2}$ ， $A = \frac{\pi }{4}$ ，则角 $C$ 的大小为( )

A. $\frac{\pi }{6}$ . B. $\frac{\pi }{6}$ 或 $\frac{5\pi }{6}$ . C. $\frac{\pi }{3}$ . D. $\frac{\pi }{3}$ 或 $\frac{2\pi }{3}$ .

12. 在 $\bigtriangleup  {ABC}$ 中，角 $A$ 、 $B$ 、 $C$ 所对应的边分别为 $a$ 、 $b$ 、 $c$ ，设 $\bigtriangleup  {ABC}$ 的面积为 $S$ ，则 $\frac{S}{{a}^{2} + {4bc}}$ 的最大值为 ()

A. $\frac{\sqrt{2}}{16}$ . B. $\frac{\sqrt{3}}{12}$ . C. $\frac{\sqrt{3}}{16}$ . D. $\frac{\sqrt{2}}{18}$ .

\section*{三、解答题}

13. 在 $\bigtriangleup {ABC}$ 中,已知 $a = 2,b = \sqrt{2},c = \sqrt{3} + 1$ .

(1)求角 $A$ 、 $B$ 、 $C$ 的大小；

(2)求 $\bigtriangleup  {ABC}$ 面积.

14. 在 $\bigtriangleup  {ABC}$ 中，内角 $A$ 、 $B$ 、 $C$ 所对的边分别为 $a$ 、 $b$ 、 $c$ ，已知 $a\sin C = \sqrt{3}c\cos A$ .

(1)求角 $A$ 的大小；

(2)若 $b = 2$ ，且 $\frac{\pi }{4} \leq  B \leq  \frac{\pi }{3}$ ，求边 $c$ 的取值范围.

15. 已知 $\bigtriangleup  {ABC}$ 的内角 $A$ 、 $B$ 、 $C$ 的对边分别记为 $a$ 、 $b$ 、 $c$ ，且 ${2c}\cos A = b - c$ .

(1)证明: $A = {2C}$ ；

(2)若 $\bigtriangleup  {ABC}$ 为锐角三角形，求 $\frac{a + b}{c}$ 的取值范围.

16. 海岸的 $A$ 哨所在凌晨 1 点 15 分发现哨所北偏东 ${30}^{ \circ  }$ 方向 20 海里处的点 $D$ 出现可疑船只，因天气恶劣能见度低,无法对船只进行识别,所以将该船雷达特征信号进行标记并上报周围哨所. 早上 5 点 15 分位于 $A$ 哨所正西方向 20 海里的 $B$ 哨所发现了该可疑船只位于 $B$ 哨所北偏西 30 \textdegree 方向 60 海里处的点 $E$ ，并识别出其为走私船,立刻命令位于 $B$ 哨所正西方向 30 海里处点 $C$ 的我方缉私船前往拦截，已知缉私船速度大小为 30 海里/时. (假设所有船只均保持匀速直线航行)

(1)求走私船的速度大小；

(2)缉私船沿什么方向行驶才能最快截获走私船，并求出截获走私船的具体时间.

\begin{center}
\includegraphics[width=0.2\textwidth]{images/bo_d7tk7oc91nqc73875fu0_107_1137_1144_322_261_0.jpg}
\end{center}
 

(第 16 题图)

\section*{§ 6.3 三角函数及其性质一}

\section*{典型例题精讲}

【例 1】已知函数 $f\left( x\right)  = \sqrt{3}\sin x\cos x - {\cos }^{2}x$ .

(1)求函数 $f\left( x\right)$ 的最小正周期和单调区间；

(2)求函数 $f\left( x\right)$ 的零点.

【例 2】已知 $f\left( x\right)  = \cos {\omega x}\left( {\omega  > 0}\right)$ .

(1) $f\left( x\right)$ 的周期是 $\pi$ ,求当 $x \in  \left\lbrack  {0,{2\pi }}\right\rbrack$ ，方程 $f\left( {x + \frac{\pi }{6}}\right)  = \frac{\sqrt{3}}{2}$ 的解集；

(2)已知 $\omega  = 1$ ， $g\left( x\right)  = {f}^{2}\left( x\right)  + \sqrt{3}f\left( {-x}\right)$ ， $f\left( {\frac{\pi }{2} - x}\right)$ ， $x \in  \left\lbrack  {0,\frac{\pi }{4}}\right\rbrack$ ，求 $g\left( x\right)$ 的值域.

【例 3】已知函数 $f\left( x\right)  = 2\cos \left( {{\omega x} + \varphi }\right) \left( {\omega  \in  \mathbf{N},\left| \varphi \right|  < \frac{\pi }{2}}\right)$ 在区间 $\left( {0,\frac{5\pi }{18}}\right)$ 上是单调函数，且 $f\left( \frac{\pi }{18}\right)  +  f\left( \frac{\pi }{6}\right)  = 0.$

(1)求 $f\left( x\right)$ 图像的一个对称中心；

(2)若 $f\left( \frac{\pi }{6}\right)  =  - 1$ ，求 $f\left( x\right)$ 的解析式.

【例 4】函数 $f\left( x\right)  = \cos \left( {{\omega x} + \varphi }\right) \left( {\omega  > 0,\left| \varphi \right|  < \frac{\pi }{2}}\right)$ 的部分图像如图所示.

(1)求 $f\left( x\right)$ 的解析式;

(2)若对任意 $x \in  \left\lbrack  {-\frac{\pi }{4},\frac{\pi }{4}}\right\rbrack$ ，都有 ${\left\lbrack  f\left( x\right) \right\rbrack  }^{2} - {mf}\left( x\right)  - 1 \leq  0$ ，求实数 $m$ 的取值范围.

\begin{center}
\includegraphics[width=0.2\textwidth]{images/bo_d7tk7oc91nqc73875fu0_109_1097_1328_322_208_0.jpg}
\end{center}
 

(例 4 图)

【例 5】已知 $f\left( x\right)  = 2{x}^{2} - {3x} + 1,g\left( x\right)  = k \cdot  \sin \left( {x - \frac{\pi }{6}}\right) \left( {k \neq  0}\right)$ .

(1)设 $f\left( x\right)$ 的定义域为 $\left\lbrack  {0,3}\right\rbrack$ ，值域为 $A$ ； $g\left( x\right)$ 的定义域为 $\left\lbrack  {0,3}\right\rbrack$ ，值域为 $B$ ，且 $A \subseteq  B$ ，求实数 $k$ 的取值范围.

(2)若方程 $f\left( {\sin x}\right)  + \sin x - a = 0$ 在 $\lbrack 0,{2\pi })$ 上恰有两个解，求实数 $a$ 的取值范围.

\section*{四基达标精练}

\section*{一、填空题}

1. $f\left( x\right)  = \left| {\sin x}\right|$ 的值域是\_\_\_\_\_.

2. 函数 $y = 2\cos x$ 的严格减区间为\_\_\_\_\_.

3. 函数 $y = \tan \left( {\frac{\pi }{6}x + \frac{\pi }{3}}\right)$ 的定义域是\_\_\_\_\_.

4. 方程 $\cos {2x} - \sin x = 0$ 在区间 $\left\lbrack  {0,{2\pi }}\right\rbrack$ 上的所有解的和为\_\_\_\_\_.

5. 若函数 $f\left( x\right)  = \cos {\omega x}\left( {\omega  > 0}\right)$ 相邻两条对称轴之间的距离为 2，则 $f\left( 2\right)  =$ \_\_\_\_\_.

6. 若函数 $y = \sin \left( {{\pi x} - \frac{\pi }{6}}\right)$ 在 $\left\lbrack  {0,m}\right\rbrack$ 上单调增,则 $m$ 的最大值为\_\_\_\_\_.

7. 若方程 $\cos x = \frac{a - 1}{2}$ 在 $x \in  \left\lbrack  {-\frac{\pi }{3},\pi }\right\rbrack$ 上有两个不同的实数根，则实数 $a$ 的取值范围为\_\_\_\_\_.

8. 已知函数 $f\left( x\right)  = a\sin x + {bx} + 1$ ，若 $f\left( {-1}\right)  = 2$ ，则 $f\left( 1\right)  =$ \_\_\_\_\_.

9. 已知 $f\left( x\right)  = {x}^{2} - \cos x$ ,对 $\left\lbrack  {-\frac{\pi }{2},\frac{\pi }{2}}\right\rbrack$ 上的任意 ${x}_{1}\text{ 、 }{x}_{2}$ ,有如下条件: $\text{  \ding{172}  }{x}_{1} > {x}_{2}$ ;  \ding{173}  ${x}_{1}^{2} > {x}_{2}^{2}$ ； \ding{174}  $\left| {x}_{1}\right|  > {x}_{2}$ . 其中能使 $f\left( {x}_{1}\right)  > f\left( {x}_{2}\right)$ 恒成立的条件序号是\_\_\_\_\_.

\section*{二、选择题}

10. 下列四个函数中,同时具有性质:  \ding{172}  最小正周期为 ${2\pi }$ ;  \ding{173}  图像关于直线 $x = \frac{\pi }{3}$ 对称的一个函数是 ( )

A. $y = \sin \left( {x - \frac{\pi }{6}}\right)$ . B. $y = \sin \left( {x + \frac{\pi }{6}}\right)$ . C. $y = \sin \left( {x + \frac{\pi }{3}}\right)$ . D. $y = \sin \left( {{2x} - \frac{\pi }{3}}\right)$ .

11. 函数 $y = x + \sin \left| x\right| \left( {x \in  \left\lbrack  {-\pi ,\pi }\right\rbrack  }\right)$ 的大致图像是 ( )

\begin{center}
\includegraphics[width=0.2\textwidth]{images/bo_d7tk7oc91nqc73875fu0_111_161_206_256_210_0.jpg}
\end{center}
 

A.

\begin{center}
\includegraphics[width=0.4\textwidth]{images/bo_d7tk7oc91nqc73875fu0_111_463_205_576_243_0.jpg}
\end{center}
 

(第 11 题图)

\begin{center}
\includegraphics[width=0.2\textwidth]{images/bo_d7tk7oc91nqc73875fu0_111_1117_199_254_210_0.jpg}
\end{center}
 

D.

12. 函数 $f\left( x\right)  = \cos \left( {{\omega x} + \varphi }\right) \left( {0 < \varphi  < \pi }\right)$ 周期为 $T = \frac{2\pi }{3}$ ，其图像关于 $\left( {\frac{\pi }{18},0}\right)$ 对称，且当 $x \in  \left\lbrack  {\frac{\pi }{6},m}\right\rbrack$ 时， $f\left( x\right)$ 的值域是 $\left\lbrack  {-1, - \frac{\sqrt{3}}{2}}\right\rbrack$ ,则 $m$ 的取值范围是 ( )

A. $\left\lbrack  {\frac{\pi }{9},\frac{7\pi }{18}}\right\rbrack$ . B. $\left\lbrack  {\frac{2\pi }{9},\frac{7\pi }{18}}\right\rbrack$ . C. $\left\lbrack  {\frac{\pi }{9},\frac{5\pi }{18}}\right\rbrack$ . D. $\left\lbrack  {\frac{2\pi }{9},\frac{5\pi }{18}}\right\rbrack$ .

\section*{三、解答题}

13. 已知函数 $f\left( x\right)  = 2\sin \left( {{2x} - \frac{\pi }{6}}\right)$ .

(1)求 $f\left( x\right)$ 的最小正周期及单调递减区间；

(2)求 $f\left( x\right)$ 在区间 $\left\lbrack  {\frac{\pi }{4},\frac{\pi }{2}}\right\rbrack$ 上的最大值与最小值.

14. 已知函数 $f\left( x\right)  = 2\sin \left( {{2\omega x} - \frac{\pi }{6}}\right) ,{x}_{1}\text{ 、 }{x}_{2}$ 是函数 $y = f\left( x\right)$ 与直线 $y = 2$ 的两个交点的横坐标，且 $\left| {{x}_{1} - {x}_{2}}\right|$ 的最小值为 $\pi$ .

(1)求 $f\left( x\right)$ 的解析式;

(2)若 $x \in  \left\lbrack  {\frac{\pi }{2},m}\right\rbrack$ ， $f\left( x\right)$ 的值域是 $\left\lbrack  {-2,1}\right\rbrack$ ，求 $m$ 的取值范围.

15. 已知函数 $f\left( x\right)  = A\sin \left( {{\omega x} + \varphi }\right) ,\left( {\omega  > 0,\left| \varphi \right|  < \frac{\pi }{2}}\right)$ 的部分图像如图所示.

(1)求 $f\left( x\right)$ 的解析式；

(2)若 $x \in  \left( {-\frac{\pi }{18},\frac{5\pi }{18}}\right)$ ，求 $f\left( x\right)$ 的取值范围.

\begin{center}
\includegraphics[width=0.2\textwidth]{images/bo_d7tk7oc91nqc73875fu0_112_1116_362_320_205_0.jpg}
\end{center}
 

(第 15 题图)

16. 如图，已知 ${OPQ}$ 是半径为 $\sqrt{7}$ ，圆心角为 $\frac{\pi }{3}$ 的扇形， $C$ 是圆弧 $\overset{\text{ ⏜ }}{PQ}$ 上的一点， ${ABCD}$ 是扇形的内接矩形,记 $\angle {BOC} = \alpha$ .

(1)若 $\alpha  = \frac{\pi }{6}$ ，求矩形 ${ABCD}$ 的周长和面积；

(2)若 ${{Rt}} \bigtriangleup  {CBO}$ 的周长为 $\frac{\sqrt{7}\left( {2\sqrt{10} + 5}\right) }{5}$ ，求 $\tan \alpha$ 的值.

\begin{center}
\includegraphics[width=0.2\textwidth]{images/bo_d7tk7oc91nqc73875fu0_112_1187_1213_253_253_0.jpg}
\end{center}
 

(第 16 题图)

\section*{§ 6.4 三角函数及其性质二}

\section*{典型例题精讲}

【例 1】把正弦函数 $f\left( x\right)  = \cos x$ 的图像向左平移 $\frac{\pi }{6}$ 个单位后,再将得到的图像上各点横坐标伸长到原来的 2 倍,纵坐标不变,得到函数 $y = g\left( x\right)$ 的图像.

(1)求函数 $g\left( x\right)$ 单调递减区间；

(2)若 $x \in  \lbrack 0,{4\pi })$ ，且 $g\left( {x}_{0}\right)  =  - \frac{\sqrt{3}}{2}$ ，求 ${x}_{0}$ 的值.

【例 2】已知函数 $f\left( x\right)  = \sin \left( {{\omega x} + \varphi }\right) \left( {\omega  > 0,\left| \varphi \right| \frac{\pi }{2}}\right)$ 的图像关于 $\left( {-\frac{\pi }{6},0}\right)$ 中心对称，且图像上相邻两个对称轴的距离为 $\frac{\pi }{2}$ .

(1)求函数 $f\left( x\right)$ 的解析式；

(2)设 ${x}_{1}$ 、 ${x}_{2} \in  \left( {-\frac{\pi }{6},\frac{\pi }{3}}\right)$ ，且 ${x}_{1} \neq  {x}_{2}$ ，若 $f\left( {x}_{1}\right)  = f\left( {x}_{2}\right)$ ，求 $f\left( {{x}_{1} + {x}_{2}}\right)$ 的值.

【例 3】若点 $\left( {{x}_{0},{y}_{0}}\right)$ 在函数 $f\left( x\right)$ 的图像上,且满足 ${y}_{0} \cdot  f\left( {y}_{0}\right)  \geq  0$ ,则称 ${x}_{0}$ 是 $f\left( x\right)$ 的 $\xi$ 点. 函数 $f\left( x\right)$ 的所有 $\xi$ 点构成的集合称为 $f\left( x\right)$ 的 $\xi$ 集.

(1)判断 $\frac{2\pi }{3}$ 是否是函数 $f\left( x\right)  = \tan x$ 的 $\xi$ 点,并说明理由;

(2)若函数 $f\left( x\right)  = \sin \left( {\pi x}\right)$ ，求 $f\left( x\right)$ 的 $\xi$ 集.

【例 4】某旅游区每年各个月份接待游客的人数近似地满足周期性规律,因而第 $n$ 个月从事旅游服务工作的人数 $f\left( n\right)$ 可近似地用函数 $f\left( n\right)  = A\cos \left( {{\omega n} + \theta }\right)  + k$ 来刻画,其中正整数 $n$ 表示月份且 $n \in  \left\lbrack  {1,{12}}\right\rbrack$ , 例如 $n = 1$ 表示 1 月份, $A$ 和 $k$ 是正整数, $\omega  > 0,\theta  \in  \left( {0,\pi }\right)$ . 统计发现,该地区每年各个月份从事旅游服务工作的人数有以下规律:

 \ding{172} 每年相同的月份，该地区从事旅游服务工作的人数基本相同；

 \ding{173} 该地区从事旅游服务工作的人数最多的 8 月份和最少的 2 月份相差 400 人；

 \ding{174} 2 月份该地区从事旅游服务工作的人数为 100 人，随后逐月递增直到 8 月份达到最多.

(1)试根据已知信息，求 $f\left( n\right)$ 的表达式；

(2)一般地，当该地区从事旅游服务工作的人数在 400 或 400 以上时，该地区也进入了一年中的旅游“旺季”，那么，一年中的哪几个月是该地区的旅游“旺季”？请说明理由.

【例 5】如图所示是某斜拉式大桥图片, 为了解桥的一些结构情况, 学校数学兴趣小组将大桥的结构进行了简化，取其部分可抽象成右图所示的模型，其中桥塔 ${AB}$ 、 ${CD}$ 与桥面 ${AC}$ 垂直，通过测量得知 ${AB} =  {50}{m},{AC} = {50}{m}$ ,当 $P$ 为 ${AC}$ 中点时, $\angle {PBD} = {45}^{ \circ  }$ .

(1)求 ${CD}$ 的长；

(2)设 ${AP} = x$ ，写出 $\tan \angle {PBD}$ 与 $x$ 的函数关系式；

(3)已知命题:函数 $y = \tan x$ 在 $x \in  \left( {0,\frac{\pi }{2}}\right)$ 内为严格增函数；求证该命题为真命题，并用该命题求解 $P$ 在线段 ${AC}$ 的何处时， $\angle {PBD}$ 达到最大，最大值为多少？

\begin{center}
\includegraphics[width=0.6\textwidth]{images/bo_d7tk7oc91nqc73875fu0_115_691_520_765_271_0.jpg}
\end{center}
 

(例 5 图)

\section*{四基达标精练}

\section*{一、填空题}

1. 将曲线 $y = \sin x$ 上所有点向左平移 $\varphi \left( {\varphi  > 0}\right)$ 个单位，得到函数 $y =  - \sin x$ 的图像，则 $\varphi$ 的最小值为\_\_\_\_\_.

2. 若函数 $f\left( x\right)  = \cos \left( {{2x} + \varphi }\right) \left( {-\pi  \leq  \varphi  < \pi }\right)$ 为奇函数，则 $\varphi  =$ \_\_\_\_\_.

3. 函数 $y = \tan x + x$ 在坐标原点处的切线的斜率为\_\_\_\_\_.

4. 已知 $\cos x = \frac{3}{5}$ ，则 $x$ 取值集合为\_\_\_\_\_. (答案用反余弦表示)

5. 已知函数 $f\left( x\right)  = \sqrt{3}\sin x - \cos x$ ，若将 $f\left( x\right)$ 的图像向左平行移动 $\frac{\pi }{6}$ 个单位长度后得到 $g\left( x\right)$ 的图像， 则 $g\left( x\right)$ 的对称中心为\_\_\_\_\_.

6. 如果函数 $y = \frac{1}{2}\sin {\omega x}$ 在区间 $\left\lbrack  {-\frac{\pi }{8},\frac{\pi }{12}}\right\rbrack$ 上单调递减，那么 $\omega$ 的取值范围是\_\_\_\_\_.

7. 函数 $f\left( x\right)  = 3\cos \left( {{2\omega x} + \frac{\pi }{3}}\right) \left( {\omega  > 0}\right)$ 在 $\left\lbrack  {-\frac{\pi }{3},\frac{5\pi }{6}}\right\rbrack$ 上是减函数,且在 $\left\lbrack  {0,{2\pi }}\right\rbrack$ 上恰好取得一次最小值 -3，则 $\omega$ 的取值范围是\_\_\_\_\_.

8. 方程 $\sin {2x} = \frac{1}{2}\tan x$ 在区间 $\left\lbrack  {0,\pi }\right\rbrack$ 上的所有解的和为\_\_\_\_\_.

9. 直线 $y = a$ 与函数 $f\left( x\right)  = \sin \left( {{\omega x} + \frac{\pi }{3}}\right)  + 1\left( {\omega  > 0}\right)$ 的图像相交,若自左至右的三个相邻交点依次为 $A\text{ 、 }B\text{ 、 }C$ ，且满足 $2\left| {AB}\right|  = \left| {BC}\right|$ ，则实数 $a =$ \_\_\_\_\_.

\section*{二、选择题}

10. 把函数 $y = f\left( x\right)$ 图像上所有点的横坐标缩短到原来的 $\frac{1}{3}$ ，纵坐标不变，再把所得曲线向左平移 $\frac{\pi }{4}$ 个单位长度,得到函数 $y = \sin \left( {x + \frac{\pi }{3}}\right)$ 的图像,则 ( )

A. $f\left( x\right)  = \sin \left( {\frac{x}{3} + \frac{\pi }{12}}\right)$ . B. $f\left( x\right)  = \sin \left( {\frac{x}{3} + \frac{7\pi }{12}}\right)$ . C. $f\left( x\right)  = \sin \left( {{3x} + \frac{\pi }{12}}\right)$ . D. $f\left( x\right)  = \sin \left( {{3x} + \frac{7\pi }{12}}\right)$ .

11. 关于函数 $f\left( x\right)  = \frac{1 + \tan x}{1 - \tan x}$ ,下列说法错误的是 ( )

A. 函数的最小正周期为 $\pi$ . B. 函数图像关于点 $\left( {\frac{\pi }{4},0}\right)$ 对称.

C. 函数图像关于直线 $x = \frac{\pi }{2}$ 对称. D. 函数的定义域为 $\{ x \mid  x \neq  \frac{\pi }{4} + {k\pi },k \in  \mathbf{Z}\}$ .

12. 已知函数 $f\left( x\right)  = \sin \left( {{\omega x} - \frac{\pi }{3}}\right) \left( {\omega  > 0}\right)$ 在 $\left( {\frac{\pi }{2},\frac{3\pi }{2}}\right)$ 上无零点,则 $\omega$ 的取值范围是 ( )

A. $\left( {0,\frac{2}{9}}\right\rbrack   \cup  \left\lbrack  {\frac{8}{9}, + \infty }\right)$ . B. $\left( {0,\frac{2}{9}}\right\rbrack   \cup  \left\lbrack  {\frac{2}{3},\frac{8}{9}}\right\rbrack$ . C. $\left( {0,\frac{2}{9}}\right\rbrack   \cup  \left\lbrack  {\frac{8}{9},1}\right\rbrack$ . D. $\left( {\frac{2}{9},\frac{8}{9}}\right\rbrack   \cup  \lbrack 1, + \infty )$ .

\section*{三、解答题}

13. 已知函数 $f\left( x\right)  = 2\sqrt{3}\sin x\cos x + 2{\cos }^{2}x$ .

(1)求函数 $f\left( x\right)$ 的最小正周期；

(2)现将函数 $f\left( x\right)$ 图像上所有的点的横坐标伸长到原来的 2 倍 (纵坐标不变)，得到函数 $g\left( x\right)$ 的图像， 求 $g\left( x\right)$ 在区间 $\left\lbrack  {0,\frac{\pi }{2}}\right\rbrack$ 上的值域.

14. 已知函数 $f\left( x\right)  = \sin x\cos x - \sqrt{3}{\sin }^{2}x + \frac{\sqrt{3}}{2}$ .

(1)若存在 $x \in  \left\lbrack  {-\frac{\pi }{3},\frac{\pi }{6}}\right\rbrack$ ，使得 $f\left( x\right)  \geq  a$ 成立，则求 $a$ 的取值范围；

(2)将函数 $f\left( x\right)$ 的图像上每个点纵坐标不变，横坐标缩短到原来的 $\frac{1}{2}$ ，得到函数 $g\left( x\right)$ 的图.

15. 某用电器电流 $I\left( {{mA}}\right)$ 随时间 $t\left( {s}\right)$ 的变化的部分图像如图所示， $I$ 和 $t$ 满足下列关系式:

$I\left( t\right)  = A\sin \left( {{\omega t} + \varphi }\right) \left( {A > 0,\omega  > 0,\left| \varphi \right|  < \frac{\pi }{2}}\right) .$

(1)求 $I\left( t\right)  = A\sin \left( {{\omega t} + \varphi }\right)$ 的解析式；

(2)若该用电器核心部件有效工作的电流 $\left| I\right|$ 必须大于 ${150}{{mA}}$ ，则在一个周期内，该用电器核心部件的有效工作时间是多少? (电流的正负表示电流的正反方向)

\begin{center}
\includegraphics[width=0.2\textwidth]{images/bo_d7tk7oc91nqc73875fu0_117_1142_1010_317_233_0.jpg}
\end{center}
 

(第 15 题图)

16. 已知函数 $f\left( x\right)  = 2\sqrt{3}\sin \left( {\frac{\pi }{4} + \frac{x}{2}}\right)  \cdot  \cos \left( {\frac{\pi }{4} + \frac{x}{2}}\right)  - \sin \left( {\pi  + x}\right)$ ,且函数 $y = g\left( x\right)$ 的图像与函数 $y =  f\left( x\right)$ 的图像关于直线 $x = \frac{\pi }{4}$ 对称.

(1)求函数 $g\left( x\right)$ 的解析式；

(2)若存在 $x \in  \left\lbrack  {0,\frac{\pi }{2}}\right)$ ，使等式 ${\left\lbrack  g\left( x\right) \right\rbrack  }^{2} - {mg}\left( x\right)  + 2 = 0$ 成立，求实数 $m$ 的最大值和最小值.

\section*{第 7 章 数列(选择性必修 1)}

\section*{知识梳理}

\section*{一、数列的有关概念}

1.数列的概念:数列是有序排列的一列数. 数列是定义在自然数集或其有限子集上的函数,当自变量由小到大依次取值时对应的一列函数值.

 \ding{172} 数列是特殊的函数:数列的学习是函数学习的延续; 可以借助函数的知识和方法研究数列;

 \ding{173} 数列是离散型函数:数列的图像是一系列孤立的点；

 \ding{174} 数列与函数比较更具有特殊性，数列的学习更多关注的是其特殊性.

2. 数列的单调性:

数列 $\left\{  {a}_{n}\right\}$ 是递增数列: 对任意 $n \in  {\mathbf{N}}^{ * },{a}_{n + 1} > {a}_{n}$ 成立;

数列 $\left\{  {a}_{n}\right\}$ 是递减数列: 对任意 $n \in  {\mathbf{N}}^{ * },{a}_{n + 1} < {a}_{n}$ 成立.

 \ding{172} 数列的单调性定义是判断数列单调性的基本依据.

 \ding{173} 数列的单调性与函数的单调性不同.

一般地，若函数 $f\left( x\right) \left( {x > 0}\right)$ 单调增 (减)，则数列 $\left\{  {f\left( n\right) }\right\}$ 单调增 (减)；

但,若数列 $\left\{  {f\left( n\right) }\right\}$ 单调增 (减),不一定函数 $f\left( x\right) \left( {x > 0}\right)$ 单调增 (减).

 \ding{174} 数列的单调性与数列的项的最值密切相关.

3. 求数列最大项或最小项的方法:

用不等式 ${a}_{n} \leq  {a}_{n + 1}$ (或 ${a}_{n} \geq  {a}_{n + 1}$ ) 寻找使得 $\left\{  {a}_{n}\right\}$ 的单调递增和单调递减的 $n$ 的范围确定数列的最大项或最小项.

注意: \ding{172}  通常当 ${a}_{n}$ 由等差数列和等比数列相应项的和(或差)构成时利用 ${a}_{n + 1} - {a}_{n} \geq  0$ ；

 \ding{173}  当 ${a}_{n}$ 由等差数列和等比数列相应项的积(或商)构成时利用 $\frac{{a}_{n + 1}}{{a}_{n}} \geq  1$ (注意各项为正)；

 \ding{174}  当数列 $\left\{  {a}_{n}\right\}$ 的各项中有正有负时,常先考虑其各项的绝对值的最大项,再还原.

\section*{二、等差数列与等比数列}

\begin{center}
\begin{tabular}{|c|c|c|}
\hline
\phantom{X} & 等差数列 & 等比数列 \\
\cline{1-3}
定义 & $\left\{  \begin{array}{l} {a}_{1}, \\  {a}_{n + 1} - {a}_{n} = d \end{array}\right.$ (常数) & $\left\{  \begin{array}{l} {a}_{1} \neq  0, \\  \frac{{a}_{n + 1}}{{a}_{n}} = q\text{ (非零常数) } \end{array}\right.$ \\
\cline{1-3}
通项公式 & ${a}_{n} = {a}_{1} + \left( {n - 1}\right) d = {a}_{k} + \left( {n - k}\right) d$ & ${a}_{n} = {a}_{1}{q}^{n - 1} = {a}_{k}{q}^{n - k}$ \\
\cline{1-3}
中项 & $A = \frac{a + b}{2}\left( {a,b \in  \mathbf{R},A\text{ 唯一 }}\right)$ & ${G}^{2} = {ab}\left( {a\text{ 、 }b \in  \mathbf{R},{ab} > 0,G\text{ 不唯一) }}\right)$ \\
\cline{1-3}
前 $n$ 项和 & ${S}_{n} = \frac{n\left( {{a}_{1} + {a}_{n}}\right) }{2} = \frac{n\left( {{a}_{k} + {a}_{n - k + 1}}\right) }{2}  = {a}_{1}n + \frac{n\left( {n - 1}\right) d}{2} = {a}_{n}n - \frac{n\left( {n - 1}\right) d}{2}$ & ${S}_{n} = \left\{  \begin{array}{l} n{a}_{1},q = 1, \\  \frac{{a}_{1}\left( {1 - {q}^{n}}\right) }{1 - q} = \frac{{a}_{1} - {a}_{n}q}{1 - q},q \neq  1 \end{array}\right.$ \\
\cline{1-3}
\phantom{X} & \phantom{X} & 各项和: $\mathop{\lim }\limits_{{n \rightarrow   + \infty }}{S}_{n} = \mathop{\sum }\limits_{{i = 1}}^{{+\infty }}{a}_{i} = \frac{{a}_{1}}{1 - q}\left( {0 < \left| q\right|  < 1}\right)$ \\
\cline{1-3}
\hline
\end{tabular}
\end{center}

\HRule

* 注:本章节所属知识模块是 “函数与分析”.

\HRule

1. 若 $m + n = k + l\left( {m\text{ 、 }n\text{ 、 }k\text{ 、 }l \in  {\mathbf{N}}^{ * }}\right)$ ,则在等差数列中: ${a}_{m} + {a}_{n} = {a}_{k} + {a}_{l}$ ; 在等比数列中: ${a}_{m}{a}_{n} = {a}_{k}{a}_{l}$ . 注意:以上两个性质可推广到若干个，但等式两端的项数必须相等，脚标之和必须相等.

2. $A = \frac{a + b}{2}$ 是 $a,A,b$ 成等差数列的充要条件.

3. ${G}^{2} = {ab}$ 是 $a,G,b$ 成等比数列的必要非充分条件.

4. 用等比数列求和公式求数列的和时,勿忘 $q \neq  1$ .

5. 分析等差数列的前 $n$ 项和,有 ${S}_{{2k} - 1} = \left( {{2k} - 1}\right) {a}_{k}$ .

6. 等差数列 $\left\{  {a}_{n}\right\}$ 的前 $n$ 项和,次 $n$ 项和,再 $n$ 项和, $\cdots$ ,第 $k$ 个 $n$ 项和成等差数列,公差为 ${n}^{2}d$ .

7. 等比数列 $\left\{  {a}_{n}\right\}$ 的前 $n$ 项和,次 $n$ 项和,再 $n$ 项和, $\cdots$ ,第 $k$ 个 $n$ 项和成等比数列,公比 ${q}^{n}$ .

8. 已知等差数列的前 $n$ 项和为 ${S}_{n}$ ,则 $\frac{{S}_{n}}{n} = {dn} + \left( {{a}_{1} - \frac{d}{2}}\right)$ ,表明 $\left\{  \frac{{S}_{n}}{n}\right\}$ 也是等差数列; 即点 $\left( {n,\frac{{S}_{n}}{n}}\right) \left( {n \in  {\mathbf{N}}^{ * }}\right)$ 共线 (也就是当 $n$ 取不同值时,任意两点连线的斜率相等).

9. 项数有限的等差数列奇数项的和 ${S}_{\text{ 奇 }}$ ,偶数项的和 ${S}_{\text{ 偶 }}$ 满足 ${S}_{\text{ 奇 }} + {S}_{\text{ 偶 }} = {S}_{n}$ ; 当 $n$ 为偶数时, ${S}_{\text{ 奇 }} - {S}_{\text{ 偶 }} =  - \frac{nd}{2}$ ; 当 $n$ 为奇数时, ${S}_{\text{ 奇 }} - {S}_{\text{ 偶 }} = \frac{{S}_{n}}{n} = {a}_{\frac{n + 1}{2}}$ .

10. 运算性质:

 \ding{172}  若数列 $\left\{  {a}_{n}\right\}$ 、 $\left\{  {b}_{n}\right\}$ 都是等差数列，则 $\left\{  {A \cdot  {a}_{n} - B \cdot  {b}_{n}}\right\}$ ( $A$ 、 $B$ 为常数)也是等差数列；

 \ding{173} 若数列 $\left\{  {a}_{n}\right\}$ 、 $\left\{  {b}_{n}\right\}$ 都是等比数列，则 $\left\{  {{a}_{n}^{A} \cdot  {b}_{n}^{B}}\right\}$ ( $A$ 、 $B$ 为常数) 也是等比数列；

 \ding{174} 若数列 $\left\{  {a}_{n}\right\}$ 是等差数列， ${b}_{n} = {a}^{{a}_{n}}\left( {a > 0,a \neq  1}\right)$ ，则 $\left\{  {b}_{n}\right\}$ 是等比数列；

 \ding{175} 若数列 $\left\{  {a}_{n}\right\}  \left( {{a}_{n} > 0}\right)$ 是等比数列， ${b}_{n} = {\log }_{a}{a}_{n}\left( {a > 0,a \neq  1}\right)$ ，则 $\left\{  {b}_{n}\right\}$ 是等差数列.

\begin{center}
\begin{tabular}{|c|c|}
\hline
证明数列 $\left\{  {a}_{n}\right\}$ 为等差数列的方法 & 证明数列 $\left\{  {a}_{n}\right\}$ 为等比数列的方法 (注意 ${a}_{1} \neq  0$ ) \\
\cline{1-2}
\makecell{ 据定义: ${a}_{n + 1} - {a}_{n} = d$ ( $d$ 为常数)； \\ 根据通项公式: ${a}_{n} = {a}_{1} + \left( {n - 1}\right) d$ ; \\ 根据性质: $2{a}_{n + 1} = {a}_{n} + {a}_{n + 2}$ } & \makecell{  \ding{172} 据定义: ${a}_{n + 1} = q{a}_{n}$ ( $q$ 为非零常数)； \\  \ding{173}  根据通项公式: ${a}_{n} = {a}_{1}{q}^{n - 1}$ ； \\  \ding{174}  根据性质: ${a}_{n + 1}^{2} = {a}_{n}{a}_{n + 2}$ } \\
\cline{1-2}
\hline
\end{tabular}
\end{center}

11. 设 ${S}_{n}$ 是数列 $\left\{  {a}_{n}\right\}$ 的前 $n$ 项和, $\left\{  {a}_{n}\right\}$ 为等差数列的充要条件是 ${S}_{n} = a{n}^{2} + {bn}\left( {a,b\text{ 为常数 }}\right)$ ,其公差是 ${2a};\left\{  {a}_{n}\right\}$ 为等比数列的充要条件是 ${S}_{n} = a{q}^{n} - a\left( {a\text{ 、 }q\left( {q \neq  1}\right) }\right.$ 为非零常数 $)$ ,其公比是 $q$ .

\section*{三、数列的通项公式}

1. 已知前 $n$ 项和 ${S}_{n}$ ,则 ${a}_{n} = \left\{  \begin{array}{l} {S}_{1},n = 1, \\  {S}_{n} - {S}_{n - 1},n \geq  2, \end{array}\right.$ 只有当 $n = 1$ 适合 $n \geq  2$ 的表达式时才能只用第二个式子表示.

2. 一阶线性递推数列的通项公式: $\left\{  \begin{array}{l} {a}_{1}, \\  {a}_{n + 1} = p{a}_{n} + q\text{ ( }p,q\text{ 为满足条件的常数). } \end{array}\right.$

求法一(迭代法):

${a}_{n} = p{a}_{n - 1} + q = p\left( {p{a}_{n - 2} + q}\right)  + q$

$= {p}^{2}{a}_{n - 2} + {qp} + q = {p}^{2}\left( {p{a}_{n - 3} + q}\right)  + {qp} + q$

$= {p}^{3}{a}_{n - 3} + q{p}^{2} + {qp} + q = \cdots$

$= {p}^{n - 1}{a}_{1} + q{p}^{n - 2} + \cdots  + q{p}^{2} + {qp} + q$

$= {p}^{n - 1}{a}_{1} + \frac{q\left( {1 - {p}^{n - 1}}\right) }{1 - p}$

求法二(构造法 \ding{172} ):

当 $n \geq  2$ 时,由 $\left\{  \begin{array}{l} {a}_{n - 1} = p{a}_{n} + q, \\  {a}_{n} = p{a}_{n - 1} + q \end{array}\right.$ 得, ${a}_{n + 1} - {a}_{n} = p\left( {{a}_{n} - {a}_{n - 1}}\right) ,{a}_{2} = p{a}_{1} + q$ , $\therefore \left\{  {{a}_{n + 1} - {a}_{n}}\right\}$ 是以 ${a}_{2} - {a}_{1} = \left( {p - 1}\right) {a}_{1} + q$ 为首项, $p$ 为公比的等比数列,

$\therefore {a}_{n + 1} - {a}_{n} = \left\lbrack  {\left( {p - 1}\right) {a}_{1} + q}\right\rbrack  {p}^{n - 1}$ ,于是

$\left. \begin{array}{l} {a}_{n} - {a}_{n - 1} = \left\lbrack  {\left( {p - 1}\right) {a}_{1} + q}\right\rbrack  {p}^{n - 2}, \\  {a}_{n - 1} - {a}_{n - 2} = \left\lbrack  {\left( {p - 1}\right) {a}_{1} + q}\right\rbrack  {p}^{n - 3}, \\  \cdots \\  {a}_{2} - {a}_{1} = \left( {p - 1}\right) {a}_{1} + q, \\  {a}_{1} = {a}_{1}, \end{array}\right\}$ 相加得 ${a}_{n} = {a}_{1} + \frac{\left\lbrack  {\left( {p - 1}\right) {a}_{1} + q}\right\rbrack  \left( {1 - {p}^{n - 1}}\right\rbrack  }{1 - p}$ .

求法三(构造法 \ding{173} )

设 ${a}_{n + 1} + x = p\left( {{a}_{n} + x}\right)$ ,则 ${a}_{n + 1} = p{a}_{n} + \left( {p - 1}\right) x$ ,于是 $x = \frac{q}{p - 1}$ .

$\therefore {a}_{n + 1} + \frac{q}{p - 1} = p\left( {{a}_{n} + \frac{q}{p - 1}}\right)$ ,

即数列 $\left\{  {{a}_{n} + \frac{q}{p - 1}}\right\}$ 是以 ${a}_{1} + \frac{q}{p - 1}$ 为首项, $p$ 为公比的等比数列,

$\therefore {a}_{n} + \frac{q}{p - 1} = \left( {{a}_{1} + \frac{q}{p - 1}}\right) {p}^{n - 1}$ ,即 ${a}_{n} = \left( {{a}_{1} + \frac{q}{p - 1}}\right) {p}^{n - 1} - \frac{q}{p - 1}$ .

3. 一阶线性递推数列的特殊情形:

 \ding{172} 当 $p = 1$ 时，即为等差数列；

 \ding{173} 当 $p \neq  0,q = 0$ ，且 $p \neq  1$ 时，即为等比数列.

\section*{四、数列求和}

1. (类似于)等差数列前 $n$ 项和公式的推导方法中的“倒序相加法”，适用于与函数图像的对称性有关的若干项求和问题.

2. (类似于) 等比数列前 $n$ 项和公式的推导方法中的“错位相减法”,适用于等差数列与等比数列相应项相乘构成的数列求和问题.

3. “若 ${c}_{n} = {a}_{n}{b}_{n}$ ,其中 $\left\{  {a}_{n}\right\}$ 是等差数列, $\left\{  {b}_{n}\right\}$ 是等比数列,求 $\left\{  {c}_{n}\right\}$ 的前 $n$ 项”用“错位相减”法.

4. “若 ${b}_{n} = \frac{1}{{a}_{n}{a}_{n + 1}}$ ，其中 $\left\{  {a}_{n}\right\}$ 是等差数列，求 $\left\{  {b}_{n}\right\}$ 的前 $n$ 项的和”用“裂项相消法”(其中常用待定系数法进行裂项,方法如下: 设 ${b}_{n} = \frac{A}{{a}_{n}} + \frac{B}{{a}_{n + 1}}$ ,则 ${b}_{n} = \frac{A{a}_{n + 1} + B{a}_{n}}{{a}_{n}{a}_{n + 1}} = \frac{\left( {A + B}\right) {a}_{n} + {Ad}}{{a}_{n}{a}_{n + 1}}$ ,于是 $\left\{  {\begin{array}{l} A + B = 0, \\  {Ad} = 1, \end{array}A =  - B = }\right.  \frac{1}{d}$ ,即 ${b}_{n} = \frac{1}{d}\left( {\frac{1}{{a}_{n}} - \frac{1}{{a}_{n + 1}}}\right)$ .

注意: \ding{172}  “若 ${b}_{n} = \frac{1}{{a}_{n}{a}_{n + 1}{a}_{n + 2}}$ ，其中 $\left\{  {a}_{n}\right\}$ 是等差数列，求 $\left\{  {b}_{n}\right\}$ 的前 $n$ 项的和”也可用 “裂项相消法”；  \ding{173}  并非只有这种类型的求和问题才能用 “裂项相消法”，要根据具体题目的特征选用，如: $n \cdot  n! = \left( {n + 1}\right) !  - n$ ! 等; 又如 ${a}_{n} = \frac{{2}^{n}}{\left( {{2}^{n - 1} + 1}\right) \left( {{2}^{n} + 1}\right) } = 2\left( {\frac{1}{{2}^{n - 1} + 1} - \frac{1}{{2}^{n} + 1}}\right)$ 等.

5. ”若 ${c}_{n} = A{a}_{n} + B{b}_{n}\left( {A\text{ 、 }B\text{ 为常数 }}\right)$ ，其中 $\left\{  {a}_{n}\right\}$ 、 $\left\{  {b}_{n}\right\}$ 的前 $n$ 项和分别已知或可求，求 $\left\{  {c}_{n}\right\}$ 的前 $n$ 项的和”用“分组求和法”.

6. 分段数列求和要注意分段的要求,分类讨论要到位.

\section*{五、数列的极限与数学归纳法}

1. 数列 $\left\{  {a}_{n}\right\}$ 的极限是 $A$ ,说明随着 $n$ 的无限增大, ${a}_{n}$ 与 $A$ 越来越近,而且无限接近,该极限与数列的前有限项无关, 只与后面的所有项有关.

2.几个常用的数列极限:

 \ding{172}  $\mathop{\lim }\limits_{{n \rightarrow  \infty }}C = C\left( {C\text{ 为常数 }}\right)$ ;  \ding{173}  $\mathop{\lim }\limits_{{n \rightarrow  \infty }}\frac{1}{n} = 0$ ;  \ding{174}  $\mathop{\lim }\limits_{{n \rightarrow  \infty }}{q}^{n} = \left\{  \begin{array}{l} 0,\left| q\right|  < 1, \\  1,q = 1, \\  \text{ 不存在, }q =  - 1\text{ 或 }\left| q\right|  > 1 \end{array}\right.$

3. 若无穷等比数列的和为 $a$ ,由 $\frac{{a}_{1}}{1 - q} = a$ 和 $0 < \left| q\right|  < 1$ 可求 ${a}_{1}$ (或 $q$ )的范围,本质是由定义域求值域.

4. 数学归纳法证明与自然数有关问题的步骤:以下斜体字在证题过程中照抄，“······”是证明过程

 \ding{172}  (验证)当 $n = {n}_{0}$ 时， $\cdots \cdots$ 结论成立；

 \ding{173}  假设 $n = k\left( {k \geq  {n}_{0},n \in  \mathbf{N}}\right)$ 时结论成立，即······

则当 $n = k + 1$ 时， $\cdots \cdots$

即当 $n = k + 1$ 时，结论成立

由 \ding{172}  \ding{173} 可知，结论对任意 $n \geq  {n}_{0},n \in  \mathbf{N}$ 成立.

5. 数学归纳法中,  \ding{172} 是归纳基础;  \ding{173} 是归纳递推, 在第二步中必须使用归纳假设, 否则就不是数学归纳法.

6. 数学归纳法可用于证明一部分与正整数 $n$ 有关的等式、不等式、整除问题和几何问题.

\section*{$\text{ § }{7.1}$ 数列的概念与性质}

\section*{典型例题精讲}

【例 1】根据下列条件,求数列 $\left\{  {a}_{n}\right\}$ 的通项公式.

(1) ${a}_{1} = 2,{a}_{n + 1} = {a}_{n} + n + 1$ ；

(2) ${a}_{1} = \frac{2}{3},{a}_{n + 1} = \frac{n}{n + 1}{a}_{n}$ .

【例 2】已知数列 $\left\{  {a}_{n}\right\}$ 满足 ${a}_{1} = 1,{a}_{n} = 3{a}_{n - 1} + 4\left( {n \geq  2}\right)$ .

(1)求证:数列 $\left\{  {a}_{n}\right\}$ 既不是等差数列，也不是等比数列；

(2)求证:数列 $\left\{  {{a}_{n} + 2}\right\}$ 是等比数列,并求数列 $\left\{  {a}_{n}\right\}$ 的通项公式.

【例 3】已知数列 $\left\{  {a}_{n}\right\}$ 满足 ${a}_{1} = 1,2{a}_{n + 1} = {a}_{n}$ ,数列 $\left\{  {b}_{n}\right\}$ 满足 ${b}_{n} = 2 - {\log }_{2}{a}_{n + 1}^{2}$ .

(1)求数列 $\left\{  {a}_{n}\right\}$ 、 $\left\{  {b}_{n}\right\}$ 的通项公式；

(2)设数列 $\left\{  {b}_{n}\right\}$ 的前 $n$ 项和为 ${T}_{n}$ ，求使得 $2{T}_{n} \leq  4{n}^{2} + m$ 对任意正整数 $n$ 都成立的实数 $m$ 的取值范围.

【例 4】已知数列 $\left\{  {a}_{n}\right\}$ 满足 ${a}_{1} + 3{a}_{2} + 5{a}_{3} + \cdots  + \left( {{2n} - 1}\right) {a}_{n} = {3}^{n}$ .

(1)求数列 $\left\{  {a}_{n}\right\}$ 的通项公式；

(2)若对任意的 $n \in  {\mathbf{N}}^{ * }$ ，不等式 ${a}_{n} \geq  {\left( -1\right) }^{n}\lambda$ 恒成立，求实数 $\lambda$ 的取值范围.

【例 5】已知 $\left\{  {a}_{n}\right\}$ 为等差数列, ${b}_{n} = \left\{  \begin{array}{l} {a}_{n} - 6,n\text{ 为奇数,记 }{S}_{n}\text{ 、 }{T}_{n}\text{ 为 }\left\{  {a}_{n}\right\}  \text{ 和 }\left\{  {b}_{n}\right\}  \text{ 的前 }n\text{ 项和, }{S}_{4} = {32},{T}_{3} = {16}. \end{array}\right.$

(1)求 $\left\{  {a}_{n}\right\}$ 的通项公式；

(2)证明:当 $n > 5$ 时， ${T}_{n} > {S}_{n}$ .

\section*{四基达标精练}

\section*{一、填空题}

1. 在数列 $1,2,\sqrt{7},\sqrt{10},\sqrt{13},\cdots$ 中， $2\sqrt{19}$ 是这个数列的第\_\_\_\_\_项.

2. 在数列 $\left\{  {a}_{n}\right\}$ 中， ${a}_{1} = 2$ ，且 ${a}_{n} = {a}_{n - 1} + \lg \frac{n}{n - 1}\left( {n \geq  2}\right)$ ，则 ${a}_{100} =$ \_\_\_\_\_.

3. 已知数列 $\left\{  {a}_{n}\right\}$ 满足: ${a}_{1} = 1,{a}_{n + 1} = {\left( \sqrt{{a}_{n}} + 1\right) }^{2}$ ,则 ${a}_{5} =$ \_\_\_\_\_.

4. 已知数列 $\left\{  {a}_{n}\right\}$ 的通项公式是 ${a}_{n} = {100} - {dn}$ ,若该数列从第 10 项起以后各项均为负值,则实数 $d$ 的取值范围是\_\_\_\_\_.

5. 数列 $\left\{  {a}_{n}\right\}$ 满足 $\frac{1}{3}{a}_{1} + \frac{1}{{3}^{2}}{a}_{2} + \cdots  + \frac{1}{{3}^{n}}{a}_{n} = {3n} + 1\left( {n \in  {\mathbf{N}}^{ * }}\right)$ ，则 ${a}_{n} =$ \_\_\_\_\_.

6. 已知数列 $\left\{  {a}_{n}\right\}$ 满足 ${a}_{1} = {33},{a}_{n + 1} - {a}_{n} = {2n}$ ，则 $\frac{{a}_{n}}{n}$ 的最小值为\_\_\_\_\_.

7. 数列 $\left\{  {a}_{n}\right\}$ 满足:点 $\left( {n,{a}_{n}}\right)$ 都在抛物线 $y = {x}^{2} + {kx} + {156}$ 上，且 ${a}_{n + 1} > {a}_{n}$ 对任意的正整数 $n$ 恒成立，则实数 $k$ 的范围是\_\_\_\_\_.

8. 设数列 $\left\{  {a}_{n}\right\}$ 的通项公式为 ${a}_{n} = {n}^{2} + {bn}$ ，若数列 $\left\{  {a}_{n}\right\}$ 是单调递增数列，则实数 $b$ 的取值范围为\_\_\_\_\_.

9. 已知数列 $\left\{  {a}_{n}\right\}$ 是单调数列,其前 $n$ 项和为 ${S}_{n}$ ,若 ${S}_{n} + {S}_{n + 1} = {n}^{2} + n$ ,则首项 ${a}_{1}$ 的取值范围是\_\_\_\_\_.

\section*{二、选择题}

10. 已知数列 $\left\{  {a}_{n}\right\}$ 共有 10 项, ${a}_{n} = \frac{{2022} - {10}^{n}}{{2023} - {10}^{n}}$ ,则其中最大项、最小项的情况是 ( )

A. 最大项为 ${a}_{1}$ 、最小项为 ${a}_{10}$ . B. 最大项为 ${a}_{10}$ 、最小项为 ${a}_{1}$ .

C. 最大项为 ${a}_{6}$ 、最小项为 ${a}_{5}$ . D. 最大项为 ${a}_{4}$ 、最小项为 ${a}_{3}$ .

11. 已知数列 $\left\{  {a}_{n}\right\}$ 满足 ${a}_{n} = \left\{  {\begin{array}{l}  - {n}^{2} + {2tn},n \leq  5, \\  \left( {t - 1}\right) n,n > 5 \end{array}\left( {n \in  {\mathbf{N}}^{ * }}\right) }\right.$ ,且数列 $\left\{  {a}_{n}\right\}$ 是严格递增数列,则实数 $t$ 的取值范围是 ( )

A. $\left( {\frac{9}{2},\frac{19}{4}}\right)$ . B. $\left( {\frac{9}{2}, + \infty }\right)$ . C. $\left( {5, + \infty }\right)$ . D. $(1,4\rbrack$ .

12. 数列的发展史, 折射出许多有价值的数学思想方法, 对时代的进步起了重要作用, 比如意大利数学家列昂纳多·斐波拉契以兔子繁殖为例，引入“兔子数列”:即 $1,1,2,3,5,8,{13},{21},{34},{55},{89},{144},{233},\cdots$ ,也即 $F\left( 1\right)  = F\left( 2\right)  = 1,F\left( n\right)  = F\left( {n - 1}\right)  + F\left( {n - 2}\right) \left( {n \geq  3,n \in  \mathbf{N}}\right)$ ,若此数列被 4 整除后的余数构成一个新的数列 $\left\{  {b}_{n}\right\}$ ,则 ${b}_{1} + {b}_{2} + {b}_{3} + \cdots  + {b}_{2023}$ 等于 ( )

A. 2699. B. 3539. C. 2027. D. 2023.

\section*{三、解答题}

13. 已知数列 $\left\{  {a}_{n}\right\}$ 满足 ${a}_{1} = 1,{a}_{2} = \frac{1}{3}$ ,且 $2{a}_{n + 1}{a}_{n}{a}_{n - 1} = {a}_{n + 1}{a}_{n} + {a}_{n}{a}_{n - 1} - 2{a}_{n + 1}{a}_{n - 1}\left( {n \geq  2}\right)$ .

(1)求证: $\left\{  {\frac{1}{{a}_{n + 1}} - \frac{1}{{a}_{n}}}\right\}$ 是等差数列；

(2)求数列 $\left\{  {a}_{n}\right\}$ 的通项公式.

14. 已知数列 $\left\{  {a}_{n}\right\}$ 满足 ${a}_{1} = \frac{5}{6}$ ,若对于任意 $n \in  \mathbf{N},n \geq  2$ ,二次方程 ${a}_{n - 1}{x}^{2} - {a}_{n}x + 1 = 0$ 都有根 $\alpha$ 、 $\beta$ ,且满足 ${3\alpha } - {\alpha \beta } + {3\beta } = 1$ .

(1)求证: $\left\{  {{a}_{n} - \frac{1}{2}}\right\}$ 是等比数列；

(2)求数列 $\left\{  {a}_{n}\right\}$ 的前 $n$ 项和 ${S}_{n}$ .

15. 已知正项等比数列 $\left\{  {a}_{n}\right\}$ 满足 ${a}_{3} = 4,{a}_{1} + {a}_{5} = {17}$ ,公比 $q > 1$ .

(1)求数列 $\left\{  {a}_{n}\right\}$ 的通项公式；

(2)若 ${c}_{n} = \frac{{n}^{2} \cdot  {a}_{n + 1}}{{3}^{n}}$ ，判断数列 $\left\{  {c}_{n}\right\}$ 有没有最大项? 若有，求出第几项为最大项；若没有，请说明理由.

16. 流行性感冒是由流感病毒引起的急性呼吸道传染病. 某市去年 11 月曾发生流感，据统计，11 月 1 日该市的新感染者有 30 人，以后每天的新感染者比前一天的新感染者增加 50 人. 由于该市医疗部门采取措施，使该种病毒的传播得到控制，从 11 月 $k + 1\left( {9 \leq  k \leq  {29},k \in  \mathbf{N}}\right)$ 日起每天的新感染者比前一天的新感染者减少 20 人.

(1)若 $k = 9$ ，求 11 月 1 日至 11 月 10 日新感染者总人数；

(2)若到 11 月 30 日止，该市在这 30 天内的新感染者总人数为 11940 人，问 11 月中哪一天，该市新感染者人数最多? 并求这一天的新感染者人数.

\section*{§ 7.2 等差数列和等比数列一}

\section*{典型例题精讲}

【例 1】等差数列 $\left\{  {a}_{n}\right\}$ 满足 ${a}_{3} + {a}_{5} = {10},{a}_{2} = 9$ .

(1)求 $\left\{  {a}_{n}\right\}$ 的通项公式；

(2)若 $\left\{  {a}_{n}\right\}$ 的前 $n$ 项和为 ${S}_{n}$ ，求 ${S}_{n}$ 的最大值.

【例 2】记 ${S}_{n}$ 为数列 $\left\{  {a}_{n}\right\}$ 的前 $n$ 项和, ${b}_{n}$ 为数列 $\left\{  {S}_{n}\right\}$ 的前 $n$ 项积,已知 $\frac{2}{{S}_{n}} + \frac{1}{{b}_{n}} = 2$ .

(1)证明:数列 $\left\{  {b}_{n}\right\}$ 是等差数列；

(2)求 $\left\{  {a}_{n}\right\}$ 的通项公式.

【例 3】已知函数 $f\left( x\right)  = \frac{1}{2}\left( {{a}_{n + 2} - {a}_{n + 1}}\right) {x}^{2} - \left( {3{a}_{n + 1} - 4{a}_{n}}\right) x\left( {n \in  {\mathbf{N}}^{ * }}\right)$ 的对称轴是 $x = 1$ ,数列 $\left\{  {a}_{n}\right\}$ 满足, ${a}_{1} = 2,{a}_{2} = 8$ .

(1)证明数列 $\left\{  {{a}_{n + 1} - 2{a}_{n}}\right\}$ 是等比数列，并求数列 $\left\{  {a}_{n}\right\}$ 的通项公式；

(2)设 ${3}^{n}{b}_{n} = {\left( -1\right) }^{n}{a}_{n}$ ，且 $\left| {b}_{1}\right|  + \left| {b}_{2}\right|  + \cdots  + \left| {b}_{n}\right|  < m - {3n}{\left( \frac{2}{3}\right) }^{n + 1}$ 对于 $n \in  {\mathbf{N}}^{ * }$ 恒成立，求实数 $m$ 的取值范围.

【例 4】已知数列 $\left\{  {a}_{n}\right\}$ 的前 $n$ 项和为 ${S}_{n}$ ,若 ${a}_{2} = 4,\frac{{S}_{n} - n}{n} = \frac{1}{2}{a}_{n}$ .

(1)求证:数列 $\left\{  {a}_{n}\right\}$ 是等差数列；

(2)若 ${b}_{n} = \left| {{a}_{n} - {11}}\right|$ ，求数列 $\left\{  {b}_{n}\right\}$ 的前 $n$ 项的和 ${T}_{n}$ .

【例 5】已知正项数列 $\left\{  {a}_{n}\right\}$ 的前 $n$ 项和为 ${S}_{n}$ ,数列 $\left\{  {b}_{n}\right\}$ 为等比数列,满足 $4{S}_{n} = {a}_{n + 1}^{2} - {4n} - 4$ ,且 ${a}_{1} =  {b}_{1} + 1 = 2,{a}_{4} = {b}_{4}.$

(1)求证:数列 $\left\{  {a}_{n}\right\}$ 为等差数列；

(2)若从数列 $\left\{  {a}_{n}\right\}$ 中去掉数列 $\left\{  {b}_{n}\right\}$ 的项后余下的项按原来的顺序组成数列 $\left\{  {c}_{n}\right\}$ ，求 ${c}_{1} + {c}_{2} + {c}_{3} + \cdots  + {c}_{100}$ .

\section*{四基达标精练}

\section*{一、填空题}

1. 已知等差数列 $\left\{  {a}_{n}\right\}$ 的前 $n$ 项和为 ${S}_{n}$ ，且 ${a}_{2} = 3,{S}_{5} = {25}$ ，则公差 $d =$ \_\_\_\_\_.

2. 已知各项均为正数的等比数列 $\left\{  {a}_{n}\right\}$ 的前 4 项和为 15，并且 $4{a}_{1},2{a}_{3},{a}_{5}$ 成等差数列，则 ${a}_{1}  =$

3. 记 ${S}_{n}$ 为等差数列 $\left\{  {a}_{n}\right\}$ 的前 $n$ 项和,已知 ${S}_{5} = 0,{a}_{6} = 6$ ,则 ${S}_{n} =$ \_\_\_\_\_.

4. 从等差数列 ${84},{80},{76},\cdots$ 的第\_\_\_\_\_项开始,以后各项均为负值.

5. 已知 $\left\{  {a}_{n}\right\}$ 为等比数列， ${S}_{n}$ 为其前 $n$ 项和，若 $2{S}_{3} = {a}_{2} + {a}_{3} + {a}_{4}$ ，则公比 $q =$ \_\_\_\_\_.

6. 已知等比数列 $\left\{  {a}_{n}\right\}$ 中， ${a}_{2} = 2$ ， ${a}_{5} = \frac{1}{4}$ ，则满足 ${a}_{1}{a}_{2} + {a}_{2}{a}_{3} + \cdots  + {a}_{n}{a}_{n + 1} \leq  \frac{21}{2}$ 成立的最大正整数 $n$ 的值为\_\_\_\_\_.

7. 设 $f\left( x\right)  = \frac{{4}^{x}}{{4}^{x} + 2}$ ,则 $f\left( \frac{1}{11}\right)  + f\left( \frac{2}{11}\right)  + f\left( \frac{3}{11}\right)  + \cdots  + f\left( \frac{10}{11}\right)  =$ \_\_\_\_\_.

8. 等差数列 $\left\{  {a}_{n}\right\}$ 的公差 $d \neq  0$ ，其前 $n$ 项和为 ${S}_{n}$ . 若 ${S}_{10} = 0$ ，则 ${S}_{i}\left( {i = 1,2,\cdots ,{2003}}\right)$ 中不同的数值有\_\_\_\_\_个.

9. 已知数列 $\left\{  {a}_{n}\right\}$ 满足, ${a}_{1} = 1$ ,且对于任意正整数 $m\text{ 、 }n$ 都有 ${a}_{m + n} = {a}_{m} + {a}_{n} + {mn}$ 成立,则 $\frac{1}{{a}_{1}} + \frac{1}{{a}_{2}} + \cdots  +  \frac{1}{{a}_{2023}} =$ \_\_\_\_\_.

\section*{二、选择题}

10. 已知等差数列 $\left\{  {a}_{n}\right\}$ 的公差为 $d,{S}_{20} = {10M}$ ,则下列正确的是 ( )

A. $M = 2{a}_{5} + {a}_{10}$ . B. $M = {a}_{1} + 2{a}_{10}$ . C. $M = {a}_{5} + {a}_{15}$ . D. $M = 2{a}_{10} + d$ .

11. 某林厂现在的森林木材存量是 1800 万 ${{m}}^{3}$ ，木材以每年 25\% 的增长率生长，而每年要砍伐固定的木材量为 $x$ 万 ${{m}}^{3}$ ，为达到经两次砍伐后木材存量增加 50\% 的目标，则 $x$ 的值是 ( )

A. 40. B. 45. C. 50 . D. 55.

12. 数列 $\left\{  {a}_{n}\right\}$ 的前 $n$ 项和为 ${S}_{n},{a}_{1} = m$ ,且对任意的 $n \in  {\mathbf{N}}^{ * }$ 都有 ${a}_{n} + {a}_{n + 1} = {2n} + 1$ ,则下列三个命题中, 所有真命题的序号是 ( )

 \ding{172}  存在实数 $m$ ，使得 $\left\{  {a}_{n}\right\}$ 为等差数列； \ding{173}  存在实数 $m$ ，使得 $\left\{  {a}_{n}\right\}$ 为等比数列；

 \ding{174} 若存在 $k \in  {\mathbf{N}}^{ * }$ ，使得 ${S}_{k} = {S}_{k + 1} = {55}$ ，则实数 $m$ 唯一.

A.  \ding{172} . B.  \ding{172}  \ding{173} . C.  \ding{172}  \ding{174} . D.  \ding{172}  \ding{173}  \ding{174} .

\section*{三、解答题}

13. 已知数列 $\left\{  {a}_{n}\right\}$ 和 $\left\{  {b}_{n}\right\}$ ,其中 ${b}_{n} = {2}^{{a}_{n}}\left( {n \in  {\mathbf{N}}^{ * }}\right)$ ,数列 $\left\{  {{a}_{n} + {b}_{n}}\right\}$ 的前 $n$ 项和为 ${S}_{n}$ .

(1)若 ${a}_{n} = {2n}$ ，求 ${S}_{n}$ ；

(2)若 $\left\{  {b}_{n}\right\}$ 是各项为正的等比数列， ${S}_{n} = {3n}$ ，求数列 $\left\{  {a}_{n}\right\}$ 和 $\left\{  {b}_{n}\right\}$ 的通项公式.

14. 已知数列 $\left\{  {{\log }_{2}\left( {{a}_{n} - 1}\right) }\right\}$ 为等差数列,且 ${a}_{1} = 3,{a}_{3} = 9$ .

(1)求数列 $\left\{  {a}_{n}\right\}$ 的通项公式；

(2)证明 $\frac{1}{{a}_{2} - {a}_{1}} + \frac{1}{{a}_{3} - {a}_{2}} + \cdots  + \frac{1}{{a}_{n + 1} - {a}_{n}} < 1$ .

15. 已知数列 $\left\{  {a}_{n}\right\}$ 的前 $n$ 项和 ${S}_{n}$ 满足 $2{S}_{n} - n{a}_{n} = n\left( {n \in  \mathbf{N}}\right)$ ,且 ${a}_{2} = 3$ .

(1)求证:数列 $\left\{  \frac{{a}_{n} - 1}{n - 1}\right\}  \left( {n \geq  2}\right)$ 是常数列；

(2)若数列 $\left\{  {b}_{n}\right\}$ 通项公式 ${b}_{n} = {3n} - 2$ ，将数列 $\left\{  {a}_{n}\right\}$ 与 $\left\{  {b}_{n}\right\}$ 的公共项按从小到大的顺序排列得到数列 $\left\{  {c}_{n}\right\}$ ,求 $\left\{  {c}_{n}\right\}$ 的前 $n$ 项和.

16. 近两年, 直播带货逐渐成为一种新兴的营销模式, 带来电商行业的新增长点. 某直播平台第 1 年初的启动资金为 500 万元，由于一些知名主播加入，平台资金的年平均增长率可达 40\%，每年年底扣除运营成本 $a$ 万元，再将剩余资金继续投入直播平台.

(1)若 $a = {100}$ ，在第 3 年年底扣除运营成本后，直播平台的资金有多少万元？

(2)每年的运营成本最多控制在多少万元，才能使得直播平台在第 6 年年底扣除运营成本后资金达到 3000 万元?(结果精确到 0.1 万元)

\section*{$\text{ § }{7.3}$ 等差数列和等比数列二}

\section*{典型例题精讲}

【例 1】已知 ${S}_{n}$ 是数列 $\left\{  {a}_{n}\right\}$ 的前 $n$ 项和,且 ${a}_{1} = 1,{a}_{n} + {a}_{n + 1} = {2n} + 1$ . (1)求数列 $\left\{  {a}_{n}\right\}$ 的通项公式；

(2)记 ${b}_{n} = \left\{  \begin{array}{l} {2}^{{a}_{n}},n\text{ 为奇数,求数列 }\left\{  {b}_{n}\right\}  \text{ 的前 }{2n}\text{ 项和 }{T}_{n}. \end{array}\right.$

【例 2】已知数列 $\left\{  {a}_{n}\right\}$ 前 $n$ 项积为 ${T}_{n}$ ,且 ${a}_{n} + {T}_{n} = 1$ .

(1)求证:数列 $\left\{  \frac{1}{1 - {a}_{n}}\right\}$ 为等差数列；

(2)设 ${S}_{n} = {T}_{1}^{2} + {T}_{2}^{2} + \cdots  + {T}_{n}^{2}$ ，求证: ${S}_{n} > {a}_{n + 1} - \frac{1}{2}$ .

【例 3】已知正项等差数列 $\left\{  {a}_{n}\right\}$ 满足: ${a}_{3n} = 3{a}_{n}\left( {n \in  {\mathbf{N}}^{ * }}\right)$ ,且 $2{a}_{1},{a}_{3} + 1,{a}_{8}$ 成等比数列.

(1)求 $\left\{  {a}_{n}\right\}$ 的通项公式；

(2)设 ${c}_{n} = \frac{{2}^{{a}_{n + 1}}}{\left( {1 + {2}^{{a}_{n}}}\right) \left( {1 + {2}^{{a}_{n + 1}}}\right) }$ ， ${R}_{n}$ 是数列 $\left\{  {c}_{n}\right\}$ 的前 $n$ 项和，若对任意正整数 $n$ 均有 ${R}_{n} < \lambda$ 恒成立，求 $\lambda$ 的最小值.

【例 4】已知数列 $\left\{  {a}_{n}\right\}$ 的前 $n$ 项和为 ${S}_{n}$ ,且 ${S}_{n} = 2{a}_{n} - 2\left( {n \in  {\mathbf{N}}^{ * }}\right)$ .

(1)求 $\left\{  {a}_{n}\right\}$ 的通项公式；

(2)设 ${b}_{n + 1} = 2{b}_{n} - {2}^{n + 1},{b}_{1} = 8,{T}_{n}$ 是数列 $\left\{  {b}_{n}\right\}$ 的前 $n$ 项和，求正整数 $k$ ，使得对任意 $n \in  {\mathbf{N}}^{ * }$ 均有 ${T}_{k} \geq  {T}_{n}$ 恒成立.

【例 5】小张计划贷款 60 万元用于购买一套商品住房, 银行给小张提供了两种贷款方式:  \ding{172}  等额本金: 每月的还款额呈递减趋势，且从第二个还款月开始，每月还款额与上月还款额的差均相同； \ding{173}  等额本息:每个月的还款额均相同. 银行规定，在贷款到账日的次月当天开始首次还款(若 2023 年 7 月 7 日贷款到账，则 2023 年 8 月 7 日首次还款).

已知小张该笔贷款年限为 20 年, 月利率为 0.004 .

(1)若小张采取等额本金的还款方式，现已得知第一个还款月应还 4900 元，最后一个还款月应还 2510 元，试计算小张该笔贷款的总利息；

(2)若小张采取等额本息的还款方式，银行规定，每月还款额不得超过家庭平均月收入的一半，已知小张家庭平均月收入为 1 万元，判断小张该笔贷款是否能够获批(不考虑其他因素)；

(3)对比两种还款方式，从经济利益的角度来考虑，小张应选择哪种还款方式.

\section*{四基达标精练}

\section*{一、填空题}

1. 已知数列 $\left\{  {a}_{n}\right\}$ 为等差数列,若 ${a}_{2} + {a}_{8} = \frac{2\pi }{3}$ ,则 $\tan \left( {{a}_{3} + {a}_{7}}\right)  =$ \_\_\_\_\_.

2. 已知等差数列 $\left\{  {a}_{n}\right\}$ 的前 $n$ 项和是 ${S}_{n}$ ，且 ${a}_{9} = 8$ ，则 ${S}_{17} =$ \_\_\_\_\_.

3. 设等差数列 $\left\{  {a}_{n}\right\}$ 的前 $n$ 项和为 ${S}_{n}$ ,若 $\frac{{S}_{2023}}{2023} - \frac{{S}_{2022}}{2022} = 2$ ,则数列 $\left\{  {a}_{n}\right\}$ 的公差为\_\_\_\_\_.

4. 已知等差数列 $\left\{  {a}_{n}\right\}$ 的项数为奇数,其中所有奇数项之和为 319,所有偶数项之和为 290,则该数列的中间项为\_\_\_\_\_.

5. 已知正数等比数列 $\left\{  {a}_{n}\right\}$ 中，公比 $q = 2,{a}_{1}{a}_{2}{a}_{3}\cdots {a}_{30} = {2}^{45}$ ,则 ${a}_{1}{a}_{4}{a}_{7}\cdots {a}_{28} =$ \_\_\_\_\_.

6. 已知实数 $x,{a}_{1},{a}_{2},y$ 成等差数列， $x,{b}_{1},{b}_{2},y$ 成等比数列，则 $\frac{{\left( {a}_{1} + {a}_{2}\right) }^{2}}{{b}_{1}{b}_{2}}$ 的取值范围是\_\_\_\_\_.

7. 数列 $\left\{  {a}_{n}\right\}$ 中， ${a}_{1} = \frac{1}{5},{a}_{n} + {a}_{n + 1} = \frac{6}{{5}^{n + 1}}\left( {n \in  {\mathbf{N}}^{ * }}\right)$ ，则 $\mathop{\sum }\limits_{{i = 1}}^{{+\infty }}{a}_{i} =$ \_\_\_\_\_.

8. 数列 $\left\{  {a}_{n}\right\}$ 满足 ${a}_{n + 2} + {\left( -1\right) }^{n}{a}_{n} = {3n} - 1$ ，前 16 项和为 540，则 ${a}_{1} =$ \_\_\_\_\_.

9. 已知数列 $\left\{  {b}_{n}\right\}$ ， ${b}_{n + 1} - {b}_{n} = 2 = {b}_{1}\left( {n \in  {\mathbf{N}}^{ * }}\right)$ ，等比数列 $\left\{  {a}_{n}\right\}$ 中， ${a}_{1} = {b}_{1},{a}_{4} = {b}_{8}$ ，若数列 $\left\{  {b}_{n}\right\}$ 中去掉与数列 $\left\{  {a}_{n}\right\}$ 相同的项后余下的项按原顺序组成数列 $\left\{  {c}_{n}\right\}$ ,则 $\left\{  {c}_{n}\right\}$ 前 200 项的和为\_\_\_\_\_.

\section*{二、选择题}

10. 已知数列 $\left\{  {a}_{n}\right\}$ 是公差不为零的等差数列, $\left\{  {b}_{n}\right\}$ 是正项等比数列,若 ${a}_{1} = {b}_{1},{a}_{7} = {b}_{7}$ ,则 ( )

A. ${a}_{4} = {b}_{4}$ . B. ${a}_{5} < {b}_{5}$ . C. ${a}_{8} > {b}_{8}$ . D. ${a}_{9} < {b}_{9}$ .

11. 已知函数 $f\left( x\right)  = \frac{1}{3}{x}^{3} + {4x}$ ,记等差数列 $\left\{  {a}_{n}\right\}$ 的前 $n$ 项和为 ${S}_{n}$ ,若 $f\left( {{a}_{1} + 2}\right)  = {100},f\left( {{a}_{2022} + 2}\right)  =$ -100,则 ${S}_{2022}$ 等于 ( )

A. -4044 . B. -2022. C. 2022. D.4044.

12. 各项均为正数的数列 $\left\{  {a}_{n}\right\}$ 满足 ${a}_{n} = 2\sqrt{{S}_{n}} - 1$ ，其中 ${S}_{n}$ 是数列 $\left\{  {a}_{n}\right\}$ 的前 $n$ 项和，若对任意 $n \in  \mathbf{N}$ ， $n \geq  2$ ,总有 $\frac{1}{{S}_{2} - 1} + \frac{1}{{S}_{3} - 1}\cdots  + \frac{1}{{S}_{n} - 1} \leq  \lambda$ 恒成立,则实数 $\lambda$ 的最小值为 ( )

A.1. B. $\frac{3}{4}$ . C. $\frac{2}{3}$ . D. $\frac{1}{3}$ .

\section*{三、解答题}

13. 设数列 $\left\{  {a}_{n}\right\}$ 的前 $n$ 项和为 ${S}_{n}$ ,满足 ${S}_{n} = 1 - n{a}_{n}$ .

(1)求数列 $\left\{  {a}_{n}\right\}$ 的通项公式；

(2)设数列 $\left\{  \frac{{\left( -1\right) }^{n}}{{a}_{n}}\right\}$ 的前 $n$ 项和为 ${T}_{n}$ ，求 ${T}_{2n}$ 的表达式.

14. 赵先生准备通过某银行贷款 5000 元, 然后通过分期付款的方式还款. 银行与赵先生约定: 每个月还款一次，分12次还清所有欠款，且每个月还款的金额都相等，贷款的月利率为 0.5\%，则赵先生每个月所要还款的金额是多少元. (精确到 0.01 元)

15. 已知 ${S}_{n}$ 是公差不为零的等差数列 $\left\{  {a}_{n}\right\}$ 的前 $n$ 项和, ${S}_{3} = 6,{a}_{3}^{2} = {a}_{1} \cdot  {a}_{9}$ .

(1)求数列 $\left\{  {a}_{n}\right\}$ 的通项公式；

(2)设 ${b}_{n} = {\left( -1\right) }^{n}\frac{4{a}_{n}}{4{n}^{2} - 1}$ ，数列 $\left\{  {b}_{n}\right\}$ 的前 ${2n}$ 项和为 ${T}_{2n}$ ，若 $\left| {{T}_{2n} + 1}\right|  < \frac{1}{2022}$ ，求正整数 $n$ 的最小值.

16. 已知数列 $\left\{  {b}_{n}\right\}$ 的前 $n$ 项和 ${S}_{n}$ 满足 $\frac{1}{{S}_{1}} + \frac{1}{{S}_{2}} + \cdots  + \frac{1}{{S}_{n}} = \frac{2n}{n + 1}$ .

(1)求数列 $\left\{  {b}_{n}\right\}$ 的通项公式；

(2)设数列 $\left\{  {c}_{n}\right\}$ 满足 ${c}_{n} = \frac{{b}_{n}}{{2}^{n}}$ ，其前 $n$ 项和为 ${T}_{n}$ ，是否存在正整数 $n$ ，使得 ${T}_{n} = 4 - n$ 成立？若存在，求出所有 $n$ 的值；若不存在，请说明理由.

\section*{$\text{ § }{7.4}$ 数列综合}

\section*{典型例题精讲}

【例 1】已知数列 $\left\{  {a}_{n}\right\}$ 的前 $n$ 项和为 ${S}_{n}$ ,且 ${a}_{1} = 1,{a}_{n} \neq  0,4{S}_{n} = {a}_{n} \cdot  {a}_{n + 1} + 1\left( {n \in  {\mathbf{N}}^{ * }}\right)$ .

(1)求证: $\left\{  {a}_{{2k} - 1}\right\}$ 、 $\left\{  {a}_{2k}\right\}  \left( {k \in  {\mathbf{N}}^{ * }}\right)$ 都是等差数列；

(2)求数列 $\left\{  {a}_{n}\right\}$ 的通项公式.

【例 2】已知数列 $\left\{  {a}_{n}\right\}$ 的前 $n$ 项和 ${S}_{n}$ ,且 ${a}_{n} \in  \mathbf{Z},{a}_{1} = 1$ ,记 $D\left( {a}_{n}\right)  = \frac{1}{n - 1}\left( {\left| {{a}_{2} - {a}_{1}}\right|  + \left| {{a}_{3} - {a}_{2}}\right|  + \cdots  + }\right.  \left. {\mid {a}_{n} - {a}_{n - 1} \mid  }\right) \left( {n \geq  2,n \in  \mathbf{N}}\right)$ .

(1)若 ${a}_{n} = {2n} - 1$ ，求 $D\left( {a}_{3}\right)  + D\left( {S}_{3}\right)$ ；

(2)若对任意正整数 $n$ ， ${3D}\left( {S}_{n}\right)  = {2D}\left( {a}_{n}\right)$ ，求 ${a}_{n}$ .

【例 3】观察下列等式:

\[
1 - \frac{1}{2} = \frac{1}{2},1 - \frac{1}{2} + \frac{1}{3} - \frac{1}{4} = \frac{1}{3} + \frac{1}{4},1 - \frac{1}{2} + \frac{1}{3} - \frac{1}{4} + \frac{1}{5} - \frac{1}{6} = \frac{1}{4} + \frac{1}{5} + \frac{1}{6},\cdots
\]

据此规律,请你猜想出第 $n$ 个等式并证明你的结论.

【例 4】设数列 $\left\{  {a}_{n}\right\}  \text{ 、 }\left\{  {b}_{n}\right\}$ 满足 ${a}_{1} = {b}_{1} = 6,{a}_{2} = {b}_{2} = 4,{a}_{3} = {b}_{3} = 3$ ,且数列 $\left\{  {{a}_{n + 1} - {a}_{n}}\right\}$ 是等差数列,数列 $\left\{  {{b}_{n} - 2}\right\}$ 是等比数列.

(1)求数列 $\left\{  {a}_{n}\right\}$ 、 $\left\{  {b}_{n}\right\}$ 的通项公式；

(2)是否存在正整数 $k$ ，使 ${a}_{k} - {b}_{k} \in  \left( {0,\frac{1}{2}}\right)$ ？若存在，求出 $k$ ；若不存在，说明理由.

【例 5】在平面直角坐标系上,有一点列 ${P}_{0},{P}_{1},{P}_{2},{P}_{3},\cdots ,{P}_{n - 1},{P}_{n}$ ,设点 ${P}_{k}$ 的坐标 $\left( {{x}_{k},{y}_{k}}\right) \left( {k \in  \mathbf{N}}\right)$ ,其中 ${x}_{k} \in  \mathbf{Z}$ 、 ${y}_{k} \in  \mathbf{Z}$ . 记 $\Delta {x}_{k} = {x}_{k} - {x}_{k - 1},\Delta {y}_{k} = {y}_{k} - {y}_{k - 1}$ ,且满足 $\left| {\Delta {x}_{k}}\right|  \cdot  \left| {\Delta {y}_{k}}\right|  = 2\left( {k = 1,2,3,\cdots ,n}\right)$ .

(1)点 ${P}_{0}\left( {0,1}\right)$ ，点 ${P}_{1}$ 坐标满足 ${x}_{0} < {x}_{1},{y}_{0} < {y}_{1}$ ，且点 ${P}_{k}\left( {k = 0,1}\right)$ 落在直线 $l$ 上，求直线 $l$ 的方程；

(2)若 ${P}_{0}\left( {0,0}\right)$ ， ${y}_{n} = {100}$ ，且 $n \geq  {2020}$ ，求 ${T}_{n} = {x}_{0} + {x}_{1} + {x}_{2} + \cdots  + {x}_{n}$ 的最大值.

\section*{四基达标精练}

\section*{一、填空题}

1. 若首项为 1、公比为 $\frac{1}{3}$ 的无穷等比数列的各项和为 $S$ ，则 $S =$ \_\_\_\_\_.

2. 在等差数列 $\left\{  {a}_{n}\right\}$ 中,已知公差 $d \neq  0,{a}_{2}^{2} = {a}_{1}{a}_{4}$ ,若 ${a}_{1},{a}_{3},{a}_{{k}_{1}},{a}_{{k}_{2}},\cdots ,{a}_{{k}_{n}},\cdots$ 成等比数列,则 ${k}_{n} =$ \_\_\_\_\_.

3. 设数列 $\left\{  {a}_{n}\right\}   : {2}^{1},{2}^{2},{2}^{2},{2}^{3},{2}^{3},{2}^{3},{2}^{4},{2}^{4},{2}^{4},{2}^{4},{2}^{5},{2}^{5},{2}^{5},{2}^{5},{2}^{5},{2}^{5},\cdots$ . 则 ${a}_{50} =$ \_\_\_\_\_.

4. 用数学归纳法证明 “ $1 + \frac{1}{2} + \frac{1}{3} + \cdots  + \frac{1}{{2}^{n} - 1} < n\left( {n \geq  2}\right)$ ” 时,由 $n = k$ 的假设证明 $n = k + 1$ 时,不等式左边需增加的项数为\_\_\_\_\_.

5. 已知数列 $\left\{  {a}_{n}\right\}$ 是以 1 为首项、2 为公差的等差数列,且 $\left\{  {{\left( -1\right) }^{n + 1}{a}_{n}^{2}}\right\}$ 的前 $n$ 项和为 ${S}_{n}$ ,则数列 $\left\{  \frac{{S}_{2n}}{n}\right\}$ 的前 10 项和为\_\_\_\_\_.

6. 若 $a$ 、 $b$ 分别是正数 $p$ 、 $q$ 的算术平均数和几何平均数，且 $a$ 、 $b$ 、-2这三个数可适当排序后成等差数列，也可适当排序后成等比数列，则 $p + q + {pq}$ 的值构成的集合是\_\_\_\_\_.

7. 若数列 $\left\{  {a}_{n}\right\}$ 满足:对任意正整数 $n$ ， $\left\{  {{a}_{n + 1} - {a}_{n}}\right\}$ 为递减数列，则称数列 $\left\{  {a}_{n}\right\}$ 为“差递减数列”. 给出下列数列 $\left\{  {a}_{n}\right\}   :$  \ding{172}  ${a}_{n} = {3n}$ ；  \ding{173}  ${a}_{n} = {n}^{2} + 1$ ； \ding{174}  ${a}_{n} = \sqrt{n}$ ； \ding{175}  ${a}_{n} = \frac{n}{n + 1}$ ，其中是“差递减数列”的有\_\_\_\_\_.

8. 若数列 $\left\{  {a}_{n}\right\}$ 满足 ${a}_{1} = {10},{a}_{n + 1} = {a}_{n} + {18n} + {10}\left( {n \in  \mathbf{N},n \geq  1}\right)$ ，记 $\left\lbrack  x\right\rbrack$ 表示不超过实数 $x$ 的最大整数，则 $\mathop{\lim }\limits_{{n \rightarrow  \infty }}\left( {\sqrt{{a}_{n}} - \left\lbrack  \sqrt{{a}_{n}}\right\rbrack  }\right)  =$ \_\_\_\_\_.

9. 已知 $n \in  \mathbf{N},n \geq  2$ ，函数 $y = \frac{n}{{n}^{2} + 3}x + \frac{3n}{n + 3}$ 的图像与 $y$ 轴相交于点 ${A}_{n}$ ，与 $y = {\log }_{\frac{1}{n}}\left( {x - 4}\right)$ 的图像相交于点 ${B}_{n}$ ， $\bigtriangleup  O{A}_{n}{B}_{n}$ 的面积为 ${S}_{n}$ ( $O$ 为坐标原点)，则 $\mathop{\lim }\limits_{{n \rightarrow  \infty }}{S}_{n} =$ \_\_\_\_\_.

\section*{二、选择题}

10. 已知 ${a}_{n} = {n}^{2} + {\lambda n}$ ,则 “ ${a}_{1} < {a}_{2}$ ” 是 “ $\left\{  {a}_{n}\right\}$ 是严格递增数列” 的 ( )

A. 充分非必要条件. B. 必要非充分条件. C. 充要条件. D. 既非充分也非必要条件.

11. 定义在 $\lbrack 0, + \infty )$ 上的函数 $f\left( x\right)$ 满足: 当 $0 \leq  x < 2$ 时, $f\left( x\right)  = {2x} - {x}^{2}$ ; 当 $x \geq  2$ 时, $f\left( x\right)  = {3f}\left( {x - 2}\right)$ . 记函数 $f\left( x\right)$ 的极大值点从小到大依次为 ${a}_{1},{a}_{2},\cdots ,{a}_{n},\cdots$ ,并记相应的极大值依次为 ${b}_{1},{b}_{2},\cdots ,{b}_{n},\cdots$ ,则 ${S}_{20} = {a}_{1}{b}_{1} + {a}_{2}{b}_{2} + \cdots  + {a}_{20}{b}_{20}$ 的值为 ( )

A. ${19} \times  {3}^{20} + 1$ . B. ${19} \times  {3}^{19} + 1$ . C. ${20} \times  {3}^{19} + 1$ . D. ${20} \times  {3}^{20} + {1711}$ .

12. 已知 $\overrightarrow{e}\text{ 、 }\overrightarrow{f}$ 是互相垂直的单位向量,向量 $\overrightarrow{{a}_{n}}$ 满足: $\overrightarrow{e} \cdot  \overrightarrow{{a}_{n}} = n,\overrightarrow{f} \cdot  \overrightarrow{{a}_{n}} = {2n} + 1,{b}_{n}$ 是向量 $\overrightarrow{f}$ 与 $\overrightarrow{{a}_{n}}$ 夹角的正切值,则数列 $\left\{  {b}_{n}\right\}$ 是 ( )

A. 单调递增数列且 $\mathop{\lim }\limits_{{n \rightarrow  \infty }}{b}_{n} = \frac{1}{2}$ . B. 单调递减数列且 $\mathop{\lim }\limits_{{n \rightarrow  \infty }}{b}_{n} = \frac{1}{2}$ .

C. 单调递增数列且 $\mathop{\lim }\limits_{{n \rightarrow  \infty }}{b}_{n} = 2$ . D. 单调递减数列且 $\mathop{\lim }\limits_{{n \rightarrow  \infty }}{b}_{n} = 2$ .

\section*{三、解答题}

13. 设 $\left\{  {a}_{n}\right\}$ 是等差数列,公差为 $d$ ,前 $n$ 项和为 ${S}_{n}$ .

(1)设 ${a}_{1} = {40},{a}_{6} = {38}$ ，求 ${S}_{n}$ 的最大值；

(2)设 ${a}_{1} = 1,{b}_{n} = {2}^{{a}_{n}}\left( {n \in  {\mathbf{N}}^{ * }}\right)$ ，数列 $\left\{  {b}_{n}\right\}$ 的前 $n$ 项和为 ${T}_{n}$ ，且对任意的 $n \in  {\mathbf{N}}^{ * }$ ，都有 ${T}_{n} \leq  {20}$ ，求 $d$ 的取值范围.

14. 已知定理: 设数列 $\left\{  {a}_{n}\right\}$ 、 $\left\{  {b}_{n}\right\}$ 的前 $n$ 项和分别为 ${A}_{n}$ 、 ${B}_{n}$ ,则 ${A}_{n} = {B}_{n} + p$ ( $p$ 为非零常数) 的充要条件是 ${a}_{1} = {b}_{1} + p,{a}_{n} = {b}_{n}\left( {n \geq  2}\right)$ .

(1)若 ${a}_{1} = {10}$ ， ${a}_{n} = {2}^{n}\left( {n \geq  2}\right)$ ，试运用上述定理求数列 $\left\{  {a}_{n}\right\}$ 的前 $n$ 项和 ${A}_{n}$ ；

(2)已知数列 $\left\{  {a}_{n}\right\}$ 的前 $n$ 项和减去数列 $\left\{  {b}_{n}\right\}$ 的前 $n$ 项和的差为 4，数列 $\left\{  {b}_{n}\right\}$ 是首项为 2 的等差数列，且数列 $\left\{  {b}_{n}\right\}$ 的前 5 项和为 40,试运用上述定理求 ${a}_{1}$ 和 ${a}_{5}$ .

15. 定义在 $\mathbf{R}$ 上的偶函数 $f\left( x\right)$ 的图像关于 $x = 1$ 对称,且 $f\left( {{x}_{1} + {x}_{2}}\right)  = f\left( {x}_{1}\right)  \cdot  f\left( {x}_{2}\right)$ 对任意 ${x}_{1}\text{ 、 }{x}_{2} \in  \left\lbrack  {0,\frac{1}{2}}\right\rbrack$ 都成立.

(1)若 $f\left( 1\right)  = a\left( {a > 0}\right)$ ， 探求 $f\left( \frac{1}{2}\right)$ ， $f\left( \frac{1}{4}\right)$ ， $\cdots$ ， $f\left( \frac{1}{{2}^{n}}\right)$ 的值；

(2)若 $f\left( 1\right)  = a\left( {a > 0}\right) ,{a}_{n} = f\left( {{2n} + \frac{1}{2n}}\right)$ ，求 $\mathop{\lim }\limits_{{n \rightarrow  \infty }}{a}_{n}$ 的值.

16. 如果存在常数 $a$ ,使得数列 $\left\{  {a}_{n}\right\}$ 满足: 若 $x$ 是数列 $\left\{  {a}_{n}\right\}$ 中的一项,则 $\frac{a}{x}$ 也是数列 $\left\{  {a}_{n}\right\}$ 中的一项,称数列 $\left\{  {a}_{n}\right\}$ 具有倒置性质 $P$ ,常数 $a$ 是它的 “倒置系数”.

(1)若数列: $1,2,4,m\left( {m > 4}\right)$ 具有倒置性质 $P$ ，求 $m$ 和倒置系数 $a$ ；

(2)已知各项为正的有穷等比数列 $\left\{  {b}_{n}\right\}$ 的项数是 ${n}_{0}\left( {{n}_{0} \geq  3}\right)$ ，公比为 $q$ ，所有项之积是 $B$ ，求证:数列 $\left\{  {b}_{n}\right\}$ 具有倒置性质 $P$ ,并用 ${n}_{0}$ 和 $B$ 表示它的倒置系数 $a$ .


\section*{第 8 章 平面向量与复数(必修 2)}

\section*{知识梳理*}

\section*{一、向量}

1. 向量 (平面向量和空间向量): 既有大小,又有方向的量.

2. 向量的运算:加法、减法(平行四边形、三角形法则).

实数与向量的乘积 $k\overrightarrow{a} : k > 0,k\overrightarrow{a}$ 与 $\overrightarrow{a}$ 同向; $k < 0,k\overrightarrow{a}$ 与 $\overrightarrow{a}$ 反向; $k = 0$ ,零向量; $k\overrightarrow{a}$ 的模为 $\left| k\right| \left| \overrightarrow{a}\right|$

向量的数量积: $\overrightarrow{a} \cdot  \overrightarrow{b} = \left| \overrightarrow{a}\right| \left| \overrightarrow{b}\right| \cos \theta$ ( $\theta$ 为 $\overrightarrow{a}$ 和 $\overrightarrow{b}$ 的夹角); 变式: $\cos \theta  = \frac{\overrightarrow{a} \cdot  \overrightarrow{b}}{\left| \overrightarrow{a}\right|  \cdot  \left| \overrightarrow{b}\right| }$ .

数量积的几何意义: $\left| \overrightarrow{a}\right|$ 与 $\overrightarrow{b}$ 在 $\overrightarrow{a}$ 方向上的数量投影的乘积 (该投影可正可负可为零).

3. 向量的坐标运算: $\overrightarrow{a} = \left( {{x}_{1},{y}_{1}}\right) ,\overrightarrow{b} = \left( {{x}_{2},{y}_{2}}\right)$ ,则 $\overrightarrow{a} \pm  \overrightarrow{b} = \left( {{x}_{1} \pm  {x}_{2},{y}_{1} \pm  {y}_{2}}\right) ;k\overrightarrow{a} = \left( {k{x}_{1},k{y}_{1}}\right) ;\overrightarrow{a} \cdot  \overrightarrow{b} =  {x}_{1} \cdot  {x}_{2} + {y}_{1} \cdot  {y}_{2}$

4. $\overrightarrow{a}//\overrightarrow{b} \Leftrightarrow$ 存在 $k \in  \mathbf{R}$ ,使 $\overrightarrow{a} = k\overrightarrow{b}\left( {\overrightarrow{b} \neq  \overrightarrow{0}}\right)  \Leftrightarrow  {x}_{1}{y}_{2} = {x}_{2}{y}_{1} \Leftrightarrow  {\left( \overrightarrow{a} \cdot  \overrightarrow{b}\right) }^{2} = {\left| \overrightarrow{a}\right| }^{2}{\left| \overrightarrow{b}\right| }^{2}$ .

5. $\overrightarrow{a} \bot  \overrightarrow{b} \Leftrightarrow  \overrightarrow{a} \cdot  \overrightarrow{b} = 0 \Rightarrow  \left| {\overrightarrow{a} + \overrightarrow{b}}\right|  = \left| {\overrightarrow{a} - \overrightarrow{b}}\right|  \Leftrightarrow  {x}_{1} \cdot  {x}_{2} + {y}_{1} \cdot  {y}_{2} = 0$ (空间向量类似).

6. $\overrightarrow{a},\overrightarrow{b}$ 的夹角为锐角 $\Leftrightarrow  \left\{  {\begin{array}{l} \overrightarrow{a} \cdot  \overrightarrow{b} > 0, \\  \overrightarrow{a} \times  \overrightarrow{b} \end{array};\overrightarrow{a}\text{ 、 }\overrightarrow{b}}\right.$ 的夹角为钝角 $\Leftrightarrow  \left\{  \begin{array}{l} \overrightarrow{a} \cdot  \overrightarrow{b} < 0, \\  \overrightarrow{a} \times  \overrightarrow{b}. \end{array}\right.$

7. 定比分点公式: 若 $M\left( {{x}_{1},{y}_{1}}\right) ,N\left( {{x}_{2},{y}_{2}}\right)$ 满足 $\overrightarrow{MP} = \lambda \overrightarrow{PN}$ ,则有 $P\left( {\frac{{x}_{1} + \lambda {x}_{2}}{1 + \lambda },\frac{{y}_{1} + \lambda {y}_{2}}{1 + \lambda }}\right)$ .

8. 常用结论: \ding{172}  ${\overset{ \rightarrow  }{a}}^{2} = {\left| \overset{ \rightarrow  }{a}\right| }^{2}$ ;  \ding{173}  $\overset{ \rightarrow  }{a} \cdot  \overrightarrow{b} \leq  \left| \overset{ \rightarrow  }{a}\right| \left| \overset{ \rightarrow  }{b}\right|$ ;  \ding{174}  ${\left( \overrightarrow{a} \pm  \overrightarrow{b}\right) }^{2} = {\left| \overrightarrow{a} \pm  \overrightarrow{b}\right| }^{2} = {\overrightarrow{a}}^{2} \pm  2\overrightarrow{a} \cdot  \overrightarrow{b} + {\overrightarrow{b}}^{2}$ ；

 \ding{175}  $\left| \right| \overrightarrow{a}\left| -\right| \overrightarrow{b}\left| \right|  \leq  \left| {\overrightarrow{a} \pm  \overrightarrow{b}}\right|  \leq  \left| \overrightarrow{a}\right|  + \left| \overrightarrow{b}\right|$ ;  \circled{5}  ${\left| \overrightarrow{a} + \overrightarrow{b}\right| }^{2} + {\left| \overrightarrow{a} - \overrightarrow{b}\right| }^{2} = 2\left( {{\left| \overrightarrow{a}\right| }^{2} + {\left| \overrightarrow{b}\right| }^{2}}\right)$ ;

 \circled{6} 非零向量 $\overrightarrow{a}$ 、 $\overrightarrow{b}$ 、 $\overrightarrow{c}$ 共面的充要条件是存在 $\lambda$ 、 $u \in  \mathbf{R}$ 使得 $\overrightarrow{a} = \lambda \overrightarrow{b} + u\overrightarrow{c}$ .

9. 向量的投影和数量投影:

$\left| \overrightarrow{a}\right| \cos \langle \overrightarrow{a},\overrightarrow{b}\rangle  \cdot  \frac{\overrightarrow{b}}{\left| \overrightarrow{b}\right| } = \frac{\overrightarrow{a} \cdot  \overrightarrow{b}}{\left| \overrightarrow{b}\right| } \cdot  \frac{\overrightarrow{b}}{\left| \overrightarrow{b}\right| }$ 表示向量 $\overrightarrow{a}$ 在向量 $\overrightarrow{b}$ 方向上的投影是向量.

其中系数 $\left| \overrightarrow{a}\right| \cos \langle \overrightarrow{a},\overrightarrow{b}\rangle  = \frac{\overrightarrow{a} \cdot  \overrightarrow{b}}{\left| \overrightarrow{b}\right| }$ 表示向量 $\overrightarrow{a}$ 在向量 $\overrightarrow{b}$ 方向上的数量投影.

10. 平面向量基本定理: 平面内任意一个向量都可以表示为该平面内两个不平行的向量的线性组合.

\section*{二、复数}

1. 引进虚数单位 ${i}$ 的原则: ${\left( {D}{i}\right) }^{2} =  - 1$ ;  \ding{173}  ${i}$ 可以和实数进行四则运算，且不改变原有的运算律.

2. 对于复数 $z = a + b{i}\left( {a\text{ 、 }b \in  \mathbf{R}}\right) ,a$ 叫实部,记为 $\operatorname{Re}z;b$ 叫虚部,记为 $\operatorname{Im}z$ .

3. 复数相等: $a + b{i} = c + d{i}\left( {a\text{ 、 }b\text{ 、 }c\text{ 、 }d \in  \mathbf{R}}\right)  \Leftrightarrow  \left\{  \begin{array}{l} a = b, \\  c = d; \end{array}\right.$ 不全是实数的两个复数不能比较大小.

4. $z \in  \mathbf{R} \Leftrightarrow  \operatorname{Im}z = 0 \Leftrightarrow  z = \bar{z} \Leftrightarrow  {z}^{2} \geq  0;z$ 为纯虚数 $\Leftrightarrow  \left\{  {\begin{array}{l} \operatorname{Re}z = 0, \\  \operatorname{Im}z \neq  0 \end{array} \Leftrightarrow  \left\{  {\begin{array}{l} z + \bar{z} = 0, \\  z \neq  0 \end{array} \Leftrightarrow  {z}^{2} < 0}\right. }\right.$ .

\HRule

*注:本章节所属知识模块是 “图形与几何”.

\HRule

5. 复数 $z = a + b{i}\left( {a\text{ 、 }b \in  \mathbf{R}}\right)$ 所对应的点即直角坐标平面内的点 $\left( {a,b}\right)$ ; 也对应于以 $\left( {a,b}\right)$ 为坐标的位置向量.

6.* 复数的三角形式:

$z = r\left( {\cos \theta  + {i}\sin \theta }\right) ,\theta  \in  \mathbf{R},\theta$ 为辐角,记作 $\operatorname{Arg}z,r$ 是复数的模,若取 $\theta  \in  \lbrack 0,{2\pi })$ ,则称 $\theta$ 为辐角主值,记作 $\arg z$ !

7. 复数的运算法则(以下 $a$ 、 $b$ 、 $c$ 、 $d \in  \mathbf{R},m$ 、 $n \in  \mathbf{N},m$ 、 $n \geq  2$ ):

 \ding{172}  $\left( {a + b{i}}\right)  \pm  \left( {c + d{i}}\right)  = \left( {a \pm  c}\right)  + \left( {b \pm  d}\right) {i}$ ；

 \ding{173}  $\left( {a + b{i}}\right)  \div  \left( {c + d{i}}\right)  = \frac{{ac} + {bd}}{{c}^{2} + {d}^{2}} + \frac{{bc} - {ad}}{{c}^{2} + {d}^{2}}{i}\left( {c + d{i} \neq  0}\right)$ ；

 \ding{174}  ${z}^{n} = \underset{n}{\underbrace{z \cdot  z \cdot  z\cdots z}}\left( {n \in  {\mathbf{N}}^{ * }}\right)$ ；

 \ding{175}  $\left( {a + b{i}}\right) \left( {c + d{i}}\right)  = \left( {{ac} - {bd}}\right)  + \left( {{bc} + {ad}}\right) {i}$ ；

 \circled{5}  ${z}^{m}{z}^{n} = {z}^{m + n},{\left( {z}^{m}\right) }^{n} = {z}^{mn},{\left( {z}_{1}{z}_{2}\right) }^{n} = {z}_{1}^{n}{z}_{2}^{n}$ .

本质:复数的加、减、乘、乘方运算仿实数的多项式运算(合并同类项)，并随时将 ${{i}}^{2}$ 换成 -1，复数的除法运算的本质是将分母实数化 (利用 $z\bar{z} = {\left| z\right| }^{2}$ ).

8. 常用的运算结论:

${{i}}^{2} =  - 1,{{i}}^{{4n} + 1} = {i},{{i}}^{{4n} + 2} =  - 1,{{i}}^{{4n} + 3} =  - {i},{{i}}^{4n} = 1,{{i}}^{n} + {{i}}^{n + 1} + {{i}}^{n + 2} + {{i}}^{n + 3} = 0\left( {n \in  \mathbf{Z}}\right)$

${\left( 1 \pm  {i}\right) }^{2} =  \pm  2{i},\frac{1 + {i}}{1 - {i}} = {i},\frac{1 - {i}}{1 + {i}} =  - {i}.$

9. 共轭复数的运算性质: $\bar{z} = z;\overline{{z}_{1} \pm  {z}_{2}} = \overline{{z}_{1}} \pm  \overline{{z}_{2}};\overline{{z}_{1} \cdot  {z}_{2}} = \overline{{z}_{1}} \cdot  \overline{{z}_{2}};\overline{\left( \frac{{z}_{1}}{{z}_{2}}\right) } = \frac{\overline{{z}_{1}}}{{z}_{2}}\left( {{z}_{2} \neq  0}\right) ;\overline{{z}^{n}} = {\left( \bar{z}\right) }^{n}$ .

10. 复数模的性质:

$z \cdot  \bar{z} = {\left| z\right| }^{2} = {\left| \bar{z}\right| }^{2} = \left| {z}^{2}\right| ;\left| {{z}_{1} \cdot  {z}_{2}}\right|  = \left| {z}_{1}\right| \left| {z}_{2}\right| ;\left| \frac{{z}_{1}}{{z}_{2}}\right|  = \frac{\left| {z}_{1}\right| }{\left| {z}_{2}\right| }\left( {\left| {z}_{2}\right|  \neq  0}\right) ;\left| {z}^{n}\right|  = {\left| z\right| }^{n}\left( {n \in  \mathbf{N},n \geq  1}\right) ;$

$\left| \right| {z}_{1}\left| -\right| {z}_{2}\left| \right|  \leq  \left| {{z}_{1} \pm  {z}_{2}}\right|  \leq  \left| {z}_{1}\right|  + \left| {z}_{2}\right| ;{\left| {z}_{1} + {z}_{2}\right| }^{2} + {\left| {z}_{1} - {z}_{2}\right| }^{2} = 2\left( {{\left| {z}_{1}\right| }^{2} + {\left| {z}_{2}\right| }^{2}}\right) .$

11. 解复数问题的两个转化:代数化，几何化. 设出一个复数后要注明实部和虚部都是实数.

12. 解复数方程往往是设 $x = a + b{i}\left( {a\text{ 、 }b \in  \mathbf{R}}\right)$ ,利用复数相等转化为实数方程组求解,其本质是两个复数相等的概念.

13. 对于实系数一元二次方程 $a{x}^{2} + {bx} + c = 0$ ,韦达定理一定成立,

 \ding{172}  当 $\Delta  = {b}^{2} - {4ac} > 0$ 时,方程有两个不相等的实根 ${x}_{1}\text{ 、 }{x}_{2} = \frac{-b \pm  \sqrt{{b}^{2} - {4ac}}}{2a}$ ;

此时 $\left| {{x}_{1} - {x}_{2}}\right|  = \sqrt{{\left( {x}_{1} - {x}_{2}\right) }^{2}} = \sqrt{{x}_{1}^{2} + {x}_{2}^{2} - 2{x}_{1}{x}_{2}} = \sqrt{{\left( {x}_{1} + {x}_{2}\right) }^{2} - 4{x}_{1}{x}_{2}} = \frac{\sqrt{{b}^{2} - {4ac}}}{\left| a\right| }$ ;

$\left| {x}_{1}\right|  + \left| {x}_{2}\right|  = \sqrt{{x}_{1}^{2} + {x}_{2}^{2} + 2\left| {{x}_{1}{x}_{2}}\right| } = \sqrt{{\left( {x}_{1} + {x}_{2}\right) }^{2} - 2{x}_{1}{x}_{2} + 2\left| {{x}_{1}{x}_{2}}\right| } = \frac{\sqrt{{b}^{2} - {2ac} + 2\left| {ac}\right| }}{\left| a\right| };$

 \ding{173}  当 $\Delta  = {b}^{2} - {4ac} = 0$ 时,方程有两个相等的实根 ${x}_{1} = {x}_{2} = \frac{-b}{2a}$ ;

此时 $\left| {{x}_{1} - {x}_{2}}\right|  = 0;\left| {x}_{1}\right|  + \left| {x}_{2}\right|  = \frac{\left| b\right| }{\left| a\right| }$ ;

 \ding{174}  当 $\Delta  = {b}^{2} - {4ac} < 0$ 时,方程有两个共轭虚根 ${x}_{1}\text{ 、 }{x}_{2} = \frac{-b \pm  {i}\sqrt{{4ac} - {b}^{2}}}{2a}$ ;

此时 $\left| {{x}_{1} - {x}_{2}}\right|  = \sqrt{{\left| {x}_{1} - {x}_{2}\right| }^{2}} = \sqrt{\left( {{x}_{1} - {x}_{2}}\right) \left( {{\bar{x}}_{1} - {\bar{x}}_{2}}\right) } = \sqrt{-{\left( {x}_{1} - {x}_{2}\right) }^{2}} = \frac{\sqrt{{4ac} - {b}^{2}}}{\left| a\right| }$ ;

$\left| {x}_{1}\right|  + \left| {x}_{2}\right|  = 2\left| {x}_{1}\right|  = 2\sqrt{{\left| {x}_{1}\right| }^{2}} = 2\sqrt{{x}_{1}{\bar{x}}_{1}} = 2\sqrt{{x}_{1}{x}_{2}} = 2\sqrt{\frac{c}{a}}.$

\section*{§ 8.1 平面向量及其运算}

\section*{典型例题精讲}

【例 1】在三角形 ${ABC}$ 中, $\overrightarrow{AB} \bot  \overrightarrow{AC}$ ,点 $F$ 为边 ${AC}$ 中点,点 $E$ 在边 ${AB}$ 上,且 $\overrightarrow{AE} = 2\overrightarrow{EB},{CE}$ 与 ${BF}$ 相交于点 $P$ .

(1)将向量 $\overrightarrow{AP}$ 用向量 $\overrightarrow{AB}$ 、 $\overrightarrow{AC}$ 表示；

(2)若 ${AB} = 3,{AC} = 4$ ，求 $\overrightarrow{AP} \cdot  \overrightarrow{BC}$ .

\begin{center}
\includegraphics[width=0.2\textwidth]{images/bo_d7tk7oc91nqc73875fu0_143_1207_653_233_273_0.jpg}
\end{center}
 

(例 1 图)

【例 2】已知非零向量 $\overrightarrow{a}$ 、 $\overrightarrow{b}$ 满足 $\left| \overrightarrow{a}\right|  = \left| \overrightarrow{b}\right|  = 2\left| {\overrightarrow{a} + \overrightarrow{b}}\right|$ .

(1)求 $\overrightarrow{a}$ 、 $\overrightarrow{b}$ 的夹角的余弦值；

(2)若 $\left| {2\overrightarrow{a} + \overrightarrow{b}}\right|  = 2\sqrt{6}$ ，求 $\left| \overrightarrow{a}\right|$ .

【例 3】如图所示,在 $\bigtriangleup {ABO}$ 中, $\overrightarrow{OC} = \frac{1}{4}\overrightarrow{OA},\overrightarrow{OD} = \frac{1}{2}\overrightarrow{OB},{AD}$ 与 ${BC}$ 相交于点 $M$ ,设 $\overrightarrow{OA} = \overrightarrow{a},\overrightarrow{OB} = \overrightarrow{b}$ .

(1)试用向量 $\overrightarrow{a}$ 、 $\overrightarrow{b}$ 表示 $\overrightarrow{OM}$ ；

(2)过点 $M$ 作直线 ${EF}$ 分别交线段 ${AC}$ 、 ${BD}$ 于点 $E$ 、 $F$ ，记 $\overrightarrow{OE} = \lambda \overrightarrow{OA},\overrightarrow{OF} = \mu \overrightarrow{OB}$ ，求证:不论点 $E$ 、 $F$ 在线段 ${AC}\text{ 、 }{BD}$ 上如何移动, $\frac{1}{\lambda } + \frac{3}{\mu }$ 为定值.

\begin{center}
\includegraphics[width=0.2\textwidth]{images/bo_d7tk7oc91nqc73875fu0_144_1215_400_267_269_0.jpg}
\end{center}
 

(例 3 图)

【例 4】已知向量 $\overrightarrow{a}$ 、 $\overrightarrow{b}$ 的夹角为 ${135}^{ \circ  }$ ，且 $\left| \overrightarrow{a}\right|  = 2$ ， $\left| \overrightarrow{b}\right|  = 2\sqrt{2}$ .

(1)求 $\left| {\overrightarrow{a} + 2\overrightarrow{b}}\right|$ ；

(2) $\overrightarrow{c} = \overrightarrow{a} + x\overrightarrow{b}$ (其中 $x \in  \mathbf{R}$ )，当 $\left| \overrightarrow{c}\right|$ 取最小值时，求 $\overrightarrow{c}$ 与 $\overrightarrow{b}$ 的夹角的大小.

【例 5】已知在平面直角坐标系中,有点 $A\left( {a,0}\right)$ 、点 $B\left( {0,b}\right)$ (其中 $a\text{ 、 }b$ 为常数,且 ${ab} \neq  0$ ),点 $O$ 为坐标原点.

(1)设点 $P$ 为线段 ${AB}$ 靠近点 $A$ 的三等分点， $\overrightarrow{OP} = \lambda \overrightarrow{OA} + \left( {1 - \lambda }\right) \overrightarrow{OB}$ ，求实数 $\lambda$ 的值；

(2)如图，设点 ${P}_{1}$ 、 ${P}_{2}$ 、 $\cdots$ 、 ${P}_{n - 1}$ 是线段 ${AB}$ 的 $n$ 等分点， $\overrightarrow{O{P}_{k}} = \mu \overrightarrow{OA} + \left( {1 - \mu }\right) \overrightarrow{OB}$ ，其中 $1 \leq  k \leq  n - 1$ ( $n$ 、 $k \in  {\mathbf{N}}^{ * },n \geq  2)$ ,当 $n = {2023}$ 时,求 $\left| {\overrightarrow{OA} + \overrightarrow{O{P}_{1}} + \overrightarrow{O{P}_{2}} + \cdots  + \overrightarrow{O{P}_{n - 1}} + \overrightarrow{OB}}\right|$ 的值(用含 $a$ 、 $b$ 的式子表示).

\begin{center}
\includegraphics[width=0.3\textwidth]{images/bo_d7tk7oc91nqc73875fu0_145_975_387_467_253_0.jpg}
\end{center}
 

(例 5 图)

\section*{四基达标精练}

\section*{一、填空题}

1. 已知向量 $\overrightarrow{a} = \left( {2,1}\right) ,\overrightarrow{b} = \left( {-2,4}\right)$ ，则 $\left| {\overrightarrow{a} - \overrightarrow{b}}\right|  =$ \_\_\_\_\_.

2. 已知向量 $\overrightarrow{a} = \left( {\lambda , - 3}\right) ,\overrightarrow{b} = \left( {1,\lambda }\right)$ ，若 $\left( {\overrightarrow{a} + \overrightarrow{b}}\right)  \bot  \overrightarrow{b}$ ，则 $\lambda  =$ \_\_\_\_\_.

3. 已知向量 $\overrightarrow{a}$ 、 $\overrightarrow{b}$ 不平行，若向量 $\overrightarrow{p} = \overrightarrow{a} + \frac{4}{3}m\overrightarrow{b}$ 与向量 $\overrightarrow{q} = \overrightarrow{b} + {3m}\overrightarrow{a}$ 平行，则 $m =$ \_\_\_\_\_.

4. 若空间向量 $\overrightarrow{a} = \left( {1,0,1}\right) ,\overrightarrow{b} = \left( {0,1, - 1}\right)$ ，则向量 $\overrightarrow{a}$ 与 $\overrightarrow{b}$ 的夹角为\_\_\_\_\_.

5. 在 $\bigtriangleup {ABC}$ 中, ${AB} = 2,{AC} = 1,D$ 是 ${BC}$ 边上的中点,则 $\overrightarrow{AD} \cdot  \overrightarrow{BC} =$ \_\_\_\_\_.

6. 若点 $M$ 是 $\bigtriangleup  {ABC}$ 所在平面内一点，且满足 $\overrightarrow{AM} = \frac{3}{5}\overrightarrow{AB} + \frac{2}{5}\overrightarrow{AC}$ . 则 $\bigtriangleup  {ABM}$ 与 $\bigtriangleup  {ABC}$ 的面积之比为

7. 已知不共线向量 $\overrightarrow{OA}\text{ 、 }\overrightarrow{OB}$ 夹角为 $\alpha ,\left| \overrightarrow{OA}\right|  = 1,\left| \overrightarrow{OB}\right|  = 2,\overrightarrow{OP} = \left( {1 - t}\right) \overrightarrow{OA},\overrightarrow{OQ} = t\overrightarrow{OB}$ ，其中 $0 \leq  t \leq  1$ ， $\left| \overrightarrow{PQ}\right|$ 在 $t = {t}_{0}$ 处取小值，则当 $0 < {t}_{0} < \frac{1}{5}$ 时， $\alpha$ 的取值范围是\_\_\_\_\_.

8. 已知 $P\left( {0,\sqrt{3}}\right)$ ,将函数 $f\left( x\right)  = 4\cos \frac{\pi }{2}x$ 和直线 $g\left( x\right)  = x - 1$ 的所有交点从左到右依次记为 ${A}_{1}$ 、 ${A}_{2}\text{ 、 }\cdots \text{ 、 }{A}_{n}$ ,则 $\left| {\overrightarrow{P{A}_{1}} + \overrightarrow{P{A}_{2}} + \cdots  + \overrightarrow{P{A}_{n}}}\right|  =$ \_\_\_\_\_.

9. 在平面直角坐标系 ${xOy}$ 中,角 $\alpha$ 的顶点在原点,始边与 $x$ 轴非负半轴重合,终边与单位圆交于点 $A$ , 角 $\alpha  + \beta$ 的终边交单位圆于点 $B$ ，且 $\alpha$ 、 $\beta  \in  \left( {0,\frac{\pi }{2}}\right)$ . 记 $A\left( {{x}_{1},{y}_{1}}\right)$ ， $B\left( {{x}_{2},{y}_{2}}\right)$ . 若 ${x}_{2} =  - \frac{1}{4}$ ，且 $\overrightarrow{OA} \cdot  \overrightarrow{OB} = \frac{1}{2}$ ， 那么 $\sqrt{3}{x}_{1} + {y}_{1} =$ \_\_\_\_\_.

\section*{二、选择题}

10. 若 $\overrightarrow{a}\text{ 、 }\overrightarrow{b}$ 为不平行的非零向量, $\overrightarrow{AB} = \overrightarrow{a} + 5\overrightarrow{b},\overrightarrow{BC} =  - 2\overrightarrow{a} + 8\overrightarrow{b},\overrightarrow{CD} = 3\overrightarrow{a} - 3\overrightarrow{b}$ ,则 ( )

A. $A\text{ 、 }B\text{ 、 }C$ 三点共线. B. $A\text{ 、 }B\text{ 、 }D$ 三点共线. C. $B\text{ 、 }C\text{ 、 }D$ 三点共线. D. $A\text{ 、 }C\text{ 、 }D$ 三点共线.

11. 已知椭圆 $\frac{{x}^{2}}{16} + \frac{{y}^{2}}{12} = 1$ 的左顶点为 $A$ ,右焦点为 $F,M$ 是椭圆上任意一点,则 $\overrightarrow{MA} \cdot  \overrightarrow{MF}$ 的取值范围为 ( )

A. $\left\lbrack  {-{16},0}\right\rbrack$ . B. $\left\lbrack  {-8,0}\right\rbrack$ . C. $\left\lbrack  {0,8}\right\rbrack$ . D. $\left\lbrack  {0,{16}}\right\rbrack$ .

12. 设向量 $\overrightarrow{a}$ 与 $\overrightarrow{b}$ 的夹角为 $\theta$ ，定义 $\overrightarrow{a} \oplus  \overrightarrow{b} = \left| {\overrightarrow{a}\sin \theta  - \overrightarrow{b}\cos \theta }\right|$ ，已知 $\left| \overrightarrow{b}\right|  = \left| {\overrightarrow{a} - \overrightarrow{b}}\right|  = 1$ ， $\left| \overrightarrow{a}\right|  = \sqrt{2}$ ，则 $\overrightarrow{a} \oplus  \overrightarrow{b} =$ ( )

A. $\frac{\sqrt{2}}{2}$ . B. $\sqrt{2}$ . C. $\frac{\sqrt{3}}{2}$ . D. $\sqrt{3}$ .

\section*{三、解答题}

13. 已知 $\left| \overrightarrow{a}\right|  = 2,\left| \overrightarrow{b}\right|  = 3,\left( {2\overrightarrow{a} + \overrightarrow{b}}\right)  \cdot  \left( {\overrightarrow{a} - 2\overrightarrow{b}}\right)  =  - {19}$ .

(1)求向量 $\overrightarrow{a}$ 与 $\overrightarrow{b}$ 的夹角；

(2)求 $\left| {\overrightarrow{a} + 2\overrightarrow{b}}\right|$ .

14. 在平面直角坐标系中, $O$ 为坐标原点, $A\left( {1,1}\right) ,B\left( {2, - 3}\right)$ .

(1)若 $\overrightarrow{{OA} \bot  }$ 上( $\overrightarrow{OA} + \lambda \overrightarrow{AB}$ )，求实数 $\lambda$ 的值；

(2)设 $C\left( {-6,k}\right)$ ，若 $\overrightarrow{AB}$ 、 $\overrightarrow{BC}$ 的夹角为钝角，求实数 $k$ 的取值范围.

15. 在平面直角坐标系中,已知 $\overrightarrow{a} = \left( {1, - 2}\right) ,\overrightarrow{b} = \left( {3,4}\right)$ .

(1)若 $\left( {3\overrightarrow{a} - \overrightarrow{b}}\right) //\left( {\overrightarrow{a} + k\overrightarrow{b}}\right)$ ，求实数 $k$ 的值；

(2) 求 $3\overrightarrow{a} - \overrightarrow{b}$ 与 $\overrightarrow{a} + 3\overrightarrow{b}$ 的夹角.

16. 在 $\bigtriangleup {ABC}$ 中,设 $\overrightarrow{AC} = \overrightarrow{a},\overrightarrow{AB} = \overrightarrow{b}$ ,若 $\left| \overrightarrow{a}\right|  = 2,\left| \overrightarrow{b}\right|  = 3,\angle {BAC} = {60}^{ \circ  },\overrightarrow{AD} = 2\overrightarrow{DB},E$ 为 ${BC}$ 的中点， ${CD}$ 与 ${AE}$ 交于点 $O$ .

(1)求 $\overrightarrow{AE} \cdot  \overrightarrow{CD}$ 的值；

(2)若 $\overrightarrow{OD} = \lambda \overrightarrow{CD}$ ，求 $\lambda$ 的值.

\section*{§ 8.2 平面向量及其应用一}

\section*{典型例题精讲}

【例 1】已知点 $A\left( {1,4}\right) ,B\left( {5,2}\right) ,C\left( {2,2}\right)$ ,点 $P$ 是直线 ${OC}$ 上的动点 $(O$ 为坐标原点), $\overrightarrow{PA}//\overrightarrow{PB}$ . (1)求 $\overrightarrow{OP}$ 的坐标；

(2)求 $\overrightarrow{OP}$ 在 $\overrightarrow{AB}$ 方向上的向量投影.

【例 2】在 $\bigtriangleup {ABC}$ 中, ${AB} = 2,{AC} = 4,\cos \angle {BAC} = \frac{1}{4},D$ 是线段 ${BC}$ 的靠近点 $B$ 的三等分点.

(1)用 $\overrightarrow{AB}$ 、 $\overrightarrow{AC}$ 表示 $\overrightarrow{AD}$ ；

(2)求 ${AD}$ 的长度.

\begin{center}
\includegraphics[width=0.3\textwidth]{images/bo_d7tk7oc91nqc73875fu0_148_1100_1297_387_166_0.jpg}
\end{center}
 

(例 2 图)

【例 3】已知 $\bigtriangleup {ABC}$ 中, $\angle {BAC} = {90}^{ \circ  }$ . 以 ${AB}$ 为一边向外作等边三角形 ${ABD}$ (如图所示),且 $\overrightarrow{CD} =  \lambda \overrightarrow{AB} + \mu \overrightarrow{AC}\left( {\lambda ,\mu  \in  \mathbf{R}}\right)$ .

(1)当 ${AB} = {AC}$ 时，求 $\lambda  + \mu$ 的值；

(2)当 $\angle {BCD} = 2\angle {ACD}$ 时，求 $\lambda  - \mu$ 的值.

\begin{center}
\includegraphics[width=0.2\textwidth]{images/bo_d7tk7oc91nqc73875fu0_149_1114_325_338_231_0.jpg}
\end{center}
 

(例 3 图)

【例 4】已知向量 $\overrightarrow{a} = \left( {2\sin \theta ,\sin \theta  + \cos \theta }\right) ,\overrightarrow{b} = \left( {\cos \theta ,2 - m}\right)$ ,函数 $f\left( \theta \right)  = \overrightarrow{a} \cdot  \overrightarrow{b}$ .

(1)当 $m = 0$ 时，求 $f\left( \frac{\pi }{6}\right)$ 的值；

(2)若不等式 $f\left( \theta \right)  + \frac{4}{\sin \theta  + \cos \theta } > {2m} - 3$ 对所有 $\theta  \in  \left\lbrack  {0,\frac{\pi }{2}}\right\rbrack$ 恒成立. 求实数 $m$ 的范围.

【例 5】在 $\bigtriangleup {ABC}$ 中,点 $P$ 为 $\bigtriangleup {ABC}$ 内一点.

(1)记 $\bigtriangleup {PBC}$ 、 $\bigtriangleup {PAC}$ 、 $\bigtriangleup {PAB}$ 的面积分别为 ${S}_{A}$ 、 ${S}_{B}$ 、 ${S}_{C}$ ，求证: ${S}_{A}\overrightarrow{PA} + {S}_{B}\overrightarrow{PB} + {S}_{C}\overrightarrow{PC} = \overrightarrow{0}$ ；

(2)若点 $P$ 为 $\bigtriangleup {ABC}$ 的垂心，且 $\overrightarrow{PA} + 2\overrightarrow{PB} + 3\overrightarrow{PC} = \overrightarrow{0}$ ，求 $\cos \angle {APB}$ .

\begin{center}
\includegraphics[width=0.2\textwidth]{images/bo_d7tk7oc91nqc73875fu0_150_1191_290_289_237_0.jpg}
\end{center}
 

(例 5 图)

\section*{四基达标精练}

\section*{一、填空题}

1. 如图，在 $\bigtriangleup  {ABC}$ 中， $D$ 为 ${AB}$ 的中点， $E$ 为 ${CD}$ 的中点，设 $\overrightarrow{AB} = \overrightarrow{a}$ ， $\overrightarrow{AC} = \overrightarrow{b}$ ，以向量 $\overrightarrow{a}$ 、 $\overrightarrow{b}$ 为基底，则向量 $\overrightarrow{AE} =$

\begin{center}
\includegraphics[width=0.2\textwidth]{images/bo_d7tk7oc91nqc73875fu0_150_185_1185_327_239_0.jpg}
\end{center}
 

(第 1 题图)

\begin{center}
\includegraphics[width=0.3\textwidth]{images/bo_d7tk7oc91nqc73875fu0_150_636_1223_352_197_0.jpg}
\end{center}
 

(第 2 题图)

\begin{center}
\includegraphics[width=0.2\textwidth]{images/bo_d7tk7oc91nqc73875fu0_150_1113_1232_316_182_0.jpg}
\end{center}
 

(第 3 题图)

2. 如图，在 $\bigtriangleup  {ABC}$ 中， ${BD} = {3DC}$ ，若 $\overrightarrow{AD} = \lambda \overrightarrow{AB} + \mu \overrightarrow{AC}$ ，则 $\frac{\lambda }{\mu } = {\_\_\_\_\_ }$ \_\_\_\_\_.

3. 如图，在平行四边形 ${ABCD}$ 中， ${AB} = 2,\cos \angle {BAD} = \frac{1}{2}$ ， $E$ 是边 ${BC}$ 的中点， $F$ 是 ${CD}$ 上靠近 $D$ 的三等分点,若 $\overrightarrow{AE} \cdot  \overrightarrow{BF} = 8$ ,则 $\left| \overrightarrow{AD}\right|  =$ \_\_\_\_\_.

4. 已知向量 $\overrightarrow{a}$ 、 $\overrightarrow{b}$ 的夹角为 ${120}^{ \circ  }$ ， $\left| \overrightarrow{a}\right|  = 4$ ， $\left| \overrightarrow{b}\right|  = 2$ ，则向量 $\overrightarrow{b}$ 在向量 $\overrightarrow{a}$ 方向上的数量投影为\_\_\_\_\_.

5. 已知 $\left| \overrightarrow{OM}\right|  = 2,\left| \overrightarrow{ON}\right|  = 2$ ,且 $\overrightarrow{OM}$ 、 $\overrightarrow{ON}$ 的夹角为 $\frac{\pi }{3}$ ,点 $P$ 在以 $O$ 为圆心的圆弧 $\overset{\text{ ⏜ }}{MN}$ 上运动,若 $\overrightarrow{OP} =  x\overrightarrow{OM} + y\overrightarrow{ON}\left( {x\text{ 、 }y \geq  0}\right)$ ,则 $x + y$ 的取值范围为\_\_\_\_\_.

6. 已知 $O$ 是 $\bigtriangleup {ABC}$ 内一点, $2\overrightarrow{OA} + 3\overrightarrow{OB} + m\overrightarrow{OC} = \overrightarrow{0}$ ,若 $\bigtriangleup {AOB}$ 与 $\bigtriangleup {ABC}$ 的面积之比为 $\frac{4}{7}$ ,则实数 $m$ 的值为\_\_\_\_\_.

7. 在 $\bigtriangleup {ABC}$ 中,内角 $A\text{ 、 }B\text{ 、 }C$ 的对边分别是 $a\text{ 、 }b\text{ 、 }c$ . 若 ${20a}\overrightarrow{BC} + {15b}\overrightarrow{CA} + {12c}\overrightarrow{AB} = \overrightarrow{0}$ ,则 $\bigtriangleup {ABC}$ 中最小角的余弦值等于\_\_\_\_\_.

8. 如图，四个边长为 1 的等边三角形有一条边在同一条直线上，边 ${B}_{4}{C}_{4}$ 上有 10 个不同的点 ${P}_{1}$ 、 ${P}_{2}$ 、...、 ${P}_{10}$ ，记 ${m}_{i} = \overrightarrow{A{B}_{2}} \cdot  \overrightarrow{A{P}_{i}}\left( {i = 1,2,\cdots ,{10}}\right)$ ，则 ${m}_{1} + {m}_{2} + \cdots  + {m}_{10} =$ \_\_\_\_\_.

\begin{center}
\includegraphics[width=0.5\textwidth]{images/bo_d7tk7oc91nqc73875fu0_151_220_281_603_191_0.jpg}
\end{center}
 

(第 8 题图)

\begin{center}
\includegraphics[width=0.3\textwidth]{images/bo_d7tk7oc91nqc73875fu0_151_950_246_372_221_0.jpg}
\end{center}
 

(第 13 题图)

9. 窗花是贴在窗纸或窗户玻璃上的剪纸，它是中国古老的传统民间艺术之一. 在 2022 年虎年新春来临之际，人们设计了一种由外围四个大小相等的半圆和中间正方形所构成的剪纸窗花(如图 1). 已知正方形 ${ABCD}$ 的边长为 2,中心为 $O$ ,四个半圆的圆心均为正方形 ${ABCD}$ 各边的中点 (如图 2),若 $P$ 在 $\overset{\text{ ⏜ }}{BC}$ 的中点, 则 $\left( {\overrightarrow{PA} + \overrightarrow{PB}}\right)  \cdot  \overrightarrow{PO} =$ \_\_\_\_\_.

\section*{二、选择题}

10. 已知平面向量 $\overrightarrow{a} = \left( {1,3}\right) ,\overrightarrow{b} = \left( {-2,4}\right)$ ,则 $\overrightarrow{a}$ 在 $\overrightarrow{b}$ 上的投影为 ( )

A. $\left( {1,3}\right)$ . B. $\left( {-1,2}\right)$ . C. $\sqrt{10}$ . D. $\sqrt{5}$ .

11. 已知单位向量 $\overrightarrow{a}$ 、 $\overrightarrow{b}$ ，若对任意实数 $x$ ， $\left| {x\overrightarrow{a} + \overrightarrow{b}}\right|  \geq  \frac{\sqrt{3}}{2}$ 恒成立，则向量 $\overrightarrow{a}$ 、 $\overrightarrow{b}$ 的夹角的取值范围为\_\_\_\_\_ \ding{115} \_\_\_\_\_.

A. $\left\lbrack  {\frac{\pi }{4},\frac{3\pi }{4}}\right\rbrack$ . B. $\left\lbrack  {\frac{\pi }{3},\frac{2\pi }{3}}\right\rbrack$ . C. $\left\lbrack  {\frac{\pi }{4},\frac{\pi }{2}}\right\rbrack$ . D. $\left\lbrack  {\frac{\pi }{3},\frac{\pi }{2}}\right\rbrack$ .

\begin{center}
\includegraphics[width=0.2\textwidth]{images/bo_d7tk7oc91nqc73875fu0_151_1190_1010_257_271_0.jpg}
\end{center}
 

(第 12 题图)

12. 如图，在 $\bigtriangleup  {ABC}$ 中， $M$ 为线段 ${BC}$ 的中点， $G$ 为线段 ${AM}$ 上一点，且 $\overrightarrow{AG} =  2\overrightarrow{GM}$ ,过点 $G$ 的直线分别交直线 ${AB}\text{ 、 }{AC}$ 于 $P\text{ 、 }Q$ 两点, $\overrightarrow{AB} = x\overrightarrow{AP}\left( {x > 0}\right) ,\overrightarrow{AC} =  y\overrightarrow{AQ}\left( {y > 0}\right)$ ,则 $\frac{1}{x} + \frac{1}{y + 1}$ 的最小值为 ( )

A. $\frac{3}{4}$ . B.1. C. $\frac{4}{3}$ . D. 4.

\section*{三、解答题}

13. 在直角梯形 ${ABCD}$ 中,已知 ${AB}//{CD},\angle {DAB} = {90}^{ \circ  },{AB} = {2AD} = {2CD} = 2$ ,点 $F$ 是 ${BC}$ 边上的中点,点 $E$ 是 ${CD}$ 边上一个动点.

(1)若 $\overrightarrow{DE} = \frac{1}{2}\overrightarrow{DC}$ ，求 $\overrightarrow{AC} \cdot  \overrightarrow{EF}$ 的值；

(2)求 $\overrightarrow{EA}$ 、 $\overrightarrow{EF}$ 的取值范围.

\begin{center}
\includegraphics[width=0.2\textwidth]{images/bo_d7tk7oc91nqc73875fu0_151_1138_1529_315_205_0.jpg}
\end{center}
 

(第 13 题图)

14. 如图,在平面四边形 ${ABCD}$ 中, ${BC}//{AD},{AB} = {BC} = 2,{AD} = 4,\angle {BAD} = {120}^{ \circ  },E\text{ 、 }F$ 分别是 ${AD}$ 、 ${DC}$ 的中点, $G$ 为线段 ${BC}$ 上一点,且 $\overrightarrow{BG} = \lambda \overrightarrow{BC}$ . 设 $\overrightarrow{AB} = \overrightarrow{a},\overrightarrow{AD} = \overrightarrow{b}$ .

(1)若 $\lambda  = \frac{1}{3}$ ，以 $\overrightarrow{a}$ 、 $\overrightarrow{b}$ 为基底表示向量 $\overrightarrow{AF}$ 与 $\overrightarrow{EG}$ ；

( 2 )若 $\lambda  \in  \left( {0,1}\right)$ ，求 $\overrightarrow{AF}$ 。 $\overrightarrow{EG}$ 的取值范围.

\begin{center}
\includegraphics[width=0.3\textwidth]{images/bo_d7tk7oc91nqc73875fu0_152_1083_359_405_175_0.jpg}
\end{center}
 

(第 14 题图)

15. 已知 $\bigtriangleup  {ABC}$ 中，角 $A$ 、 $B$ 、 $C$ 所对的边分别是 $a$ 、 $b$ 、 $c$ ，向量 $\overrightarrow{m} = \left( {b + a + c\text{ ， }\sqrt{3}b}\right)$ ， $\overrightarrow{n} = \left( {\sqrt{3}c,b - a + c}\right)$ ， 且 $\overrightarrow{m}//\overrightarrow{n}$ .

(1)若 $\overrightarrow{BD} = 3\overrightarrow{DC},{AD} = 1$ ，求 $\bigtriangleup  {ABC}$ 面积的最大值；

(2)若 $\bigtriangleup  {ABC}$ 为锐角三角形,且 $a = 3$ ，求 $\bigtriangleup  {ABC}$ 周长的取值范围.

16. 如图，在 $\bigtriangleup  {ABC}$ 中，已知 ${AC} = 1,{AB} = 3,\angle {BAC} = {60}^{ \circ  }$ ，且 $\overrightarrow{PA} + \overrightarrow{PB} + \overrightarrow{PC} = \overrightarrow{0}$ .

(1)求 ${PA}$ 的长；

(2)求 $\cos \angle {APC}$ .

\begin{center}
\includegraphics[width=0.3\textwidth]{images/bo_d7tk7oc91nqc73875fu0_152_1111_1622_377_151_0.jpg}
\end{center}
 

(第 16 题图)

$\text{ § }{8.3}$ 平面向量及其应用二

\section*{典型例题精讲}

【例 1】已知向量 $\overrightarrow{OA} = \left( {0, - 6}\right) ,\overrightarrow{OB} = \left( {8,0}\right) ,\overrightarrow{OC} = \left( {x,y}\right)$ .

(1)若点 $A$ 、 $B$ 、 $C$ 三点共线，求实数 $x$ 、 $y$ 满足的关系；

(2)若 $x = 1$ ，且 $\angle {ACB}$ 为钝角，求实数 $y$ 的取值范围.

【例 2】如图,已知 $\bigtriangleup {ABC}$ 的外接圆 $O$ 的半径为 $4,\overrightarrow{AO} = 2\overrightarrow{AC} + \overrightarrow{AB}$ .

(1)求 $\bigtriangleup  {ABC}$ 中 ${AC}$ 边的长；

(2)求 $\overrightarrow{OC} \cdot  \overrightarrow{AB}$ .

\begin{center}
\includegraphics[width=0.2\textwidth]{images/bo_d7tk7oc91nqc73875fu0_153_1158_1256_296_295_0.jpg}
\end{center}
 

(例 2 图)

【例 3】如图,一条河两岸平行,河的宽度 ${AC} = \sqrt{3}{{km}}$ ,一艘船从河边的点 $A$ 出发到达对岸的点 $B$ ,船只在河内行驶的路程 ${AB} = 2{{km}}$ ，行驶时间为 0.2 h. 已知船在静水中的速度 ${v}_{1}$ 的大小为 $\left| {v}_{1}\right|$ ，水流的速度 ${v}_{2}$ 的大小为 $\left| {v}_{2}\right|  = 2{{km}}/{h}$ . 求:

(1) $\left| {v}_{1}\right|$ ；

(2)船在静水中速度 ${v}_{1}$ 与水流速度 ${v}_{2}$ 夹角的余弦值.

\begin{center}
\includegraphics[width=0.2\textwidth]{images/bo_d7tk7oc91nqc73875fu0_154_1163_365_328_207_0.jpg}
\end{center}
 

(例 3 图)

【例 4】在平面直角坐标系 ${xOy}$ 中,已知点 $A\left( {-1,0}\right) ,B\left( {0,\sqrt{3}}\right) ,C\left( {\cos \theta ,\sin \theta }\right)$ ,其中 $\theta  \in  \left\lbrack  {0,\frac{\pi }{2}}\right\rbrack$ . (1)求 $\overrightarrow{AC} \cdot  \overrightarrow{BC}$ 的最大值；

(2)是否存在 $\theta  \in  \left\lbrack  {0,\frac{\pi }{2}}\right\rbrack$ ，使得 $\bigtriangleup  {ABC}$ 为钝角三角形？若存在，求出 $\theta$ 的取值范围；若不存在，说明理由.

【例 5】已知在平面直角坐标系中， $O$ 为坐标原点，定义非零向量 $\overrightarrow{OM} = \left( {a,b}\right)$ 的“相伴函数”为 $y =  a\sin x + b\cos x\left( {x \in  \mathbf{R}}\right)$ ,向量 $\overrightarrow{OM} = \left( {a,b}\right)$ 称为函数 $y = a\sin x + b\cos x\left( {x \in  \mathbf{R}}\right)$ 的“相伴向量”; 记平面内所有向量的“相伴函数”构成的集合为 $S$ .

(1)已知 $\alpha  \in  \mathbf{R}$ ， $h\left( x\right)  = \cos \left( {x + \alpha }\right)  + 2\cos x$ ，若函数 $h\left( x\right)$ 为集合 $S$ 中的元素，求其 “相伴向量” 的模的取值范围;

(2)当向量 $\overrightarrow{OM} = \left( {\sqrt{3},1}\right)$ 时，“相伴函数”为 $f\left( x\right)$ ，若 $x \in  \left\lbrack  {0,\frac{11\pi }{6}}\right\rbrack$ ，方程 ${f}^{2}\left( x\right)  + \left( {2 - a}\right) f\left( x\right)  + a - 3 = 0$ 存在 4 个不相等的实数根,求实数 $a$ 的取值范围.

\section*{四基达标精练}

\section*{一、填空题}

1. 向量 $\overrightarrow{a} = \left( {3,4}\right)$ 在向量 $\overrightarrow{b} = \left( {1,0}\right)$ 方向上的数量投影为\_\_\_\_\_.

2. 已知 $G$ 为 $\bigtriangleup  {ABC}$ 的重心,用向量 $\overrightarrow{AB}$ 、 $\overrightarrow{BC}$ 表示向量 $\overrightarrow{AG} =$ \_\_\_\_\_.

3. 已知平面单位向量 $\overrightarrow{a}\text{ 、 }\overrightarrow{b}\text{ 、 }\overrightarrow{c}$ 满足 $\langle \overrightarrow{a},\overrightarrow{b}\rangle  = \langle \overrightarrow{b},\overrightarrow{c}\rangle  = \langle \overrightarrow{c},\overrightarrow{a}\rangle  = \frac{2\pi }{3}$ ，则 $\left| {3\overrightarrow{a} + 2\overrightarrow{b} + \overrightarrow{c}}\right|  =$ \_\_\_\_\_.

4. 如图所示，两块斜边长均等于 $\sqrt{2}$ 的直角三角板拼在一起，则 $\overrightarrow{OD}$ . $\overrightarrow{AB} =$ \_\_\_\_\_.

\begin{center}
\includegraphics[width=0.2\textwidth]{images/bo_d7tk7oc91nqc73875fu0_155_259_604_343_177_0.jpg}
\end{center}
 

(第 4 题图)

\begin{center}
\includegraphics[width=0.3\textwidth]{images/bo_d7tk7oc91nqc73875fu0_155_815_549_447_226_0.jpg}
\end{center}
 

(第 8 题图)

5. 已知非零向量 $\overrightarrow{a}$ 、 $\overrightarrow{b}$ 、 $\overrightarrow{c}$ 两两不平行，且 $\overrightarrow{a}//\left( {\overrightarrow{b} + \overrightarrow{c}}\right)$ ，方 $//\left( {\overrightarrow{a} + \overrightarrow{c}}\right)$ ，设 $\overrightarrow{c} = {x\overrightarrow{a}} + {y\overrightarrow{b}}\left( {x\text{ 、 }y \in  \mathbf{R}}\right)$ ，则 $x + {2y}  =$ \_\_\_\_\_.

6. 一条东西方向的河流两岸平行，河宽 ${250}\sqrt{3}{m}$ ，河水的速度为向东 $2{{km}}/{h}$ . 一艘小货船准备从河南岸的码头 $A$ 处出发,航行到位于河对岸 $B$ ( ${AB}$ 与河的方向垂直)的正西方向并且与 $B$ 相距 ${250}{m}$ 的码头 $C$ 处卸货. 若流水的速度与小货船航行的速度的合速度为 $6{{km}}/{h}$ ，则当小货船的航程最短时，小货船航行的速度大小是\_\_\_\_\_km/h.

7. 已知点 $A$ 、 $B$ 在圆 $O : {x}^{2} + {y}^{2} = {16}$ 上,且 $\left| {AB}\right|  = 4,P$ 为圆 $O$ 上任意一点,则 $\overrightarrow{AB} \cdot  \overrightarrow{BP}$ 的最小值为\_\_\_\_\_.

8. 青花瓷，又称白地青花瓷，常简称青花，是中国瓷器的知名品种之一. 如图 1，这是一个青花瓷圆盘. 该圆盘中的两个圆的圆心重合,如图 2,其中大圆半径 $R = 3$ ,小圆半径 $r = 2$ ,点 $P$ 在大圆上,过点 $P$ 作小圆的切线，切点分别是 $E$ 、 $F$ ，则 $\overrightarrow{PE}$ . $\overrightarrow{PF} =$ \_\_\_\_\_.

9. 已知平面向量 $\overrightarrow{a}$ 、 $\overrightarrow{b}$ 、 $\overrightarrow{c}$ 满足 $\left| \overrightarrow{a}\right|  = \left| \overrightarrow{b}\right|  = \overrightarrow{a} \cdot  \overrightarrow{b} = 2$ ，且 $\left( {\overrightarrow{b} - \overrightarrow{c}}\right)  \cdot  \left( {2\overrightarrow{b} - \overrightarrow{c}}\right)  = 0$ ，则 $\left| {\overrightarrow{a} - 2\overrightarrow{c}}\right|$ 的最大值为\_\_\_\_\_.

\section*{二、选择题}

10. 设非零向量 $\overrightarrow{a}$ 、 $\overrightarrow{b}$ 的夹角为 $\theta$ ，定义运算 $\overrightarrow{a} * \overrightarrow{b} = \left| \overrightarrow{a}\right|  \cdot  \left| \overrightarrow{b}\right|  \cdot  \sin \theta$ . 下列叙述正确的是 ( )

A. 若 $\overrightarrow{a} * \overrightarrow{b} = 0$ ，则 $\overrightarrow{a} \bot  \overrightarrow{b}$ .

B. 若 $\left| \overrightarrow{a}\right|  \cdot  \left| \overrightarrow{b}\right|  = 1$ ，则 ${\left( \overrightarrow{a} * \overrightarrow{b}\right) }_{\min } =  - 1$ .

C. 设在 $\bigtriangleup {ABC}$ 中, $\overrightarrow{AB} = \overrightarrow{a},\overrightarrow{AC} = \overrightarrow{b}$ ,则 $2{S}_{\bigtriangleup {ABC}} = \overrightarrow{a} * \overrightarrow{b}$ .

D. $\overrightarrow{a} * \left( {\overrightarrow{b} + \overrightarrow{c}}\right)  = \overrightarrow{a} * \overrightarrow{b} + \overrightarrow{a} * \overrightarrow{c}\left( \overrightarrow{c}\right)$ 为任意非零向量).

11. 已知点 ${A}_{1},{A}_{2},\cdots ,{A}_{n},\cdots$ 和数列 $\left\{  {a}_{n}\right\}$ 、 $\left\{  {b}_{n}\right\}$ 满足 $\overrightarrow{{A}_{n}{A}_{n + 1}} = \left( {\cos \frac{2n\pi }{3},\sin \frac{2n\pi }{3}}\right) \left( {n \in  {\mathbf{N}}^{ * }}\right)$ ,若 ${a}_{1} = 1$ , ${S}_{n}\text{ 、 }{T}_{n}$ 分别为数列 $\left\{  {a}_{n}\right\}  \text{ 、 }\left\{  {b}_{n}\right\}$ 的前 $n$ 项和,则 ${S}_{60} + 2{T}_{60}$ 等于 ( )

A. -20 . B. ${24}\sqrt{3}$ . C. ${48}\sqrt{3} - {20}$ . D.O.

12. 如图， ${Ox}$ 、 ${Oy}$ 是平面内相交成 $\theta$ ( $\theta  \neq  \frac{\pi }{2}$ )角的两条数轴， $\overrightarrow{{e}_{1}}$ 、 $\overrightarrow{{e}_{2}}$ 分别是与 $x\text{ 、 }y$ 轴正方向同向的单位向量，则称 ${xOy}$ 为 $\theta$ 斜坐标系，若 $\overrightarrow{OM} = x\overrightarrow{{e}_{1}} + y\overrightarrow{{e}_{2}}$ ，则把有序数对 $\left( {x,y}\right)$ 叫做向量 $\overrightarrow{OM}$ 的斜坐标,记为 $\overrightarrow{OM} = \left( {x,y}\right)$ . 在 $\frac{\pi }{4}$ 斜坐标系中, $\overrightarrow{a} = \left( {\frac{1}{2},\frac{\sqrt{3}}{2}}\right) ,\overrightarrow{b} = \left( {\sqrt{3}, - 1}\right)$ . 则下列结论中,正确的是 ( )

\begin{center}
\includegraphics[width=0.2\textwidth]{images/bo_d7tk7oc91nqc73875fu0_156_1177_141_310_243_0.jpg}
\end{center}
 

(第 12 题图)

A. $\overrightarrow{a} - \overrightarrow{b} = \left( {\frac{1}{2} - \sqrt{3},\frac{\sqrt{3}}{2} + 1}\right)$ . B. $\left| \overrightarrow{a}\right|  = 1$ .

C. $\overrightarrow{a} \bot  \overrightarrow{b}$ . D. $\overrightarrow{b}$ 在 $\overrightarrow{a}$ 上的投影向量为 $\left( {\frac{2\sqrt{2} + \sqrt{3}}{5},\frac{2\sqrt{6} + 2}{5}}\right)$ .

\section*{三、解答题}

13. 在 $\bigtriangleup {ABC}$ 中,已知 $\overrightarrow{AB} \bot  \overrightarrow{AC},M$ 是 ${BC}$ 的中点.

(1)若 $\left| \overrightarrow{AB}\right|  = \left| \overrightarrow{AC}\right|$ ，求向量 $\overrightarrow{AB} + 2\overrightarrow{AC}$ 与向量 $2\overrightarrow{AB} + \overrightarrow{AC}$ 的夹角的余弦值；

(2)若 $O$ 是线段 ${AM}$ 上任意一点，且 $\left| \overrightarrow{AB}\right|  = \left| \overrightarrow{AC}\right|  = \sqrt{2}$ ，求 $\overrightarrow{OA} \cdot  \overrightarrow{OB} + \overrightarrow{OC} \cdot  \overrightarrow{OA}$ 的最小值.

14. 已知 $A$ 、 $B$ 、 $C$ 为 $\bigtriangleup {ABC}$ 的三个内角,向量 $\overrightarrow{m} = \left( {2 - 2\sin A,\sin A + \cos A}\right)$ 平行于 $\overrightarrow{n} = (\sin A - \cos A$ ， $1 + \sin A)$ ,且 $\overrightarrow{AB} \cdot  \overrightarrow{AC} > 0$ .

(1)求角 $A$ ；

(2)求函数 $y = 2{\sin }^{2}\frac{B}{2} + \cos \frac{C - B}{2}$ 的值域.


15. 假设在一个以 ${m}$ 为单位的空间直角坐标系 $O - {xyz}$ 中,平面 ${xOy}$ 内有一跟踪和控制飞行机器人 $T$ 的控制台 $A,A$ 的位置为 $\left( {{170},{200},0}\right)$ ,上午 10 时 07 分测得飞行机器人 $T$ 在 $P\left( {{150},{80},{120}}\right)$ 处,并对飞行机器人 $T$ 发出指令:以速度 ${v}_{1} = {13}{m}/{s}$ 沿单位向量 ${\overrightarrow{d}}_{1} = \left( {\frac{3}{13},\frac{12}{13}, - \frac{4}{13}}\right)$ 作匀速直线飞行 (飞行中无障碍物)， ${10}{s}$ 后到达点 $Q$ ，再发出指令让机器人在点 $Q$ 原地盘旋 $2{s}$ ，在原地盘旋过程中逐步减速并降速到 $8{m}/{s}$ ，然后保持 $8{m}/{s}$ ，再沿单位向量 ${\overrightarrow{d}}_{2} = \left( {\frac{1}{2}, - \frac{\sqrt{2}}{2}, - \frac{1}{2}}\right)$ 作匀速直线飞行 (飞行中无障碍物),当飞行机器人 $T$ 最终落在平面 ${xOy}$ 内发出指令让它停止运动,机器人 $T$ 近似看成一个点.

(1)求从 $P$ 点开始出发 ${20}{s}$ 后飞行机器人 $T$ 的位置；

(2)求在整个飞行过程中飞行机器人 $T$ 与控制台 $A$ 的最近距离(精确到米).

\begin{center}
\includegraphics[width=0.2\textwidth]{images/bo_d7tk7oc91nqc73875fu0_157_1184_572_259_260_0.jpg}
\end{center}
 

(第 15 题图)

16. 如图所示，在 $\bigtriangleup  {ABO}$ 中， $\overrightarrow{OC} = \frac{1}{3}\overrightarrow{OA}$ ， $\overrightarrow{OD} = \frac{1}{2}\overrightarrow{OB}$ ， ${AD}$ 与 ${BC}$ 相交于点 $M$ . 设 $\overrightarrow{OA} = \overrightarrow{a}$ ， $\overrightarrow{OB} = \overrightarrow{b}$ .

(1)试用向量 $\overrightarrow{a}$ 、 $\overrightarrow{b}$ 表示 $\overrightarrow{OM}$ ；

(2)在线段 ${AC}$ 上取点 $E$ ，在线段 ${BD}$ 上取点 $F$ ，使 ${EF}$ ， 过点 $M$ . 设 $\overrightarrow{OE} = \lambda \overrightarrow{OA},\overrightarrow{OF} = \mu \overrightarrow{OB}$ ，其中 $\lambda$ 、 $\mu  \in  \mathbf{R}$ .

当 ${EF}$ 与 ${AD}$ 重合时, $\lambda  = 1,\mu  = \frac{1}{2}$ ,此时 $\frac{1}{\lambda } + \frac{2}{\mu } = 5$ ;

当 ${EF}$ 与 ${BC}$ 重合时, $\lambda  = \frac{1}{3},\mu  = 1$ ,此时 $\frac{1}{\lambda } + \frac{2}{\mu } = 5$ ;

能否由此得出一般结论: 不论 $E\text{ 、 }F$ 在线段 ${AC}\text{ 、 }{BD}$ 上如何变动,等式 $\frac{1}{\lambda } + \frac{2}{\mu } = 5$ 恒成立,请说明理由.

\begin{center}
\includegraphics[width=0.3\textwidth]{images/bo_d7tk7oc91nqc73875fu0_157_1087_1563_363_261_0.jpg}
\end{center}
 

(第 16 题图)

\section*{$\text{ § }{8.4}$ 复数及其运算}

(本小节若无特别说明, i 表示虚数单位)

\section*{典型例题精讲}

【例 1】已知 $m \in  \mathbf{R}$ ,复数 $z = \frac{m\left( {m + 2}\right) }{m - 1} + \left( {{m}^{2} + {2m} - 3}\right) {i}$ ,当 $m$ 为何值时,

(1) $z \in  \mathbf{R}$ ；(2) $z$ 是虚数；(3) $z$ 是纯虚数；(4) $\bar{z} = \frac{1}{2} + {4i}$ .

【例 2】在复数范围内,下列命题中为真命题的序号是\_\_\_\_\_.

 \ding{172}  ${\left| \bar{z}\right| }^{2} = \left| {z}^{2}\right|$ ； \ding{173}  若 ${z}_{1} - {z}_{2} > 0$ ，则 ${z}_{1} > {z}_{2}$ ； \ding{174}  若 ${\left( {z}_{1} - {z}_{2}\right) }^{2} + {\left( {z}_{2} - {z}_{3}\right) }^{2} = 0$ ，则 ${z}_{1} = {z}_{2} = {z}_{3}$ ；

 \ding{175}  $\left| {{z}_{1} - {z}_{2}}\right|  = \sqrt{{\left( {z}_{1} + {z}_{2}\right) }^{2} - 4{z}_{1}{z}_{2}}$ ;  \circled{5}  ${z}_{1}^{2} = {z}_{2}^{2}$ ，则 ${z}_{1} \cdot  {\bar{z}}_{1} = {z}_{2} \cdot  {\bar{z}}_{2}$ ;  \circled{6}  $2{z}_{1}{z}_{2} \leq  {z}_{1}^{2} + {z}_{2}^{2}$ ；

 \circled{7} 两个共轭复数的差是纯虚数； \circled{8} 若 $\left| {z + {i}}\right|  = \left| {z - {i}}\right|$ ，则 $z$ 必为实数.

【例 3】设复数 $z = a + b{i}\left( {a\text{ 、 }b \in  \mathbf{R}}\right) ,{z}_{1} = z + k{i},{z}_{2} = \bar{z} \cdot  k{i}$ (其中 $k \in  \mathbf{R}$ ).

(1)设 $a = b = \frac{1}{2}$ ，若 $\left| {z}_{1}\right|  = \left| {z}_{2}\right|$ ，求出实数 $k$ 的值；

(2)若复数 $z$ 满足条件:存在实数 $k$ ，使得 ${z}_{1}$ 与 ${z}_{2}$ 是某个实系数一元二次方程的两个虚数根，求符合条件的复数 $z$ 的模的取值范围.

【例 4】设 $\alpha \text{ 、 }\beta$ 为方程 ${x}^{2} + {2x} + t = 0\left( {t \in  \mathbf{R}}\right)$ 的两个根, $f\left( t\right)  = \left| \alpha \right|  + \left| \beta \right|$ .

(1)求 $f\left( t\right)$ 的解析式；

(2)证明关于 $t$ 的方程 $f\left( t\right)  = m$ ，当 $m > 2$ 时恰有两个不等的根，且两根之和为定值.

【例 5】已知 $z = m + 3 + 3\sqrt{3}{i}$ ,其中 $m \in  \mathbf{C}$ ,且 $\frac{m + 3}{m - 3}$ 为纯虚数.

(1)求 $m$ 对应点的轨迹；(2)求 $\left| z\right|$ 的最大值、最小值.

\section*{四基达标精练}

\section*{一、填空题}

1. 已知 $\frac{z - 3}{1 + {2z}} = {i}$ ,则 $z =$ \_\_\_\_\_.

2. 若 $1 + \sqrt{2}{i}$ 是关于 $x$ 的实系数方程 ${x}^{2} + {bx} + c = 0$ 的一个虚根,则 $c =$ \_\_\_\_\_.

3. 已知复数 $a + b{i} = \frac{3 + {i}}{1 - {i}}\left( {a\text{ 、 }b \in  \mathbf{R}}\right)$ ,则 $\left| {a + b{i}}\right|  =$ \_\_\_\_\_.

4. 在复数范围内分解因式: $4{x}^{3} - 8{x}^{2} + {10x} =$ \_\_\_\_\_.

5. 若复数 $z$ 满足 $\left| {z - 1 + \sqrt{3}{i}}\right|  = 3$ ，则 $\left| z\right|$ 的最大值为\_\_\_\_\_.

6. 已知关于 $x$ 的实系数方程 ${x}^{2} + \left( {1 - k}\right) x + {k}^{2} - 1 = 0$ 有两个虚根 ${x}_{1}\text{ 、 }{x}_{2}$ ,且 $\left| {x}_{1}\right|  + \left| {x}_{2}\right|  = 2$ ,则满足条件的实数 $k =$ \_\_\_\_\_.

7. 复数 ${z}_{1}$ 和 ${z}_{2}$ 满足 $\left| {{z}_{1} + 1}\right|  = \left| {{z}_{1} - 1}\right|$ ，且 $\left| {{z}_{2} - 2 - {3{i}}}\right|  = 1$ ，则 $\left| {{z}_{1} - {z}_{2}}\right|$ 的最小值是\_\_\_\_\_.

8. 已知关于 $x$ 的方程 ${x}^{2} + {zx} + 4 + {3{i}} = 0$ 有实数根，则复数 $z$ 的模的最小值为\_\_\_\_\_.

9. 已知 $k + 2$ 个两两互不相等的复数 ${z}_{1}\text{ 、 }{z}_{2}\text{ 、 }\cdots \text{ 、 }{z}_{k}\text{ 、 }{\omega }_{1}\text{ 、 }{\omega }_{2}$ ,满足,且 $\left| {{\omega }_{m} - {z}_{n}}\right|  \in  \{ 1,3\}$ (其中 $m = 1,2$ , $n = 1,2,\cdots ,k)$ ，则 $k$ 的最大值为\_\_\_\_\_.

\section*{二、选择题}

10. 若复数 $z = \left( {m + 2}\right)  + \left( {m - 4}\right) {i}$ 是虚数,则实数 $m$ 取值的集合是 ( )

A. $\{ m \mid  m > 4\}$ . B. $\{ m \mid  m < 4\}$ . C. $\{ m \mid  m \neq  4,m \in  \mathbf{R}\}$ . D.R.

D. $- 3 + {i}$ .

11. 在复平面内,平行四边形 ${ABCD}$ 的三个顶点, $A\text{ 、 }B\text{ 、 }C$ 对应的复数分别为 $- 1 + 2{i}\text{ 、 }3 - {i}\text{ 、 }1 + 2{i}$ ,则点 $D$ 对应的复数为 ( )

A. $- 3 + 5{i}$ . B. $1 - {i}$ . C. $1 + 3{i}$ .


12. 数学家欧拉发现了复指数函数和三角函数的关系,并给出公式 ${{e}}^{{i}x} = \cos x + {i}\sin x$ ,(其中 ${e}$ 是自然对数的底数, $x \in  \mathbf{R}$ ),这个公式在复变论中有非常重要的地位,被称为 “数学中的天桥”,根据此公式,有下列四个结论, 其中正确的是 ( )

A. ${{e}}^{{i}\pi } - 1 = 0$ . B. $2\cos x = {{e}}^{-{i}x} + {{e}}^{{i}x}$ . C. $2\sin x = {{e}}^{{i}x} - {{e}}^{-{i}x}$ . D. ${\left( \frac{\sqrt{2}}{2} + \frac{\sqrt{2}}{2}{i}\right) }^{2022} =  - 1$ .

\section*{三、解答题}

13. 已知 ${i}$ 为虚数单位.

(1)计算: $\frac{-2\sqrt{3} + {i}}{1 + 2\sqrt{3}{i}} - {\left( \frac{1 + {i}}{\sqrt{2}}\right) }^{2022}$ ;

(2)若 ${\left| z\right| }^{2} + \left( {z + \bar{z}}\right) {i} = 1 - {i}$ ，求复数 $z$ .

14. 已知复数 $z = \left( {{m}^{2} + m - 6}\right)  + \left( {{m}^{2} - {3m} + 2}\right) {i}$ ，其中 $m \in  \mathbf{R}$ .

(1)当 $m$ 取何值时， $z$ 为纯虚数；

(2)如果复数 $z$ 在复平面上对应的点位于第二象限，求实数 $m$ 的取值范围.

15. 已知复数 ${z}_{1} = a + 2 + \left( {{a}^{2} - 3}\right) {i},{z}_{2} = 2 - \left( {{3a} + 1}\right) {i}\left( {a \in  \mathbf{R}}\right)$ .

(1)若复数 ${z}_{1} - \overline{{z}_{2}}$ 在复平面上对应点落在第一象限，求实数 $a$ 的取值范围；

(2)若虚数 ${z}_{1}$ 是实系数一元二次方程 ${x}^{2} - {6x} + m = 0$ 的根，求实数 $m$ 的值.

16. 对于一组复数 ${z}_{1}\text{ 、 }{z}_{2}\text{ 、 }\cdots \text{ 、 }{z}_{n}\left( {n \in  \mathbf{N},n \geq  3}\right)$ ,令 ${S}_{n} = {z}_{1} + {z}_{2} + {z}_{3} + \cdots  + {z}_{n}$ ,如果存在 ${z}_{p}(p \in  \{ 1,2$ , $3,\cdots ,n\} )$ ,使得 $\left| {z}_{p}\right|  \geq  \left| {{S}_{n} - {z}_{p}}\right|$ ,那么称 ${z}_{p}$ 是该复数组的 “ $M$ 复数”.

(1)已知 ${z}_{1} = {i},{z}_{2} = 1 + {i}$ ，是否存在复数 ${z}_{3}$ 使得 ${z}_{1}$ 、 ${z}_{2}$ 、 ${z}_{3}$ 均是复数组 ${z}_{1}$ 、 ${z}_{2}$ 、 ${z}_{3}$ 的“ $M$ 复数”? 若存在， 求出所有的 ${z}_{3}$ ; 若不存在,说明理由;

(2)若 ${z}_{n} = {\left( \frac{5}{9}\right) }^{n - 1} + {i} \cdot  {\left( -1\right) }^{n}\left( {n \in  {\mathbf{N}}^{ * }}\right)$ ，复数组 ${z}_{1}$ 、 ${z}_{2}$ 、 $\cdots$ 、 ${z}_{n}\left( {n \in  \mathbf{N},n \geq  3}\right)$ 是否存在 “ $M$ 复数”? 给出你的结论, 并说明理由.

